% !TeX program = xelatex
% Standalone English volume. Use XeLaTeX until the references stabilize.
\documentclass[a4paper,11pt]{article}
\usepackage[T1]{fontenc}
\usepackage{lmodern}
\usepackage[english]{babel}
\usepackage[left=26mm,right=26mm,top=25mm,bottom=25mm,headheight=15pt,headsep=7mm,footskip=11mm]{geometry}
\usepackage{amsmath,amssymb,mathtools,amsthm}
\newtheorem{proposition}{Proposition}[section]
\newtheorem{lemma}[proposition]{Lemma}
\newtheorem{theorem}[proposition]{Theorem}
\newtheorem{corollary}[proposition]{Corollary}
\newtheorem{maintheorem}[proposition]{Main Theorem}
\theoremstyle{definition}
\newtheorem{definition}[proposition]{Definition}
\theoremstyle{remark}
\newtheorem{remark}[proposition]{Remark}
\makeatletter
\renewcommand*\l@section{\@dottedtocline{1}{0em}{2.4em}}
\renewcommand*\l@subsection{\@dottedtocline{2}{1.45em}{3.45em}}
\renewcommand*\l@subsubsection{\@dottedtocline{3}{3.8em}{4.55em}}
\makeatother
\usepackage{xcolor,fancyhdr,titlesec,needspace,xurl}
\usepackage{longtable,multicol}
\usepackage{hyperref}
\definecolor{ink}{HTML}{1A3542}
\definecolor{muted}{HTML}{4F606A}
\hypersetup{unicode=true,colorlinks=true,linkcolor=ink,filecolor=ink,urlcolor=ink,citecolor=ink,bookmarksnumbered=true,
 pdfdisplaydoctitle=true,pdflang={en},
 pdftitle={The generalized Riemann hypothesis for Dirichlet L-functions. Edition 930, Volume XXIII},
 pdfsubject={A natural-language reconstruction of the proposed proof},
 pdfauthor={A mathematical reconstruction}}
\titleformat{\section}{\Large\bfseries\color{ink}}{\thesection}{.65em}{}
\titlespacing*{\section}{0pt}{.8em}{.65em}
\setlength{\parindent}{1.2em}
\linespread{1.045}
\setlength{\parskip}{.22em plus .05em minus .03em}
\setlength{\emergencystretch}{2em}
\setcounter{topnumber}{2}
\setcounter{tocdepth}{3}\setcounter{secnumdepth}{3}
\widowpenalty=4500
\clubpenalty=4500
\pagestyle{fancy}\fancyhf{}
\fancyhead[L]{\small\color{muted}Dirichlet GRH \textperiodcentered\ Volume I}
\fancyhead[R]{\small\color{muted}Dirichlet zeros and power intervals}
\fancyfoot[C]{\small\thepage}
\renewcommand{\headrulewidth}{.3pt}
\newcommand{\lead}[1]{\par\smallskip\noindent{\bfseries\color{ink}#1.}\enspace}
\newcommand{\sour}[2]{\unskip\penalty150\hspace{.35em}{\small\hyperlink{ref#1}{[#2]}}}
\titleformat{\subsection}{\large\bfseries}{\thesubsection}{.6em}{}
\titlespacing*{\subsection}{0pt}{1.15em plus .15em}{.4em}
\setlength{\abovedisplayskip}{9pt plus 2pt minus 3pt}
\setlength{\belowdisplayskip}{9pt plus 2pt minus 3pt}
\newcommand{\RR}{\mathbb R}\newcommand{\CC}{\mathbb C}\newcommand{\QQ}{\mathbb Q}
\newcommand{\ZZ}{\mathbb Z}\newcommand{\ord}{\operatorname{ord}}
\newcommand{\Frob}{\operatorname{Frob}}\newcommand{\Gal}{\operatorname{Gal}}
\newcommand{\len}{\operatorname{length}}\newcommand{\Res}{\operatorname*{Res}}
\newcommand{\wmean}{\operatorname{av}}\newcommand{\Ker}{\operatorname{ker}}
\newcommand{\Lh}{\mathcal L}\newcommand{\Gg}{\mathcal G}
\newcommand{\real}{\operatorname{Re}}\newcommand{\imag}{\operatorname{Im}}
\newcommand{\abs}[1]{\left|#1\right|}\newcommand{\eps}{\varepsilon}
\newcommand{\anchoritem}[2]{\par\Needspace{2.8\baselineskip}\hypertarget{ref#1}{}\noindent{\bfseries\color{ink}#2.}\enspace}
\newcommand{\srcloc}[2]{{\small\nolinkurl{#1.lean}}, line~#2}
\numberwithin{equation}{section}
\raggedbottom

\newcommand{\Sel}{\mathcal S}
\newcommand{\Cl}{\operatorname{Cl}}
\newcommand{\Aut}{\operatorname{Aut}}
\newcommand{\GL}{\operatorname{GL}}
\newcommand{\OO}{\mathcal O}
\hypersetup{bookmarksdepth=3}


\addto\captionsenglish{\renewcommand{\partname}{Volume}}


\makeatletter
\expandafter\def\csname r@cor:allcost454\endcsname{{141.22}{22}{}{proposition.141.22}{Dirichlet-GRH-v930-Volume-XXVII.pdf}}
\expandafter\def\csname r@cor:move55\endcsname{{65.4}{23}{}{proposition.65.4}{Dirichlet-GRH-v930-Volume-XIII.pdf}}
\expandafter\def\csname r@cor:normaldegree347\endcsname{{127.2}{27}{}{proposition.127.2}{Dirichlet-GRH-v930-Volume-XXIV.pdf}}
\expandafter\def\csname r@cor:original200\endcsname{{83.7}{57}{}{proposition.83.7}{Dirichlet-GRH-v930-Volume-XVI.pdf}}
\expandafter\def\csname r@cor:orth\endcsname{{70.9}{10}{}{proposition.70.9}{Dirichlet-GRH-v930-Volume-XIV.pdf}}
\expandafter\def\csname r@cor:rank111\endcsname{{9.14}{119}{}{proposition.9.14}{Dirichlet-GRH-v930-Volume-I.pdf}}
\expandafter\def\csname r@cor:twomain185\endcsname{{104.15}{13}{}{proposition.104.15}{Dirichlet-GRH-v930-Volume-XX.pdf}}
\expandafter\def\csname r@eq:cov25\endcsname{{70.10}{12}{}{equation.70.10}{Dirichlet-GRH-v930-Volume-XIV.pdf}}
\expandafter\def\csname r@eq:gap\endcsname{{6.3}{26}{}{equation.6.3}{Dirichlet-GRH-v930-Volume-I.pdf}}
\expandafter\def\csname r@eq:gram25\endcsname{{70.9}{12}{}{equation.70.9}{Dirichlet-GRH-v930-Volume-XIV.pdf}}
\expandafter\def\csname r@eq:iter\endcsname{{70.6}{8}{}{equation.70.6}{Dirichlet-GRH-v930-Volume-XIV.pdf}}
\expandafter\def\csname r@eq:letters111\endcsname{{9.1}{117}{}{equation.9.1}{Dirichlet-GRH-v930-Volume-I.pdf}}
\expandafter\def\csname r@eq:sieve\endcsname{{70.2}{4}{}{equation.70.2}{Dirichlet-GRH-v930-Volume-XIV.pdf}}
\expandafter\def\csname r@eq:sub81\endcsname{{66.22}{67}{}{equation.66.22}{Dirichlet-GRH-v930-Volume-XIII.pdf}}
\expandafter\def\csname r@lem:basis59\endcsname{{65.8}{27}{}{proposition.65.8}{Dirichlet-GRH-v930-Volume-XIII.pdf}}
\expandafter\def\csname r@lem:basis63\endcsname{{65.14}{34}{}{proposition.65.14}{Dirichlet-GRH-v930-Volume-XIII.pdf}}
\expandafter\def\csname r@lem:degree181\endcsname{{104.6}{6}{}{proposition.104.6}{Dirichlet-GRH-v930-Volume-XX.pdf}}
\expandafter\def\csname r@lem:dimension239\endcsname{{94.22}{55}{}{proposition.94.22}{Dirichlet-GRH-v930-Volume-XVIII.pdf}}
\expandafter\def\csname r@lem:forms203\endcsname{{110.10}{46}{}{proposition.110.10}{Dirichlet-GRH-v930-Volume-XXI.pdf}}
\expandafter\def\csname r@lem:highestline347\endcsname{{83.27}{74}{}{proposition.83.27}{Dirichlet-GRH-v930-Volume-XVI.pdf}}
\expandafter\def\csname r@lem:hook86\endcsname{{92.4}{5}{}{proposition.92.4}{Dirichlet-GRH-v930-Volume-XVIII.pdf}}
\expandafter\def\csname r@lem:jetweyl453\endcsname{{141.17}{18}{}{proposition.141.17}{Dirichlet-GRH-v930-Volume-XXVII.pdf}}
\expandafter\def\csname r@lem:partner76\endcsname{{66.23}{63}{}{proposition.66.23}{Dirichlet-GRH-v930-Volume-XIII.pdf}}
\expandafter\def\csname r@lem:sphere60\endcsname{{66.8}{51}{}{proposition.66.8}{Dirichlet-GRH-v930-Volume-XIII.pdf}}
\expandafter\def\csname r@lem:split61\endcsname{{65.12}{32}{}{proposition.65.12}{Dirichlet-GRH-v930-Volume-XIII.pdf}}
\expandafter\def\csname r@lem:terminalclass259\endcsname{{111.2}{58}{}{proposition.111.2}{Dirichlet-GRH-v930-Volume-XXI.pdf}}
\expandafter\def\csname r@lem:terminalfourier454\endcsname{{141.20}{20}{}{proposition.141.20}{Dirichlet-GRH-v930-Volume-XXVII.pdf}}
\expandafter\def\csname r@nav:I-28\endcsname{{6.6}{32}{}{subsection.6.6}{Dirichlet-GRH-v930-Volume-I.pdf}}
\expandafter\def\csname r@prop:adjtheta234\endcsname{{110.5}{41}{}{proposition.110.5}{Dirichlet-GRH-v930-Volume-XXI.pdf}}
\expandafter\def\csname r@prop:branch275\endcsname{{94.29}{62}{}{proposition.94.29}{Dirichlet-GRH-v930-Volume-XVIII.pdf}}
\expandafter\def\csname r@prop:cartan244\endcsname{{95.2}{66}{}{proposition.95.2}{Dirichlet-GRH-v930-Volume-XVIII.pdf}}
\expandafter\def\csname r@prop:der95\endcsname{{93.16}{33}{}{proposition.93.16}{Dirichlet-GRH-v930-Volume-XVIII.pdf}}
\expandafter\def\csname r@prop:fibres211\endcsname{{98.20}{20}{}{proposition.98.20}{Dirichlet-GRH-v930-Volume-XIX.pdf}}
\expandafter\def\csname r@prop:flagtrace304\endcsname{{112.10}{83}{}{proposition.112.10}{Dirichlet-GRH-v930-Volume-XXI.pdf}}
\expandafter\def\csname r@prop:index90\endcsname{{94.6}{41}{}{proposition.94.6}{Dirichlet-GRH-v930-Volume-XVIII.pdf}}
\expandafter\def\csname r@prop:minor181\endcsname{{104.8}{7}{}{proposition.104.8}{Dirichlet-GRH-v930-Volume-XX.pdf}}
\expandafter\def\csname r@prop:moments242\endcsname{{94.23}{55}{}{proposition.94.23}{Dirichlet-GRH-v930-Volume-XVIII.pdf}}
\expandafter\def\csname r@prop:normalframe349\endcsname{{104.34}{29}{}{proposition.104.34}{Dirichlet-GRH-v930-Volume-XX.pdf}}
\expandafter\def\csname r@prop:phase301\endcsname{{126.5}{14}{}{proposition.126.5}{Dirichlet-GRH-v930-Volume-XXIV.pdf}}
\expandafter\def\csname r@prop:rank11\endcsname{{6.9}{30}{}{proposition.6.9}{Dirichlet-GRH-v930-Volume-I.pdf}}
\expandafter\def\csname r@prop:ranks47\endcsname{{92.1}{3}{}{proposition.92.1}{Dirichlet-GRH-v930-Volume-XVIII.pdf}}
\expandafter\def\csname r@prop:schurframe259\endcsname{{75.7}{120}{}{proposition.75.7}{Dirichlet-GRH-v930-Volume-XIV.pdf}}
\expandafter\def\csname r@prop:shell263\endcsname{{75.5}{117}{}{proposition.75.5}{Dirichlet-GRH-v930-Volume-XIV.pdf}}
\expandafter\def\csname r@prop:source305\endcsname{{112.3}{76}{}{proposition.112.3}{Dirichlet-GRH-v930-Volume-XXI.pdf}}
\expandafter\def\csname r@prop:split173\endcsname{{72.15}{58}{}{proposition.72.15}{Dirichlet-GRH-v930-Volume-XIV.pdf}}
\expandafter\def\csname r@prop:split233\endcsname{{94.17}{51}{}{proposition.94.17}{Dirichlet-GRH-v930-Volume-XVIII.pdf}}
\expandafter\def\csname r@prop:terminalcore255\endcsname{{95.4}{68}{}{proposition.95.4}{Dirichlet-GRH-v930-Volume-XVIII.pdf}}
\expandafter\def\csname r@prop:torus\endcsname{{70.8}{10}{}{proposition.70.8}{Dirichlet-GRH-v930-Volume-XIV.pdf}}
\expandafter\def\csname r@prop:tracemass301\endcsname{{112.4}{77}{}{proposition.112.4}{Dirichlet-GRH-v930-Volume-XXI.pdf}}
\expandafter\def\csname r@sec:lowerlattice347\endcsname{{83.13}{73}{}{subsection.83.13}{Dirichlet-GRH-v930-Volume-XVI.pdf}}
\expandafter\def\csname r@sec:missing\endcsname{{6.6}{32}{}{subsection.6.6}{Dirichlet-GRH-v930-Volume-I.pdf}}
\expandafter\def\csname r@sec:outsideprimeperiods453\endcsname{{141.6}{17}{}{subsection.141.6}{Dirichlet-GRH-v930-Volume-XXVII.pdf}}
\expandafter\def\csname r@seriesguide\endcsname{{1.1}{4}{}{subsection.1.1}{Dirichlet-GRH-v930-Volume-I.pdf}}
\expandafter\def\csname r@thm:abslow349\endcsname{{104.32}{28}{}{proposition.104.32}{Dirichlet-GRH-v930-Volume-XX.pdf}}
\expandafter\def\csname r@thm:allflags298\endcsname{{111.14}{69}{}{proposition.111.14}{Dirichlet-GRH-v930-Volume-XXI.pdf}}
\expandafter\def\csname r@thm:block47\endcsname{{92.2}{3}{}{proposition.92.2}{Dirichlet-GRH-v930-Volume-XVIII.pdf}}
\expandafter\def\csname r@thm:bothlast311\endcsname{{125.2}{4}{}{proposition.125.2}{Dirichlet-GRH-v930-Volume-XXIV.pdf}}
\expandafter\def\csname r@thm:center\endcsname{{70.7}{9}{}{proposition.70.7}{Dirichlet-GRH-v930-Volume-XIV.pdf}}
\expandafter\def\csname r@thm:class61\endcsname{{65.11}{31}{}{proposition.65.11}{Dirichlet-GRH-v930-Volume-XIII.pdf}}
\expandafter\def\csname r@thm:columns235\endcsname{{109.38}{32}{}{proposition.109.38}{Dirichlet-GRH-v930-Volume-XXI.pdf}}
\expandafter\def\csname r@thm:common111\endcsname{{9.13}{118}{}{proposition.9.13}{Dirichlet-GRH-v930-Volume-I.pdf}}
\expandafter\def\csname r@thm:compact313\endcsname{{95.6}{71}{}{proposition.95.6}{Dirichlet-GRH-v930-Volume-XVIII.pdf}}
\expandafter\def\csname r@thm:coreeuler252\endcsname{{71.2}{14}{}{proposition.71.2}{Dirichlet-GRH-v930-Volume-XIV.pdf}}
\expandafter\def\csname r@thm:covectorpaths352\endcsname{{73.1}{73}{}{proposition.73.1}{Dirichlet-GRH-v930-Volume-XIV.pdf}}
\expandafter\def\csname r@thm:direct76\endcsname{{66.24}{63}{}{proposition.66.24}{Dirichlet-GRH-v930-Volume-XIII.pdf}}
\expandafter\def\csname r@thm:dyad234\endcsname{{110.4}{40}{}{proposition.110.4}{Dirichlet-GRH-v930-Volume-XXI.pdf}}
\expandafter\def\csname r@thm:edge242\endcsname{{67.9}{82}{}{proposition.67.9}{Dirichlet-GRH-v930-Volume-XIII.pdf}}
\expandafter\def\csname r@thm:elliptic312\endcsname{{67.12}{86}{}{proposition.67.12}{Dirichlet-GRH-v930-Volume-XIII.pdf}}
\expandafter\def\csname r@thm:entireinterval310\endcsname{{21.14}{52}{}{proposition.21.14}{Dirichlet-GRH-v930-Volume-IV.pdf}}
\expandafter\def\csname r@thm:first250\endcsname{{71.10}{25}{}{proposition.71.10}{Dirichlet-GRH-v930-Volume-XIV.pdf}}
\expandafter\def\csname r@thm:fixedadic332\endcsname{{83.15}{64}{}{proposition.83.15}{Dirichlet-GRH-v930-Volume-XVI.pdf}}
\expandafter\def\csname r@thm:frob244\endcsname{{102.5}{92}{}{proposition.102.5}{Dirichlet-GRH-v930-Volume-XIX.pdf}}
\expandafter\def\csname r@thm:global175\endcsname{{72.19}{61}{}{proposition.72.19}{Dirichlet-GRH-v930-Volume-XIV.pdf}}
\expandafter\def\csname r@thm:gram66\endcsname{{65.27}{45}{}{proposition.65.27}{Dirichlet-GRH-v930-Volume-XIII.pdf}}
\expandafter\def\csname r@thm:gramheight314\endcsname{{125.4}{7}{}{proposition.125.4}{Dirichlet-GRH-v930-Volume-XXIV.pdf}}
\expandafter\def\csname r@thm:last244\endcsname{{67.11}{84}{}{proposition.67.11}{Dirichlet-GRH-v930-Volume-XIII.pdf}}
\expandafter\def\csname r@thm:lowerlattice347\endcsname{{83.28}{74}{}{proposition.83.28}{Dirichlet-GRH-v930-Volume-XVI.pdf}}
\expandafter\def\csname r@thm:markedcore253\endcsname{{102.7}{95}{}{proposition.102.7}{Dirichlet-GRH-v930-Volume-XIX.pdf}}
\expandafter\def\csname r@thm:markedstrict337\endcsname{{22.11}{70}{}{proposition.22.11}{Dirichlet-GRH-v930-Volume-IV.pdf}}
\expandafter\def\csname r@thm:minors239\endcsname{{83.9}{58}{}{proposition.83.9}{Dirichlet-GRH-v930-Volume-XVI.pdf}}
\expandafter\def\csname r@thm:norm177\endcsname{{72.21}{64}{}{proposition.72.21}{Dirichlet-GRH-v930-Volume-XIV.pdf}}
\expandafter\def\csname r@thm:norm212\endcsname{{98.41}{36}{}{proposition.98.41}{Dirichlet-GRH-v930-Volume-XIX.pdf}}
\expandafter\def\csname r@thm:norm91\endcsname{{94.9}{44}{}{proposition.94.9}{Dirichlet-GRH-v930-Volume-XVIII.pdf}}
\expandafter\def\csname r@thm:one61\endcsname{{65.13}{33}{}{proposition.65.13}{Dirichlet-GRH-v930-Volume-XIII.pdf}}
\expandafter\def\csname r@thm:ordinarymoment304\endcsname{{112.8}{81}{}{proposition.112.8}{Dirichlet-GRH-v930-Volume-XXI.pdf}}
\expandafter\def\csname r@thm:originalgram315\endcsname{{126.1}{8}{}{proposition.126.1}{Dirichlet-GRH-v930-Volume-XXIV.pdf}}
\expandafter\def\csname r@thm:outsidedifferent453\endcsname{{141.18}{18}{}{proposition.141.18}{Dirichlet-GRH-v930-Volume-XXVII.pdf}}
\expandafter\def\csname r@thm:pen185\endcsname{{104.12}{10}{}{proposition.104.12}{Dirichlet-GRH-v930-Volume-XX.pdf}}
\expandafter\def\csname r@thm:periodden312\endcsname{{125.3}{6}{}{proposition.125.3}{Dirichlet-GRH-v930-Volume-XXIV.pdf}}
\expandafter\def\csname r@thm:pointlattice353\endcsname{{99.20}{60}{}{proposition.99.20}{Dirichlet-GRH-v930-Volume-XIX.pdf}}
\expandafter\def\csname r@thm:pol49\endcsname{{64.10}{12}{}{proposition.64.10}{Dirichlet-GRH-v930-Volume-XIII.pdf}}
\expandafter\def\csname r@thm:poscolumn316\endcsname{{126.2}{10}{}{proposition.126.2}{Dirichlet-GRH-v930-Volume-XXIV.pdf}}
\expandafter\def\csname r@thm:proj\endcsname{{70.1}{3}{}{proposition.70.1}{Dirichlet-GRH-v930-Volume-XIV.pdf}}
\expandafter\def\csname r@thm:range239\endcsname{{88.1}{52}{}{proposition.88.1}{Dirichlet-GRH-v930-Volume-XVII.pdf}}
\expandafter\def\csname r@thm:ranges301\endcsname{{112.6}{79}{}{proposition.112.6}{Dirichlet-GRH-v930-Volume-XXI.pdf}}
\expandafter\def\csname r@thm:ranktrace302\endcsname{{112.7}{80}{}{proposition.112.7}{Dirichlet-GRH-v930-Volume-XXI.pdf}}
\expandafter\def\csname r@thm:signedfibers331\endcsname{{99.12}{53}{}{proposition.99.12}{Dirichlet-GRH-v930-Volume-XIX.pdf}}
\expandafter\def\csname r@thm:splitmass323\endcsname{{126.4}{13}{}{proposition.126.4}{Dirichlet-GRH-v930-Volume-XXIV.pdf}}
\expandafter\def\csname r@thm:terminaldegrees454\endcsname{{141.21}{21}{}{proposition.141.21}{Dirichlet-GRH-v930-Volume-XXVII.pdf}}
\expandafter\def\csname r@thm:terminalformula347\endcsname{{127.1}{26}{}{proposition.127.1}{Dirichlet-GRH-v930-Volume-XXIV.pdf}}
\expandafter\def\csname r@thm:third263\endcsname{{74.6}{107}{}{proposition.74.6}{Dirichlet-GRH-v930-Volume-XIV.pdf}}
\expandafter\def\csname r@thm:torsion312\endcsname{{66.35}{74}{}{proposition.66.35}{Dirichlet-GRH-v930-Volume-XIII.pdf}}
\expandafter\def\csname r@thm:transport233\endcsname{{94.16}{50}{}{proposition.94.16}{Dirichlet-GRH-v930-Volume-XVIII.pdf}}
\expandafter\def\csname r@thm:transport59\endcsname{{65.9}{28}{}{proposition.65.9}{Dirichlet-GRH-v930-Volume-XIII.pdf}}
\expandafter\def\csname r@thm:unique221\endcsname{{109.15}{13}{}{proposition.109.15}{Dirichlet-GRH-v930-Volume-XXI.pdf}}
\expandafter\def\csname r@thm:vac242\endcsname{{67.8}{81}{}{proposition.67.8}{Dirichlet-GRH-v930-Volume-XIII.pdf}}
\expandafter\def\csname r@thm:vectorchain353\endcsname{{73.2}{74}{}{proposition.73.2}{Dirichlet-GRH-v930-Volume-XIV.pdf}}
\expandafter\def\csname r@thm:zerotrace303\endcsname{{112.2}{75}{}{proposition.112.2}{Dirichlet-GRH-v930-Volume-XXI.pdf}}
\makeatother
\begin{document}
\setcounter{part}{22}
\fancyhead[L]{\small\color{muted}Dirichlet GRH \textperiodcentered\ Volume XXIII}
\fancyhead[R]{\small\color{muted}Coefficient lattices and period tests}
\thispagestyle{plain}\begin{center}
{\LARGE\bfseries The generalized Riemann hypothesis\\[2pt]for Dirichlet \(L\)-functions}\\[6pt]
{\small 27 November 2026 \quad\textperiodcentered\quad Edition 930}\end{center}
\part{Integral coefficient lattices and tests for periods}\label{part:XXIII}
\begin{abstract}\noindent We construct integral tests for the projected period. Branching between Weyl lattices and elementary divisors of the coefficient tables control the finite projections on the marked rows.\end{abstract}
\setcounter{section}{118}
\noindent The \hyperref[seriesguide]{guide in Volume~I} lists all twenty-seven volumes.
\begingroup\small\setlength{\parskip}{0pt}\tableofcontents\endgroup
\clearpage

\section{Refined columns and integral class pairings}\label{nav:chapter-refined-rational-columns}\label{nav:volume-XXIII}
\subsection{When two refined columns contribute to the trace}\label{nav:IX-59}
Fix a supported transverse label with \(\mathcal G>0\), a
coarse rational row cell \(tK^*\), and a supporting fine
representative \(t\), as in Proposition~\ref{prop:fibres211}.
Index the two flags by \(i=0,1\). Their coarse column
representatives are \(y_i\in H_f\). At the places \(v\mid Q\)
write \(u_{i,v}\) for the middle blocks of the changes of
frame in that proposition and \(q_v\) for the residue
cardinality. Here \(U_k\) denotes the upper unitriangular
group over the residue field.

Let \(K_{\rm line}^*\) be the coarse column level obtained
from \(K_H^*\) by allowing the nonzero covector at each
\(v\mid T\) to be multiplied by a residue unit. It preserves
the parent lattice and the kernel of each terminal covector.
Put
\[
 {\cal A}_t=\{\alpha\in H(F):y_1^{-1}\alpha y_0\in K_{\rm line}^*\},
 \qquad A_{\alpha,v}=u_{1,v}^{-1}
       \overline{(y_1^{-1}\alpha y_0)_v}\,u_{0,v}\quad(v\mid Q).
\]
The bar here means reduction in the first split factor.
For \(A\in\GL_k(\mathbb F_q)\) let \(P_k(A;q)\) be the fraction
of the cosets \(zU_k\) for which \(z^{-1}Az\in U_k\).

\begin{proposition}[An exact count of scalar-orbit matches]\leavevmode
\label{prop:loops302}
The fraction of common fine refinements whose column classes
belong to the same terminal scalar orbit is
\[
                    \sum_{\alpha\in{\cal A}_t}
                            \prod_{v\mid Q}P_k(A_{\alpha,v};q_v).
\]
Every matching refinement has a unique rational comparison
\(\alpha\) in the chosen coordinates. One may restrict the
sum by any condition on the conjugacy classes of the
\(A_{\alpha,v}\). The resulting fraction is also the
fraction of the finite covolume over this cell satisfying
the condition. Uniformly in the cell,
\(\log\max(1,|{\cal A}_t|)\ll fk\log^2(2fk)\).
\end{proposition}
\begin{proof}
The terminal scalar action multiplies the covectors independently
at the places above \(T\). Its inverse image in the parent compact
is exactly the stabilizer of their lines. Indeed the stabilizer of
a nonzero covector has quotient of order \(T-1\) in the stabilizer
of its line, with every unit realized by a scalar matrix. Hence
two fine column classes have the same scalar orbit precisely
when a rational isometry carries the first to the second
at the fine level above \(Q\), at the line level above \(T\),
and at the given level elsewhere.

By Proposition~\ref{prop:fibres211} the common parameters
are \(z=(z_vU_k)_{v\mid Q}\), and the two column representatives
are obtained from \(y_i\) by right multiplication by integral
lifts of the \(u_{i,v}z_v\) at \(Q\). Thus the condition at
\(Q\) for a given \(\alpha\in{\cal A}_t\) is exactly
\(z_v^{-1}A_{\alpha,v}z_v\in U_k\) at every marked place.
The fraction for this \(\alpha\) is the product in the assertion.

If \({\cal A}_t\) is nonempty, composition with one of its
members identifies it with the rational stabilizer of a coarse
projective column. This group embeds in the automorphism group
of the positive parent lattice and hence is finite. The
underlying integral trace lattice has rank \(2fk\). The finite
linear group bound used in Theorem~\ref{thm:norm212} therefore
gives the asserted order bound also at this projective level.
Two comparisons for the same \(z\) would differ by a rational
automorphism of its fine projective column. At each marked
prime \(g\) this automorphism belongs, up to conjugacy, to the
inverse image of \(U_k(\mathbb F_g)\) in \(\GL_k(\ZZ_g)\).
Every finite-order element of this inverse image has order a
power of \(g\). Indeed all its reductions modulo \(g^n\)
lie in finite \(g\)-groups and their kernels intersect in the
identity. The embeddings into the local linear groups are
faithful. Applying this at the two distinct primes \(g_1,g_2\)
forces the automorphism to be the identity. The sets of
parameters counted for distinct \(\alpha\) are disjoint.

For clarity, the class conditions used here are intrinsic.
The two middle blocks identify a common residue quotient of the
row with the two column coordinates. A change of column frames
transports the rational maps with those frames. A change of the
common row frame conjugates their residue endomorphisms.
In particular a rational row stabilizer permutes the parameters
and transports each comparison with its conjugacy data.
By Proposition~\ref{prop:fibres211} the quotient of the coarse
row stabilizer by the common unipotent stabilizer acts freely
on the parameters. Every fine orbit has the same size and the
same finite covolume. Its size cancels in the fraction, which
also proves the assertion with a class condition.
\end{proof}

\subsection{Integral cross classes in the remaining height range}\label{nav:IX-60}
Fix a supported transverse rational label with \(\mathcal G>0\)
and \(H_*<b^{Mfd}\) as in Theorem~\ref{thm:ranges301}. Write
\(M_0\) for its monomial cross matrix from
Lemma~\ref{lem:forms203}. Its blocks have sizes \(h_i\),
with a single pivot scalar \(c_i\) times a signed reversal
in each block. An integral cross representative \(Z\) is the
cross matrix of two primitive integral frames for this label.
Changing the frames acts on the left by integral lower
unitriangular matrices and on the right by integral upper
unitriangular matrices.

\begin{lemma}[Denominators and the number of integral cross classes]\leavevmode
\label{lem:classes304}
There exists a nonzero \(a\in\OO_E\) depending only on the
rational label, with
\(\log |N_{E/\QQ}a|\ll fd^3\log b\), such that every
representative has the form \(Z=L M_0 V\). Here \(L\) is
lower block unitriangular with identity diagonal blocks,
\(V\) is upper unitriangular, and \(aL,aV\) are integral.
The number of integral cross classes for this label is at
most \(\exp(Cfd^6\log b)\).
\end{lemma}
\begin{proof}
Rational changes of the frames give a lower and an upper
unitriangular factor on either side of \(M_0\). Factor
the lower matrix into a lower block matrix with identity
diagonal blocks times its block diagonal part. Conjugating
the latter past the reversals turns it into an upper
unitriangular matrix, which can be absorbed into \(V\).
Let \(P\) reverse the rows within each block. Then
\(P Z=(P L P^{-1})(P M_0)V\) is an ordinary triangular
decomposition with diagonal middle matrix \(P M_0\).
For \(1\le p\le d\), the leading \(p\)-minor \(\omega_p\)
of \(P Z\) is the product of its first \(p\) pivots. It
is determined by the label. It is a nonzero algebraic
integer whenever an integral representative exists. We
may assume one exists.

For \(i>j\), the entry at \((i,j)\) of the lower factor
equals the minor in rows \(1,\ldots,j-1,i\) and columns
\(1,\ldots,j\) divided by \(\omega_j\). The corresponding
formula for an entry at \((j,i)\) of the upper factor
uses rows \(1,\ldots,j\) and columns \(1,\ldots,j-1,i\).
These formulas follow by taking determinants of the
triangular factors in the indicated rows and columns.
Thus \(a=\prod_{p=1}^d\omega_p\) clears both factors.
Within a block every pivot has modulus \(|\tau c_i|\)
at any one complex embedding. Hence \(\log|\tau\omega_p|\)
interpolates linearly between its values at the two
endpoints, starting at zero. At the proper endpoints
the original determinants are \(Q^{j_i}a_i\), up to sign.
At the last endpoint \(D_d=Q^d a_t-TD_s\), with \(D_s=0\)
or \(D_s=Q^s a_{t-1}\). From \(H_i\le H_*^{j_i}\),
\(\log T=O(d\log b)\) and \(\log Q=O(\log b)\), the sum
of the positive parts of the logarithms at any endpoint
over the \(f\) complex embeddings is \(O(fd^2\log b)\).
For the last endpoint use \(|v+w|\le
2\max(1,|v|)\max(1,|w|)\). Convexity of the positive part
gives the same bound at every \(p\). Taking both conjugate
embeddings and summing over \(p\) proves the norm bound.

Record each free entry of \(L,V\) in
\(a^{-1}\OO_E/a^{d-1}\OO_E\), a set of size \(|N_{E/\QQ}a|^d\).
If two records agree, the two lower factors differ by
left multiplication by an integral lower unitriangular
matrix. Indeed the entrywise difference is in
\(a^{d-1}\OO_E\), and inversion of a \(d\)-square
unitriangular matrix with denominator \(a\) costs at
most \(a^{d-1}\). The two upper factors likewise differ
by an integral right multiplication. The cross matrices
are therefore in one integral class. There are at most
\(d^2\) recorded entries, which proves the count.
\end{proof}

\subsection{The conjugates of a prescribed signed class sum}\label{sec:fixedcolumn332}
Fix \(x,E_x,\xi\) as in Section~\ref{sec:missing}. At each
\(v\in Q(\xi)\) prescribe a nontrivial residue character \(\chi_v\),
a row \(l_v\), and a primitive frequency \(\nu_v\) for the signed
integral of Theorem~\ref{thm:signedfibers331}. A rational class may
determine positions at several of these places at once. A common
cyclotomic action also changes frequencies simultaneously at places
above the same rational prime. Both constraints enter the moment
at a prescribed tuple.

\begin{proposition}[The joint character at the prescribed tuple]\leavevmode
\label{prop:jointcharacter332}
Put \(G=\prod_{v\in Q(\xi)}(\mathbb F_{p_v}^\times\times\ZZ/p_v^s\ZZ)\).
On a product of the prescribed frequency balls consider an integrable
weight, or an absolutely summable discrete class weight on the
supported positions. Let \(H(g)\) be its total weight in the joint
fiber \(g=((b_v,y_v))_v\), using the original measure or the original
discrete weights. Its integral or sum against the product of local
kernels is
\[
   \left(\prod_v c_{\chi_v}\chi_v(\bar\nu_v)^{-1}\right)
            \sum_{t\bmod L} C_t\zeta^t ,
 \qquad
 L=\mathop{\rm lcm}_{v}\{p_v^s,\operatorname{ord}(\chi_v)\},
\]
where \(\zeta\) is a fixed primitive \(L\)-th root. More precisely,
\[
 \lambda(g)=\prod_v\chi_v(b_v)\psi_v(p_v^{-s}\widetilde\nu_v\widetilde y_v),
 \qquad
 C_t=\sum_{\lambda(g)=\zeta^t}H(g).
\]
For the empty tuple take \(L=1\). For every tuple each fiber of
\(\lambda\) has \(|G|/L\) elements.
\end{proposition}
\begin{proof}
On each joint fiber the local scalar has the value of
Theorem~\ref{thm:signedfibers331}. Integration over a finite
partition, or regrouping an absolutely convergent series, gives the
formula with \(H(g)\). Haar fiber volumes, class multiplicities and
stabilizers are already included in \(H\). The character on each
additive factor has order \(p_v^s\) because \(\nu_v\) is a unit.
The restrictions to the other factors have orders
\(\operatorname{ord}(\chi_v)\). The complete product group contains
each factor separately, so its image under \(\lambda\) is the finite
cyclic group of the asserted order. All of its fibers are cosets
of the kernel. This argument imposes no invariance on the weights
within those cosets.
\end{proof}

For example, at two places over the same odd \(p\), choose the
additive characters to give the same primitive \(p^s\)-th root.
Give weight one to \(b_1=b_2=1\), \(y_1=y_2\), and zero elsewhere.
After removing the displayed row scalar, the value at
\((\nu_1,\nu_2)=(1,-1)\) is \(p^s\), whereas at \((1,1)\) it is
\(\sum_y\psi_p(2p^{-s}y)=0\). These are two different rows on the
same joint array.

For the arithmetic column of Theorem~\ref{thm:column242} take its
normalized scalar \(Z=\mathcal A\), optionally multiplied by the
denominator of Theorem~\ref{thm:finitecolumn313}. Under its
arithmetic hypotheses the finite class expansion, including the
character projector, has algebraic coefficients and CM derivatives.
Theorem~\ref{thm:classden312} gives the roots of unity in its
finite translations. Take a common multiple \(L\) of the orders of
the chosen finite character and additive phases, or \(L=1\) for an
empty list. In each term write their product as \(\zeta^t\) and sum
the remaining signed coefficients into \(C_t\), retaining all original
divisors. This gives \(Z=\sum_t C_t\zeta^t\) directly on the arithmetic
side of the comparison in Section~\ref{sec:missing}.

\begin{theorem}[Second moment and nonvanishing for a fixed algebraic scalar]\leavevmode
\label{thm:scalarorbits332}
Let \(Z=\sum_t C_t\zeta^t\) be either of the preceding sums with
algebraic \(C_t\), with any common algebraic scalar included in the
coefficients if desired. Let \(F=\QQ(\zeta)\), let \(B\) be a number
field containing \(F\) and all the \(C_t\), let \(u_1,\ldots,u_d\)
be an \(F\)-basis of \(B\), and put \(n=\varphi(L)\). Write uniquely
\[
    C_t=\sum_{i=1}^d u_i\sum_{j=0}^{n-1}q_{ijt}\zeta^j,\qquad
    H_i(k)=\sum_{\substack{j,t\\t+j\equiv k\bmod L}} q_{ijt},
               \qquad q_{ijt}\in\QQ .
\]
For a rational prime \(p\mid L\) let \(E_p\) average the translates
of an array on \(\ZZ/L\ZZ\) by the subgroup generated by \(L/p\).
Put \(P=\prod_{p\mid L}(1-E_p)\), with \(P=1\) for \(L=1\).
Then \(Z=0\) if and only if \(P H_i=0\) for every \(i\).

For \(a\in(\ZZ/L\ZZ)^\times\) put \(b_i(a)=\sum_k H_i(k)\zeta^{ak}\).
Define the positive definite Hermitian matrix
\[
   G_{ih}(a)=d^{-1}\sum_{\sigma:\,\sigma(\zeta)=\zeta^a}
                         \overline{\sigma(u_i)}\,\sigma(u_h),
       \qquad \sigma:B\hookrightarrow\CC .
\]
The conjugate second moment of this one scalar is
\begin{equation}\label{eq:fixedmoment332}
  \frac{1}{[B:\QQ]}\sum_\sigma |\sigma Z|^2
        =\frac1n\sum_a b(a)^*G(a)b(a).
\end{equation}
Writing \(\lambda_-,\lambda_+\) for the least and greatest
eigenvalues over all these matrices, the moment lies between
\(\lambda_- (L/n)\sum_i\|PH_i\|_2^2\) and
\(\lambda_+ (L/n)\sum_i\|PH_i\|_2^2\), with counting measure on
\(\ZZ/L\ZZ\). For rational \(C_t\) take \(B=F,u_1=1\).
In this case the moment is exactly \((L/n)\|PC\|_2^2\).
\end{theorem}
\begin{proof}
The powers of \(\zeta\) up to degree \(n-1\) are a rational basis
of \(F\), by irreducibility of the cyclotomic polynomial
\(\Phi_L\). Substitution gives \(Z=\sum_i u_i b_i(1)\).
Independence over \(F\) shows that \(Z\) vanishes exactly when
all \(b_i(1)\) vanish. As the \(H_i(k)\) are rational, their values
at the primitive frequencies are precisely the Galois conjugates
of \(b_i(1)\). Hence vanishing at one primitive frequency is
equivalent to vanishing at all of them. On a character of frequency
\(a\) the multiplier of \(E_p\) is \(p^{-1}\sum_{j=0}^{p-1}
\exp(2\pi\mathrm{i}aj/p)\), up to replacing the primitive root by
a conjugate. This equals one for \(p\mid a\) and zero otherwise.
The commuting projections therefore leave in \(P\) exactly the
frequencies prime to \(L\). This proves the vanishing test and,
by character orthogonality on \(\ZZ/L\ZZ\), gives
\(\sum_a |b_i(a)|^2=L\|PH_i\|_2^2\), the sum being over units.

Every embedding of \(F\) has \(d\) extensions to \(B\). At such
an extension \(\sigma Z=\sum_i\sigma(u_i)b_i(a)\), since all
coordinates in the \(H_i\) are rational. Expanding the squared
modulus proves \eqref{eq:fixedmoment332}. For a fixed \(a\), the
matrix with entries \(\sigma(u_i)\) for its \(d\) extensions is
invertible. For a power basis of a primitive element of \(B/F\)
this is a Vandermonde determinant on distinct images, since the
extension is separable. A change to the given basis preserves
invertibility. This proves positive definiteness of \(G(a)\).
Its two eigenvalue bounds and the orthogonality identity give
the inequalities. When \(B=F,u_1=1\), every \(G(a)\) equals one.
For \(L=1\) the unit group and frequency set each contain a single
element and the same proof applies.
\end{proof}

The test uses rational sums and differences on the original class
coefficients after expansion in the stated basis. For rational
\(C_t\) in particular, if every term occurs with equal total weight
in each translate by \(L/p\) for one \(p\mid L\), the designated
sum vanishes. For general algebraic coefficients the rational arrays
\(H_i\) incorporate the action on the coefficients as well as on
their displayed phases. Theorem~\ref{thm:fixedadic332} gives the
finite valuations of the same \(Z\).

\begin{corollary}[A lower bound from the orders of the occupied phases]\leavevmode
\label{cor:separatedphases332}
For each rational array \(H_i\) in Theorem~\ref{thm:scalarorbits332}
write \(S_i\) for its support. For \(t\ne u\bmod L\) write
\(d_{tu}=L/\gcd(L,t-u)\). Let \(\delta\) be the maximum, over all indices
\(i\) and all \(t\in S_i\), of
\[
   \sum_{\substack{u\in S_i\\u\ne t}}
               \frac{|\mu(d_{tu})|}{\varphi(d_{tu})},
\]
with maximum zero for empty supports. If \(\delta<1\), the moment
in \eqref{eq:fixedmoment332} is bounded below by
\(\lambda_-(1-\delta)\sum_{i,t}|H_i(t)|^2\).
In particular it is positive whenever one array is nonzero.
For rational \(C_t\) and \(L=p^s\), without a sparsity hypothesis,
the moment is exactly
\[
 \frac{p}{p-1}\sum_{v\bmod p^{s-1}}\sum_{u=0}^{p-1}
       |C_{v+up^{s-1}}-\overline C_v|^2,\qquad
 \overline C_v=p^{-1}\sum_{u=0}^{p-1}C_{v+up^{s-1}} .
\]
\end{corollary}
\begin{proof}
For fixed \(t,u\) the mean of \(\zeta^{a(t-u)}\) over units \(a\bmod L\)
equals \(\mu(d)/\varphi(d)\), where \(d=d_{tu}\), also taking
\(d=1\) on the diagonal. Indeed reduction of units modulo \(d\) is
onto by the Chinese remainder theorem, with equal fiber sizes. The
sum of the primitive \(d\)-th roots is \(\mu(d)\), by M\"obius
inversion of the sum of all \(d\)-th roots. Expand the mean of
\(|b_i(a)|^2\) into pairs \(t,u\). Its diagonal is
\(\sum_t H_i(t)^2\). Using
\(2|H_i(t)H_i(u)|\le H_i(t)^2+H_i(u)^2\) on each unordered
off-diagonal pair bounds the absolute value of the remaining sum
by \(\delta\sum_t H_i(t)^2\). The lower bound now follows from the
positive matrices in \eqref{eq:fixedmoment332}. For \(L=p^s\) the
operator \(P=1-E_p\) subtracts exactly the indicated progression
mean and \(L/\varphi(L)=p/(p-1)\). This gives the formula even
if entire progressions occur in the support.
\end{proof}


\subsection{The unit normal at the two marked levels}\label{sec:markedline350}

Fix a period field \(E=E_x\), a word, a translate and the prime
\(T=\pi\bar\pi\) of Theorem~\ref{thm:markedstrict337}. Use the
base auxiliary tuple, whose extra set of frames is empty. Write
\(n=r+mt\) for the smaller terminal rank and \(N=n+1\).
Let \(S_T\) consist of the primes \(P\) of \(E\) dividing \(\pi\).
Write \(k_v=\mathcal O_E/v\) and \(A_T=\prod_{P\in S_T}k_P^\times\).
The two marked rational primes \(g_1,g_2\) are different from \(T\).
In particular \(\pi\) is a unit at every marked frame \(w\) over
either \(g_i\).

In the small rank, an object is a unimodular Hermitian lattice
\(M\) in the prescribed definite rational space together with
residue frames \(\beta_w:M/wM\longrightarrow k_w^n\), taken
modulo the specified groups \(J_w\). In the ordinary endpoint,
\(J_w\) is block upper triangular with blocks \(r,t,\ldots,t\).
Its first diagonal block is upper unitriangular. Its level
character is the given additive character of the sum of the first
superdiagonal entries in this block times the assigned determinant
characters of the \(t\)-blocks. More generally take a subgroup
of \(\mathrm{GL}_r(k_w)\) with its specified character on the
first block. Taking the prescribed symmetric subgroup in its chosen
order gives the symmetric pair of endpoints. We use the following
arguments for each pair with matching first block and character.

Write \(e\) for the last standard unit vector. The lift \(i(M)\)
has lattice \(M\perp\mathcal O_Ee\), marked frames
\(\beta_w^+(u+ye)=(\beta_w u,\pi^{-1}y)\) after reduction, and
covectors \(f_P(u+ye)=\bar\pi y\bmod P\) on its residues.
The terminal subgroup at \(w\) preserves the extended block flag
and fixes its last basis column. It is exactly
\(\operatorname{diag}(J_w,1)\). Indeed block upper triangularity
annuls the last row to the left of its final entry, and the
fixed-column condition gives that last entry and all entries above
it. The character restricts to the small character, since its
formula uses only the original blocks. These descriptions can
also be found in \srcloc{Formalism/LocalToGlobal/Columns/StdLatUnimodular}{168},
\srcloc{Formalism/LocalToGlobal/Columns/ColForms}{92} and
\srcloc{Formalism/LocalToGlobal/Columns/LevGSym}{81}.
A morphism between marked lattices is a rational unitary isometry
carrying one lattice onto the other, respecting the prescribed
orbits of frames and carrying each given covector to the other.
Elements of \(A_T\) act by multiplying the individual covectors.

\begin{lemma}[Recovery of a unit normal from a residue frame]\leavevmode
\label{lem:markedline350}
Let \(M'\perp\mathcal O_Ee\) be such a lattice. A lattice vector
\(v\) of Hermitian length one whose last coordinate agrees with
that of \(e\) modulo one \(w\mid g_i\) equals \(e\), provided
\(g_i-1>[E:\mathbb Q]\). This inequality holds on the intervals
of Theorem~\ref{thm:terminalformula347} for all sufficiently large \(x\).
\end{lemma}
\begin{proof}
Write \(v=u+\alpha e\) with \(u\in M'\) and \(\alpha\in\mathcal O_E\).
The congruence \(\alpha\equiv1\bmod w\) gives \(\alpha\ne0\).
At every complex embedding, orthogonality and positive definiteness
give \(|\sigma\alpha|\le1\). The absolute norm of \(\alpha\) is a
positive integer, so the product of these absolute values is one
and each equals one. The orthogonal length identity at one embedding
then forces \(u=0\). The same bound holds for every power of
\(\alpha\). The characteristic polynomials of these powers in
\(E\) have bounded integer coefficients by the elementary
symmetric function formula. Finitely many roots occur among this
finite set of polynomials, so two powers agree and \(\alpha\) is
a root of unity. Its order \(h\) is prime to \(g_i\), since
otherwise \(E\) would contain a primitive \(g_i\)-th root, of
degree \(g_i-1\). If \(\alpha\ne1\), the geometric sum
\(1+\alpha+\cdots+\alpha^{h-1}\) is zero. Its residue at \(w\)
is \(h\ne0\), a contradiction. This proves the assertion.
\end{proof}

\begin{proposition}[All scalar translates of the lifted classes]\leavevmode
\label{prop:scalarlifts350}
For small objects \(a,b\) and \(c,d\in A_T\), a marked isometry
from \(c\,i(a)\) to \(d\,i(b)\) exists only if \(c=d\).
For \(c=d\), extension by the identity on \(e\) gives a bijection
from the isometries \(a\longrightarrow b\) to those
\(c\,i(a)\longrightarrow c\,i(b)\). In particular the large
class set contains \( |A_T| \) disjoint copies of the full small
class set with its stabilizer orders.
\end{proposition}
\begin{proof}
Apply a putative isometry \(\gamma\) to \(e\). Its marked residue
in the target frame equals \(\pi^{-1}e_N\) at every \(w\), since
the change of frame fixes \(e_N\). Its last lattice coordinate
is therefore congruent to one. Lemma~\ref{lem:markedline350}
gives \(\gamma e=e\). The orthogonal complement of \(e\) and
its intersection with the lattice must now go to the corresponding
complement and intersection. On the small lattice the marked
relations are exactly the changes in \(J_w\), by the displayed
formula for the frames. Evaluating the terminal covector relation
at \(e\) gives \(c_P\bar\pi=d_P\bar\pi\bmod P\) for every \(P\).
Here \(\bar\pi\) is a unit, proving \(c=d\). Conversely every
small isometry extends by the identity and respects all the
indicated data. The residue changes, if recorded as part of a
morphism, are determined uniquely because each frame is surjective.
Applying this bijection with \(a=b\) identifies the stabilizers.
\end{proof}

\subsection{Simultaneous restriction of terminal class values}\label{sec:allclasses350}

Let \(\mathcal H_N\) be the space of all large coefficient tables
at the matched level above, and \(\mathcal H_n\) the small table on
the output side of the restriction. Thus each value transforms by
the coefficient representation and level character, and the
bilinear testing space uses the dual action and inverse character.
At each complex embedding use the normalized Fischer forms of
Section~\ref{sec:lowerlattice347}. On each space use the sum of
the squared norms of its class values with weight \(1/|\Gamma_a|\),
where \(\Gamma_a\) is the corresponding stabilizer.
These give the finite class measures of Theorem~\ref{thm:classtrace326}.
Write \(L=\bigotimes_\tau L_\tau\) for standard lowering,
\(s=\prod_\tau c_\tau^2\kappa_\tau>0\), and \(L_z=ZL\) in the
rational frames, where \(Z\) is given in Corollary~\ref{cor:normaldegree347}.
For \(c\in A_T\) define \((R_cP)(a)=L_z P(c\,i(a))\).
The smaller block equivariance and the matching characters make this
an element of \(\mathcal H_n\). We write \(D=|Z|^2s\) at the
chosen complex embedding.

\begin{theorem}[Restriction of all class values at once]\leavevmode
\label{thm:allclasses350}
For \(c,d\in A_T\) one has \(R_cR_d^\dagger=\delta_{cd}D\,I\)
on \(\mathcal H_n\), with \(\dagger\) denoting Hermitian adjoint.
Consequently the map \(P\mapsto(R_cP)_{c\in A_T}\) is onto the
direct sum of these small tables. Its minimum-norm right inverse
is defined over the coefficient field and preserves the specified
integral coefficient lattices at all primes above \(g_1,g_2,T\)
simultaneously.
\end{theorem}
\begin{proof}
In rational polynomial frames put \(B=L^*/s\), where the star
denotes the rational Fischer adjoint. Its complexification equals
the Hermitian adjoint divided by \(s\), since the coefficients
of \(L\) and the two Gram forms are rational and real. By
Theorem~\ref{thm:lowerlattice347}, \(LB=I\). To lift a family
\((Y_c)_c\), assign \(Z^{-1}B Y_c(a)\) to \(c\,i(a)\) and zero
on all the other classes. The operator \(B\) intertwines the
small block action by that same theorem. Proposition~\ref{prop:scalarlifts350}
identifies every stabilizer with its small block, including its
character. Thus the assignment respects stabilizers and extends
by equivariance to the entire large table. The proposition gives
disjoint supports for the different \(c\). Equality of mass
weights shows that \(R_c^\dagger\) inserts \(\bar ZL^*\) on the
\(c\)-copy and zero elsewhere. Hence the asserted operator product
follows from \(LL^*=sI\). The constructed inverse is the adjoint
of the joint map divided by \(D\). Its image is orthogonal to the
kernel, which proves minimality.

For integrality, put \(z_\tau=\tau(\bar\pi)\) and
\(\delta_\tau=|\lambda_\tau|-|\mu_\tau|\), using the sums of the
actual weights with their determinant terms. In the image of
\(B_\tau=L_\tau^*/(c_\tau^2\kappa_\tau)\) the normal factor
\(\operatorname{diag}(I_n,v)\) acts by \(v^{\delta_\tau}\).
Indeed the smaller block acts by its prescribed type and the
large center has weight \(|\lambda_\tau|\). Applying the small
center with its inverse gives this normal action.
Choose an integral frame \(G'\) of \(M\) at a coefficient place,
at each embedding. In the large dual lattice the frame is
\(G=\operatorname{diag}(G',z_\tau^{-1})\), by
Proposition~\ref{prop:normalframe349}. Scaling \(f_P\) by a unit
preserves its kernel and hence preserves this lattice for every
\(c\). With the transpose action on polynomial covectors, small
block equivariance gives
\[
 \rho_{\lambda_\tau}(G^{\mathsf t})\,z_\tau^{\delta_\tau}B_\tau
       = B_\tau\,\rho_{\mu_\tau}((G')^{\mathsf t}).
\]
In integral frames the right inverse is therefore the tensor
product of the standard \(B_\tau\). Each preserves the two flag
lattices at the stated primes by Theorem~\ref{thm:lowerlattice347}
and the cutoff in Theorem~\ref{thm:terminalformula347}. This
proves simultaneous integrality for the very same inverse,
including primes at which \(Z^{-1}\) has a denominator in the
rational frames.
\end{proof}

\begin{theorem}[The last restriction on the separately zero-sum table]\leavevmode
\label{thm:augmented350}
For \(P\in S_T\) let \(U_h\) denote the action of \(h\in k_P^\times\)
by translating covectors on large tables. In the complex table let
\[
 E^\circ=\prod_{P\in S_T}
     \left(I-\frac{1}{|k_P^\times|}\sum_{h\in k_P^\times}U_h\right),
 \qquad d_T=\prod_{P\in S_T}(1-|k_P^\times|^{-1}).
\]
Then \(R_1E^\circ R_1^\dagger=d_TD\,I\). On its separately
zero-sum domain every nonzero singular value of \(R_1\) is
\(\sqrt{d_TD}\), and its image is the full small table.
For the integral lattice defined by vanishing of the unnormalized
sum in each coordinate, \(R_1\) has a single right inverse
integral above \(g_1,g_2,T\). With \(f=|S_T|\) the archimedean
loss satisfies \(0\le-\log d_T\le f/(T-2)\).
\end{theorem}
\begin{proof}
Translation permutes the large classes with their stabilizers and
preserves their coefficient frames. It is unitary. The factors of
\(E^\circ\) are commuting orthogonal projections, since each group
average projects onto its invariants. On the union of the copies
of Proposition~\ref{prop:scalarlifts350} they act by translation
in the separate coordinates of the finite group \(A_T\). A vector
supported on the identity copy returns there after projection with
coefficient \(d_T\), the product of the identity coefficients in
the factors. The formula follows from
Theorem~\ref{thm:allclasses350}. The same identity gives the singular
values and surjectivity since \(T>2\). All the residue fields
at \(S_T\) have order \(T\), and
\(-\log(1-1/(T-1))=\log(1+1/(T-2))\le1/(T-2)\).

At each \(P\) choose \(h_P\ne1\). For an integral small table
\(Y\), first use the section of Theorem~\ref{thm:allclasses350}
on the identity copy alone. Apply \(\prod_P(I-U_{h_P})\) to this
lift. Summing over any one coordinate kills its difference factor.
On the identity copy the original value survives with coefficient
one by disjointness of the copies. Every translation preserves the
integral lattice and every coefficient used is an integer. This
gives the claimed section even at primes dividing \(T-1\).
\end{proof}

Theorems~\ref{thm:allclasses350} and~\ref{thm:augmented350}
compute the last map on all admissible class values. In the
period of Section~\ref{nav:I-28} the input values are supplied
jointly by the earlier signed theta transforms. A lower bound on
that image after the specified restriction and comparison with
the analytic rows is still needed there.


\subsection{The integral terminal pairing at all primes}\label{sec:weylperiod351}

Retain the period fields, intervals, matched levels and base tuple
of Section~\ref{sec:allclasses350}. Fix a coefficient number
field \(J\) containing all the embeddings and finite character
values in use. On the large side take the integral class table
with the column Weyl lattice \(\Delta\) in its specified lattice
frames. On the small side let \(S\) be the integral table of
tests with the dual Weyl lattice \(\nabla\), the dual action and
the inverse finite character. These are the lattices specified
in \srcloc{Formalism/LocalToGlobal/Columns/ColForms}{105}
and \srcloc{Formalism/LocalToGlobal/Columns/RowForms}{280}.
Use normalized bilinear Fischer contraction and weight
\(1/|\Gamma_a|\) on each small class \(a\). Thus the functional
from a large table \(P\) sends \(F\) to
\[
            \sum_a |\Gamma_a|^{-1}
                          B_m\big((R_1P)(a),F(a)\big),
\]
where \(B_m\) here denotes the tensor of the small forms.
This is the finite normalization of Theorem~\ref{thm:terminalformula347}.
For adelic measure multiply by the original small level volume.

Write \(e=e_{\mathcal X}\) for the prescribed central projection
on the two finite marked blocks of Theorem~\ref{thm:terminaltypes327}.
The formulas below also allow \(e=1\).
Put \(S_0=S\cap e(S\otimes J)\) and \(S'=eS\), in the same
projected space over \(J\). For each of the three choices
\(T_0=S,S_0,S'\), let \(\mathcal B(T_0)\) consist of the
restrictions of the displayed functionals to its \(J\)-span as
the integral large table varies. Write \(T_0^\vee\) for the
lattice of functionals integral on \(T_0\).

\begin{theorem}[A common integral estimate on the projected test lattice]\leavevmode
\label{thm:weylperiod351}
At a fixed \(x\) let \(D_x\) be the integer in
Theorem~\ref{thm:abslow349}. Let \(d\) bound \(n[E_x:\mathbb Q]\)
for the terminal ranks under consideration, put
\(M=|\mathrm{GL}_d(\mathbb F_3)|\), and write \(m_G\) for the
order of the two-prime block group of
Theorem~\ref{thm:terminaltypes327}. For every word and translate
at this \(x\) one has the lattice inclusions
\[
 \begin{aligned}
 D_x^2 T_0^\vee&\ \subseteq\ \mathcal B(T_0)\
                      \subseteq\ M^{-1}T_0^\vee
                          &&(T_0=S,\ S_0),\\
 D_x^2 (S')^\vee&\ \subseteq\ \mathcal B(S')\
                      \subseteq\ (m_GM)^{-1}(S')^\vee .
 \end{aligned}
\]
Both lines also hold if the large table is required to have
separate scalar sums zero at every \(P\mid\pi\). For every fixed
\(A>0\) the single integer \(D_x^2m_GM\) has logarithm
\(o(b/x^A)\) on the common intervals, where \(b=\exp(x^\eps)\).
It bounds the total finite logarithmic loss per exterior degree
of this terminal step.
\end{theorem}
\begin{proof}
Fix the word and put \(C=\prod_\tau c_\tau\).
Work first in the local ring of \(J\) at an arbitrary finite
prime. Denote by \(Y\) the small integral output table with
its Weyl lattice and positive finite character. Theorem~\ref{thm:weyllower351}
and the change of frame in Proposition~\ref{prop:normalframe349}
put \(R_1P\) in \(Y\). Indeed the matrix on each pair of
adapted integral frames is the tensor of the standard \(L_\tau\).
The formula for the inverse frames in the proof of
Theorem~\ref{thm:allclasses350} holds over the coefficient field
at every finite prime. In those frames \(C^2\) times the
section at the identity copy is the tensor of the integral
\(c_\tau^2 B_\tau\) of Theorem~\ref{thm:weyllower351}.
Extension by zero commutes with the class stabilizers by
Proposition~\ref{prop:scalarlifts350}. The section followed by
the product of the integral differences in
Theorem~\ref{thm:augmented350} has the same restriction on this
copy. In either large domain the image therefore lies between
\(C^2Y\) and \(Y\).

At each small class, the two ambient coefficient lattices are
exactly dual by Lemma~\ref{lem:weyldual351} and the opposite
changes of frame. The invariant test submodule is saturated,
being an intersection of kernels of integral endomorphisms of a
free module. Its quotient in the ambient test lattice is free
over the local ring. Any integral functional on the invariant
submodule thus extends to the ambient test lattice and is given
by an integral Weyl vector \(z\). Take its trace
\(\sum_{\gamma\in\Gamma_a}\gamma z\), using the coefficient
action with the positive finite character. This is a permissible
value in \(Y\). Against invariant tests its pairing divided by
\(|\Gamma_a|\) reproduces the given functional, by invariance of
the contraction. Doing this separately on all the classes realizes
every element of \(S^\vee\) using \(Y\), even at primes dividing
a stabilizer order.

The order of every \(\Gamma_a\) divides \(M\). Its rational
isometry acts faithfully on the underlying integral lattice of
rank at most \(d\), since its frame changes are determined by
that isometry. Add an identity block to obtain rank \(d\). The
kernel of reduction of \(\mathrm{GL}_d(\mathbb Z)\) modulo three
contains no nontrivial finite-order element. Otherwise it has
an element of prime order \(q\) of the form \(I+3^jU\),
\(j\ge1\), with \(U\) nonzero modulo three. For \(q\ne3\),
the first binomial term in its \(q\)-th power is nonzero modulo
\(3^{j+1}\). For \(q=3\) it is nonzero modulo \(3^{j+2}\)
while the higher terms vanish there. This proves the asserted
divisibility. Integral vectors in \(Y\) therefore give
functionals integral after multiplication by \(M\).
Together the two trace assertions give the first line for \(S\),
since \(|C|\) divides \(D_x\).

The submodule \(S_0\) is saturated in \(S\). Extend each
integral functional from \(S_0\) to \(S\) locally, as above.
Restriction proves the first line for \(S_0\).
The finite projection satisfies \(m_G eS\subset S_0\) by
Theorem~\ref{thm:terminaltypes327}, while \(e\) is the
identity on \(S_0\). Thus
\(S_0\subseteq S'\subseteq m_G^{-1} S_0\).
Duality on their shared rational span puts \((S')^\vee\)
between \(m_G S_0^\vee\) and \(S_0^\vee\), proving the
second line. The class tables, the kernels imposing invariance
or the zero sums, and the images of the period pairings commute
with localization, by flatness and finiteness of the class
tables. The inclusions at all finite primes give the stated
inclusions over the ring of integers.

The ranks and field degrees are polynomial in \(x\) on the given
family. One can take \(d\) to be their largest indicated product.
Then \(\log M\le d^2\log 3\), and
Theorem~\ref{thm:terminaltypes327} bounds \(\log m_G\)
by a polynomial in \(x\) for the very same marked levels.
Theorem~\ref{thm:abslow349} proves the asserted size.
Explicitly, at each prime over \(p\), normalize \(v(p)=1\).
Every elementary divisor valuation of the image relative to the
indicated dual test lattice belongs to
\([-v(m_GM),2v(D_x)]\). Valuations in exterior degree \(j\)
are sums of \(j\) such numbers. After dividing by \(j>0\),
the sum over \(p\) of the worst absolute bound times \(\log p\)
is at most \(\log(D_x^2m_GM)\), independently of the word.
\end{proof}

The estimate controls the cost of pairing arbitrary integral
large class values with the prescribed projected tests. For the
signed theta period those large values belong to the image of
the preceding theta transforms. Their restriction to the lifted
classes, and agreement with the incoming analytic row, are the
remaining period questions in Section~\ref{nav:I-28}.


\subsection{Which rational paths interact on the initial column}\label{sec:pathnorm352}

Use the actual finite path terms of Theorem~\ref{thm:covectorpaths352}
with a fixed testing form \(F\). For one unrefined initial class \(a\),
write \(\Xi_a=\prod_{P\mid\pi}(V_{0,P}^\vee\setminus\{0\})\). For
\(\xi\in\Xi_a\), let
\[
              u_{a,\xi}=
                 \sum_{\gamma:\,a_0=a,\ \eta_\gamma=\xi} W_{\gamma,F},
\]
with value zero for an empty sum. Thus \(u_{a,\xi}\) is a covector
on the initial coefficient module in the fixed class frame.
Use the finite class norm and normalized compact coefficient norms.
The unrefined ordinary initial classes have trivial stabilizers as in
Corollary~\ref{cor:rank111}. Indeed an automorphism has finite order
and acts on an integral lattice of rank \(k[E_x:\mathbb Q]\).
Its order is prime to \(g_1\) for \(g_1-1>k[E_x:\mathbb Q]\),
since an element of order \(g_1\) would require a primitive
\(g_1\)-th root as a rational matrix eigenvalue. Its reduction
at a marked prime over \(g_1\) belongs to the upper unitriangular
group there, which has \(g_1\)-power order. It therefore reduces to
the identity. For order \(h\) prime to \(g_1\), the matrix
\(1+\gamma+\cdots+\gamma^{h-1}\) is invertible locally since it
reduces to \(hI\). Multiplication by \(\gamma-1\) gives zero,
so \(\gamma=1\). Consequently initial refinements in these
coordinates carry unit weights and their coefficient values can be
chosen independently in the unrestricted refined table.

\begin{theorem}[The squared norm and integral values of the signed path sum]\leavevmode
\label{thm:pathnorm352}
Let \(\mathcal P_F^\circ\) be the functional of
Theorem~\ref{thm:covectorpaths352} restricted to the ordinary
initial table with separate scalar sums zero, in the arithmetic
class frames. For two nonzero covectors \(v,w\) over \(k_P\)
write \([v]\) for the line spanned by \(v\) and put
\[
 k_P(v,w)=\mathbf1_{v=w}
                -\frac{\mathbf1_{[v]=[w]}}{|k_P^\times|}.
\]
In the dual Hilbert norm its square is
\[
 \|\mathcal P_F^\circ\|^2
 =\sum_a\ \sum_{\gamma,\gamma':\,a_0=a_0'=a}
     \langle W_{\gamma,F},W_{\gamma',F}\rangle
          \prod_{P\mid\pi}k_P(\eta_{\gamma,P},\eta_{\gamma',P}).
\]
In particular interference can occur only within a common tuple
of projective residue lines at the same initial class.

There is also an exact formula over the coefficient ring of
integers. On each product of scalar orbits choose \(v_P,w_P\)
in its \(P\)-th orbit, and in its common integral initial
coefficient lattice choose \(z\). Put
\(\xi_{B,P}=w_P\) for \(P\in B\) and \(\xi_{B,P}=v_P\) otherwise.
The fractional ideal of values on the separately zero-sum
integral initial table is generated by
\[
        \sum_{B\subseteq S_T}(-1)^{|B|}\,u_{a,\xi_B}(z)
\]
over all such choices and all classes \(a\), for algebraic \(F\)
and the arithmetic coefficient maps. The zero ideal is allowed.
\end{theorem}
\begin{proof}
The matrix of the identity minus averaging within a scalar orbit
has entry \(k_P(v,w)\) on its index set of nonzero covectors.
The separate averages commute and are self-adjoint. Their
complementary projections have tensor product the orthogonal
projection onto the indicated zero-sum space. The squared norm
of a functional restricted to an orthogonal summand is the
inner product of its full coefficient array with the projection
of that array. Apply this to the array \(u_{a,\xi}\), using the
unit mass of each initial class, and expand both entries as their
path sums. All sums are finite, giving the formula.

For integrality, the initial lattice depends on the kernels of
the terminal covectors and is identical on each product of
scalar orbits. Choose one base point in each factor orbit.
Its free coordinate module splits integrally into the line
generated by the base point and the augmentation kernel,
with basis the differences of all other basis coordinates from
the base point. Expanding the tensor product of these
decompositions shows that the kernel of summation separately
in every coordinate is the tensor of the augmentation kernels:
summation in a given coordinate is an isomorphism on its
base-point line and zero on its kernel. Tensoring the differences
and the common coefficient lattice, then summing over all
orbits and classes, therefore generates precisely the integral
zero-sum table. Evaluating the period on these generators gives
the asserted alternating sums. This calculation uses no inverse
of \(|k_P^\times|\) in the coefficient ring.
\end{proof}

\begin{proposition}[Separation of short composite covectors at the chosen prime]\leavevmode
\label{prop:shortpaths352}
At an initial representative fix any \(E_x\)-basis
\(e_1,\ldots,e_k\) contained in its lattice. For a surviving path
choose rational representatives for its matrices and set
\(\ell_\gamma=e_N^\vee x_{m-1}\cdots x_0\).
Choose a nonzero \(b_\gamma\in\mathcal O_{E_x}\), a unit at all
\(P\mid\pi\), making each \(b_\gamma\ell_\gamma(e_j)\) integral.
Define
\[
      H_\gamma=\max_{\sigma:E_x\hookrightarrow\mathbb C,\ j}
                         |\sigma(b_\gamma\ell_\gamma(e_j))|.
\]
If \(\ell_\gamma,\ell_{\gamma'}\) are linearly independent and
\(2H_\gamma H_{\gamma'}<\sqrt T\), their cross term in
Theorem~\ref{thm:pathnorm352} is zero. Thus for any collection
with \(H_\gamma<T^{1/4}/\sqrt2\), its contribution taken alone
is an orthogonal sum over its rational initial covector lines.
\end{proposition}
\begin{proof}
The rational representatives are integral between the local
lattices at every \(P\mid\pi\). Hence the ideal of simultaneous
denominators of the coordinates under consideration localizes
to the full ring at each such \(P\). It contains an element
equal to one modulo their product by the Chinese remainder
theorem. This supplies \(b_\gamma\).

Suppose the two tuples have the same projective residue lines.
The factor \(\bar\pi\) is a unit on these places and is common
to both. Thus for any \(i,j\) the algebraic integer
\[
 a_{ij}=b_\gamma b_{\gamma'}
  \big(\ell_\gamma(e_i)\ell_{\gamma'}(e_j)
                -\ell_\gamma(e_j)\ell_{\gamma'}(e_i)\big)
\]
belongs to every \(P\mid\pi\). Independence gives at least one
nonzero \(a_{ij}\). For this pair its principal ideal is
divisible by \((\pi)\), the product of those distinct primes.
Writing \(d=[E_x:\mathbb Q]\), the absolute ideal norm of
\((\pi)\) is \(T^{d/2}\), since \(\pi\bar\pi=T\) and conjugation
preserves ideal norm. At each complex embedding the triangle
inequality bounds the absolute value of \(a_{ij}\) by
\(2H_\gamma H_{\gamma'}\). Taking norms gives
\[
                  T^{d/2}
                  \le |N_{E_x/\mathbb Q}(a_{ij})|
                  \le (2H_\gamma H_{\gamma'})^d ,
\]
a contradiction. Therefore some local projection entry vanishes.
This proves both claims upon expanding the squared norm.
\end{proof}

\begin{corollary}[Separation when choosing the common terminal prime]\leavevmode
\label{cor:pathprimes352}
At a given \(x\) fix a finite list of rational matrix paths on
the period field, before choosing \(T\) from the candidate primes
of Theorem~\ref{thm:markedstrict337}. The list may include paths
for several words. Retain the nonzero composite normal rows.
Take one positive integer \(D_0\) clearing the denominators of
all the matrix maps between their lattices and of these rows
on their initial lattices. In the preceding height definition
use \(b_\gamma=D_0\) for the entire list. Let \(\mathcal J\)
consist of the unordered pairs of independent rows in matching
initial coordinates. Write \(d=[E_x:\mathbb Q]\) and set
\[
 \mathcal L=\log D_0+d\sum_\gamma\log H_\gamma+
              2\sum_{\{\gamma,\gamma'\}\in\mathcal J}
                         \log(2H_\gamma H_{\gamma'}).
\]
Among candidate primes \(T>Y>1\), at most
\(\mathcal L/\log Y\) have to be excluded in order that all
the matrices reduce, that every listed row survives at every
terminal place, and that each pair in \(\mathcal J\) gives zero
cross term. If \(\mathcal L\le Y^{1-\rho}\) for fixed \(\rho>0\),
there remain \(Y^{1-o(1)}\) choices on the cohomological intervals
with the original word and marked-prime conditions.
\end{corollary}
\begin{proof}
Exclude the primes dividing \(D_0\). For each row choose a
nonzero integral coordinate \(D_0\ell_\gamma(e_j)\). If the
row vanishes at any one terminal place away from \(D_0\), the
underlying rational prime divides the absolute norm of this
coordinate. The sum of logarithms of these primes for this
row is at most \(d\log H_\gamma\).

For each pair in \(\mathcal J\) fix a nonzero integral minor
\(a_{ij}\) as in Proposition~\ref{prop:shortpaths352}. At every
colliding prime its ideal is divisible by the full chosen
\((\pi)\). The ideals obtained for distinct rational \(T\) are
pairwise coprime. Taking the product of their norms bounds
the sum of their \(\log T\) by
\((2/d)\log|N_{E_x/\mathbb Q}(a_{ij})|
\le2\log(2H_\gamma H_{\gamma'})\).
Adding these bounds and \(\log D_0\), then dividing by \(\log Y\),
proves the stated count. The original common selection leaves
\(Y^{1-o(1)}\) primes by Theorem~\ref{thm:markedstrict337}.
For fixed \(\rho>0\) the exclusions are a vanishing fraction
of this number. All the stated cohomological conditions persist
on the remaining primes.
\end{proof}

The height and prime selection arguments control the indicated
cross terms in the actual finite array. A lower bound for the
complete signed norm must also control the other paths and the
signed sums inside a projective class, including rationally
proportional paths.

\section{Branching and elementary divisors of coefficient lattices}\label{nav:chapter-coefficient-weyl-lattices}

\subsection{Relative Fischer norms of products of flag minors}\label{sec:commonfischer348}

Use polynomials in an \(N\)-square variable matrix with \(N=n+1\)
and \(n\ge1\). For \(1\le i\le n\), let \(D_i,D_i'\) have first
\(i\) columns and rows \((1,\ldots,i)\), respectively
\((1,\ldots,i-1,N)\). For nonnegative exponent arrays write
\[
       X(a,b)=\prod_iD_i^{a_i}(D_i')^{b_i},\qquad
       q(X)=\sum_\alpha \alpha! X_\alpha^2 .
\]
All these polynomials have integral coefficients. Fix a box of
exponent arrays bounded coordinatewise by a lower array and an
upper array. Write \(X_-\) and \(X_+\) for the two extremal
products, \(D=\deg X_+\), and \(L=\deg X_+-\deg X_-\).
The degree \(D\) can be much larger than the number of boxes in
the Young diagrams that change, measured here by \(L\).

\begin{lemma}[An integral multiple of the ratio of two squared norms]\leavevmode
\label{lem:contractwindow348}
There is one positive integer \(A\) depending on the box with
\[
       \log A\ll N^2(1+L)^2\log(2+N+D).
\]
If the arrays of \(X,Y\) in the box are ordered coordinatewise
in that order, then \(Aq(Y)/q(X)\) is a positive integer and
\[
             q(Y)/q(X)\le (N^2\max(1,D))^{\deg Y-\deg X}.
\]
\end{lemma}
\begin{proof}
The variables carry the defining column action of \(\mathrm{GL}_N\)
and, on the rows, the defining action of \(\mathrm{GL}_n\) on the
first \(n\) rows and the identity on the last. Choose the row and
column triangular subgroups for which initial minors are highest
vectors. Every indicated product is a joint highest vector. Adding an earlier
column to a later one leaves each used minor unchanged. The same applies to the first \(n\) rows because each
minor uses an initial segment of them. Its joint diagonal weight
\(\eta=(\eta^{\rm col},\eta^{\rm row})\) counts the occurrences
of these columns and rows. Each such highest line in its joint
type is unique. In degree \(d\), the Cauchy decomposition pairs
one column module of partition \(\lambda\) with one large row
module of that partition. Restricting the latter to the first
\(n\) rows has multiplicity one, by the interlacing rule. Both
statements are also proved by the character arguments used in
Lemma~\ref{lem:degree181} and Proposition~\ref{prop:branch275}.
The polynomial spanning the line is primitive. Extract the terms
in the last row in decreasing column order, as in
Lemma~\ref{lem:highestline347}. This gives
\(\prod_iD_i^{a_i}D_{i-1}^{b_i}\), whose diagonal monomial has
coefficient one. This argument works for all the arrays in the box.

Write \(Y=SX\), where \(S\) is the product of the extra minors
and \(e=\deg S\le L\). The Fischer adjoint of multiplication by
\(S\) is the constant-coefficient differential operator
\(S(\partial)\). Indeed multiplication by a variable is adjoint
to its derivative on the factorial monomial basis, and the
coefficients here are real. The polynomial \(v=S(\partial)Y\)
is integral. Its weight for the two diagonal tori is the weight
of \(X\). It belongs to the image of the contraction map from
the module generated by \(Y\) tensored with the dual of all
degree-\(e\) polynomials. This map is equivariant because a
linear change of variables acts dually on constant derivatives.

For either factor of rank \(j=N,n\), write \(\eta(Y)\) for its
part of the highest weight. Every irreducible type \(\gamma\)
in that contraction satisfies
\[
                  \eta(Y)_i-e\le\gamma_i\le\eta(Y)_i .
\]
In fact the dual of degree-\(e\) polynomials embeds in the
\(e\)-fold dual tensor power by symmetrization in characteristic
zero. For columns each factor is a sum of dual defining modules.
For the smaller row group trivial summands also occur. Tensoring
an irreducible module by a dual defining module subtracts one
from one coordinate in all ways that keep the entries dominant.
To verify this, multiply its Weyl character numerator
\(\det(t_j^{\eta_i+r-i})\) in rank \(r\) by \(\sum_jt_j^{-1}\).
Expansion of the determinant groups the terms by the exponent
lowered. Equal resulting exponents give zero, and every other
determinant remains in decreasing order. Division by the same
Weyl denominator proves the rule, including determinant twists.
Iteration proves the two bounds.

Let \(\eta_-\) be the weight of \(X_-\) in the factor of rank
\(j\). Every coordinate of \(\eta(Y)-\eta_-\) lies between zero
and \(L\), by counting the added variables. All shifted entries
\(\gamma_i+j-i\) therefore belong to
\[
 I_j=\bigcup_{i=1}^j
 [\eta_{-,i}+j-i-L,\eta_{-,i}+j-i+L]\cap\mathbb Z .
\]
For the desired type \(\eta(X)\) let
\(S_j=\{\eta(X)_i+j-i:1\le i\le j\}\). Apply the single
Capelli projection of \eqref{eq:sieve}, taking its product over
\(b\in I_j\setminus S_j\). On a type \(\gamma\) its numerator
has eigenvalue \(\prod_{b\notin S_j}\prod_i(\gamma_i+j-i-b)\)
where \(b\) runs through \(I_j\). Dividing by the value on
\(\eta(X)\) makes the operator the orthogonal isotypic
projection on the entire contraction space. Every other type
has a shifted entry outside \(S_j\), as the strictly decreasing
shifted entries determine the type. The numerator is integral on
the full polynomial ring in view of the integral matrix
differential operators in the Capelli formula of
Theorem~\ref{thm:proj}. A common multiple of the absolute
denominators for both projections is
\[
             A=\prod_{j\in\{n,N\}}\ 
                \prod_{a,b\in I_j,\ a\ne b}|a-b| .
\]
There are at most \(j(2L+1)\) entries in \(I_j\), and every
absolute difference is at most \(2L+D+N\). This gives the
claimed logarithmic bound.

The projections for columns and rows commute. Applied to \(v\)
their product is \(tX\), since the torus weight already is the
joint highest weight and that line occurs once. Pair with \(X\)
in the Fischer form. Orthogonality gives
\[
           t q(X)=(S(\partial)Y,X)_{\rm F}
                         =(Y,SX)_{\rm F}=q(Y).
\]
Since \(AtX\) is integral and \(X\) has a coefficient one,
\(At\) is a positive integer. Finally, on degree \(j\) the norm
of multiplication by a coordinate is at most \(\sqrt{j+1}\)
by the monomial basis. The sum of the absolute coefficients of
an order-\(i\) determinant is \(i!\le N^i\). Applying the
triangle inequality to all factors of \(S\) gives
\(\|SX\|_{\rm F}\le N^e\max(1,D)^{e/2}\|X\|_{\rm F}\),
as asserted.
\end{proof}

\begin{theorem}[One integer for all the products in a box]\leavevmode
\label{thm:commonfischer348}
For all products in the box there are positive divisors \(a_X\)
of a single positive integer \(J\) such that
\[
       \frac{q(X)}{q(X')}=\frac{a_X}{a_{X'}},\qquad
       \log J\ll N^2(1+L)^2\log(2+N+D) .
\]
\end{theorem}
\begin{proof}
Apply Lemma~\ref{lem:contractwindow348} in both directions along
the ordered triple \(X_-,X,X_+\). The positive integers
\(a_X=Aq(X)/q(X_-)\) and \(Aq(X_+)/q(X)\) have product
\(J=A^2 q(X_+)/q(X_-)\), independent of \(X\). Applying the
upper bound of the lemma to the two extremal products gives the
estimate for \(\log J\). Dividing the two formulas for
\(a_X,a_{X'}\) proves the identity.
\end{proof}


\subsection{The three large powers at the last hyperplane}\label{sec:threeminors349}

Keep the initial minors \(D_i,D'_i\) of Section~\ref{sec:commonfischer348}
in rank \(N=n+1\). For the actual weights, only three of their
exponents can grow faster than every fixed power of \(x\).
The following calculation keeps their full multiplicities. Write
\(q\) for the squared Fischer norm of a real polynomial.

\begin{lemma}[The normalizing ratio for three minor powers]\leavevmode
\label{lem:threeminors349}
Let \(3\le k\le n\) and \(A,B,C\) be nonnegative integers.
Define two polynomials
\[
     X_0=(D'_1)^B D_{k-1}^A(D'_k)^C,\qquad
     P_0=D_1^B D_{k-1}^A D_k^C.
\]
Then
\[
 \frac{q(X_0)}{q(P_0)}
     =\frac{(A+k-2)!(B+k-2)!}{(k-2)!(A+B+k-2)!}
     =\prod_{j=1}^A \frac{k-2+j}{B+k-2+j}.
\]
The numerator and denominator of this reduced fraction both divide
an integer of logarithm at most \(2A\log(2+A+B+k)\). Its value
is one whenever \(A=0\).
\end{lemma}
\begin{proof}
The variables used by \(X_0\) are in the \(k\)-square matrix
\(Y\) on rows \(1,\ldots,k-1,N\) and the first \(k\) columns.
Rename the last of these rows as row \(k\), and write \(D\)
for the leading minor of order \(k-1\) in \(Y\).
Then \(X_0=(\det Y)^C Q\) with \(Q=Y_{k1}^B D^A\).
The variable \(Y_{k1}\) is disjoint from those in \(D\), so
monomial factorials factor and
\(q(Q)=B!H_{(A^{k-1})}\), using Lemma~\ref{lem:hook86}.
Here \(H_v\) denotes the hook product of the partition \(v\).

For column operations \(Q\) is highest of weight
\(v=(A+B,A^{k-2},0)\). Its column degrees give these entries,
and adding earlier columns to later ones fixes both factors.
By the Cauchy decomposition the full space of highest column
polynomials of this weight is one irreducible row module.
Multiplication by \((\det Y)^C\) changes its row action by
a determinant character, which has absolute value one on the
unitary group. Consequently the squared norm after multiplication
is a scalar times the old squared norm on that module.
Evaluate on its leading-minor polynomial \(p_v\). This gives
the scalar \(H_{v+C\mathbf1}/H_v\), because
\((\det Y)^Cp_v=p_{v+C\mathbf1}\). The polynomial \(P_0\),
on the first \(k\) original rows, is also this leading-minor
product. Thus the ratio in the statement is
\(B!H_{(A^{k-1})}/H_v\).

In rank \(k-1\) the hook formula is
\[
 H_\nu=\frac{\prod_{i=1}^{k-1}(\nu_i+k-1-i)!}
                   {\prod_{i<j}(\nu_i-\nu_j+j-i)} .
\]
Increasing the first entry of the rectangle by \(B\) changes
only its first factorial and the differences involving the first
row. Their ratios give the two factorials in the numerator and
the two in the denominator of the asserted expression. In its
product over \(j\), each of the \(2A\) integer factors is at
most \(2+A+B+k\). Their product is a common multiple of both
parts of the reduced fraction, which proves the bound.
\end{proof}


\subsection{Integral branching on the two polynomial lattices}\label{sec:weylrestriction351}

Let \(O\) be the valuation ring at any finite place of a number
field and \(K\) its fraction field. In a rank \(q\) polynomial
model of partition \(v\), write \(V_v\) for the flag space over
\(K\), \(p_v\) for the product of leading minors, and \(B_v\)
for the bilinear Fischer form normalized by \(B_v(p_v,p_v)=1\).
Put \(\nabla_v=V_v\cap O[X]\). Write \(\Delta_v\) for the
\(O\)-span of all iterated coefficients of elementary row
substitutions on \(p_v\). In a substitution one replaces a
row \(X_j\) by \(X_j+tX_i\) and takes any coefficient in \(t\).
This includes \(i=j\). These are the dual Weyl and Weyl lattices
used in the terminal test and column, respectively. Their
definitions in the supplied polynomial model are in
\srcloc{Formalism/GradedWeylDivisibility/StatementA}{24}
and lines 63--74 there. For general integral highest weights we
use an integral frame for the separate determinant line.

\begin{lemma}[Exact duality in the highest-vector normalization]\leavevmode
\label{lem:weyldual351}
The \(B_v\)-dual of \(\Delta_v\) is exactly \(\nabla_v\), and
the \(B_v\)-dual of \(\nabla_v\) is exactly \(\Delta_v\), in
every residue characteristic.
\end{lemma}
\begin{proof}
For \(P\in V_v\) one has \(B_v(P,p_v)=P(I_q)\).
Both functionals vanish on row torus weights different from \(v\).
For evaluation, a surviving monomial at the identity uses only
diagonal entries, whose exponents are its prescribed column
degrees. For the pairing, distinct monomials are orthogonal.
The row weight \(v\) is the one-dimensional highest line and
both functionals equal one on \(p_v\). If the matrix \(G\)
substitutes \(X\mapsto GX\), its Fischer adjoint substitutes
\(X\mapsto G^{\mathsf t}X\). This follows on a product of linear
forms by contracting each factor with one in the second product
and summing over permutations. Thus, as a polynomial identity in
the entries of \(G\),
\[
                    B_v(P,p_v(GX))=P(G^{\mathsf t}).
\]
In particular every integral-coefficient \(P\) pairs integrally
with every iterated row coefficient of \(p_v\).

Conversely assume that all these pairings belong to \(O\).
Use \(G\) equal to a product of elementary row substitutions,
with independent parameters. Every coefficient of
\(P(G^{\mathsf t})\) then belongs to \(O\) by the identity.
Choose a product parametrizing a lower unitriangular matrix,
a diagonal matrix with entries \(1+t_i\), and an upper
unitriangular matrix. Such products of elementary matrices
give all the indicated triangular factors by row elimination.
Over the algebraic closure of the residue field their image
contains every matrix with nonzero leading principal minors,
by triangular factorization. This is a nonempty Zariski open
subset of the matrix space, and its transpose is dense as well.
If \(P\) has a nonintegral coefficient, multiply \(P\) by a
positive power of a uniformizer making its coefficients integral
and at least one a unit. Its reduction is a nonzero polynomial.
In the displayed substitution the reduction is identically zero,
by integrality of the original substituted coefficients. Density
gives a contradiction. This proves the first duality. Both
modules are full lattices. For \(\Delta_v\) the same argument
over \(K\) shows that a vector orthogonal to all its generators
vanishes on a dense set, so is zero. Finite generation follows
from containment in the finite free monomial module over \(O\).
Lattice biduality over a discrete valuation ring proves the
second equality.
\end{proof}

Use the interlacing weights \(\lambda,\mu\) of
Section~\ref{sec:lowerlattice347} in ranks \(N=n+1,n\). Their
last entries agree. Let \(l,m\) denote their partitions after
subtracting this common entry. Write \(C_b\) for extraction
of the last-row monomial of exponents \(b_i=\lambda_i-\mu_i\),
with the final column set to zero. Write \(A=P_m C_b\), where
\(P_m\) is the projection onto the small flag space in its
column multidegree. Let \(I:V_m\longrightarrow V_l\) be the
small-block intertwiner carrying \(p_m\) to the interlacing
highest vector \(X\) of Lemma~\ref{lem:highestline347}.
The branching rule gives uniqueness of its copy, and
\(A I=1\) by \(C_b X=p_m\). Set \(c=c_\tau\) and
\(\kappa=B_l(X,X)\) in this single-embedding calculation.

\begin{proposition}[Integral maps on the horizontal strip]\leavevmode
\label{prop:tableaulift351}
There is an \(O\)-linear map \(H:\nabla_m\longrightarrow\nabla_l\)
with \(C_bH=1\), taking values in last-row degree \(\sum_i b_i\).
The two equivariant maps \(c A\) and \(c I\) preserve the dual
Weyl lattices in their respective directions.
\end{proposition}
\begin{proof}
For a semistandard tableau of shape \(m\) on \(1,\ldots,n\),
form the product of its column minors, using the entries of
each diagram column as the row list and the initial matrix
columns of the same length. These products form an \(O\)-basis
of \(\nabla_m\). Indeed, in lexicographic order read the
matrix columns successively, putting smaller row indices first
in each. A minor with increasing row list has its diagonal
product as leading monomial with coefficient one. In a tableau
product this monomial records the multiset of entries in each
diagram row. Since the row is sorted, it recovers the tableau.
The products are independent and their number equals the
dimension by the Schur character formula for semistandard
tableaux.~\cite[Chapters 1 and 8]{XXIII-XVIII-fulton351}
They therefore span over \(K\). For a polynomial with integral
coefficients, successively subtract its largest remaining basis
monomial coefficient times its tableau product. At each stage
that coefficient lies in \(O\), giving the asserted \(O\)-basis.

Fill the boxes of \(l/m\) with the new letter \(N\). They are
at row ends, and at most one lies in any diagram column by
interlacing. The result is semistandard. Define \(H\) on the
basis by the new product of minors. There are exactly \(b_i\)
modified minors of height \(i\), each containing the last row
once. To obtain all the required \(X_{Nn}\) factors one must
choose the final-column term in every modified minor of height
\(n\). Proceed in decreasing height. At height \(i\) the
same argument forces the choice of \(X_{Ni}\) in every remaining
modified minor of that height. Each sign is positive because
the last row occurs last in the row list. The cofactors give
the original tableau product. This proves \(C_bH=1\).

On the large module keep just the indicated last-row degree,
by the integral monomial projection \(Q\). Among the small
types \(\nu\) on this slice, interlacing and equality of the
small central weight give
\(\nu_i\in I_i\), where \(I_i\) is the interval in
Section~\ref{sec:lowerlattice347}. In detail the nonnegative
differences \(\nu_i-\lambda_{i+1}\) have sum
\(\sum_i(\mu_i-\lambda_{i+1})\). For integral \(y\), use the
central block operator \(K_n(y)=(-1)^n C_n(-y-n)\), in the
Capelli notation of Theorem~\ref{thm:proj}. Its eigenvalue on
type \(\nu\) is \(\prod_j(y-\nu_j+j)\). Its entries are integral
polynomial differential operators, including the integral
diagonal shift from the determinant line. Thus the product
\[
            \prod_{i=1}^n\ \prod_{v\in I_i,\ v\ne\mu_i} K_n(v-i)
\]
preserves \(\nabla_l\). On the slice it acts as \(c\) on type
\(\mu\) and zero on every other type, using any differing
coordinate of \(\nu\). If \(P\) denotes the large projection
onto this copy on the slice, it follows that \(cPQ\) is integral.

The map \(A\) kills every other type by equivariance, and
vanishes outside the slice by extraction. On that copy
\(C_b\) carries \(X\) to \(p_m\) and intertwines the smaller
group, so it already takes its values in \(V_m\). Consequently
\(A=C_b P Q\) and \(P Q=I A\). Extraction is integral, proving
that \(cA\) preserves integral coefficients. For \(Y\in\nabla_m\),
\(cIY=cPQH(Y)\) since \(A H(Y)=Y\). This proves integrality
of \(cI\).
\end{proof}

\begin{theorem}[The column Weyl lattice at every finite prime]\leavevmode
\label{thm:weyllower351}
For the prescribed standard last restriction \(L=c\kappa A\),
put \(B=L^*/(c^2\kappa)\). In integral frames at every finite
place one has
\[
       L(\Delta_l)\subseteq\Delta_m,\qquad
       c^2 B(\Delta_m)\subseteq\Delta_l,\qquad LB=1 .
\]
In particular each elementary divisor of the cokernel on these
lattices divides \(c^2\). Under tensor product over the CM
embeddings this holds with \(c^2\) replaced by
\((\prod_\tau c_\tau)^2\).
\end{theorem}
\begin{proof}
The calculation on the highest line in
Theorem~\ref{thm:lowerlattice347} gives \(A^*p_m=X/\kappa\).
Since the adjoint intertwines the small block, this proves
\(I=\kappa A^*\). Hence \(L=(cI)^*\) and \(c^2B=(cA)^*\).
Dualizing the two integral maps of Proposition~\ref{prop:tableaulift351}
with Lemma~\ref{lem:weyldual351} gives both inclusions.
The identity \(LB=1\) is the adjoint identity of
Theorem~\ref{thm:lowerlattice347}. Multiplication by \(c^2\)
therefore annihilates the cokernel. Tensoring these explicit
maps proves the final assertion, including after extension of
coefficient rings.
\end{proof}


\subsection{Coefficient lattices over residue hyperplanes}\label{sec:openflags353}

Let \(L\) be a free \(\mathbb Z_T\)-lattice of rank \(k\),
where \(T>2\) is prime. Write \(W_\lambda(L)\) for the Weyl
lattice \(\Delta\) of Section~\ref{sec:weylrestriction351} with
its determinant power, for an integral decreasing highest
weight \(\lambda\). For a nonzero residue covector \(a\) put
\(L_a=\{x\in L:a(\bar x)=0\}\). This is independent of scaling
\(a\). For a finite list \(\mathcal U\) of nonzero vectors in
\(L/TL\), let \(A\) consist of the \(a\) with \(a(u)\ne0\)
for every \(u\in\mathcal U\).

\begin{lemma}[Sum and intersection over the allowed hyperplanes]\leavevmode
\label{lem:openflags353}
Put \(s=\lambda_k\) and \(h=\lambda_1-\lambda_k\).
Assume \(T>h+k\) and \(T-|\mathcal U|>h\), with \(A\) restricted
to nonzero covectors also for an empty list. Then
\[
  \sum_{a\in A} W_\lambda(L_a)=T^{\lambda_k}W_\lambda(L),
  \qquad
  \bigcap_{a\in A} W_\lambda(L_a)=T^{\lambda_1}W_\lambda(L).
\]
The identities persist over the valuation ring of any finite
extension of \(\mathbb Q_T\).
\end{lemma}
\begin{proof}
Since \(\det L_a=T\det L\), remove the factor \(T^s\) and
write \(\nu_i=\lambda_i-s\). Its Young diagram has \(h\)
columns and at most \(k-1\) rows. In a basis of \(L\) adapted
to \(L_a\), one basis vector is multiplied by \(T\).
Every normal-coordinate weight of \(W_\nu\) is between zero
and \(h\): in a product of its \(h\) column minors the vector
occurs at most once in each minor. The integral Weyl lattice
is the direct sum of its torus weight lattices, since each
iterated row coefficient of the highest vector is homogeneous
for the diagonal torus. These facts give the two immediate
inclusions of each lattice between \(T^hW_\nu(L)\) and \(W_\nu(L)\).

The translates of the highest vector by residue bases span the
reduction of \(W_\nu(L)\). Indeed each elementary row substitution
applied to any Weyl vector has degree at most \(h\) in its
parameter. This follows from the minor products over the fraction
field, since a rank-one change of a minor is at most linear.
Its coefficients lie in the Weyl lattice by definition of the
divided row operators. Modulo \(T\) each coefficient is a linear
combination of values at invertible elementary matrices, by
interpolation at \(h+1\) distinct parameters. In a diagonal
substitution omit the parameter \(-1\). The hypothesis on \(T\)
leaves enough choices. Iteration proves the spanning assertion.

In such a translate the last basis vector is unused. Let \(a_0\)
vanish on the span of the first \(k-1\). If \(a_0\in A\) the
translate is already in the image of \(W_\nu(L_{a_0})\) modulo
\(T\), by lifting these basis vectors into \(L_{a_0}\).
Otherwise choose a covector \(b\) nonzero on every \(u\in\mathcal U\)
with \(a_0(u)=0\). Fewer than \(T\) hyperplanes are excluded,
so their union has fewer than \(T^k\) elements. The covectors
\(a_0,b\) are independent since \(a_0\) was excluded.
Choose \(y\) with \(a_0(y)=1,\ b(y)=0\).
The transvection \(g_t=1-t yb\) sends \(\ker a_0\) onto
\(\ker(a_0+tb)\). At most \(|\mathcal U|\) parameters \(t\)
are forbidden. The translate of our highest vector under \(g_t\)
is polynomial of degree at most \(h\) in \(t\). Interpolation
from the allowed values recovers its value at zero. Thus the
sum of the allowed lattices surjects modulo \(T\) onto \(W_\nu(L)/T\).
Nakayama's lemma gives equality.

For the intersection use the normalized Fischer duality.
All hooks of \(\nu\) are \(<T\), so the polynomial Weyl and dual
Weyl lattices agree in a standard frame by
Lemma~\ref{lem:weyldual351} and Theorem~\ref{thm:terminaldual325}.
Their pairing is perfect there. The transpose adjoint for a
change of frame identifies the dual lattice of \(W_\nu(L_a)\)
with \(W_\nu(L_a^\vee)\) in the dual frame. Here
\(L_a^\vee=L^\vee+T^{-1}\mathbb Z_T\tilde a\), with \(\tilde a\)
any lift of \(a\). Thus \(T^hW_\nu(L_a^\vee)\subset W_\nu(L^\vee)\).
Its reduction contains every highest-vector translate whose first
basis vector is \(a\): lift that residue basis and use
\(T^{-1}\tilde a\) as the first vector in the enlarged lattice.
Each of the \(h\) minors uses this vector once.

These translates again span even with \(a\) restricted to \(A\).
Starting with an excluded first vector \(a_0\), use \(b,y\) as
above and the transvection on \(L^\vee/T\) given by
\(z\mapsto z+t z(y)b\). The first vector becomes \(a_0+tb\).
Interpolation at the same allowed parameters proves the assertion.
Nakayama gives
\(\sum_{a\in A}T^hW_\nu(L_a^\vee)=W_\nu(L^\vee)\).
Taking lattice duals gives the claimed intersection. Reinstating
the determinant gives \(s+h=\lambda_1\). Finite flat scalar
extension preserves the sums and intersections, proving the last
assertion.
\end{proof}


\subsection{Elementary divisors of the transformed coefficient table}
\label{sec:fourierslopes354}

At a fixed ordinary initial class \(a\), let \(S_{\rm cov}\)
and \(S_{\rm vec}\) be the separately zero-sum tables on
covectors and vectors, respectively, taking all values in the
single unrefined lattice \(D_a\) of Theorem~\ref{thm:pointlattice353}.
Use \(S_a^\circ\) for the actual lattice on the covector side.
These three lattices have the same rank \(d\). We compute the
elementary divisors of \(S_a^\circ\) relative to \(S_{\rm cov}\)
and bound the change in each of them when the target becomes
\(S_{\rm vec}\) under Fourier transformation. Elementary-divisor
exponents are allowed to be negative for fractional lattices.

\begin{theorem}[The full list of weight exponents and the Fourier contribution]\leavevmode
\label{thm:fourierslopes354}
Fix a coefficient prime above \(T\), with valuation \(v(T)=1\).
For each \(\tau\) with \(v(\tau(\bar\pi))=1\), list the weights
of a single diagonal coordinate of \(W_{\lambda_\tau}\) with
multiplicity. For every other \(\tau\), use zero with multiplicity
\(\dim W_{\lambda_\tau}\). Form all sums of one entry per list
and repeat the multiset once per basis element of the scalar
separately zero-sum table. Order the result as
\(\alpha_1\le\cdots\le\alpha_d\). These are the exponents of the
identity \(S_a^\circ\longrightarrow S_{\rm cov}\).
If \(\beta_1\le\cdots\le\beta_d\) are the exponents of
\(\mathcal F_0:S_a^\circ\longrightarrow S_{\rm vec}\), then
\[
        \alpha_i\le\beta_i\le\alpha_i+fk
                        \qquad(1\le i\le d).
\]
In every exterior degree \(j\), the valuations of the two ideals
generated by the \(j\)-rowed minors differ by a number in
\([0,jfk]\), with the larger valuation for the transform.
At each finite prime away from \(T\), both lists of exponents
consist of zeros.
\end{theorem}
\begin{proof}
Fix a tuple of covector lines and choose bases adapted to their
kernels as in Theorem~\ref{thm:pointlattice353}. In each
affected coefficient factor a single basis vector has been
multiplied by \(T\). Its matrix on the direct sum of torus
weight lattices is diagonal with entries \(T^w\), one for
each weight \(w\) of this coordinate including the determinant
twist. The existence of the integral weight decomposition was
proved in Lemma~\ref{lem:openflags353}. A change of the chosen
line or of its integral adapted frame leaves the elementary
divisors unchanged. Tensoring adds the exponents. On a product
of scalar orbits the table of coefficients is a tensor of
augmentation modules with this single coefficient lattice.
The augmentation of an orbit of size \(T-1\) has a basis of
the \(T-2\) differences against a fixed orbit member.
Consequently this tuple contributes \((T-2)^f\) copies of
the listed coefficient weights. The different line tuples
are independent direct summands. This proves the assertion
for the identity and justifies every multiplicity.

In the two unrefined lattices the Fourier matrix and its
opposite are integral, including on the augmentation modules.
The two compositions are \(Q I\) for \(Q=T^{fk}\) by
Theorem~\ref{thm:vectorchain353}. For a lattice \(L\) over
the local valuation ring define \(\operatorname{ord}_L(z)\),
for \(z\ne0\), as the minimum valuation of its coordinates
in an integral basis. The preceding two inclusions give
\[
  \operatorname{ord}_{S_{\rm cov}}(z)
       \le \operatorname{ord}_{S_{\rm vec}}(\mathcal F_0z)
       \le \operatorname{ord}_{S_{\rm cov}}(z)+fk .
\]
For an invertible linear map \(A\) between two lattices of
rank \(d\), its \(i\)-th increasing exponent equals
\[
 \max_{\dim U=d-i+1}\
 \min_{\substack{z\in U\\ \operatorname{ord}_{\rm source}(z)=0}}
                   \operatorname{ord}_{\rm target}(Az),
\]
where \(U\) runs over subspaces of the fraction-field source.
Indeed, in Smith bases the span of the last \(d-i+1\) vectors
attains the \(i\)-th exponent. Every competing subspace
meets the first \(i\) coordinates nontrivially. Scaling a vector
of this intersection to be primitive gives an image of order
at most the \(i\)-th exponent. Applying this formula with the
same source \(S_a^\circ\) to the two orders gives the claimed
pointwise inequalities. The valuation of the \(j\)-minor ideal
is the sum of the \(j\) smallest exponents, by diagonalization.
Away from \(T\) the three coefficient frames agree and the
opposite transform divided by \(Q\) is integral.
\end{proof}

\begin{corollary}[The mean exponent and its global normalization]\leavevmode
\label{cor:fouriermean354}
At each prime over \(T\),
\[
       \frac1d\sum_{i=1}^d\beta_i
            =\frac{fk}{2}
                +\frac1k
                  \sum_{\tau:\,v(\tau(\bar\pi))=1}|\lambda_\tau|,
     \qquad |\lambda_\tau|=\sum_{j=1}^k\lambda_{\tau,j}.
\]
The vertical bars here denote the signed sum of the entries of
a weight. In the class Hilbert metrics at every complex
embedding, every singular value of \(\mathcal F_0\) is
\(T^{fk/2}\). Consequently the normalized sum of the mean
finite exponents times \(\log T\), less the normalized
archimedean log determinant per rank, is
\[
                 \frac{\log T}{2k}\sum_\tau|\lambda_\tau|.
\]
For each fixed translate and terminal prime, this quantity
and each local mean in the first formula agree for all the
equal-weight words of Theorem~\ref{thm:common111}.
\end{corollary}
\begin{proof}
In each representation the coordinate permutations interchange
the diagonal coordinates, preserving multiplicities. Scalar
matrices act through the power \(|\lambda_\tau|\), so the
sum of all \(k\) coordinates of every torus weight is
\(|\lambda_\tau|\). Thus the average of the chosen coordinate
is \(|\lambda_\tau|/k\). Taking the mean of the list in the
theorem gives the second summand in the local formula.

For the determinant between the two unrefined lattices,
identify each residue space with its dual using a chosen
basis and the associated symmetric dot product. This
identification permutes the coordinates of the augmented
tables integrally. The resulting Fourier endomorphism
squares to \(Q\) times reflection in zero by character
orthogonality. Reflection is an integral automorphism of
the augmented table. The valuation of the determinant of
the Fourier endomorphism is therefore \(d fk/2\). Composition
with the identity of lattices accounts for both terms of
the displayed formula. The archimedean assertion follows
from Theorem~\ref{thm:vectorchain353}, since Galois conjugation
preserves nontriviality of every residue additive character.

In the finite mean give a prime \(w\mid T\) weight
\([J_w:\mathbb Q_T]/[J:\mathbb Q]\), and average over all
embeddings of \(J\) for the archimedean term. The constant
Fourier contributions cancel. For each \(\tau\),
\(\tau(\bar\pi)\) is supported above \(T\) and all its
archimedean absolute values are \(\sqrt T\), since
\(\pi\bar\pi=T\) in the imaginary quadratic subfield.
The norm formula shows that the weighted finite mean of
\(v_w(\tau(\bar\pi))\) is \(1/2\). This gives the global
formula, with signed exponents also for negative weights.
Finally \eqref{eq:letters111} shows that the variation of the
weight sum at each embedding depends on the word only through
its Hamming weight. This proves the last assertion even when
the ranks of the different word spaces differ.
\end{proof}


\subsection{The distribution of one coordinate of a Weyl weight}
\label{sec:oneweight355}

The elementary divisors in Theorem~\ref{thm:fourierslopes354}
are controlled by torus weights counted with their actual
multiplicities. For a decreasing integral weight \(\lambda\)
in rank \(k\ge2\), choose one weight of \(W_\lambda\) at random,
giving it probability equal to its multiplicity divided by the
dimension. Let \(w\) be its first coordinate. A permutation of
the coordinates gives the same distribution for any other one.

\begin{lemma}[An exponential moment for a coordinate of a Weyl weight]\leavevmode
\label{lem:oneweight355}
Put \(l_j=\lambda_j+k-j\). If \(p=(p_1,\ldots,p_k)\) is
uniform on the simplex \(p_i\ge0,\ \sum_i p_i=1\), then
for \(t>0\)
\[
  \mathbb E(e^{tw})
       \left(\frac{e^t-1}{t}\right)^{k-1}
       =\mathbb E\exp\!\left(t\sum_j p_j l_j\right).
\]
Suppose \(l_1\le A\), \(l_j\le C\) for \(j\ge2\), with
\(A>0,\ C\ge0\). Then for \(t=(k-1)/(2A)\) one has
\(\mathbb E(e^{tw})\le2e^{tC}\). When also \(l_1\le C\)
the bound improves to \(e^{tC}\).
\end{lemma}
\begin{proof}
The left character evaluated at \((e^t,1,\ldots,1)\), divided
by its dimension, is \(\mathbb E(e^{tw})\). In the Weyl character
formula let the last \(k-1\) variables tend to one, writing the
variables as exponentials. The rows for these variables tend,
after divided differences, to the successive derivatives of
\(e^{yl_j}\) at \(y=0\) of orders \(0,\ldots,k-2\), divided by
their factorials. Expanding the determinant along its first
row gives the Vandermonde minors of those \(k-1\) rows. Relative
to the full Vandermonde in the \(l_i\), the signed factor at
column \(j\) is \(1/\prod_{i\ne j}(l_j-l_i)\) up to a sign
common to the columns. The denominator alternant in the
original variables leaves the factor \((e^t-1)^{k-1}\); all
other factors are independent of \(t\). Thus the normalized
character is
\[
   \frac{(k-1)!}{(e^t-1)^{k-1}}
             \sum_{j=1}^k\frac{e^{tl_j}}{\prod_{i\ne j}(l_j-l_i)}.
\]
To verify its constant, the partial fraction expansion of
\(\prod_i(z-l_i)^{-1}\) has simple-pole coefficients
\(1/\prod_{i\ne j}(l_j-l_i)\). Expansion at infinity shows
that the sums with numerator \(l_j^m\) vanish for \(m<k-1\)
and equal one for \(m=k-1\). The displayed character therefore
tends to one as \(t\) tends to zero. The \(l_j\) are distinct
even when some highest weight entries coincide. A determinant
power also allows arbitrary negative entries in \(\lambda\).

For nonnegative integers \(d_j\) of sum \(n\), integration
of the monomial \(\prod p_j^{d_j}\) against uniform probability
on the simplex gives \((k-1)!\prod d_j!/(n+k-1)!\). This follows
by integrating one coordinate at a time, rescaling the remaining
simplex after each integration, and using the beta integral for
integer exponents. Thus the coefficient of \(t^n\) in the
simplex expectation is \((k-1)!h_n(l)/(n+k-1)!\), where
\(h_n(l)=\sum_{\sum d_j=n}\prod l_j^{d_j}\). In the same
partial fraction expansion at infinity the coefficient of
\(z^{-k-n}\) is precisely \(h_n(l)\). Expanding the exponentials
in the preceding display proves the first identity coefficient
by coefficient after multiplying by its stated factor.

Now \(\sum p_jl_j\le C+A p_1\), and
\(\Pr\{p_1>v\}=(1-v)^{k-1}\) for \(0\le v\le1\), by
rescaling the simplex. For \(s=(k-1)/2\), integrating the tail
gives
\[
      \mathbb E(e^{s p_1})
       =1+s\int_0^1 e^{sv}(1-v)^{k-1}\,dv
       \le 1+s\int_0^\infty e^{-(k-1-s)v}\,dv=2 .
\]
Use \((e^t-1)/t\ge1\) in the identity. When every \(l_j\le C\)
the simplex average is at most \(e^{tC}\) directly.
\end{proof}



\subsection{Divisibility of a theta operator in the Weyl lattice}
\label{sec:rowjetcontent449}

Fix \(p\) and a class section \(F=F_{ch}\) as in
Theorem~\ref{thm:direct76}. Assume that \(\pi=\sigma_pc_E\)
interchanges two sets \(H,L\) partitioning the embedding labels.
This occurs for \(H=\Phi_0,\ L=\Phi_1\) at the good primes inert
in \(K_d\) and split in \(K_q\). Work at any coefficient place
above \(p\), in its valuation ring \(O\), with \(v(p)=1\).
There are \(s\) rows and \(r\) columns of matrix variables per
embedding. Our aim is to retain the individual row lengths of a
highest weight when a derivative is paired with an integral test.

Put \(f(m)=\sum_{a\ge1}\lfloor m/p^{2a}\rfloor\).
For a list \(Z\) of high variables, let \(\mathcal A(Z)\) be the
ring of series whose coefficient at \(Z^\alpha\) has valuation
at least \(-\sum_i f(\alpha_i)\). We also use polynomial
parameters. The same bound is then required for every coefficient
in those parameters. Superadditivity of \(f\) proves closure
under multiplication. Integral linear substitutions preserve this
ring. Indeed for \(\sum_i b_i=m\) the multinomial valuation is
\(v(m!)-\sum_i v(b_i!)\), and \(v(m!)-f(m)\) is the sum of
the floors at odd powers of \(p\), hence superadditive. Applying
this to each factor proves the substitution assertion.

\begin{lemma}[The exceptional part after a change of high frame]\leavevmode
\label{lem:rankoneshift449}
Let \(A_h=(a_{h,li})\) be \(r\)-by-\(r\) matrices over an
integral polynomial ring over \(O\). Replace the high variables by
\[
                 Z_{h,bi}=\sum_{l=1}^r a_{h,li}X_{h,bl}.
\]
Write \(y_{\mu,bi}\) for the low variables, and
\(g_M=[z^M](\vartheta'/\vartheta)(z)\) for the one-variable
marked theta. To any fixed total degree in \(X\) define matrices
\[
 (U_{h l j})_{bi}=X_{h,bl}^{M}a_{h,li}^{M},\qquad M=p^j,\quad
 j\ge1\ \text{odd},\qquad
 S_\mu=\sum_{\pi^j h=\mu}\ \sum_{l=1}^r g_{p^j}U_{h l j}.
\]
The sum uses the odd \(j\) up to that degree. Every undivided
low derivative at \(y=0\) of
\[
 F(Z(X),y)\exp\!\left(-\sum_{\mu,b,i}y_{\mu,bi}(S_\mu)_{bi}\right)
\]
belongs through that degree to \(\mathcal A(X)\). Each
\(U_{h l j}\) is an outer product.
\end{lemma}
\begin{proof}
Factor the section as in Theorem~\ref{thm:direct76}. The initial
integral elliptic germ and all its low derivatives have bounded
coefficients on the high variables. Each vacuum restricted to
its chosen high half has the same bound by isotropy and
Theorem~\ref{thm:gram66}. For a low derivative in direction
\((\mu,b,i)\), the first logarithmic derivative of the product of
vacua on this half is \(E_{\mu,bi}\) of
Lemma~\ref{lem:partner76}. All logarithmic derivatives of order
at least two satisfy the bound. In fact the first ordinary
derivative of \(\vartheta'/\vartheta\) is an integral elliptic
germ, and the invariant elliptic derivation preserves integrality.
Passing from elliptic to additive coordinates uses the substitution
estimate above.

In \(E_{\mu,bi}\) the coefficient at \(Z_{h,bi}^{M}\) is
\(g_M\sum_d \overline V_{\mu bi,d}\,\overline V_{hbi,d}^{\,M}\).
For odd \(j\) this sum is congruent modulo \(pO\) to one when
\(\mu=\pi^jh\) and to zero otherwise. Indeed Frobenius sends
the second factor modulo \(p\) to \(V_{\pi^j hbi,d}\), and the
rows of \(V\) are orthonormal. All entries before the congruence
lie in an unramified ring. This identifies the same scalar
\(g_M\) for each matching matrix entry. Since
\(v(g_M)\ge-f(M)-1\), the error from this congruence belongs to
\(\mathcal A(Z)\). All the other coefficients satisfy its bound
by Lemma~\ref{lem:partner76}.

For a mixed term of \((\sum_l a_l X_l)^M\) with exponent vector
\(\beta\), each \(\beta_l<M\). In
\[
 v\binom{M}{\beta_1,\ldots,\beta_r}
       - f(M)+\sum_l f(\beta_l)
   =\sum_{k\ {\rm odd}}
       \left(\lfloor M/p^k\rfloor-\sum_l\lfloor\beta_l/p^k\rfloor\right)
\]
the term \(k=j\) equals one and all terms are nonnegative.
Consequently multiplication by \(g_M\) leaves every mixed term
in \(\mathcal A(X)\). Removing the pure terms from each first
logarithmic derivative therefore leaves a member of this ring.
Those removed terms form precisely \(S_\mu\). The exponential
in the statement subtracts them from the first logarithmic
derivatives and has zero higher logarithmic derivatives. For
any undivided derivative of a vacuum, differentiating a product
successively gives the vacuum times a sum of products of its
logarithmic derivatives with integer multiplicities. Apply this
to the corrected vacua, and use the ordinary product rule on the
initial integral germ. Closure under products proves the claim,
including sections which vanish at the origin. All calculations
in truncated series take place in characteristic zero.
\end{proof}

For \(\nu_\tau=(\lambda_{\tau,i}+n_\tau)_{1\le i\le r}\) in
\eqref{eq:sub81}, let \(\nu_\tau^+\) denote its zero padding to
length \(s\). Write \(\Delta_\lambda=\bigotimes_\tau
\Delta_{\lambda_\tau}\) for the Weyl lattice in the chosen
integral frames. Put \([u]_+=\max(u,0)\) and
\(d_M(t,s)=\max(0,\lceil(t-s(M-1))/M\rceil)\).

\begin{theorem}[A valuation bound in terms of the rows of the partition]\leavevmode
\label{thm:weylrowjet449}
For the class section and the two halves fixed above set
\[
 G_p(\nu)=
 \sum_{h\in H}\sum_{l=1}^r\sum_{\substack{j\ge1\\j\ {\rm odd}}}
 \big[\,d_{p^j}(\nu_{h,l},s)-\nu_{\pi^jh,1}\,\big]_+ .
\]
The class operator of \eqref{eq:sub81} satisfies
\[
             {\cal T}_{F,n}(\Delta_\lambda)
                     \ \subseteq\ p^{G_p(\nu)}\Delta_{\nu^+}.
\]
\end{theorem}
\begin{proof}
Write \(p_\lambda\) for the product of leading minors. Generators
of the source in one factor are coefficients of \(p_\lambda(AX)\)
for products \(A\) of elementary off-diagonal row additions with
independent parameters. To check this from the definition of
\(\Delta\), start with the row-homogeneous highest vector.
An off-diagonal coefficient of order \(d\) moves \(d\) units of
row degree from one row to another. A diagonal substitution on a
row of degree \(t\) scales by \((1+z)^t\), so its coefficients
give integral multiples of the previous vector. This proves the
generator assertion by induction. Since \(\det A=1\), multiplying
by \(\det(X)^n\) gives \(p_\nu(AX)\).

Pair at each output embedding with an arbitrary
\(P_\tau\in\nabla_{\nu_\tau^+}\). Such tensor tests span the
dual lattice by Lemma~\ref{lem:weyldual351}. For the normalized
bilinear Fischer form the resulting differential symbol is
\[
 K_\tau(D)=P_\tau\big([D A_\tau^{\mathsf t},0]\big).
\]
Indeed, place the \(r\) rows of \(A_\tau D^{\mathsf t}\) in
the top of an \(s\)-square matrix \(B\), filling the rest with
zeros. Then the polynomial in the output frame is
\(p_{\nu_\tau^+}(BY)\), and its pairing with \(P_\tau\) is
\(P_\tau(B^{\mathsf t})\) by the polynomial identity in the
proof of Lemma~\ref{lem:weyldual351}. Thus \(K_\tau\) has
integral coefficients even when \(p\) is small relative to the
weight. The scalar to be bounded is
\((\prod_\tau K_\tau(\partial))F(0)\).

At \(h\in H\) perform the substitution of
Lemma~\ref{lem:rankoneshift449} with this \(A_h\). By the chain
rule the high symbol becomes \(P_h([D,0])\) in the new variables.
Every monomial \(\partial_X^\alpha\) in these high symbols
satisfies \(\sum_b\alpha_{h,bl}=\nu_{h,l}\) by column homogeneity.
At \(\mu\in L\), the degree of \(K_\mu(D+tuv^{\mathsf t})\) in
\(t\) is at most \(\nu_{\mu,1}\) for arbitrary column vectors
\(u,v\) of the indicated matrix sizes. Indeed \(P_\mu\) is a
linear combination of tableau products of flag minors, as proved
in Proposition~\ref{prop:tableaulift351}. A minor is affine in
\(t\), since any term using two perturbed columns has proportional
columns. The right multiplication by \(A_\mu^{\mathsf t}\) and
zero padding preserve the outer-product form of the perturbation.
The number of minor factors is the number \(\nu_{\mu,1}\) of
columns of the Young diagram.

Let \(F_0\) be the corrected germ in Lemma~\ref{lem:rankoneshift449}.
Conjugation by its removed exponential replaces the low operator
by \(\prod_{\mu\in L}K_\mu(\partial_y+S_\mu)\) on \(F_0\) at
\(y=0\). This identity holds monomial by monomial because
\(S_\mu\) is independent of \(y\). Expand first with a separate
indeterminate in place of every \(g_{p^j}\) multiplying
\(U_{hlj}\). The degree in each such indeterminate is at most
\(\nu_{\pi^jh,1}\) by the outer-product calculation, even with
all the other perturbations present. Coefficients in this
expansion are integral polynomials in the matrix entries, the
frame parameters and the commuting low derivative symbols.
If one indeterminate occurs \(k_{hlj}\) times, each monomial in
its coefficient uses \(k_{hlj}\) factors \(X_{h,bl}^{p^j}\),
distributed over \(b\). Thus a term contributing to the high
index \(\alpha\) satisfies both bounds
\[
 k_{hlj}\le\nu_{\pi^jh,1},\qquad
 k_{hlj}\le\sum_{b=1}^s\lfloor\alpha_{h,bl}/p^j\rfloor .
\]

Restore the scalars \(g_{p^j}\) and differentiate \(F_0\) in
the low variables. Its resulting high coefficients satisfy the
elliptic factorial bound. Superadditivity of \(f\) absorbs the
cost \(f(p^j)\) of each exceptional factor together with the
factorial cost of the remaining coefficient into
\(\sum_{h,b,l} f(\alpha_{h,bl})\). The additional valuation
loss is at most \(\sum k_{hlj}\). Multiplication by the
undivided high factorials leaves at least
\[
 \sum_{h,l}\sum_{j\ {\rm odd}}
 \left[\,\sum_b\lfloor\alpha_{h,bl}/p^j\rfloor-\nu_{\pi^jh,1}\,\right]_+ .
\]
For a column of total degree \(\nu_{h,l}\), division with
remainder bounds its floor sum below by \(d_{p^j}(\nu_{h,l},s)\).
This gives \(G_p(\nu)\). A truncation at the total high degree
suffices throughout. Every coefficient in the independent
parameters obeys this valuation bound. Extract the source
generators and use the exact integral duality with the output
tests to obtain the lattice inclusion.
\end{proof}

The estimate applies separately to each original section in the
local frames of Theorem~\ref{thm:direct76}. A sum of sections with
integral scalar coefficients inherits it by linearity. The first
row in a low factor occurs here because every column minor is
affine along a matrix of rank one.


\subsection{A rank-one terminal frequency on the integral Weyl lattice}
\label{sec:terminalrankonetheta455}

At a split terminal prime \(T\), the Fourier sum of
Lemma~\ref{lem:terminalfourier454} involves characters
\(X\mapsto\psi_T(\xi'\bar X u)\) at the position place \(P\).
Their frequency matrices have rank at most one. To use this
in estimating an individual summand, first compare two choices
of elliptic product basis at an individual coefficient place.

\begin{lemma}[The phase between two locally integral product bases]\leavevmode
\label{lem:unitframe455}
Let \(p\) be a rational prime. Suppose \(B,B'\) are rational
orthonormal product bases both generating the fixed trace
lattice \(M\otimes\mathbb Z_p\) over
\(\mathcal O_K\otimes\mathbb Z_p\). In the fixed marked
differential the phase \(u(B',B)\) of
Theorem~\ref{thm:framechange314} is a unit at every place
above \(p\).
\end{lemma}
\begin{proof}
In the product basis \(B\) apply the reverse normalized
lattice sum to the standard vacuum to obtain the finite test
\(\phi_B\) used at the end of the proof of
Theorem~\ref{thm:vac242}. Its section \(G^B_{\phi_B}\) is the
marked product section, which has value one at the origin.
Indeed the forward and reverse intertwiners compose to the
identity on the vacuum line by finite character orthogonality
and equality of the two indices for self-dual lattices.
The finite vector \(\phi_B\) is a sum of translates of the
standard vacuum with root-of-unity coefficients times the
inverse square root of a lattice index. That index is prime
to \(p\), since the lattices coincide there, and the shifts
have order prime to \(p\). Expanding that same standard
vacuum in the \(B'\) product model has the same two properties.
Composing the sums thus expresses \(G^{B'}_{\phi_B}(0)\)
as an integral linear combination at \(p\) of marked elliptic
products translated by prime-to-\(p\) torsion. Each has
integral value there by Theorem~\ref{thm:torsion312} in order
zero. Since \(G^{B'}_{\phi_B}=u(B',B)G^B_{\phi_B}\), it follows
that \(u(B',B)\) is integral there. Using \(\phi_{B'}\) in the
opposite order proves integrality of its inverse.
\end{proof}

For one Fourier summand of one edge write \(\eta_v\) for the
residue frequency functional at \(P\in S_T\) over \(v\mid T\);
it is the map \(\bar X\mapsto\xi'\bar Xu\). At a coefficient
place \(w\mid T\) set
\[
 H_w(\eta)=\{\tau:P_{\tau,w}\in S_T,\
             \eta_{P_{\tau,w}\cap F}\ne0\}.
\]
Here \(P_{\tau,w}\) is the place pulled back by \(\tau\) as
in Theorem~\ref{thm:terminaldegrees454}. On a conjugated edge
use the conjugate coefficient place.

\begin{theorem}[The cost of a terminal Fourier character]\leavevmode
\label{thm:rankonecost455}
Use unrefined integral coefficient frames at \(T\). Remove
the explicit coefficient of the Fourier summand. In a product
basis generating \(M\otimes\mathbb Z_T\), each matrix
coordinate of the class operator between the Weyl lattices
has valuation at least
\[
                         -\sum_{\tau\in H_w(\eta)}\nu_{\tau,1}.
\]
Prime-to-\(T\) shifts with coefficients integral at \(w\)
may be included in the class section.
\end{theorem}
\begin{proof}
At a place \(v\) with \(\eta_v\ne0\), lift \(\xi'\) and \(u\)
to primitive integral vectors at the position place. Their
product gives the primitive functional
\(\widetilde\eta_v(X)=\widetilde\xi' X\widetilde u\) of rank
one on the position lattice. Extend this element to a basis of
the dual position lattice, choosing it as the first basis
vector. This is possible since one of its coefficients is a
unit. In the position lattice use the exact dual basis.
The sums of corresponding dual and position vectors are an
orthonormal basis in the split Hermitian trace model. A
transpose in identifying a dual matrix with the opposite
matrix component preserves rank. At \(v\) with \(\eta_v=0\)
use any such pair of dual bases. Approximate these bases at
\(T\) by a single rational orthonormal basis to precision
\(T\), as in Theorem~\ref{thm:terminaldegrees454}. This choice
can depend on the Fourier summand. Lemma~\ref{lem:unitframe455}
shows that changing to it multiplies the section by a unit
at the given place, for every differential order.

In the product coordinates of this basis, the \(T\)-primary
part of the shift over each \(v\) is concentrated in its
first coordinate if \(\eta_v\ne0\). In fact the class is
represented in the dual position lattice by
\(T^{-1}\widetilde\eta_v\) modulo that lattice, and reducing
the change of basis modulo \(T\) suffices to transport a
class killed by \(T\). By the conjugate annihilator formula
of Theorem~\ref{thm:torsion312}, the only coordinates at
which one may lose one unit of valuation per differentiation
are the first coordinates belonging to \(v=P_{\tau,w}\cap F\)
with \(\tau\in H_w(\eta)\). The map
\(\tau\mapsto P_{\tau,w}\cap F\) is bijective since the CM
type induces all the embeddings of \(F\) and \(T\) splits
completely. Call this list of first coordinates \(I_w\).
All other product coordinates have integral undivided
derivatives in every order. For a multi-index in product
coordinates this bounds the loss by the total order on
\(I_w\), because the translated section is a product of
one-variable sections. Finite prime-to-\(T\) intertwiners
produce sums of such products with the same bounds.

Let \(D_\tau\) denote the matrix of physical differential
symbols at \(\tau\). In the chain rule a product differential
\(\partial_a\) has a coefficient matrix \(M_{\tau a}\) in
\(D_\tau\), consisting of conjugate embedded basis entries.
For \(a\in I_w\) in a different \(v\)-block this matrix is
zero modulo \(T\). For the first coordinate \(a(\tau)\) of
the matching block, when \(\tau\in H_w(\eta)\), the matrix
is congruent modulo \(T\) to an integral rank-one matrix
\(U_\tau\). Indeed conjugation selects the dual position
component for this embedding, and in the exact local basis
its entries form \(\widetilde\eta_v\), up to the transpose
identification. The rational basis agrees with it modulo \(T\).
Consequently \(D_\tau\) is the sum of a matrix \(D_\tau^0\)
and, just when \(\tau\in H_w(\eta)\), the term
\(U_\tau\partial_{a(\tau)}\). All entries of \(D_\tau^0\)
are integral linear combinations of the commuting operators
\[
                 \partial_a\quad(a\notin I_w),\qquad
                 T\partial_a\quad(a\in I_w).
\]

For an elementary-parameter generator of the source Weyl
lattice and a dual target test, the differential polynomial
on this block is \(K_\tau(D)=P_\tau([DA_\tau^{\mathsf t},0])\)
from Lemma~\ref{lem:jetweyl453}. Its degree along
\(D+yU_\tau\) in the indeterminate \(y\) is at most
\(\nu_{\tau,1}\). Indeed \(P_\tau\) is an integral linear
combination of products of flag minors, with one minor for
each column of the Young diagram. Right multiplication and
zero padding preserve rank at most one. Each minor of a
rank-one perturbation is affine in \(y\) by multilinearity.
This proof applies with commuting differential operators for
the entries of \(D\).

Introduce a separate \(y_\tau\) in every affected block
and expand the product of all \(K_\tau\) around \(D_\tau^0\).
Every coefficient is an integral polynomial in the regular
operators displayed above and the source parameters. Substitute
\(\partial_{a(\tau)}\) for \(y_\tau\). In each term every
occurrence of a derivative on \(I_w\) coming from a regular
operator has its compensating factor \(T\). The remaining
number of such derivatives is at most
\(\sum_{\tau\in H_w(\eta)}\nu_{\tau,1}\). Applying the product
valuation bound and extracting the source parameters gives
the claimed result for every integral dual test and hence the
full Weyl lattice. Finally one can return to the same locally
integral product basis for all Fourier summands by the lemma.
\end{proof}

\begin{corollary}[First-row lengths in the centered terminal cost]\leavevmode
\label{cor:terminalrow455}
In Theorem~\ref{thm:terminaldegrees454} replace \(d_{ij}/2\)
by \(\tfrac12\sum_\tau \nu_{ij,\tau,1}\). In the centered
inequality of Corollary~\ref{cor:allcost454} its terminal term
can consequently be replaced by
\[
                 \frac{\log T}{q}\sum_{i,j}
                       \left(2fr_j+\sum_\tau\nu_{ij,\tau,1}\right).
\]
The same rescaled Gram matrices, or the same explicit correction
to the original Gram matrices, occur in that inequality.
\end{corollary}
\begin{proof}
For every summand enlarge \(H_w(\eta)\) to all
\(\tau\) with \(P_{\tau,w}\in S_T\). This makes the loss
independent of the summand. Restore the \(f\) factors
\(T^{-r_j}\) per edge and use the integral coefficients at
the other marked places. Summing the Fourier terms, composing
the edges with their fixed covectors, and summing the original
paths thus preserve the lower bound. The initial and terminal
scalars, class masses and test projections are exactly those
of Theorem~\ref{thm:terminaldegrees454}. Each \(\tau\) occurs
with total local degree weight \(1/2\) on averaging over
\(w\mid T\). This gives the first assertion for both test
lattices. Adding their two bounds in the product formula
gives the displayed contribution.
\end{proof}


\subsection{Elementary divisors of split matrix shifts}
\label{sec:splitmatrixtorsion456}

Fix a prime \(p\) splitting completely in the period field
\(E\). For a class of \(s\)-by-\(r\) matrices
\(b\) over \(\mathbb Q_p\) modulo the integral matrices, let
\[
                 a_1(b)\ge\cdots\ge a_r(b)\ge0
\]
be the exponents of the elementary divisors of the finite
image of \(\mathbb Z_p^r\) under \(b\) in
\((\mathbb Q_p/\mathbb Z_p)^s\), padded by zeros. We use
\(r\le s\); transposition gives the same nonzero exponents.
There are integral row and column bases in which \(b\) has a
representative diagonal with entries \(p^{-a_t(b)}\) for
\(a_t(b)>0\) and zero elsewhere. In fact one can clear the
denominator \(p^A\), put the resulting integral matrix in Smith
normal form, divide by \(p^A\), and discard the integral entries.
This also proves that the exponents are independent of the lift.

\begin{lemma}[Degrees of a flag polynomial in several rank-one directions]\leavevmode
\label{lem:flagdegree456}
Let \(K(D)=P([D A^{\mathsf t},0])\) be the integral differential
symbol for a partition \(\nu\) as in Lemma~\ref{lem:jetweyl453}.
If \(U_1,\ldots,U_r\) are rank-at-most-one matrices and the
coefficient of \(y_1^{m_1}\cdots y_r^{m_r}\) in
\(K(D+\sum_t y_t U_t)\) is nonzero, then
\[
       \sum_{t=1}^k m_t\le\sum_{t=1}^k\nu_t
                        \qquad(1\le k\le r).
\]
Consequently \(\sum_t a_t m_t\le\sum_t a_t\nu_t\) for all
\(a_1\ge\cdots\ge a_r\ge0\).
\end{lemma}
\begin{proof}
Hold \(y_{k+1},\ldots,y_r\) fixed. Over the field of rational
functions, the column space of \(\sum_{t\le k}y_t U_t\) has
dimension at most \(k\). The same holds after the fixed right
substitution and padding. In a minor of height \(h\), any term
using more than \(k\) perturbed columns vanishes by linear
dependence. The degree in \(y_1,\ldots,y_k\) is therefore at
most \(\min(k,h)\). A tableau monomial of shape \(\nu\) is
a product of one minor per diagram column. The sum of
\(\min(k,h)\) over their heights is the number of boxes in
the first \(k\) rows. Passing to a linear combination proves
the degree bounds, as polynomial identities also after
integral specialization. With \(a_{r+1}=0\), multiply the
\(k\)-th inequality by \(a_k-a_{k+1}\) and sum.
The weighted degree is attained by the highest polynomial on a
diagonal matrix with entries \(y_t\), where the product of its
leading minors equals \(\prod_t y_t^{\nu_t}\).
\end{proof}

Write a split Weyl shift on a given edge as a tuple of matrix
classes \(b_Q\), one for each \(Q\mid p\) in \(E\). At the
chosen position place of a split pair, \(b_Q\) is a position
class. At its conjugate place it is the frequency class in
the dual matrix lattice, using the trace pairing. Since \(E\)
splits completely at \(p\), these lattices are the ordinary
integral matrix lattices in dual local frames. We work with
the marked CM model and its unrefined physical frames.

\begin{theorem}[The valuation of a translated Weyl operator at a split prime]\leavevmode
\label{thm:matrixshift456}
For a finite shift \(b\), at every coefficient place \(w\mid p\)
every coordinate of its class operator on the Weyl lattices
has valuation at least
\[
                 -\sum_\tau\sum_{t=1}^r
                    \nu_{\tau,t}\,a_t(b_{\overline{P_{\tau,w}}}).
\]
This holds in a product basis generating the local trace
lattice. A conjugated class operator uses the conjugate
place. A scalar coefficient \(c\) adds \(v_w(c)\) to the
lower bound.
\end{theorem}
\begin{proof}
Fix \(w\). For each \(v\mid p\) in \(F\), precisely one
embedding \(\tau\) in the CM type induces \(v\) under \(w\).
At the opposite component \(\overline{P_{\tau,w}}\) take
Smith row and column bases for the shift as above. The matrix
units in these bases form a basis of this component of the
trace lattice, each member of rank one in the physical
coordinates. The corresponding diagonal matrix units carry
the components of positive order \(p^{a_t}\). Use the exact
trace-dual basis at \(P_{\tau,w}\) and take the sums of dual
partners. This is a local orthonormal basis. If \(A\ge1\)
bounds all exponents of both components at every \(v\), choose
a rational orthonormal basis agreeing with it modulo \(p^A\)
at both places of \(K\), by the Cayley approximation in
Theorem~\ref{thm:adapted314}. The inverses agree to the same
precision. Multiplying the difference by a shift of order
dividing \(p^A\) therefore changes its product coordinates
by integral elements. The choices may be made for this one
shift and this one place, by Lemma~\ref{lem:unitframe455}.

In the resulting elliptic coordinates, label a diagonal
coordinate of positive order in the opposite component
by \((v,t)\), and give it cost \(a_t\); give all other
coordinates cost zero. The conjugate annihilator bound of
Theorem~\ref{thm:torsion312} says exactly that an undivided
derivative of order \(n\) on a coordinate of cost \(a\)
has valuation at least \(-na\). Indeed the annihilator over
the opposite split component contains \(p^a\), and
\(\kappa_p=0\). The component on the other side of the same
elliptic coordinate imposes no condition on this local
exponent. The normalized intertwiners for the local trace
lattice add shifts of order prime to \(p\), with integral
coefficients. Each resulting elliptic product obeys the
additive bound on all its separate derivative orders.

Call \(p^a\partial_{v,t}\) the scaled variable on a coordinate
of cost \(a\), and use the ordinary derivative on a zero-cost
coordinate. Each matrix \(D_\tau\) of physical differential
symbols is a matrix \(D_\tau^0\) of integral combinations
of these variables, plus
\[
         \sum_{t:\,a_t>0} U_{\tau t}\,\partial_{v,t},
                     \qquad v=P_{\tau,w}\cap F ,
\]
where \(U_{\tau t}\) has rank one. To check this, in the exact
local basis the coefficient matrices of the indicated
derivatives are the embedded conjugate entries of the Smith
matrix units, up to transposition; entries on other \(v\)-blocks
vanish. In the rational basis used for differentiation the
coefficients differ by multiples of \(p^A\). Since \(A\ge a\)
for every coordinate, all these differences multiply scaled
variables with integral coefficients. Thus the asserted
decomposition is exact, with the difference absorbed into
\(D_\tau^0\).

Now start from a Weyl generator and pair with any integral
dual test as in Lemma~\ref{lem:jetweyl453}. Expand its
\(K_\tau(D_\tau)\) in independent symbols \(y_t\) for the
unscaled \(\partial_{v,t}\) in the displayed sum. The
coefficients are polynomials in the scaled variables with
integral coefficients. Lemma~\ref{lem:flagdegree456} bounds
the valuation loss on substituting \(\partial_{v,t}\) for
\(y_t\) by \(\sum_t a_t\nu_{\tau,t}\) in that block.
The compensating powers in all scaled variables cover their
own differentiations, even when they land on the same
elliptic factor as an unscaled variable. Multiplication over
the embeddings, evaluation on the product sections, and
extraction of the source-generator coefficients prove the
bound for all pairings. Duality gives it on the integral
target lattice. Returning to any common locally integral
basis only multiplies the section by a unit at \(w\).
\end{proof}

For a prescribed column \(i\), edge \(j\), and nonterminal
row prime \(p\), use the normalized split functions of
Section~\ref{sec:outsideprimeperiods453}. Let \(L\) be the
product of their position lattices over \(v\mid p\) and choose
a common period sublattice \(L'\subset L\). For each function
\(\phi\) in the list, each supported coset \(z+L\) and each
\(\eta\in(L/L')^\vee\), use the coefficient
\[
         c_{\phi,z,\eta}=\frac1{[L:L']}
                    \sum_{x\in L/L'}\phi(z+x)\overline{\eta(x)} .
\]
The bar on a finite character denotes its inverse. In the
symmetric Weyl convention there can also be a root-of-unity
phase, which has valuation zero. Write \(b(z,\eta)\) for
the position and frequency pair. Its two local matrix
classes are evaluated in the indicated split frames.
In each fixed \(w\) maximize over the nonzero coefficients
and include \(0\) among the candidates; set
\[
 J^{\rm mat}_{ij,p}=
 \sum_{w\mid p}\omega_w
   \max\left(\{0\}\cup
    \left\{\sum_{\tau,t}\nu_{ij,\tau,t}\,
           a_t\!\left(b(z,\eta)_{\overline{P_{\tau,w}}}\right)
              -v_w(c_{\phi,z,\eta})\right\}\right).
\]
Conjugate the local calculation on a conjugated edge.

\begin{corollary}[The original row-prime costs from finite Fourier coefficients]\leavevmode
\label{cor:splitcost456}
In Corollary~\ref{cor:allcost454}, with the terminal improvement
of Corollary~\ref{cor:terminalrow455}, one can replace every
\(J^{\rm row}_{ij,p}\) by \(J^{\rm mat}_{ij,p}\). These
quantities satisfy \(J^{\rm mat}_{ij,p}\le J^{\rm row}_{ij,p}\).
\end{corollary}
\begin{proof}
Finite Fourier inversion expresses \(\phi\) exactly as the
specified linear combination of shifted unit vacua. Taking a
smaller period lattice adds characters with zero coefficient
by orthogonality. Changing a lift of \(z\) changes coefficients
only by phases and changes its matrix by an integral matrix.
Use Theorem~\ref{thm:matrixshift456} on each nonzero summand.
Maximizing its loss and then averaging over \(w\) is valid
in the proof of Theorem~\ref{thm:outsidedifferent453}; all
path operations there stay linear and keep the specified
coupling coset. This proves the replacement. For the estimate
against the previous bound take \(L'=\prod_v p^{b_v}L_v\).
Every \(c_{\phi,z,\eta}\) has valuation at least
\(-c_{ij,p}-\sum_v b_v D_{jv}\) by the defining finite sum.
Over either place \(\epsilon\) of \(K\), all the relevant
matrix exponents at the opposite place are at most
\(h_{-\epsilon}\). Thus their weighted partition sum is at
most \(d_{ij}h_{-\epsilon}\). Each \(\epsilon\) has total
local degree weight \(1/2\). This gives the claimed inequality.
\end{proof}

\section{Arithmetic functionals on the prescribed rows}\label{nav:g-XII-0}

\subsection{The effect of theta lifting on the Gram estimates}\label{nav:I-49}
The diagonal and mixed estimates preceding Section~\ref{sec:missing}
use different dimension factors. We compare them with the same
factors on the period and theta sides, on the entire interval.

\begin{corollary}[Transport of the two normalized Gram estimates]
\label{cor:blocks239}
Use the original data and the maps \(S,\widehat\Theta\) of
Theorem~\ref{thm:minors239}. For each block choose a finite
Hilbert space \(V_i\) of dimension \(1\le R_i\le\dim D_{\rm all}\)
and a map \(X_i:V_i\to D_{\rm all}\). These maps can include
the column restrictions and chosen scalar normalizations.
Put \(B_i=SX_i,\ C_i=\widehat\Theta B_i\) and
\(G_{ij}=B_i^*B_j,\ H_{ij}=C_i^*C_j\). Use \(D_*\) from
Lemma~\ref{lem:dimension239}. Suppose for one of the two families
\(K=G,H\) the diagonal quantities
\(\|K_{ii}-1\|_{\rm HS}/\sqrt{R_i}\) are at most one.
Then on every sufficiently distant common interval,
\[
 \frac{\|H_{ii}-G_{ii}\|_{\rm HS}}{\sqrt{R_i}}\le\eta_x,
 \qquad
 \sqrt{\frac{D_*}{R_iR_j}}\,\|H_{ij}-G_{ij}\|_{\rm HS}
       \le\eta_x,\qquad
 \eta_x=\exp\!\big(-(\log 2)x^{377}/3\big).
\]
Hence the normalized diagonal errors, and the normalized mixed
Hilbert--Schmidt norms, differ by at most \(\eta_x\) between the
two families. The estimate applies to all pairs with the same
choice of theta tests.
\end{corollary}
\begin{proof}
On the span of the images write \(L\) for the compressed
normalized theta Gram operator from Theorem~\ref{thm:minors239}.
Let \(g_i=R_i^{-1}\operatorname{tr}G_{ii}\).
Cauchy--Schwarz for the matrix trace shows
\(|\operatorname{tr}(K_{ii}-1)|\le\sqrt{R_i}\|K_{ii}-1\|_{\rm HS}\).
The hypothesis gives \(\operatorname{tr}K_{ii}\le2R_i\).
The quadratic comparison in that theorem consequently gives
\(g_i\le2e^{\Delta_x}\) in either case.

Each entry of \(H_{ij}-G_{ij}\) in orthonormal bases is a
pairing of one column of \(B_i\) with \(L-1\) applied to a
column of \(B_j\). Summing the squared Cauchy--Schwarz estimates
for all entries gives
\[
   \|H_{ij}-G_{ij}\|_{\rm HS}
       \le (e^{\Delta_x}-1)\sqrt{R_iR_jg_ig_j}.
\]
Divide by \(\sqrt{R_i}\) for the diagonal bound and multiply
by \(\sqrt{D_*/(R_iR_j)}\) for the mixed bound.
The remaining dimension factors are at most \(\exp(x^{40}/2)\)
by Lemma~\ref{lem:dimension239}. The estimate for \(\Delta_x\)
in the proof of Theorem~\ref{thm:minors239} absorbs these factors
and \(2e^{\Delta_x}\) in the stated \(\eta_x\).
The reverse triangle inequality gives the final assertion.
\end{proof}

In particular an estimate for the signed orbital Gram matrices
at \(K^0,\varpi\) would control these complex theta Gram matrices
at the precision used in Section~\ref{sec:missing}. The signed
orbital estimate itself, and the lower bound for the pairing
with the designated integral row in Corollary~\ref{cor:original200},
are still required there.

\subsection{The arithmetic period of the definite column tower}\label{nav:I-51}
Use the fields and the two towers of Proposition~\ref{prop:ranks47}
throughout the interval of Section~\ref{sec:missing}. All the
column groups \(H_i=U(X_i)\) are definite. The last row pair
\(U(Y_{m-1}),U(Y_m)\) is definite as well. The column ranks
are even, and the last two row ranks differ by the odd integer
\(t_{\rm r}\). Thus the rank product in each of these pairs is
even, as required in Theorem~\ref{thm:vac242}. For the signed
column operator use \eqref{eq:iter} with its fixed splittings and
finite-order intermediate determinant characters.

We use a single marked elliptic differential on each \(E_x\).
Theorem~\ref{thm:edge242} applies in the oriented matrix coordinates.
In the tensor convention the action on a coefficient matrix is
\(Z\mapsto h Zg^t\), which is the oriented action evaluated at
\(((g^{-1})^t,h)\) in the standard definite bases. If the
alternating form has the opposite sign, \(Z\mapsto\bar Z\) on
rational matrices is an isometry to the chosen form, since it
negates \(B\). Carrying both the group action and the finite and
real tests by this isometry gives the same theta sum by the
bijection on rational summation variables. Both descriptions use
the original integration variables and determinant characters.

\begin{lemma}[The real central characters fix the surviving degrees]
\label{lem:degrees242}
Suppose the initial function in the definite column integral has
one character on the real center and each real Schwartz vector is
a Fock polynomial homogeneous at each place. With the splittings
fixed, every column giving a nonzero integral has exactly the same
degree \(d_{i,\tau}\) at edge \(i\) and real place \(\tau\).
This applies throughout the determinant-twist family of
Theorem~\ref{thm:block47}.
\end{lemma}
\begin{proof}
Fix \(\tau\). On the matrix coordinates the central circle at
either end of an edge acts by \(z\) or \(z^{-1}\). Rotation
preserves the vacuum line. Hence on homogeneous degree \(d_i\)
that circle acts by \(z^{a_{i,e}+\epsilon_{i,e}d_i}\), where
\(e=L,R\) denotes the end, \(a_{i,e}\) is the vacuum weight,
and \(\epsilon_{i,e}\in\{-1,1\}\) is determined by the
action on the degree-one variables. The vacuum weight is an
integer because the splitting gives a circle character. The
argument also works after transporting the action as above.
Write \(b_0\) for the input weight at \(\tau\). Translating
the first compact integration variable by this central circle
shows that nonvanishing requires
\(b_0+a_{1,L}+\epsilon_{1,L}d_1=0\).
At each interior variable the determinant twist is trivial on
the connected circle, because it has finite order. Translation
there requires
\[
 a_{i,R}+\epsilon_{i,R}d_i+
 a_{i+1,L}+\epsilon_{i+1,L}d_{i+1}=0 .
\]
Indeed any nontrivial power of \(z\) has circle average zero.
The initial equation fixes \(d_1\) uniquely, and each succeeding
one fixes \(d_{i+1}\) uniquely. A negative or nonintegral
required degree gives zero for every column. Applying the
argument separately at each real place proves the lemma.
\end{proof}

\begin{theorem}[CM derivatives for the projected definite column period]
\label{thm:column242}
Pair the restriction of \eqref{eq:iter} to
\(H_*=U(Y_m)\subset U(X_m)\) with a fixed vector \(b\) of the
chosen terminal test space. An orthogonal projection onto that
space may be taken before pairing. Allow any finite signed sum
of the determinant twists and retain all their prescribed finite
tests. Write \(a\) for the initial vector. Use the polynomial
coordinates of Theorem~\ref{thm:vac242}. For an arithmetic
statement suppose the finite tests have algebraic coefficients
in the finite lattice model, the determinant factors and split
Weil actions are as in Theorem~\ref{thm:one61}, and the finite
class values of \(a,b\) are algebraic matrix coefficients at
infinity. Take algebraic coefficients in the real test polynomials
and in the signed sum. Under these hypotheses the actual period
has the form
\[
 \mathcal P=V_J\left(\prod_{i=1}^m C_i\right)c^{-t/2}\mathcal A,
 \qquad
 t=\sum_{i,\tau} d_{i,\tau},\qquad
 |C_i|=\nu^{-[E:K]\,N_{i-1}N_i/2}.
\]
Here \(N_i=\dim_E X_i\), \(c=\pi/A\) and \(\nu\) are those of
Theorem~\ref{thm:vac242}. The factor \(V_J\) is the product of the volumes
of the real group times the finite integration level at the
vertices \(0,\ldots,m-1,*\). The scalar \(\mathcal A\) is an
explicit sum of products of ordinary CM derivatives, with the
original signs. If no degree profile survives then
\(\mathcal P=0\).

More precisely, at common levels fixing the integrands let
\(A_j\) be the least common multiple of the rational stabilizer
orders at vertex \(j\). For each \(j,\tau\) let \(L_{j,\tau}\)
bound the degrees in the ordinary and the conjugate entries
of that compact matrix variable in the expansion below.
Write \(n_j\) for the dimension of the group being integrated,
so \(n_*=N_m-1\). Let \(I\) be the fractional ideal generated
by the nonzero coefficients of this expansion, including its
CM derivatives, in any common number field. Use the zero ideal
if all the coefficients vanish. Then
\[
 \left(\prod_j A_j\right)
 \left(\prod_{j,\tau}\mathcal D(n_j,L_{j,\tau})\right)
                  \mathcal A\ \in I .
\]
In particular the compact integrations introduce no denominator at
a prime exceeding every \(n_j+L_{j,\tau}-1\) and prime to
\(\prod_j A_j\).
\end{theorem}
\begin{proof}
Fix the endpoint argument \(g_*\oplus1\) in the last kernel,
multiply by \(\overline{b(g_*)}\), and integrate over \([H_*]\).
For an orthogonal projection \(p\) and \(b\) in its range,
\(\langle p u,b\rangle=\langle u,b\rangle\).
If a fixed test outside its range is used, replace \(b\) by
\(p b\) and use its actual coefficients. Intersect the finite
stabilizers of the data on all incident edges, the character
kernels, and the levels of \(a,b\). At the endpoint take the
inverse image of the last test's stabilizer. This gives compact
open subgroups inside the prescribed ones because the number of
tests is finite. Splitting the original quotients into classes
at these levels evaluates the same integrals. In particular a
finite character projector prescribed in the data keeps every
term of its defining finite average and its divisor.

For each tuple \(\mathbf u\) of finite representatives, choose
the representatives in Theorem~\ref{thm:class61} away from the
finitely many bad places. Denote by \(\phi_{i,\mathbf u}\)
the actual finite translate of the test on edge \(i\). Write
\(k_j\) for the compact real matrix variables, using \(k_*\oplus1\)
at the endpoint. On their product the integrand after removal of the common scalar is
\[
 a(k_0,u_0)\,\overline{b(k_*,u_*)}
 \prod_{j=1}^{m-1}\chi_j(\det u_j)
 \prod_{i=1}^m
   \bigl[\omega_{i,\infty}(k_{i-1},k_i)P_i\bigr](\partial)
                                 G_{\phi_{i,\mathbf u}}(0).
\]
The symbol \(k_m\) here means \(k_*\oplus1\).
Form this expression for every summand with its coefficient,
omitting the degree profiles that integrate to zero by the
lemma. Expand it as a polynomial in all the entries \(k_j,\bar k_j\).
Define \(\mathcal A\) by taking the probability Haar average
of each such polynomial, dividing by the product of its vertex
stabilizer orders, and summing over tuples and signed terms.
This is the expansion defining \(I\), before averaging and
division by stabilizers.

To verify the period formula apply Theorem~\ref{thm:vac242}
pointwise to every kernel, then use the class integration
argument of Theorem~\ref{thm:edge242} at every vertex. Each
kernel contributes \(C_i c^{-\sum_\tau d_{i,\tau}/2}\). The
lemma allows their product to be taken outside the signed sum.
Finite sums, compact integrals, and the uniformly convergent
theta kernels admit these interchanges. The same calculation
for tests written as \(P_i(\overline{\Phi(a')})\) in raw
matrix coordinates uses \((c\bar\delta)^{-t}\) in place of
\(c^{-t/2}\), since \(\bar w=\bar\delta\,\overline{\Phi(a')}\).
If the original theta character is used, first transport its
vectors and splittings by Theorem~\ref{thm:transport233}. That
identity holds pointwise and retains the given terminal character.

It remains to justify the coefficient assertion. At good split
places a class translate acts on a unit lattice vector by normalized
change of variables and its algebraic character factor. The change
of lattice is a finite sum with coefficients that are roots of
unity times inverse square roots of indices, by
Lemma~\ref{lem:split61} and Theorem~\ref{thm:one61}. The same
argument applies to a translated vector by conjugating its
translation. Any finite adelic translate of the new vacuum has
a rational lift modulo its periods, and changing the lift only
introduces a root of unity. At the other places the chosen
representatives act trivially. Thus the resulting sections
have algebraic Taylor coefficients by Theorem~\ref{thm:transport59}.
At finer levels a rational change of representative acts on
endpoint matrix coefficients by an algebraic matrix. The
old level factors on algebraic endpoint data are retained in
their class coordinates.

On the compact real groups the polynomial action is substitution
of unitary matrices, times the vacuum character. This character
is a product of integral powers of determinants, since every
character of \(U(n)\) factors through its determinant. Inverses
of determinants on a unitary group are conjugate determinants.
The same description holds for the stipulated algebraic matrix
coefficients, since every algebraic representation of
\(\mathrm{GL}_n\) has regular matrix coefficients in the entries
and the inverse determinant. Complex conjugation preserves
algebraicity. All coefficients just described, and all ordinary
derivatives occurring here, therefore belong to a single number
field for the fixed finite family. Proposition~\ref{prop:moments242}
puts the compact average in \(I/\prod_{j,\tau}\mathcal
D(n_j,L_{j,\tau})\). Clearing the finitely many stabilizer
denominators proves the assertion for \(\mathcal A\).
\end{proof}

At the last pair of the row tower Theorem~\ref{thm:edge242}
also evaluates the integral against the incoming form, with its
actual coefficients. The earlier row passes through indefinite
groups. The size and nonvanishing of its pairing with the
designated column in Corollary~\ref{cor:original200} still
require a quantitative argument on the common interval.

\subsection{The two integral positions along the row}\label{nav:I-56}
Fix a field and interval of Theorem~\ref{thm:third263}. At a
horizontal edge put \(S=V_{i,0}\), \(T=V_{i+1,0}\), with
\(2\le i<m\). The staircase of \(S\) has \(c=i-1\) columns
\(b_j\), and that of \(T\) has \(c+1\) columns \(a_j\), numbered
from one. Write \(r=\lfloor c/2\rfloor\) and
\(\ell=\lfloor(c+1)/2\rfloor\). Let \(S_0,T_0\) be the orthogonal
complements of the \(r,\ell\) hyperbolic planes, respectively.
Their lattices are the intersections with the given full lattices.
In either staircase opposite columns pair to one, and the middle
column in odd length has norm \(-1\). Put
\[
 d=-\xi_0,\qquad A=\mathfrak A/\bar{\mathfrak m},\quad
 B=\bar{\mathfrak A}/\mathfrak m,\quad
 I_j^{(c)}=A^{c+1-j}B^{-j},\quad
 \mathfrak j=\mathfrak m\bar{\mathfrak m}.
\]
Thus \(\Lambda_S\) contains the direct summands \(I_j^{(c)}b_j\);
the expression with \(c+1\) gives the analogous summands in
\(\Lambda_T\). These formulas describe fractional ideals and
require no global choice of generators.

Use two successive polarizations of \(\operatorname{Hom}(-T,S)\).
First retain the positions \(x_j=X(a_j)\), \(1\le j\le\ell\).
In the remaining map from \(T_0\) to \(S\) retain the coefficient
\(p_k\in T_0^\vee\) of \(b_{c+1-k}\), \(1\le k\le r\). Here
\(T_0^\vee\) is the space of \(E\)-linear functionals. The
remaining oscillator is that on \(\operatorname{Hom}(-T_0,S_0)\).
In the scaled frame used for the Fourier character the first
coordinates are evaluated on \(d^j a_j\); their base values are
\(d^j b_j\). Denote the corresponding unscaled row values by
\(p_{*,k}\).

\begin{proposition}[Integral tests in the two polarizations]\leavevmode
\label{prop:positions264}
Put \(L_j=(I_j^{(c)})^{-1}\Lambda_S\) and
\(M_k=\mathfrak A I_{c+1-k}^{(c)}\Lambda_{T_0}^{\vee}\), where
\(\Lambda_{T_0}^{\vee}\) is the integral linear dual. At all finite
places take the product of the indicators
\[
  x_j\in b_j+\mathfrak j\widehat L_j,\qquad
  p_k\in p_{*,k}+\mathfrak m\widehat M_k .
\]
For \(c=2r+1\) all \(p_{*,k}\) are zero. For \(c=2r\) the only
nonzero row, \(p_{*,r}\), vanishes on the core in \(T_0\) and
takes a value in \(\{1,-1\}\) on the middle column \(a_{r+1}\).
These tests tensored with an original residual lattice test have
the prescribed depth. If the residual test has one marked
coset per place, the change to the original split polarizations
has exactly
\[
              (N\mathfrak m)^{\ell\dim_E S+r\dim_E T_0}
\]
distinct products of marked integral coset indicators. In the
measures of Proposition~\ref{prop:shell263} their coefficients
are roots of unity. At every unmarked place the resulting test
is the original lattice vacuum, with the lattice normalization.
\end{proposition}
\begin{proof}
The successive frame scalars are \(d^j\) on both sides, including
the last vector \(d^\ell b_\ell\) when it has norm of nonzero
absolute value. They cancel in \(X(a_j)\). For \(c=2r\) the last
row of the second position uses the vector \(-d^{r+1}a_{r+1}\)
and the adjoint multiplier \(-\xi_0^{-1}\). The physical opposite
column paired to \(d^r b_r\) is
\(\bar d^{-r}b_{r+1}\). Thus the scalar multiplying the row
\(h_{T_0}(a_{r+1},\,\cdot\,)\) is
\(\bar d^{-r}(-\xi_0^{-1})(-\bar d^{r+1})=1\).
Reversing the input Hermitian form reverses the row's sign.
This gives a unit at every finite prime, including primes
dividing the absolute different. For \(c=2r+1\) the grid
prescribes the zero second position.
\srcloc{Development/AbelianVarietyTheory/Definitions/TwistedInnerProducts/IsAuxiliaryVectorWittDec}{160--194}.
\srcloc{Development/AbelianVarietyTheory/Definitions/TwistedInnerProducts/WittDecompSymplInclIsotropicInclAnisotropic}{183--201}.

The undeepened map lattice on the \(j\)-th input column is
\(\mathfrak A(I_j^{(c+1)})^{-1}\Lambda_S=\bar{\mathfrak m}L_j\),
because \(I_j^{(c+1)}=A I_j^{(c)}\). On the retained output row
it is \(M_k\). Each is paired with its opposite column or row by
exact trace duality, by Proposition~\ref{prop:tracelat253}.
Different hyperbolic planes are mutually orthogonal, and the
second change uses only the residual map space after the first.
For a retained lattice \(D_1\) of either kind the allowed
translations are \(\mathfrak m D_1\). The allowed support is
\(\bar{\mathfrak m}^{-1}D_1\), the annihilator of the deepened
opposite lattice. We have \(b_j\in L_j\). In the case of the
nonzero row its values on \(a_{r+1}\), as elements of \(M_r\),
range over the ideal
\[
  \mathfrak A I_{r+1}^{(2r)}(I_{r+1}^{(2r+1)})^{-1}
      =\bar{\mathfrak m}.
\]
The middle line is a direct summand and the row vanishes on its
complement, so \(p_{*,r}\in\bar{\mathfrak m}^{-1}M_r\).
All the indicated cosets therefore satisfy both depth conditions.

At a marked pair \(w,\bar w\) of depth \(a\), retain the given
coset on the \(w\)-half of each position. On the other half the
test is the translate of the full unshifted position lattice,
since \(\mathfrak m_{\bar w}\) is the unit ideal. Its inverse
Fourier transform is the full opposite lattice at \(w\) times
the character of the center of that translate. This follows by
integrating the character over the lattice in dual measures
with that lattice of volume one. The character is trivial on
the sublattice of depth \(a\) by the support calculation above.
Splitting the opposite lattice into these cosets gives \((Nv)^{a e}\)
terms for a position of \(E\)-dimension \(e\). They have disjoint
supports. Multiplication over the independent positions and
marked places proves the count. Away from them each center is
integral, and Fourier transform interchanges full dual lattice
indicators. This proves the vacuum assertion also at split
places above two.
\end{proof}

\subsection{A finite column orbit for all the terminal characters}\label{nav:I-73}
Fix the field, interval, framed finite row, terminal frequency
and finite column level of Theorem~\ref{thm:ordinarymoment304}.
Let \({\cal O}\) run over the scalar orbits of the fine column
classes, each of size \(|Z|=(T-1)^f\). For a prescribed
terminal space let \(J_{\cal O}\) be its subspace supported
on this single orbit. Define \(M_{J,{\cal O}}\) as the mean
squared norm of the ordinary periodization when the signed
column of Theorem~\ref{thm:dyad234} is sampled inside
\({\cal O}\).

\begin{theorem}[One set of orbits for the prescribed real words and characters]\leavevmode
\label{thm:commonorbit305}
Apart from a fraction at most \(\exp(-fkd\log k/3)\) of
the scalar orbits, one has simultaneously
\[
             |M_{J,{\cal O}}-1|
                   \le \exp\bigl(-(1-o(1))fk\log k\bigr)
\]
for all prescribed words sharing the fixed framed data and for
every terminal character nontrivial at each place, as well as
for the simultaneous zero-sum space. These orbits can also
be intersected with the primitive determinant set of
Theorem~\ref{thm:columns235}.
\end{theorem}
\begin{proof}
For any single orbit the covariance of the signed column
is \(P_{J,{\cal O}}/\dim J_{\cal O}\), by the finite averaging
in Theorem~\ref{thm:dyad234} restricted to this orbit.
The normalized trace of the identity is one, and
Theorem~\ref{thm:allflags298} bounds the proper-intersection
and transverse full-column error on every orbit by
\(((k+1)!)^{-f}\exp(Crf)\) through its operator norm.

We exhibit one nonnegative bound \(e({\cal O})\) for the
remaining trace error, independent of the character and
the word. In the full space, a matching smaller-column
fine cell has absolute normalized trace at most its
unipotent covolume times
\[
   \frac{\operatorname{vol}_f K^0}
          {n_{\rm cl}\,d_\rho c_{\rho,K^0}}\,
                      \prod_\tau{\cal P}_\tau
\]
by Proposition~\ref{prop:tracemass301}. The product
\(d_\rho c_{\rho,K^0}\) depends only on the fixed framed
data. Replace the real product by the infimum of the
explicit upper bounds obtained by multiplying
Proposition~\ref{prop:phase301} over the places with any
permitted endpoint powers. All these bounds are uniform
in the word. This gives a common numerical majorant
for the cell. The matching test is intrinsic to the two
finite columns, independent of the character. The
underlying relevant labels for a single word depend
only on the two equal unipotent characters and the
fixed frames, by Lemma~\ref{lem:forms203}.

The sum of these majorants over all matching smaller-column
cells is at most \(\exp(-(\tfrac12-o(1))fkd\log k)\).
To verify this point, on the two ranges of
Theorem~\ref{thm:ranktrace302} its proof sums finite
measures and the real majorants with the indicated
endpoint powers. The rank-zero sum in
Theorem~\ref{thm:zerotrace303} uses only the index of
the row and the same positive measures. On the complement,
Proposition~\ref{prop:flagtrace304} counts the measures
after the uniform covolume bound, and
Theorem~\ref{thm:ordinarymoment304} uses endpoint powers
zero for its real factors. Thus each range admits the
claimed estimate for these positive majorants before
summing any signed contributions.

Set \(e({\cal O})\) equal to the sum of the majorants of
cells whose matching columns lie in \({\cal O}\), multiplied
by the total number \(n_{\rm cl}/|Z|\) of orbits. This
factor accounts for normalization on one orbit. Indeed
the projector coefficient after division by dimension
is bounded in modulus by \(1/(|Z|d_\rho)\) there, for
either terminal space. Each matching cell lies on just one
orbit. Hence the uniform mean of \(e({\cal O})\) is the
full sum just bounded. Markov's inequality shows that
\(e({\cal O})>\exp(-fk\log k)\) on a fraction at most
\(\exp(-fkd\log k/3)\) for sufficiently large \(x\).
Add the other operator-norm bound on every remaining
orbit to obtain the theorem. The earlier determinant
set is a union of whole orbits of fraction tending to one,
so the intersection also has fraction tending to one.
\end{proof}

\subsection{The arithmetic functional on the actual incoming row}\label{nav:I-76}
On the interval of Section~\ref{sec:missing} use the same word,
levels, and measures as in Theorem~\ref{thm:lastrow311}.
For a nonzero actual incoming row \(X\), take the detecting
endpoint \(V\) given there. At the fine source classes write
\(X(k,u)=\langle\rho_-(k)t_-,Y_u\rangle\), with the conjugations
of Theorem~\ref{thm:last244}. Use the ordered coordinate wedge
tensors \(e_I\) of the ambient source tensor product.

\begin{theorem}[The last arithmetic functional on the preceding lift]\leavevmode
\label{thm:pullback311}
There is an array \(a_{u,I}(V)\) in the arithmetic coefficient
field, independent of \(X\) once \(V\) is fixed, such that
\[
 \frac{D_-D_+c^{d/2}}{w_-w_+ C}\langle TX,V\rangle
                  =\sum_{u,I}(Y_u)_I a_{u,I}(V).
\]
At every \(\mathfrak p\mid p\) in Theorem~\ref{thm:bothlast311}
its entries have valuation at least \(e(\mathfrak p/p)g_p\).
At every complex embedding they have absolute value at most
\(i_+ B_+B_\sigma\) in that theorem's notation. On the source
Cartan spaces the array satisfies
\[
 \left(\sum_u
   \big\|a_{u,\cdot}(V)|_{V_n}\big\|_{\rm op}^2\right)^{1/2}
 \ \ge\ \frac{i_+c^{d/2}\sqrt{D_-a_*}}
                 {2|C|\sqrt{w_-v_+}} .
\]
Here \(\|\cdot\|_{\rm op}\) is the norm of a linear functional
in the unitary tensor norm. These statements use the same
endpoint for the actual Bessel pairing and for the arithmetic
array.
\end{theorem}
\begin{proof}
Expand only the endpoint character sum in the class formula
of Theorem~\ref{thm:last244}. At a source class \(u\) evaluate
it on \(e_I\), using the minor formula on the entire exterior
tensor product. This defines \(a_{u,I}(V)\) as a sum of
\(i_+\) coordinate CM derivatives with the endpoint tensor
inserted. Expanding the actual \(Y_u\) in these coordinates
gives the identity. The same integral polynomial calculation
as in Theorem~\ref{thm:bothlast311} gives the finite valuations.

For the estimate at infinity, fix an endpoint Cartan tensor
\(Z'\). On the ambient source tensor space the map in
Proposition~\ref{prop:cartan244} satisfies
\(\|M(y\otimes Z')\|_{\rm F}\le
\sqrt{h_\lambda}\|y\|\|Z'\|\). Indeed the Cauchy decomposition
of the homogeneous matrix polynomials pairs the target
partition only with the same source partition. Its multiplicity
in the source is one because a diagonal one-parameter subgroup
with strictly decreasing weights
has a unique maximal weight line there, obtained by choosing
the leading wedge in every slot. It is the highest line of
the Cartan component. Equivariance consequently kills the
orthogonal complement of that component when the target is
in its Cartan component. On the Cartan component the asserted
norm is the isometry formula, including its scalar
\(\sqrt{h_\lambda}\). Apply this argument at each embedding
with \(y=e_I\) of norm one and then the sphere estimate of
Theorem~\ref{thm:bothlast311}.

Finally \(\sum_u\|Y_u\|^2=D_-\|X\|^2/w_-\) by compact
orthogonality. Apply Cauchy--Schwarz to the displayed identity
and insert the lower bound on \(|\langle TX,V\rangle|\) from
Theorem~\ref{thm:lastrow311}. Cancelling \(\|X\|\) gives the
claimed lower bound.
\end{proof}

There is also an exact conversion to the intrinsic integral
class lattices if those lattices are used for the tests.
On each fine chart let \(\Lambda_u\) be the lattice of matrix
coefficients generated by the Schur and determinant models of
Proposition~\ref{prop:schurframe259} and its dual. Let
\(H_u\) be the module generated by the coefficients, as a
polynomial in \(\bar Z\), of
\(\langle\rho(k)t_R,v_R(Z)\rangle\). Both are taken over the
ring of integers of a common coefficient field. The quotient
\((H_u+\Lambda_u)/\Lambda_u\) is finite. Indeed the two
realizations of the partition are defined over the coefficient
field by the exterior and Young tensor constructions. The
matrix coefficients of the Schur lists span the isotypic
space over \(\CC\) by their frame bounds. Linear equations
over the coefficient field with a complex solution already
have a solution over that field, by row reduction. Clearing
denominators for the finitely many generators of \(H_u\)
proves finiteness. Thus the least common multiple of these
quotient exponents over the source charts, respectively over
the target charts, gives integers \(b_-,b_+\) directly from
the two lattices. Multiplying the integral-family tests by
these integers puts them in the intrinsic class lattices.
The character sums preserve the transported lattices.
The minors and both of their size bounds then acquire
the same factor \((b_-b_+)^j\).

For the full scalar, the displayed identity still contains
the class coefficients \(Y_u\) of the preceding lifts.
The remaining period estimate in Section~\ref{sec:missing}
requires their arithmetic normalization and the prescribed
signed contraction. The local cohomology remainder is given
in Theorem~\ref{thm:entireinterval310}.

\subsection{Comparison for one signed definite column period}\label{nav:I-77}
Use the single scalar \(\mathcal A\), its data and its interval
from Theorem~\ref{thm:finitecolumn313}. Index a summand by the
signed term \(b\) and its tuple of finite representatives
\(\boldsymbol u\). Write \(s_{b,\boldsymbol u}\) for the scalar
coefficient outside the polynomials, endpoints and class sections,
including finite twists and prescribed character averages in that
position. Let \(\gamma_{\boldsymbol u}\) be the product of its
vertex stabilizer orders. Write \(a_{b,\boldsymbol u}\) and
\(\widetilde b_{b,\boldsymbol u}\) for its two endpoint
polynomials in the compact variables. Here \(\widetilde b\) is
the polynomial actually occurring after the convention for
conjugating the endpoint has been applied.

For an embedding \(\sigma:J\to\CC\) apply \(\sigma\) to the
algebraic coefficients in this expansion. Set
\[
 F_{b,\boldsymbol u,\sigma}
 =\sup_{k_0} |a_{b,\boldsymbol u}^{\sigma}(k_0)|\,
       \sup_{k_*}|\widetilde b_{b,\boldsymbol u}^{\sigma}(k_*)|.
\]
The suprema run over the two compact endpoint groups and can also
be bounded by the sums of the absolute coefficients. Write
\(P_{i,b}\) for the real Schwartz polynomial in the dimensionless
Fock coordinates of Theorem~\ref{thm:column242}.
Let \(b_{i,b,\boldsymbol u,\sigma}\) bound the sum of the absolute
coefficients in the actual finite test expanded in translates of
its normalized lattice vacua after conjugation by \(\sigma\).
These coefficients come from the finite changes of lattice and
finite Fourier sums of Theorem~\ref{thm:bothlast311}, applied to
the indicated edge. With \(D_i=[E_x:K]\ell_i\), write
\[
 \begin{split}
 W_{i,b,\boldsymbol u,\sigma}^{\,2}
  ={}& b_{i,b,\boldsymbol u,\sigma}^{\,2}
   [e\max(1,2\ell_i\operatorname{rd}(E_x)/\pi)\nu_\sigma]^{D_i}
          c_\sigma^{d_i}\prod_\tau H_{\ell_i}(d_{i,\tau}),\\
 U_\sigma={}&\sum_{b,\boldsymbol u}
   \frac{|\sigma(s_{b,\boldsymbol u})|F_{b,\boldsymbol u,\sigma}}
        {\gamma_{\boldsymbol u}}\,
   \prod_i \big(\|P_{i,b}^{\sigma}\|_{\rm F}
                         W_{i,b,\boldsymbol u,\sigma}\big).
 \end{split}
\]
Here \(c_\sigma=\pi/A_\sigma\), \(\nu_\sigma\) is the squared norm
of the marked elliptic vacuum, and \(H_\ell\) is the explicit
sphere factor in Lemma~\ref{lem:sphere60}. The embeddings of a CM
block are permuted with \(\sigma\).

\begin{theorem}[Both estimates for the full signed column]\leavevmode
\label{thm:archcolumn313}
The original projected column satisfies
\(|\sigma\mathcal A|\le U_\sigma\) at every embedding. Put
\[
 h_i=\frac1{[J:\QQ]}\sum_{p\notin S}\sum_{\mathfrak p\mid p}
  [J_{\mathfrak p}:\QQ_p]\Gamma_{i,\mathfrak p}\log p,\qquad
 G_S=\sum_{p\in S}\sum_i g_{i,p}\log p .
\]
For any nonzero \(\mathcal A\) from these prescribed data,
\[
 G_S\le \log D_0+\sum_i h_i+\frac1{[J:\QQ]}\sum_\sigma\log U_\sigma .
\]
At the original embedding \(\sigma_0\), its lower bound is
\[
 \log|\sigma_0\mathcal A|\ \ge\
 [J:\QQ]\bigl(G_S-\log D_0-\sum_i h_i\bigr)
                -\sum_{\sigma\ne\sigma_0}\log U_\sigma .
\]
\end{theorem}
\begin{proof}
After the Haar integrals have been expressed by their rational
moments, applying \(\sigma\) applies it to the class derivatives
and polynomial coefficients and keeps each moment fixed.
Consequently one may take the same compact averages of the
coefficientwise transported integrands. In particular the
coefficients of \(\widetilde b\) are transported after its original
conjugation. The integrals use probability Haar at every real
vertex because their volume factors are already in \(V_J\).

For a fixed edge and tuple, Theorem~\ref{thm:transport59} bounds
the squared metric supremum of the normalized lattice vacuum
at \(\sigma\) by
\([e\max(1,2\ell_i\operatorname{rd}(E_x)/\pi)\nu_\sigma]^{D_i}\).
Hermitian translations preserve the metric supremum, so the
coefficient sum gives the additional factor
\(b_{i,b,\boldsymbol u,\sigma}^{\,2}\).
Lemma~\ref{lem:sphere60} bounds the Fischer norm of its specified
homogeneous Taylor component by \(W_{i,b,\boldsymbol u,\sigma}\),
after taking this finite sum. For a polynomial \(P\) of those
block degrees, \(P(\partial)G(0)\) is the bilinear Fischer
contraction with that Taylor component. Indeed the contribution
of \(Z^\alpha\) in each factor is weighted by \(\alpha!\).
Cauchy--Schwarz bounds it by \(\|P\|_{\rm F}W_{i,b,\boldsymbol u,\sigma}\).

At every compact argument the polynomial in the integrand is
obtained from \(P_{i,b}^\sigma\) by unitary matrix substitution
and a determinant character of absolute value one. This is the
compact action proved in Theorem~\ref{thm:column242}, expressed
by substitutions with integer coefficients in the matrix variables.
Transposing a matrix and embedding \(k_*\) as \(k_*\oplus1\)
retain unitarity. The Fischer norm is invariant by its Gaussian
integral formula. Therefore the last bound holds uniformly over
the compact variables, including the common variables on adjacent
edges. Bound the endpoint factors by \(F_{b,\boldsymbol u,\sigma}\),
integrate and sum the absolute values. This proves the first assertion
on the scalar defined with all the original signs.

If \(\mathcal A\ne0\), apply the product formula to \(D_0\mathcal A\).
Theorem~\ref{thm:finitecolumn313} gives a finite-place sum at least
\(G_S-\sum_i h_i\) in the normalization by \([J:\QQ]\).
Thus the mean of \(\log|\sigma\mathcal A|\) is at least
\(G_S-\sum_i h_i-\log D_0\). Bound every embedding above, or all
of them except \(\sigma_0\), to obtain the two inequalities.
Each \(h_i\) is finite and bounded explicitly by the last
inequality of Theorem~\ref{thm:periodden312} at that edge.
\end{proof}

Returning to the complex period multiplies
\(|\sigma_0\mathcal A|\) by \(V_J|\prod_i C_i|c^{-t/2}\).
The comparison uses the entire definite column. Section~\ref{sec:missing}
also requires a lower bound on its prescribed pairing with the
arithmetic normalization of the Bessel row, and control of the
errors on the scale of the weighted prime count.

\subsection{The discriminant bound for a detected terminal minor}\label{nav:I-78}
Use the original field, word, levels and the two integral families
in Theorems~\ref{thm:lastrow311} and~\ref{thm:minors311}.
For their selected nonzero \(j\)-row minor \(z_j\), take \(T,Q\)
from the fixed finite vacuum model of Theorem~\ref{thm:adapted314}
for every class section entering this matrix. Write \(d\) for
the last-edge degree and
\[
 B_{\rm fin}=\log T+d\log Q+\tfrac d2\log\operatorname{rd}(E)
 -\left(\tfrac d4-\tfrac12\lfloor d/2\rfloor
                    \mathbf1_{d_0\nmid Q}\right)\log d_0 .
\]

\begin{theorem}[Two-sided control for the selected arithmetic minor]\leavevmode
\label{thm:terminalheight314}
At the original embedding the selected minor satisfies
\(|z_j|\ge j!(i_-i_+L_*)^j\), with the common terminal test of
Theorem~\ref{thm:lastrow311} when \(j=1\). For \(S\) as in
Theorem~\ref{thm:bothlast311} away from \(TQ\), the very same
minor gives
\[
 \sum_{p\in S}g_p\log p\le B_{\rm fin}+
       \tfrac12\log j+\log(i_-i_+)+\tfrac d2\log(ns)
                     +\frac1{[J:\QQ]}\sum_\sigma\log B_\sigma .
\]
Both \(T\) and \(Q\) are determined before any elliptic product
basis is chosen.
\end{theorem}
\begin{proof}
Keep the physical derivative coordinates and both test families
fixed. For each \(p\) change just the product basis at the last
edge to the one in Theorem~\ref{thm:adapted314}. Every matrix
entry acquires one common scalar by
Theorem~\ref{thm:framechange314}. The determinant acquires its
\(j\)-th power. Therefore \(z_j\overline{z_j}\) is identical
in the two evaluations. Every term of the determinant contains
exactly \(j\) derivatives of the given block degrees with integral
exterior tensor and finite character coefficients. The normalized
matrix of Theorem~\ref{thm:bothlast311} already uses the fine
levels and its unnormalized character sums, so there is no
compact integration or stabilizer denominator left in this
matrix.

At \(p\in S\) use the full-column minor weights in
Theorem~\ref{thm:frob244} for each factor, as in
Theorem~\ref{thm:bothlast311}. At the remaining primes use
Theorem~\ref{thm:adapted314} in both conjugate factors.
As in Theorem~\ref{thm:gramheight314}, the sum of the two
coordinate losses per derivative is \(d\rho_p\). Thus at
every \(\mathfrak p\mid p\) the valuation of the diagonal
product is at least \(2jg_p\) in \(S\) and at least
\(-2j(v_p(T)+\Phi_p(d,v_p(Q))+d\rho_p/2)\) elsewhere.
The sum of this last loss divided by \(2j\), with logarithmic
prime weights, is \(B_{\rm fin}\) by the discriminant computation.
Apply the product formula to the nonzero diagonal product, use
the archimedean bound for \(z_j\) of
Theorem~\ref{thm:bothlast311}, and divide by \(2j\). The original
lower bound follows from the choice in Theorem~\ref{thm:minors311}.
\end{proof}

The minor with \(j=1\) uses a terminal test which also detects
the preceding Bessel row. The corresponding evaluation on that
row is the contraction in Theorem~\ref{thm:pullback311}.
Determining its arithmetic normalization and controlling the
remaining strict cohomology orders of
Theorem~\ref{thm:entireinterval310} are still necessary for the
interval estimate.

\subsection{One row isometry for all the period pairings}\label{nav:I-79}
The comparison in Section~\ref{sec:missing} uses a single
isometry on the common row. The source may have a different
isometry and a different nonzero normalization for each word.
This order of quantifiers is explicit in
\nolinkurl{Formalism/SharpRestriction/GlobalTheoremInputs.lean},
in the definition \texttt{BesselId}. The realization
\texttt{besS} in
\nolinkurl{Formalism/LocalRingCardinality/RhoDetPrimeEqMuTildeMulChiPow.lean}
pairs the embedded row with the periodized column. Here
is a finite version of the comparison using those pairings
on every test.

Let \(T_u:M_u\to E_0\) and \(C_u:M'_u\to E_1\) be the
Riesz representatives of two finite lists of pairings, using
the conjugate row for a bilinear pairing as in
Theorem~\ref{thm:closedperiod318}. The complex Hilbert spaces
have \(\dim E_0=\dim E_1=n\ge1\) and
\(\dim M_u=\dim M'_u\). For a nonzero operator \(T\) put
\(N(T)=TT^\dagger/\operatorname{tr}(TT^\dagger)\).

\begin{theorem}[A simultaneous comparison on all tests]\leavevmode
\label{thm:commonunitary318}
There are a unitary \(V:E_0\to E_1\), unitaries
\(U_u:M_u\to M'_u\) and nonzero scalars \(a_u\) with
\(T_u=a_u V^\dagger C_u U_u\) for all \(u\) exactly under
the following conditions. The two lists have the same zero
words, and, on their common nonzero words, every ordered
list \(i_1,\ldots,i_k\) of length at most
\[
       L=\min\{2n^2-1,\,2d^2\},\qquad
       d=\sum_{u:T_u\ne0}\dim M_u ,
\]
satisfies
\[
 \operatorname{tr}\bigl(N(T_{i_1})\cdots N(T_{i_k})\bigr)
   =\operatorname{tr}\bigl(N(C_{i_1})\cdots N(C_{i_k})\bigr).
\]
For nonzero words one may take
\(|a_u|^2=\operatorname{tr}(T_uT_u^\dagger)/
                    \operatorname{tr}(C_uC_u^\dagger)\).
\end{theorem}
\begin{proof}
The comparison implies unitary conjugacy of the normalized
output operators and hence equality of all the traces.
Conversely write \(P_i=N(T_i), Q_i=N(C_i)\) on the nonzero
words. In the direct sum of the two full matrix algebras
let \(B_j\) be the span of the paired products of length
at most \(j\), including \((I,I)\). If \(B_j=B_{j+1}\),
right multiplication by any generator \((P_i,Q_i)\)
preserves \(B_j\), so all further spaces agree. Their
dimensions are at most \(2n^2\). They are also at most
\(1+2d^2\), since every positive-length product is supported
on the span of the output images on its side, of dimension
at most \(d\). Indeed each generator kills the orthogonal
complement of this span. Starting with dimension one,
the whole algebra is therefore spanned at length \(L\).
The difference of the two traces vanishes on the identity
and on the tested products, hence on this algebra.

Both sets of generators are self-adjoint. For any polynomial
\(f\) in noncommuting variables, including a constant term,
trace equality for \(f^*f\) implies
\(\|f(P)\|_{\rm HS}^2=\|f(Q)\|_{\rm HS}^2\).
Consequently \(f(P)\mapsto f(Q)\) is a well-defined unital
trace-preserving isomorphism \(\phi\) of the generated matrix
algebras that commutes with adjoints.

We show directly that \(\phi\) is conjugation by a unitary.
In the first algebra write the identity as a sum of mutually
orthogonal minimal projections \(p_i\). To obtain them,
split a projection whenever its corner contains a nonscalar
self-adjoint matrix, using the spectral projections in that
corner. They are polynomials in the matrix and the corner
identity. Repeated splitting terminates by finite rank,
and each resulting corner is scalar, by taking real and
imaginary parts. If \(p_i B p_j\) contains a nonzero
matrix \(x\), where \(B\) is the first algebra, \(x^*x\)
and \(xx^*\) are positive scalar multiples of \(p_j,p_i\).
Their scalars agree by associativity applied to \(xx^*x\).
Rescale \(x\) to a unitary between these two ranges, zero
on the orthogonal complement of the first domain.
Adjoints and composition show that the relation of having
a nonzero corner is an equivalence relation.

In each class choose a reference projection \(p_e\) and
such a matrix \(v_i\) from \(p_e\) to \(p_i\), with
\(v_e=p_e\). Then \(v_i v_j^*\) spans \(p_i B p_j\)
in this class, since \(v_i^* yv_j\) belongs to the scalar
corner at \(e\) for every \(y\) in that corner \(p_iBp_j\).
The other corners vanish. The ranks of \(p_e\) and
\(\phi(p_e)\) agree by trace preservation. Choose a unitary
\(W_e\) between their ranges. The rules
\(v_i z\mapsto\phi(v_i)W_e z\) on the orthogonal ranges
therefore define a unitary \(V:E_0\to E_1\). On each
spanning matrix \(v_i v_j^*\), direct evaluation gives
\(V(v_i v_j^*)V^\dagger=\phi(v_i v_j^*)\). In particular
\(VP_iV^\dagger=Q_i\) for all the words.

For each nonzero word let \(s_u\) be the positive square
root of the stated trace ratio and put
\(D=s_u V^\dagger C_u\). Then \(DD^\dagger=T_uT_u^\dagger\).
The rule \(T_u^\dagger y\mapsto D^\dagger y\) is well-defined
and isometric, since the two squared norms agree for every
\(y\), and it maps onto the range of \(D^\dagger\).
The equal source dimensions allow it to extend to a unitary
\(U_u:M_u\to M'_u\). Taking adjoints gives \(T_u=D U_u\).
For a zero word take any source unitary and scalar one.
This also handles a list consisting entirely of zero words.
\end{proof}

For the signed definite periods the left traces are exactly
those computed in Theorem~\ref{thm:closedperiod318}. If
\(R_u=C_uC_u^\dagger\) on the Bessel row, the equalities
on nonzero words can be tested without division as
\[
 \operatorname{tr}\Big(\prod_j K_{i_j}\Big)
                  \prod_j\operatorname{tr}R_{i_j}
    =\operatorname{tr}\Big(\prod_j R_{i_j}\Big)
                  \prod_j\operatorname{tr}K_{i_j}.
\]
Products inside traces retain their indicated order.
Theorem~\ref{thm:rowcycles318} expresses the right trace
as an integral with the actual row projector, and bounds
the change under a single matching complex lift on its
ordinary periodization spaces. Equality of these traces
for the designated row and the signed definite data is
still open. Even after a unitary comparison, the arithmetic
normalization requires identifying the pairing on its fixed
rational lattice with its specified values there. The trace
calculation is invariant under rescaling each whole word.

\subsection{One endpoint for the Bessel row and an arithmetic period}\label{nav:VIII-60}
Fix a field \(E_x\) on the interval of Section~\ref{sec:missing}
and one of the prescribed words. Use a good orbit of
Theorem~\ref{thm:commonorbit305} and its terminal character.
Take the last two highest lines as in Theorem~\ref{thm:pen185}
and Proposition~\ref{prop:minor181}. In particular the last
Schwartz polynomial is the coefficient-one corner polynomial
\(P_*\) of Theorem~\ref{thm:last244}. Use its marked Fock
coordinates and carry the tests with the character as in
Theorem~\ref{thm:transport233}. All finite tests and external
characters below are those of this word. Write \(T\) for its
last forward lift and \(X\) for a vector obtained by applying
the preceding lifts to a Bessel row \(Sb\).

Let \(a_*>0\) denote the lower factor in the one-step Rallis
bound for this last lift on the space of incoming rows.
More explicitly, take the product of the fixed last-step
factors in Theorem~\ref{thm:norm177}, with its finite depth
factor \(\prod_{v\in S_{\rm row}}(Nv)^{-a_v n(n+s)}\) and its
real test \(P_*\), and multiply by
\(\exp(-Cfr\,2^{-(z_0-c_r)})\) in the notation of that theorem.
Here \(n,s\) are the last source and target ranks. With compact
source Haar probability, the unnormalized real matrix-coefficient
integral in the marked, unit-variance Fock variables is
\(h_\lambda/D_-\), where \(h_\lambda=\|P_*\|_{\rm F}^2\) and
\(D_-=D_n\). Indeed \(P_*\) corresponds to the two unit highest
tensors under an isometry times \(\sqrt{h_\lambda}\).
Integrating the squared source highest coefficient gives
\(1/D_-\). Any Haar volume and oscillator dilation used in
the actual norm formula are retained there.

For the two definite groups, indicated by \(-,+\), let \(K_\pm\)
be the original finite levels and \(\chi_\pm\) their characters.
Take normal open \(J_\pm\subseteq\ker(\chi_\pm)\) inside the
levels, fixing the tests and with trivial rational stabilizers.
One obtains such subgroups by taking the normal core of fixing
levels and imposing one odd split principal congruence condition.
At that place the first congruence subgroup is torsion free.
Indeed a nonidentity matrix in its \(h\)-th layer has its
prime-th power first in layer \(h+1\) at an unramified split
rational prime, and a power prime to that prime stays in layer
\(h\). The definite rational stabilizers are finite. Write
\(i_\pm=[K_\pm:J_\pm]\), \(h_\pm\) for the numbers of fine
classes, \(w_\pm\) for the measure of one chart in the original
Haar measures, and \(v_\pm=h_\pm w_\pm\) for the quotient volumes.
Write \(D_+=D_s\), \(c=\pi/A\), and \(C\) for the scalar of
Theorem~\ref{thm:vac242} at this edge.

On one fine class take the function
\(k\mapsto\langle\rho_\pm(k)t_\pm,v\rangle\), zero on the
other classes, for \(v\) from Theorem~\ref{thm:frames311}.
Apply the unnormalized sum
\(\sum_{a\in K_\pm/J_\pm}\overline{\chi_\pm(a)}R(a)\)
of right translations. Denote the resulting tests by \(U,V\)
on the source and target respectively. They have the original
levels and characters. Define their normalized period by
\[
 z(U,V)=\frac{D_-D_+c^{d/2}}{w_-w_+ C}\langle TU,V\rangle,
 \qquad L_*=\frac{c^{d/2}\sqrt{a_*}}{4|C|\sqrt{v_-v_+}} .
\]

\begin{theorem}[A common terminal test at the prescribed character]\leavevmode
\label{thm:lastrow311}
For every nonzero incoming \(X\) there is a test \(V\) with
\[
 |\langle TX,V\rangle|
       \ge \frac{i_+\sqrt{w_+a_*}}{2D_+\sqrt{h_+}}\|X\|.
\]
For that same \(V\) there is a source test \(U\) for which
\(|z(U,V)|\ge i_- i_+ L_*\). The tests may be chosen from
either of the two families at each end. With the integral
family their coefficients in the specified fine charts are
unnormalized sums of integral exterior tensors and finite
character values. The normalized period is algebraic.
\end{theorem}
\begin{proof}
The first \(m-1\) local contractions can reach the real source
line used here. The compact projection described after
Proposition~\ref{prop:minor181} sets that line on the preceding
edge, and Theorem~\ref{thm:pen185} supplies its explicit
highest-type test on the indicated lowest source line. To reach
this latter line from the earlier lifts, translate a nonzero
contraction in its irreducible real representation and project
to the multiplicity-one compact line, as in the proof of
Corollary~\ref{cor:twomain185}. Equivariance moves
the translation and projection onto the real test. Continuity
and density of polynomial tests permit a polynomial with
nonzero projected contraction. The real choices are common to
the constituents with the prescribed real parameters. The finite
norm factors and the orbit estimate use the original finite
levels throughout. In particular there is a nonzero \(Sb\) on
the chosen orbit by Theorem~\ref{thm:commonorbit305}, and the
successive positive norm bounds give a nonzero \(X\).

For a single incoming constituent the last-step norm identity
in the proof of Theorem~\ref{thm:norm177} gives
\(\|TX\|^2\ge a_*\|X\|^2\). Its varying good-place factor is
\(L^S(z_0,B_0\otimes\mu_m^{-1})\), since all earlier parameter
strings in this last factorization are fixed. The estimate
for its absolute logarithm there gives precisely the stated
one-step error. Mixed terms between orthogonal constituents
vanish by that same mixed norm identity, including for distinct
multiplicity copies. Hence the inequality holds on the span
of the incoming vectors.

At either end choose a fine class uniformly and use the
probabilities of Theorem~\ref{thm:frames311}. Any
highest-line, original-character form \(F\) has on class \(u\)
the unique expression \(\langle\rho(k)t,y_u\rangle\).
Compact orthogonality gives
\(\|F\|^2=(w/D)\sum_u\|y_u\|^2\). Pairing with a one-class
test has absolute value \((w/D)|\langle y_u,v\rangle|\).
The finite character sum multiplies that pairing by \(i\),
since it is \(i\) times the orthogonal projector. Thus for
the projected random test \(Q\),
\[
 {\mathbb E}|\langle F,Q\rangle|^2
                 \ge \frac{i^2w}{4hD^2}\|F\|^2 .
\]
Apply this first to \(TX\) to choose \(V\). The functional
\(F\mapsto\langle TF,V\rangle\) on the original finite-level,
highest-line source has norm at least
\(|\langle TX,V\rangle|/\|X\|\). Riesz representation and the
source version of the same bound choose \(U\). A conjugate-linear
lift is treated on the conjugate source Hilbert space, which
leaves all absolute values unchanged. Multiplication by the
normalization of \(z\) gives the claimed lower bound.

Each pair of terms in the two finite sums is supported on a
single pair of fine classes. Use the actual translated
representatives for that term. The stabilizers have order
one. Theorem~\ref{thm:last244} evaluates it, after the
normalization defining \(z\), as
\(M(v_-\otimes v_+)(\partial)G_{uv}(0)\), times its finite
character factor. Here \(M\) is the integral minor map, and
\(G_{uv}\) is the CM section for the same translated kernel.
This proves algebraicity term by term. With a different fixed
representative on a translated class the tensor is transported
by the corresponding rational isometry. Using the translated
representative itself keeps its displayed integral coordinates.
\end{proof}

\subsection{Cyclic integrals on the prescribed Bessel row}\label{nav:VII-66}
For a fixed set of framed data let \(S_u:M_u\to\mathcal H\)
be the ordinary Bessel periodizations in one ambient Hilbert
space. Use the column identifications and norms occurring
in the pairing. An isometric realization
\(\iota:E\to\mathcal H\) of the row gives the period
\(\langle\iota y,S_u x\rangle\), with the inner product
linear in its second argument. If the original realization
is conjugate-linear in the row, use its conjugate space
for \(E\). Write \(\Pi=\iota\iota^\dagger\). This is the
orthogonal projector onto the realized row. The physical
column on that row is \(\iota^\dagger S_u\).

Put \(t_u=\|S_u\|_{\rm HS}^2\), where the Hilbert--Schmidt
norm squared is the sum of the squared norms on an
orthonormal basis. For \(t_u>0\) set
\(\epsilon_u=\|(I-\Pi)S_u\|_{\rm HS}^2/t_u\).
Given a cyclic list \(i_1,\ldots,i_k\) of words, let
\(\Gamma_\Pi\) be the trace of the ordered product of the
output Gram operators of \(\Pi S_{i_j}\). Denote by
\(m_\Pi\) this trace divided by the product of the individual
Gram traces when all are positive. The same definitions
with \(\Pi=I\) give \(\Gamma_I,m_I\). The row embedding
identifies the nonzero part of the product for \(\Pi\)
with the corresponding product for \(\iota^\dagger S_u\).

\begin{theorem}[Projected cyclic integrals and their error]\leavevmode
\label{thm:rowcycles318}
Let \(h_u=\dim M_u>0\). On the finite orbit spaces of
Theorem~\ref{thm:commonorbit305} take an independent signed
unit column \(b_j\) from the distribution for each
occurrence \(i_j\). Write \(v_j=\Pi S_{i_j}b_j\) and
\(v_{k+1}=v_1\). Then
\[
 \Gamma_\Pi=\left(\prod_j h_{i_j}\right)
       \mathbb E\prod_{j=1}^k\langle v_j,v_{j+1}\rangle,
 \qquad
 \|\Pi S_u\|_{\rm HS}^2=t_u(1-\epsilon_u).
\]
On a common good orbit \(t_u=h_u M_{J,\mathcal O}\) for
the corresponding word and terminal space. For
\(i_{k+1}=i_1\), and whenever \(m_\Pi\) is defined, one has
\[
 |m_\Pi-m_I|
 \le \sum_{j=1}^k\sqrt{\epsilon_{i_j}\epsilon_{i_{j+1}}}
               +1-\prod_{j=1}^k(1-\epsilon_{i_j}).
\]
On the common real line and with the single matching
complex operator \(\widehat\Theta=\alpha^{-1/2}\Theta\)
of Theorem~\ref{thm:minors239}, let \(m_\Theta\) be the
normalized trace for the family \(\widehat\Theta S_u\).
For every ordinary list of positive traces on these data,
\[
 |m_\Theta-m_I|\le 2(e^{k\Delta_x}-1).
\]
For \(k\le 2(\dim D_{\rm all})^2\) this last error is at
most \(\exp(-(\log 2)x^{377}/3)\) throughout sufficiently
distant common intervals.
\end{theorem}
\begin{proof}
For independent vectors with
\(\mathbb E(b_jb_j^*)=I_{M_{i_j}}/h_{i_j}\), the expectation
of the product of the rank-one operators \(v_jv_j^*\)
is the product of their expectations in the same order.
The trace for a fixed tuple of vectors is the cyclic
product of inner products in the assertion. The required
covariance on each orbit is the one in the proof of
Theorem~\ref{thm:commonorbit305}. Taking one factor with
\(\Pi=I\) gives the value of \(t_u\). Pythagoras on a
source orthonormal basis gives the projected trace.

By cyclicity the same unnormalized trace can be computed
on \(M_{i_1}\) as the trace of the product of the rectangular
maps \(S_{i_j}^\dagger\Pi S_{i_{j+1}}\), in order. This
identity also follows by summing their diagonal coordinates.
Put \(R_u=(I-\Pi)S_u\). The change from the ordinary
factor to the projected one subtracts \(R_i^\dagger R_l\)
for adjacent words \(i,l\). Its trace norm is at most
\(\|R_i\|_{\rm HS}\|R_l\|_{\rm HS}\). Indeed the trace norm
is the supremum of the absolute trace after multiplication
by a contraction of the opposite shape, by singular value
decomposition. Hilbert--Schmidt Cauchy--Schwarz gives the
bound under that supremum. Each unchanged factor in
either product has operator norm at most \(\sqrt{t_i t_l}\).
Replace one factor at a time and use the trace norm bound
on the changed factor. After division by \(\prod_jt_{i_j}\)
the sum of the differences is bounded in modulus by the
sum of square roots in the theorem. Every normalized output
Gram has trace norm one and operator norm at most one.
Thus the normalized trace of their product has modulus
at most one. Since
\(\Gamma_\Pi/\prod_jt_{i_j}=m_\Pi\prod_j(1-\epsilon_{i_j})\),
the triangle inequality proves the first estimate.

For the last comparison let \(H\) be the Gram operator of
the restriction of \(\widehat\Theta\) to the span of the
ordinary images. Theorems~\ref{thm:range239}
and~\ref{thm:norm177} give \(e^{-\Delta_x}I\le H\le
e^{\Delta_x}I\). Now a changed factor is
\(S_i^\dagger(H-I)S_l\), whose trace norm is at most
\((e^{\Delta_x}-1)\sqrt{t_i t_l}\) by the same inequality.
Each factor already changed contributes at most another
\(e^{\Delta_x}\) to its operator norm bound. Summing the
geometric progression shows that the difference between
the two unnormalized traces divided by \(\prod_jt_{i_j}\)
has modulus at most \(e^{k\Delta_x}-1\). The product of
the changed individual traces divided by \(\prod_jt_{i_j}\)
lies in \([e^{-k\Delta_x},e^{k\Delta_x}]\). The normalized
trace has modulus at most one, giving the second copy of
\(e^{k\Delta_x}-1\). Finally Lemma~\ref{lem:dimension239}
gives \(\log k\le\log 2+2x^{40}\), while the estimate
in Theorem~\ref{thm:minors239} gives
\(\log\Delta_x\le-(\log 2)x^{377}/2+O(x^{40})\).
Together they give the asserted uniform error.
\end{proof}

The first identity applies with the actual designated
projector whenever the columns are realized as these
periodizations. The ordinary orbit estimate controls \(t_u\).
The contribution removed by the designated projector is
the separate sum of squares \(t_u\epsilon_u\). The last
estimate concerns the common matching theta lift on the
ordinary space. Its proof uses the same \(H\) for every
factor in a cycle.

\subsection{Vanishing forced by the level of the kernels}\label{nav:I-44}
Let \(U\subseteq C\) be the image of the integral norm-one ideles supported at the places above \(P\). The lower-bound argument in Proposition~\ref{prop:torus} gives
\[
             |U|\ge (P-1)^{[E^+:\QQ]}/w_E .
\]
Here all other finite local units already belong to the congruence kernel defining \(C(P)\). The following observation applies whenever both adjacent kernels are taken at a level with these determinant units.

\begin{proposition}[Vanishing in a fixed central character]\label{prop:level25}
At every intermediate vertex \(i\) suppose there is a compact subgroup \(J_i\subseteq G_i(\mathbb A_f)\) whose determinant image in \(C\) contains \(U\). Suppose both kernels meeting that vertex are right invariant in the variable \(g_i\) under \(J_i\). In any family of \(h\) twists with a fixed output central character, at most \(|C/U|^{m-2}\) operators can be nonzero. This holds after every fixed output map as above. For any \(b=\dim V>0\) one has
\begin{equation}
  \frac{1}{hb}\sum_\eta\|A_\eta^*A_\eta-I_V\|_{\rm HS}^2
        \ \ge\ 1-|U|^{-(m-2)} .                       \label{eq:level25}
\end{equation}
For the common fields \(E_x\), \(m=\lfloor\sqrt x\rfloor\), and \(P\) as in Corollary~\ref{cor:orth}, the proportion allowed to be nonzero is at most \(\exp(-c_P x^{5/2+8\eps})\) for all sufficiently large \(x\).
\end{proposition}
\begin{proof}
In \eqref{eq:iter}, replace just \(g_i\) by \(g_i k\) with \(k\in J_i\). Right translation preserves the quotient measure, and invariance of the two incident kernels leaves as the only factor \(\chi_i(\det k)\). If this factor differs from one, the integral vanishes identically as a function of the endpoint. Thus a nonzero operator requires every \(\chi_i\) to be in \(X=\widehat{C/U}\).

The homomorphism from \(\widehat C^{m-1}\) to \(\widehat C\) given by \( (\chi_i)\mapsto\prod_i\chi_i^{s_i n_i}\) has a right inverse \(\omega\mapsto(\omega,\omega,1,\ldots,1)\) since \(-2+3=1\). Every fiber has size \(h\). The same right inverse on \(X\) shows that its restriction to \(X^{m-1}\) has either zero or \(|X|^{m-2}\) points in that fiber. A zero operator contributes \(\operatorname{tr}(I_V)=b\) to the sum in \eqref{eq:level25}. This proves both assertions about a fixed central character. Finally \(\log |U|\ge [E_x^+:\QQ]\log(P-1)-\log w_{E_x}\gg_P x^{2+8\eps}\), using Corollary~\ref{cor:orth}. Multiplying by \(m-2\gg\sqrt x\) gives the displayed exponential bound on the surviving proportion.
\end{proof}

The determinant-level hypothesis is realized, for example, for the standard definite spaces at an unramified \(P\) by using the integral stabilizers of the standard lattices with kernels invariant there. The diagonal action of an integral norm-one unit on a coordinate line lies in that stabilizer and has the given determinant. Right invariance can also be imposed on given kernels by averaging their Schwartz vectors over the corresponding compact groups at both ends of every edge. Those actions commute through the chosen splitting, and each integral on a smooth finite-place vector is a finite sum because the vector has an open stabilizer. The resulting theta kernels have exactly the invariance used in the proposition. The argument allows a vanishing average.

For this choice of levels, the error in \eqref{eq:level25} tends to one, whereas the rank argument in Volume~I uses \(e\le1/4\). Applying the rank argument to the full family therefore requires different local test data or a different selection of twists. Formulas \eqref{eq:gram25} and \eqref{eq:cov25} continue to hold with such data. The within-component correlations, the lower divisibility reflecting the prime count, and the covolumes for the particular integral periods required in Section \ref{nav:I-28} remain to be estimated.

\subsection{A common ramified test function at a split place}\label{nav:I-45}
\label{sec:flag29}

We keep the actual chain \eqref{eq:iter}, with \(n_i=i+1\), and use the same Schwartz function for every twist at a given edge. The special test function below is the indicator of a set of integral matrices. Its expansion under diagonal changes of basis selects exactly one term for each twist after integration over the intermediate groups.

\begin{lemma}[Split places on the power interval]\label{lem:split29}
Fix distinct rational primes \(P,Q_0\ge5\) prime to \(2F\). For all sufficiently distant common intervals, both are unramified in \(E=E_x\), and every place of \(k=E^+\) above \(P\) splits in \(E/k\). Put \(q_v=|\OO_k/v|\) for \(v\mid P\). Any finite-order character of the product of the integral local norm-one units at these places can be extended to a finite-order automorphic character of \(T\), allowing ramification at \(Q_0\).
\end{lemma}
\begin{proof}
Eventually the discriminant primes \(q,d\) of Volume~I are distinct from \(P,Q_0\), and \(N=2^{l+2}Fqd\) is prime to both. Let \(g\) be Frobenius at \(P\) in \(\Gal(E/\QQ)\). Its order on \(C_l\) is a power of two. For large \(l\) this order is even, since \(P\) is then neither \(1\) nor \(-1\) modulo \(2^{l+2}\). If \(g^j\) were complex conjugation, it would act trivially on \(C_l\), forcing \(j\) even. But \(g^j\) would then act trivially on \(\QQ(\sqrt{-q})\), contrary to complex conjugation. The decomposition group at \(P\) is generated by \(g\), as \(P\) is unramified. Its intersection with \(\Gal(E/k)=\{1,c\}\) is therefore trivial, giving the splitting assertion.

A prescribed continuous finite-order character on the product of unit groups factors modulo \(P^b\) for some \(b\). Map that finite quotient into \(C(P^b Q_0)\) using ideles supported at \(P\). The kernel is trivial. Indeed the rational norm-one element in a relation must be integral and a unit at every finite place, hence a root of unity by its archimedean absolute values. It is congruent to one at \(Q_0\). Since \(E\) is unramified at \(Q_0\), it cannot contain a primitive \(Q_0\)-th root of unity, so the order of that root is prime to \(Q_0\). If \(z\equiv1\bmod w\) and \(z^t=1\) with \(Q_0\nmid t\), then \(1+z+\cdots+z^{t-1}\) is a unit at \(w\mid Q_0\). This forces \(z=1\). The asserted kernel is consequently the original congruence subgroup before taking the finite quotient. A character of a subgroup of a finite abelian group extends to the full group. This follows by adjoining one generator at a time and taking an appropriate root of its prescribed power in the circle. Proposition~\ref{prop:torus} gives finiteness of \(C(P^bQ_0)\), so the extension has finite order.
\end{proof}

At a split place \(v\mid P\), write \(L=k_v\), \(\mathfrak o=\OO_L\), and \(\mathbf k_v\) for its residue field. The two idempotents of \(E\otimes_k L=L\oplus L\) give two dual summands of each unitary space. Choose the first summand at even vertices and the second at odd vertices. Thus the norm-one scalar \((z,z^{-1})\) acts on the chosen summand of vertex \(i\) by \(z^{s_i}\), \(s_i=(-1)^i\). On the edge from dimension \(a=i\) to \(b=i+1\), one Lagrangian of the symplectic tensor space is
\(\operatorname{Hom}(L^a,L^b)\). Indeed the dual of the chosen summand at vertex \(i-1\) uses the same idempotent as the chosen summand at \(i\). Their tensor product and its dual are two complementary isotropic spaces. In the Schr\"odinger model on \(M_{b,a}(L)\) the action on compact elements is therefore
\begin{equation}
       ((g,h)\phi)(X)=\phi(h^{-1}Xg),
       \quad g\in\GL_a(\mathfrak o),\quad h\in\GL_b(\mathfrak o).
       \label{eq:action29}
\end{equation}
Here we use the splitting given by pullback on that Lagrangian, with the square root of the absolute Jacobian for general elements. The Jacobian is one on the displayed compact subgroups.~\cite{XXIII-XVIII-IV-kudla29}

We explain how to arrange \eqref{eq:action29} for all edges using global splittings. On the split local general linear groups the ratio of an initial unitary splitting and the Schr\"odinger splitting is a character into the central circle. It factors through the two determinants. Indeed a character of the direct product is a product of characters, and the special linear group over the infinite field \(L\) is generated by elementary transvections. Each is a commutator in the general linear group, by conjugation by a diagonal matrix with one diagonal entry different from one. On a compact unit group each resulting smooth character has finite image. By Lemma~\ref{lem:split29} the inverses of its unit restrictions extend to global characters of the norm-one torus. Multiplying each edge splitting by the corresponding automorphic determinant characters changes it by a central scalar, preserves the rational splitting, and imposes \eqref{eq:action29} on the compact subgroups at all \(v\mid P\). Make this choice once for the entire twist family. The construction permits additional levels at \(Q_0\).

Write \(I_n=\{g\in\GL_n(\mathfrak o):\bar g\text{ is upper triangular}\}\). A bar here means reduction to \(\mathbf k_v\). Let \(\Omega_{a,b}\) consist of the integral matrices for which \(\bar X_{ij}=0\) if \(i>j\), and the first \(a\) diagonal entries are nonzero. The common vector \(\phi^0_{a,b}\) requires all these diagonal entries to be one: it is the indicator of that smaller set. For an \(a\)-tuple \(\lambda=(\lambda_1,\ldots,\lambda_a)\) of characters of \(\mathbf k_v^\times\), put
\[
 \phi_\lambda(X)=\mathbf1_{\Omega_{a,b}}(X)
                      \prod_{j=1}^{a}\lambda_j(\bar X_{jj}),
\]
with value zero off the indicated set. Then
\begin{equation}
  \phi^0_{a,b}=(q_v-1)^{-a}\sum_\lambda \phi_\lambda.
  \label{eq:common29}
\end{equation}
This is ordinary character orthogonality on each diagonal entry. Under \((g,h)\in I_a\times I_b\), the function \(\phi_\lambda\) transforms by
\[
     \prod_{j\le a}\lambda_j(\bar g_{jj})\lambda_j(\bar h_{jj})^{-1}.
\]
In fact the set \(\Omega_{a,b}\) is stable in \eqref{eq:action29}, as is seen by multiplying upper triangular residue matrices. The \(j\)-th diagonal entry of \(\bar h^{-1}\bar X\bar g\) is \(\bar h_{jj}^{-1}\bar X_{jj}\bar g_{jj}\).

\subsection{Orthogonality after Bessel restriction}\label{nav:I-46}
\label{sec:bessel29}

Use \(\phi^0_{i,i+1}\) on every edge at every \(v\mid P\), with arbitrary fixed Schwartz functions at other places. Take an input function on \([G_0]\) in a single central character, invariant under the integral units at \(P\); the constant function is an example. Fix a nondegenerate codimension-one subspace \(V'_m\subset V_m\), put \(G'_m=U(V'_m)\), and use local bases at the endpoint in which this subgroup acts in the first \(m\) coordinates. Its compact group at these places contains \(I'=\prod_{v\mid P} I_m\) embedded block diagonally. For \(W\) a finite-dimensional right \(I'\)-invariant subspace of \(L^2([G'_m])\), let \(B_{\boldsymbol\chi}\) be the operator obtained by restricting \eqref{eq:iter} to \([G'_m]\) and projecting orthogonally to \(W\). Pairing it with \(h\in W\) gives the codimension-one Bessel integral against \(\bar h\). All the restricted functions are continuous on a compact quotient, so this construction requires no boundedness assertion for restriction on the ambient \(L^2\)-space.

\begin{theorem}[A fixed central character and a fixed Bessel target]\label{thm:bessel29}
For \(C=C(P)\), let \(U\subset C\) be the image of the integral local norm-one units at \(P\). There are \(|U|^{m-2}\) tuples of determinant characters \(\chi_i\in\widehat C\), all with the same prescribed product \(\prod_i\chi_i^{s_i n_i}=1\), for which
\[
             B_{\boldsymbol\chi}^{*} B_{\boldsymbol\chi'}=0
                    \qquad(\boldsymbol\chi\ne\boldsymbol\chi').
\]
The test functions, splittings, initial central character and \(W\) are the same for the entire family. On the common intervals the number of tuples is at least \(\exp(c_P x^{5/2+8\eps})\) for sufficiently large \(x\).
\end{theorem}
\begin{proof}
Fix a place over \(P\) for the moment. Let \(\delta_i\) be the local residue character furnished by \(\chi_i\) in the chosen general linear coordinate: the internal multiplier on \(I_{n_i}\) is \(\delta_i(\prod_j\bar g_{jj})\). It is the restriction of \(\chi_i^{s_i}\) in the first norm-one coordinate. Expand each common vector by \eqref{eq:common29} and write \(\lambda_{i,j}\) for the characters on the \(i\)-th edge. Right translation of \(g_0\) by its integral units and invariance of the input force \(\lambda_{1,1}=1\) in every surviving term. Right translation of the intermediate integration variable \(g_i\) by \(I_{n_i}\) next forces
\begin{equation}
 \lambda_{i+1,j}=\delta_i^{-1}\lambda_{i,j}\ (j\le i),\qquad
 \lambda_{i+1,i+1}=\delta_i^{-1}\qquad(1\le i<m).           \label{eq:rec29}
\end{equation}
Indeed, at coordinate \(j\) the incoming edge contributes \(\lambda_{i,j}^{-1}\) when \(j\le i\), the outgoing edge contributes \(\lambda_{i+1,j}\), and the twist contributes \(\delta_i\). These diagonal entries can be chosen independently. If the product differs from one, translation invariance of Haar measure makes that summand zero. This leaves just the tuple of edge labels in \eqref{eq:rec29} at this place. The same argument applies independently at all places over \(P\).

The right character on the output \(I_{m+1}\) has diagonal coordinates \(\gamma_j=\lambda_{m,j}^{-1}\) for \(1\le j\le m\) and \(\gamma_{m+1}=1\). The recursion implies
\[
 \gamma_1=\gamma_2,\qquad
 \gamma_j=\prod_{i=\max(1,j-1)}^{m-1}\delta_i\ (1\le j\le m),\qquad
 \delta_i=\gamma_{i+1}\gamma_{i+2}^{-1}.
\]
Thus the first \(m\) characters, which are still present on the Bessel subgroup, determine every \(\delta_i\). This proves that different tuples of unit restrictions give different characters on \(I'\). Right translation is unitary on \(L^2([G'_m])\), so its two distinct character spaces are orthogonal. Since \(W\) is right \(I'\)-invariant, its orthogonal projection commutes with those translations. The projected pairings vanish for any pair of source vectors, giving the asserted operator identity.

It remains to choose global tuples. Restriction \(\widehat C\to\widehat U\) is onto, by extension of characters of a finite abelian group. For either character group the product map \((\chi_i)\mapsto\prod_i\chi_i^{s_i n_i}\) has section \(\omega\mapsto(\omega,\omega,1,\ldots,1)\), since \(-2+3=1\). Each tuple in the kernel on \(\widehat U\) lifts to one in the kernel on \(\widehat C\). In fact any lifts have product trivial on \(U\), and dividing by the image of this product under the section corrects it without changing restrictions. Choose one lift for each of the \(|U|^{m-2}\) tuples. Theorem~\ref{thm:center} gives the same global output central character for all of them. The lower bound follows from \(\log|U|\ge [k:\QQ]\log(P-1)-\log w_E\) in Proposition~\ref{prop:torus} and \(m=\lfloor\sqrt{x}\rfloor\).
\end{proof}

\begin{proposition}[Integral factors for the common local tests]\label{prop:int29}
Each common vector has values in \(\ZZ\) and takes the value one. With the integral matrix lattice assigned Haar volume one, their product has squared local \(L^2\)-norm
\[
    \prod_{v\mid P}\prod_{i=1}^m\|\phi^0_{i,i+1}\|_2^2
       =P^{-[k:\QQ]m(m+1)(m+5)/6}.
\]
The common denominator in the full expansion \eqref{eq:common29} is at most
\(D_{\rm in}=\prod_{v\mid P}(q_v-1)^{m(m+1)/2}\).
On any integral \(I'\)-stable test lattice fixed by the subgroup with all residue diagonal entries one, each endpoint character projector is integral after multiplication by
\(D_{\rm out}=\prod_{v\mid P}(q_v-1)^m\), over a coefficient field containing the character values.
\end{proposition}
\begin{proof}
In the matrix of size \((i+1)\times i\), the \(i(i-1)/2\) entries strictly above the diagonal are free in the residue field, and every other residue entry is fixed for \(\phi^0\). Its squared norm is \(q_v^{-i(i+3)/2}\). Sum the exponents over \(i\) and use \(\prod_{v\mid P}q_v=P^{[k:\QQ]}\) since \(P\) is unramified. Multiplying the denominators in \eqref{eq:common29} gives \(D_{\rm in}\). At the endpoint, \(I'\) modulo the subgroup in the statement is \(\prod_{v\mid P}(\mathbf k_v^\times)^m\). The projector to a character is its finite character average. Its numerator preserves the lattice since each coefficient is a root of unity. Its denominator is \(D_{\rm out}\). In particular \(\log D_{\rm in}\le \tfrac12 m(m+1)[k:\QQ]\log P\) and \(\log D_{\rm out}\le m[k:\QQ]\log P\).
\end{proof}

Theorem~\ref{thm:bessel29} evaluates the off-diagonal blocks for this family. A projected operator may still vanish. The norm in Proposition~\ref{prop:int29} is the norm of the test functions in the local Weil representation. A lower bound on the diagonal blocks of the global Bessel operators requires their actual period norms. Realizing the particular translates and weights of Volume~I, proving their lower divisibility at the other places, and estimating the integral indices of their common target remain necessary for inequality \eqref{eq:gap}.

\subsection{A Bessel test space of exactly the required rank}\label{nav:I-47}
\label{sec:target32}

We now specify the common test space in Theorem~\ref{thm:bessel29}. Take the standard Hermitian spaces, \(k=E^+\), and let \(H=U(E^m)\) be the subgroup preserving the first \(m\) coordinates at the endpoint. Use the fixed unramified primes \(P,Q_0\) of Lemma~\ref{lem:split29}. At \(v\mid P\) the quotient of the Iwahori by the subgroup with every residue diagonal entry one gives
\[
 A=\prod_{v\mid P}(\mathbf k_v^\times)^m .
\]
Write \(K^{\max}\) for the finite adelic stabilizer of the standard lattice. Let \(K_0\subset K^{\max}\) have Iwahori level at \(P\), principal level \(Q_0\) at every place above \(Q_0\), and full level elsewhere. Let \(K_1\triangleleft K_0\) have the same level with residue diagonals one at \(P\), so \(K_0/K_1=A\). Define
\[
 X=H(k)\backslash H(\mathbb A_{k,f})/K_1,\qquad
 h=|H(k)\backslash H(\mathbb A_{k,f})/K_0|.
\]
For the family of twists of Theorem~\ref{thm:bessel29}, let \(\Xi\) be their residue characters on \(A\), and put \(Z=\bigcap_{\gamma\in\Xi}\ker\gamma\) and \(n_0=|\Xi|\).

\begin{lemma}[The level classes and the residue characters]\label{lem:free32}
The action of \(A\) on \(X\) is free, with \(h\) orbits. The set \(\Xi\) is a subgroup of \(\widehat A\), and \(\Xi=\widehat{A/Z}\) has order \(n_0=|U|^{m-2}\), with \(U\) as in Theorem~\ref{thm:bessel29}.
\end{lemma}
\begin{proof}
For \(y\in H(\mathbb A_{k,f})\), the rational group \(H(k)\cap yK_0 y^{-1}\) preserves the Hermitian lattice with finite completions \(y\widehat{\OO_E^m}\). Its positive definite integral trace form has finite automorphism group, because there are finitely many possibilities for the images of each basis vector of specified length. Thus every element has finite order. At \(Q_0\) it is conjugate to a matrix congruent to the identity modulo \(Q_0\). The principal congruence group is torsion-free. Indeed all completions of \(E\) at \(Q_0\) are unramified and \(Q_0\ge5\). If \(1+Q_0^aB\), \(a\ge1\), has an entry of \(B\) a unit, raising it to a prime \(p\ne Q_0\) leaves a nonzero term \(pQ_0^a B\) modulo \(Q_0^{a+1}\). Raising it to \(Q_0\) gives \(Q_0^{a+1}B\) as the first nonzero term of its difference from one modulo \(Q_0^{a+2}\). No element of prime order is possible. Faithfulness after completion implies the rational stabilizer is trivial. A stabilizer for the right action of \(K_0/K_1\) on \(X\) would arise from such a rational element, so the action is free. Its orbit space is the double quotient at level \(K_0\). These double quotients are finite by definiteness, also by Lemma~\ref{lem:mass32} below.

The proof of Theorem~\ref{thm:bessel29} parametrizes the twists by the kernel of the product homomorphism on \(\widehat U^{m-1}\). Their endpoint characters are given by the products of residue coordinates in \eqref{eq:rec29}, independent of which global lifts are chosen. That formula is multiplicative and recovers all residue coordinates, hence its image is a subgroup isomorphic to the kernel. The order is \(|U|^{m-2}\). Duality for finite abelian groups gives \(\widehat{A/Z}=\Xi\): for instance the evaluation map \(A\to\widehat\Xi\) is onto since every character of the subgroup \(\Xi\subseteq\widehat A\) extends to \(\widehat A\), and its kernel is \(Z\).
\end{proof}

\begin{lemma}[The number of level classes]\label{lem:mass32}
Let \(d=[E:\QQ]\), \(n=md\), and suppose \(E\subseteq\QQ(\zeta_N)\). For an absolute constant \(C\),
\[
                  \log h\le C n^4\log(2nNPQ_0).
\]
\end{lemma}
\begin{proof}
At full level the double classes are the isometry classes of lattices in the genus of the standard Hermitian lattice. Indeed the finite adelic translates give those lattices, intersection with the rational space recovers a lattice from its completions, and two translates differ by a rational isometry precisely when they give the same class. For such a lattice use the integral symmetric form \(B(v,w)=\operatorname{Tr}_{E/\QQ}\langle v,w\rangle\). Its rank is \(n\), it is positive definite, and its Gram determinant is \(\Delta=|D_E|^m\). For the standard lattice this follows from the trace discriminant and the determinant of conjugation, a unit on \(\OO_E\). Local isometry gives the same valuation at every rational prime for all lattices in the genus; positivity gives the asserted equality. In particular \(\Delta\le N^n\) by the abelian conductor--discriminant formula.

There is an algebraic integer \(\theta\) generating \(E\) with absolute value at most \(d^3N\) at every embedding. Traces from \(\QQ(\zeta_N)\) of the powers of \(\zeta_N\) span \(E\) over \(\QQ\), since trace is onto and these powers span the larger field. Take a basis of \(d\) such traces. Each conjugate has size at most \(N\). For a pair of distinct embeddings of \(E\), the coefficient vectors giving equal values on a linear combination form a proper hyperplane. On the grid \(\{0,\ldots,d^2\}^d\) each hyperplane has at most \((1+d^2)^{d-1}\) points, by solving for a coordinate with nonzero coefficient. There are fewer than \(1+d^2\) pairs. A grid point outside their union gives the desired generator and bound (in degree one one may take \(1\)).

Every trace lattice has an integral basis with norms at most \(B_0=(4n)^{2n}\sqrt\Delta\). To see this, its shortest nonzero vector has norm at least one by integrality. Minkowski's second theorem for the Euclidean unit ball bounds its last successive minimum by \(n^{n/2}\sqrt\Delta\), using the cube of side \(2/\sqrt n\) inside the ball to bound its volume. Choose independent vectors of norm at most twice this minimum. Intersect the lattice with their successive real spans. Each consecutive quotient is free of rank one, and the orthogonal length of a generator is at most the norm of the new independent vector, since that vector gives a nonzero integer multiple in the quotient. Choose lifts of the generators consecutively. Subtract integer multiples of the previous lifts in reverse order to make each orthogonalized coefficient lie in \([-1/2,1/2]\). The squared length of each resulting basis vector is at most \((1+n/4)\) times the square of twice the last minimum. This implies the stated bound.~\cite[Thm.~6.3.3]{XXIII-XVIII-IV-minima32}

In this basis record the Gram matrix and the integral matrix of multiplication by \(\theta\). The entries of the former have magnitude at most \(B_0^2\). Those of the latter have magnitude at most \(d^3NB_0^n\), by Cramer's rule and Hadamard's inequality, since multiplication by \(\theta\) has operator norm at most \(d^3N\) and the lattice covolume is \(\sqrt\Delta\ge1\). Equal matrix pairs yield a trace-lattice isometry commuting with \(\theta\), hence linear over \(E=\QQ(\theta)\). It preserves the Hermitian form because its trace pairing after multiplication of one argument by every \(a\in E\) determines that form by trace nondegeneracy. Thus the number of genus classes is bounded by the number of matrix pairs. Taking logarithms, and using \(\log B_0\ll n\log(2nN)\), gives \(C n^4\log(2nN)\). The index of \(K_0\) in the full stabilizer is at most \((P Q_0)^{n^2}\), by reduction on the underlying rank-\(n\) integral lattice at those two rational primes; the kernels of both reductions lie in the specified levels. Each full-level class splits into at most that many classes. This proves the bound.
\end{proof}

Fix a representative \(y\in H(\mathbb A_{k,f})\) of any one of the \(h\) orbits. Put \(H_\infty=H(k\otimes_\QQ\RR)\). Choose a single continuous function \(b\) of unit \(L^2\)-norm on \(H_\infty\), for Haar probability; \(b=1\) gives integer-valued tests. Every coset \(ya\), \(a\in A\), has a chart parameterized by \(H_\infty\times K_1\) in the full adelic quotient, using a representative of \(a\) at \(P\). There are no rational identifications in this chart because the rational stabilizer just proved is trivial. Define \(W_y\) to consist of functions supported on this orbit which in the chart of \(ya\) have value \(b(h_\infty)f(aZ)\), with \(f\) an arbitrary complex function on \(A/Z\). The definition is independent of the choice of lifts in \(K_0\). It is right Iwahori invariant as a subspace. Distinct orbits have disjoint support.

\begin{theorem}[An integral Bessel target with known covolume]\leavevmode\label{thm:lattice32}
The space \(W_y\) has dimension \(n_0\) and contains exactly one line for each character \(\gamma\in\Xi\). Normalize Haar measure on the global compact quotient to have total mass one. With scalar coordinates as above the squared norm on \(W_y\) is
\[
           \|f\|^2=(h n_0)^{-1}\sum_{aZ\in A/Z}|f(aZ)|^2.
\]
Take the cyclotomic field \(J=\QQ(\mu_{\exp(A/Z)})\) and the lattice \(\mathcal L=\OO_J^{A/Z}\) in these coordinates. For \(b=1\) it is the full lattice of \(\OO_J\)-valued test functions in the space at one embedding. Over all embeddings use the metric weighting each real place by one and each pair of complex places by two. Then
\begin{align}
 \frac{-2\log\operatorname{covol}(\mathcal L)}{[J:\QQ]n_0}
       &=\log(hn_0)-\frac{\log|D_J|}{[J:\QQ]},                 \label{eq:cov32}\\
 [\mathcal L:\textstyle\sum_{\gamma\in\Xi}\OO_J\gamma]
       &=n_0^{[J:\QQ]n_0/2}.                                  \label{eq:index32}
\end{align}
On functions already in the fixed archimedean coordinate and finite level \(K_1\), projection to \(W_y\) preserves the integral lattice after multiplication by \(|Z|=|A|/n_0\).
\end{theorem}
\begin{proof}
Each chart at level \(K_1\) has the same measure: it is a translate of \(H_\infty\times K_1\) divided by a trivial rational stabilizer. There are \(h|A|\) charts. Haar probability on the two compact factors and the norm of \(b\) therefore give the stated formula by summing over the \(|Z|\) cosets above each scalar coordinate. The characters of \(A/Z\) form an orthogonal basis of its regular representation. Finite-level projection consists of keeping the chosen orbit and averaging its coefficients over \(Z\), proving the denominator. The unscaled trace metric on \(\OO_J\) has covolume \(\sqrt{|D_J|}\), by the determinant of an integral embedding matrix with the complex pairs weighted as specified. Scaling each squared norm by \(1/(h n_0)\), in real dimension \([J:\QQ]n_0\), gives \eqref{eq:cov32}. In the coset-coordinate basis the columns of the character lattice form the square Fourier matrix. Under every embedding its Hermitian Gram matrix is \(n_0 I\). The absolute norm of its determinant is consequently \(n_0^{[J:\QQ]n_0/2}\), the asserted index.

In particular \(\log |A|\le m[k:\QQ]\log P\), and \(\log|D_J|/[J:\QQ]\le\log(\exp(A/Z))\le [k:\QQ]\log P\): the exponent of \(A\) divides the product of the \(q_v-1\). Together with Lemma~\ref{lem:mass32} this gives a uniform bound for the magnitude of \eqref{eq:cov32} on the common fields and intervals. For a general choice of \(b\) the integral structure asserted is on its scalar coordinates.
\end{proof}

\subsection{The remaining diagonal period values}\label{nav:I-48}
\label{sec:diagonal32}

Fix a unit source function at the one-dimensional unitary group in Theorem~\ref{thm:bessel29}, and give its complex span dimension one. For each chosen twist denote by \(F_\gamma\) its iterated lift restricted to the Bessel subgroup before projection, with \(\gamma\in\Xi\) its right residue character. Define an actual scalar integral on one finite class by
\[
 p_\gamma(y)=\int_{H_\infty}\int_{K_1}F_\gamma((h_\infty,yk))
                      \overline{b(h_\infty)}\,dk\,dh_\infty
\]
with probability measures on the two compact factors. This includes the projection onto the chosen archimedean coordinate.

\begin{proposition}[Two moment conditions from the individual periods]\label{prop:diag32}
On \(W_y\), the column of the projected Bessel operator on the unit source has scalar coordinates \(p_\gamma(y)\gamma(a)\). Put \(d_\gamma=|p_\gamma(y)|^2/h\) and \(\bar d=n_0^{-1}\sum_\gamma d_\gamma\). Its Gram operator on the source is the scalar \(d_\gamma\). If 
\[
      n_0^{-1}\sum_{\gamma\in\Xi}(d_\gamma-1)^2\le e\le1/4,
\]
the two moment hypotheses of Proposition \ref{prop:rank11}, hold with this \(e\) and \(\eta=e/(1-\sqrt e)^2\).
\end{proposition}
\begin{proof}
Choose representatives for \(a\) supported at \(P\). Moving \(a\) past the \(K_1\) average is allowed because it normalizes \(K_1\) and preserves its Haar probability. On the right of \(F_\gamma\) it acts by \(\gamma(a)\). Thus the archimedean and finite average on the coset \(ya\) equals \(p_\gamma(y)\gamma(a)\), already constant on \(Z\). Orthogonal projection is this averaging in the charts followed by restriction to the chosen orbit and the \(Z\)-average. The norm formula gives \(d_\gamma\). The columns lie in orthogonal lines which together have dimension \(n_0\), including lines where the column vanishes. Hence the eigenvalues of the sum of the output Gram operators are exactly \(d_\gamma\) with multiplicity. The first moment hypothesis is the assumed estimate. Cauchy--Schwarz gives \(\bar d\ge1-\sqrt e\), and \(n_0^{-1}\sum(d_\gamma-\bar d)^2\le e\). Consequently
\[
  \frac{n_0\sum_\gamma d_\gamma^2}{(\sum_\gamma d_\gamma)^2}
       =1+\frac{n_0^{-1}\sum_\gamma(d_\gamma-\bar d)^2}{\bar d^2}
       \le 1+\frac e{(1-\sqrt e)^2}.
\]
This is the second moment hypothesis. No nonvanishing of an individual scalar was needed for these identities.
\end{proof}

The common rank, Haar normalization, finite projector and covolume here are determined before the periods are evaluated. Finishing Section \ref{nav:I-28} by this construction would require a bound on the squared deviations of these actual integrals on a common orbit with the required normalization, and a proof that their finite-place valuations realize the prime counts for every translate. The estimates for the lattice of test functions alone supply neither assertion.

\subsection{The terminal Bessel character and its integral projector}\label{nav:VI-08}
\label{sec:char94}

Use a common field \(E\), its real subfield \(F_+=E^+\), and the
signed rank-\(r\) row lattice of
Section \ref{nav:I-28}.
Write \(T=\varpi\bar\varpi\) for a terminal prime there, with
\(\varpi\equiv1\bmod 2M\OO_{K_d}\), and \(f=[F_+:\QQ]\).
Pair \(s=m/2-1\ge1\) positive and negative coordinate vectors.
In the \(j\)-th pair take \(u_j=e_{+,j}+e_{-,j}\),
\(v_j=(e_{+,j}-e_{-,j})/2\). Thus \(h(u_i,v_j)=\delta_{ij}\),
with \(h\) linear in the second variable. Their orthogonal complement
is \(W=W_0\perp Ee\), with \(W_0\) positive of rank \(k\) and
\(h(e,e)=-1\). Set \(u=u_s,\ v=v_s\).
Let \(G=U(h)\), \(H=U(W_0)\), the latter fixing all the displayed
vectors. Let \(N\) be the unipotent radical of the parabolic
preserving the flag of the spans of \(u_1,\ldots,u_j\).
It acts as the identity on every graded piece of the flag completed
by its orthogonal flag, including the middle quotient \(W\).

For \(n\in N\), let \(d_i(n)\) be the coefficient of \(u_i\) in
\((n-1)u_{i+1}\) for \(i<s\), and let \(d_s(n)\) be the coefficient
of \(u_s\) in \((n-1)e\). Each is additive under multiplication in
\(N\). Indeed the cross term \((n_1-1)(n_2-1)\) takes the relevant
vector into the span of the earlier \(u_j\), \(j<i\) or \(j<s\),
respectively. The group \(H\) fixes these coordinates under
conjugation. Use the standard adelic additive character \(\Psi_F\)
of \(F_+\backslash\mathbb A_{F_+}\), with positive exponential at
infinity, and put \(\Psi_E=\Psi_F\circ\operatorname{Tr}_{E/F_+}\).
For \(\beta\in E^\times\) define
\[
 \chi_\beta(n)=\Psi_E\!\left(\sum_{i<s}d_i(n)+\beta d_s(n)\right).
\]
This is the character with last vector \(\beta e\). With
\([J]=J(F_+)\backslash J(\mathbb A_{F_+})\), its global Bessel
period on a row form \(F\) and a column form \(\phi\) is
\begin{equation}
 {\cal P}_\beta(F,\phi)
   =\int_{[H]}\overline{\phi(h_0)}
       \int_{[N]} F(nh_0)\overline{\chi_\beta(n)}\,dn\,dh_0 .
 \label{eq:period94}
\end{equation}
Use quotient probability on both groups. The inner coefficient is
invariant under rational left translation of \(h_0\), by conjugating
\(N\) and using invariance of its quotient measure and character.
Both quotients are compact. This follows for \(N\) from its
filtration with additive quotients and compactness of the additive
adelic quotient, and for \(H\) from definiteness, its finite genus
class set and compactness at infinity. Smooth automorphic inputs
therefore suffice for the integrations below. A finite right
translate at the endpoint can be included in \(F\). The coordinates
and integrand occur in
\srcloc{Development/AbelianVarietyTheory/BlockTheory/HeightDefect/BesselDatum}{63}
and
\srcloc{Development/AbelianVarietyTheory/Definitions/TwistedInnerProducts/WittSmallGroupEmbeddingMemPointwiseStabilizerIso}{176}.
Complex conjugating the column input changes between the two usual
pairing conventions.

\begin{lemma}[A translation law for the global period]\label{lem:root94}
For an adelic scalar \(t\) define \(n(t)\) on \(u,e,v\) by
\[
 u\longmapsto u,\qquad e\longmapsto e+t u,\qquad
 v\longmapsto v+\bar t e+(t\bar t/2)u .
\]
It is the identity on the orthogonal complement. For
\(R(a)F(y)=F(ya)\) one has
\begin{equation}
 {\cal P}_\beta(R(n(t))F,\phi)=\Psi_E(\beta t)\,{\cal P}_\beta(F,\phi).
 \label{eq:shift94}
\end{equation}
For rational \(t\), transporting the Witt decomposition and the last
vector by \(n(t)\) gives the same functional on the same forms.
\end{lemma}
\begin{proof}
The three images have the original Gram matrix. In particular the
cross terms in \(h(n(t)e,n(t)v)\) are \(-\bar t+\bar t\), and
\(h(n(t)v,n(t)v)=t\bar t-t\bar t=0\).
The flag relations put \(n(t)\) in \(N\). It commutes with \(H\),
and its only nonzero simple coordinate is \(d_s(n(t))=t\)
when \(t\ne0\). In the inner integral change the variable from
\(n\) to \(n n(t)\). The unipotent group is unimodular, so
right translation preserves the quotient measure. The conjugate
character supplies the factor in \eqref{eq:shift94}. If \(t\)
is rational, the transported flag has the same radical and the
same embedded \(H\). Its character is the original character
composed with an inner automorphism of \(N\), hence is identical.
This proves the assertion for the full functional. In particular,
with \(t=Mc/\varpi\), the two last vectors \(z_c\) and \(z_c/\varpi\)
in the transported frame give the functionals with \(\beta=\varpi\)
and \(\beta=1\), respectively.
\end{proof}

\begin{theorem}[Vanishing at a lattice containing the normalized vectors]
\label{thm:van94}
Fix a place \(w\mid T\) of \(F_+\). The split algebra there is
\(E_w=\QQ_T\times\QQ_T\), ordered with \(\varpi\) of valuation one
in its first factor. Suppose a local integral Hermitian lattice \(L\)
contains \(z_c/\varpi\) and \(z_{c'}/\varpi\) for two distinct
integers in \([1,T-1]\). If a row form \(F\) is right invariant
under the stabilizer of \(L\) at \(w\), then
\[
             {\cal P}_\beta(F,\phi)\ne0
       \quad\Longrightarrow\quad \beta\in\varpi\OO_{E_w}.
\]
In particular the functional for the normalized vector vanishes on
every such form.
\end{theorem}
\begin{proof}
Subtracting the two vectors gives \(\varpi^{-1}u\in L\), since
\(M(c-c')\) is a local unit. Substitution then gives \(e\in L\).
For \(t\in\varpi^{-1}\OO_{E_w}\) put \(a=tu\in L\). On an arbitrary
vector \(y\) the preceding isometry reads
\[
 n(t)y=y-a h(e,y)+e h(a,y)+\tfrac12 a h(a,y).
\]
Since \(T\) is odd and \(L\) is integral, this preserves \(L\),
as does its inverse \(n(-t)\). Use \eqref{eq:shift94} with \(t\)
supported at \(w\). A nonzero period forces \(\Psi_{E,w}(\beta t)=1\)
for every such \(t\). The standard character of \(\QQ_T\) has
kernel \(\ZZ_T\), and the local trace adds the two components.
Testing one component at a time gives the dual of
\(\varpi^{-1}\OO_{E_w}\) as \(\varpi\OO_{E_w}\). The scalar \(1\)
is outside this ideal. The condition for \(\beta=\varpi\) holds on
this root subgroup.
\end{proof}

\begin{theorem}[A projector retaining the period at a smaller level]
\label{thm:project94}
At each \(w\mid T\) denote the original stabilizer by \(K_w\).
In the first split factor let \(K'_w\) stabilize the lattice
enlarged by \(T^{-1}u\), with the dual lattice in the conjugate
factor. Let \(K'_{w,1}\) be its principal subgroup modulo \(T\).
For any row form fixed by all \(K_w\) one can form a finite average
\(P F\) fixed by all \(K'_{w,1}\) with
\[
           {\cal P}_1(PF,\phi)={\cal P}_1(F,\phi).
\]
This average is an orthogonal projection on the subspace invariant
under the integral root translations below, in any unitary
realization. Its denominator divides \(T^f\) in the lattice spanned
by the translates over a coefficient ring containing \(\mu_T\).
The denominator is exactly \(T^f\) on the standard local coset test,
with the product taken over \(w\mid T\).
\end{theorem}
\begin{proof}
In the first factor write \(n_w(x)=n((x,0))\). Its matrix is
\(1+xE_{u,e}\) in the Witt basis, so these translations add.
In the second factor its action is determined by duality.
Modulo those with \(x\in\ZZ_T\), the translations with
\(x\in T^{-1}\ZZ_T\) form a group of order \(T\).
Take the product \(C\) of these quotients at all \(f\) places.
They act on the subspace fixed by the integral translations.
Write \(R(b)\) for their action, and \(\chi(b)\) for the
character induced by \(\chi_1\). Define
\[
                  PF=T^{-f}\sum_{b\in C}\overline{\chi(b)}R(b)F.
\]
Equation \eqref{eq:shift94} gives the asserted period. Reindexing
the sum gives \(R(a)P=\chi(a)P\), and the average is the identity
on this character space. In a unitary representation \(R(a)^*=R(-a)\),
so \(P=P^*\).

At \(w\), write the original and enlarged first-factor lattices as
\(L_0\subset L_1\). They satisfy \(TL_1\subset L_0\).
Consequently any element acting trivially on \(L_1/TL_1\) preserves
\(L_0\), as does its inverse. Thus \(K'_{w,1}\subset K_w\).
All translations in the average belong to \(K'_w\), which normalizes
\(K'_{w,1}\). Each translated vector is therefore fixed by
\(K'_{w,1}\). Their span over the stated coefficient ring contains
\(T^fPF\) whenever it contains the translates of \(F\).
For sharpness take the characteristic function of the product of
the \(K_w\) in the product of the local groups, with the lattice
of integral-valued compactly supported functions. Two distinct
translates in the sum have disjoint supports. Indeed their
difference in some first split factor has a nonintegral
\((u,e)\)-entry. Each nonzero value of the projected test is
a root of unity divided by \(T^f\). This proves the exact
denominator there.
\end{proof}

These formulas apply to the global Bessel integral at the specified
level. The smaller-level average retains any nonzero value it starts
with. To use it for the integral theta lift requires comparing its
translate lattice to the prescribed arithmetic lattice of row forms.

\subsection{Uniform restriction to the Bessel orbit}\label{nav:VI-09}
\label{sec:restrict96}

In the notation of Section~\ref{sec:char94}, let
\(G_\infty=\prod_{\tau:F_+\hookrightarrow\RR}U(p_+,p_-)\),
\(K_\infty=\prod_\tau(U(p_+)\times U(p_-))\), and
\(K_\Lambda=\operatorname{Stab}(\widehat\Lambda)\).
For an open subgroup \(K_f\subset K_\Lambda\) write \(I=[K_\Lambda:K_f]\).
Use the compact-dual Haar density at infinity and give \(K_\Lambda\)
finite volume one, as in Theorem \ref{thm:norm91}. The global \(L^2\)
norm below uses the resulting quotient measure on \([G]\). Put
\[
 n_0=2fr,\qquad d_0=fr^2,\qquad
 B=(8n_0)^{2n_0}|D_E|^{r/2},
\]
where \(D_E\) is the discriminant of \(E\). The positive Euclidean
norm on the sum of the archimedean coordinate spaces is the sum of
the squared coordinate moduli, followed by the square root.
Use the Frobenius inner product on each real Lie algebra, the
orthogonal sum on their product, and an orthonormal basis
\(X_1,\ldots,X_{d_0}\) taken factor by factor. For a word \(w\)
in this basis, \(D_w\) denotes successive differentiations of the
right regular representation, including the identity for the
empty word.

\begin{lemma}[Integral lattices along the period orbit]\label{lem:inj96}
Let \(g=nh_0\) with \(n\in N(\mathbb A_{F_+})\) and
\(h_0\in H(\mathbb A_{F_+})\). In the real coordinate space put
\[
 {\mathcal L}_g=g_\infty^{-1}
       \bigl(E^r\cap g_f\widehat\Lambda\bigr)
\]
using the diagonal archimedean embedding. This lattice has first
successive minimum at least \(1/4\) and contains \(n_0\) independent
vectors of norm at most \(B\). Consequently, for
\(\gamma\in G(F_+)\cap g_fK_\Lambda g_f^{-1}\),
\[
 \big\|g_\infty^{-1}\gamma g_\infty-1\big\|_{\rm op}<1/(4B)
                    \quad\Longrightarrow\quad \gamma=1 .
\]
\end{lemma}
\begin{proof}
Order the graded pieces of a rational flag as
\[
 Eu_1,\ldots,Eu_s,\quad W_0,\quad Ee,\quad Ev_s,\ldots,Ev_1.
\]
Each partial sum is stable under \(NH\). On its quotient by the
preceding sum the action of \(N\) is the identity. The action of
\(H\) is the identity except on \(W_0\), where it preserves the
positive Hermitian form. These graded summands are orthogonal for
the positive archimedean metric. Given a nonzero rational vector
in \(E^r\cap g_f\widehat\Lambda\), pass to the first partial sum
containing it. Its nonzero image on that quotient has the following
local integrality. On a coordinate line, its coefficient belongs
to \(\tfrac12\OO_E\). Indeed the coordinates in \(u_j,v_j\)
of a pair with original coordinates \(a,b\) are
\((a+b)/2,a-b\), and the nonscalar finite action on \(W_0\)
does not enter these line quotients. On \(W_0\) the image belongs
to \(h_{0,f}\widehat{\OO}_E^k\) at the finite places. Its
Hermitian norm is then integral at every finite place and totally
positive. In the line case the squared moduli of the conjugates
of twice the coefficient have product at least one, by the
absolute norm of a nonzero algebraic integer. In the \(W_0\) case
the archimedean values of the Hermitian norm also have product at
least one. The arithmetic--geometric mean inequality in either
case bounds their sum below by \(f\). The lengths of \(u_j,v_j\)
are \(\sqrt2,1/\sqrt2\). The graded component of
\(g_\infty^{-1}\) times the vector thus has squared norm at least
\(f/8\). The other components only increase its positive norm.

The covolume of \({\mathcal L}_g\) is
\(2^{-fr}|D_E|^{r/2}\), that of \(\OO_E^r\) in its complex
embeddings. At every rational prime the finite isometry preserves
the nondegenerate rational trace form of \(h\), so its determinant
has absolute value one. Hence the two finite lattices have equal
volume. Passing to their common intersection shows that their
global indices there are equal. The archimedean isometry has
real determinant of absolute value one as well. This proves the
covolume assertion. Minkowski's second theorem applied to the
unit ball bounds the product of the successive minima by
\(n_0^{n_0/2}|D_E|^{r/2}\), using the inscribed cube of side
\(2/\sqrt{n_0}\).~\cite[Thm.~6.3.3]{XXIII-XVIII-IV-minima32}
The first minimum bound shows that twice the last minimum is
at most \(B\). Finally \(g_\infty^{-1}(\gamma-1)g_\infty\)
maps \({\mathcal L}_g\) to itself. Under the displayed hypothesis
it takes each of those \(n_0\) independent vectors to a vector
of norm below \(1/4\). Each image is zero, proving \(\gamma=1\).
\end{proof}

\begin{theorem}[A uniform bound for restriction]\label{thm:rest96}
For an absolute constant \(C\), every smooth right-\(K_f\)-invariant
automorphic form \(F\) with square-integrable derivatives through
order \(d_0\) satisfies
\begin{equation}
  \sup_{n,h_0}|F(nh_0)|
  \ \le\ I^{1/2}\exp\{C d_0\log(2d_0 B)\}
       \max_{|w|\le d_0}\|D_wF\|_{L^2([G])}.                 \label{eq:rest96}
\end{equation}
The supremum uses the full adelic groups of the preceding lemma.
\end{theorem}
\begin{proof}
Write \(d=d_0\) and \(\delta=(1000d^2 B)^{-1}\). Use the chart
\[
          \Phi(t)=\exp(t_1X_1)\cdots\exp(t_dX_d),
                        \qquad |t_i|\le\delta .
\]
For any partial product \(h\), the operator norms of \(h\) and
\(h^{-1}\) are at most \(e^{d\delta}\) in the block diagonal
matrix realization. Expanding the product, or telescoping the
exponentials, gives \(\|h-1\|_{\rm op}\le e^{d\delta}-1\).
The derivative of \(\Phi\) in its \(j\)-th variable is a product
of \(X_j\) with such factors on either side. It differs from
\(X_j\) in Frobenius norm by at most \(4d\delta\). Summing the
columns by Cauchy--Schwarz, the derivative differs from the
linear isometric inclusion at zero in operator norm by at most
\(4d^{3/2}\delta<1/2\). Integration on a segment shows that
\(\Phi\) is one-to-one on the cube. The left-translated tangent
columns are \(\operatorname{Ad}(h_{>j}^{-1})X_j\), where \(h_{>j}\)
is the product after \(j\). They have the same error bound from
\(X_j\), by the two operator norm estimates. Their volume
Jacobian for the left-invariant Frobenius density is at least
\(2^{-d}\).

The compact-dual density differs from this density by the inverse
Frobenius volume of \(U(r)^f\). This compact volume is at most
\(\exp(Cd)\). Indeed the first-column map from \(U(j)\) to the
unit sphere in \(\CC^j\) has fiber \(U(j-1)\). In the horizontal
directions \((ia,z)\) the Frobenius squared norm is
\(a^2+2\|z\|^2\), and the sphere norm is \(a^2+\|z\|^2\).
The quotient density is therefore \(2^{j-1}\) times sphere
volume. The sphere has volume \(2\pi^j/(j-1)!\), giving the
volume ratio \((2\pi)^j/(j-1)!\). Multiplication over \(j,\tau\)
proves the bound. Thus the Haar Jacobian of the cube is at
least \(\exp(-Cd)\).

For every \(g=nh_0\), the image of \(g(\Phi(t),k_f)\) in \([G]\),
as \(k_f\) varies in \(K_f\), has no rational identifications.
An identification would give \(\gamma\) as in Lemma~\ref{lem:inj96},
with \(g_\infty^{-1}\gamma g_\infty=\Phi(t')\Phi(t)^{-1}\).
The distance in operator norm to the identity is at most
\(e^{2d\delta}-1<1/(4B)\). The lemma forces \(\gamma=1\), and
then the two parameters agree.

For a smooth function \(b\) on the real cube one has
\[
 |b(0)|\le
 \sum_{J\subset\{1,\ldots,d\}} (2\delta)^{|J|-d}
      \int_{[-\delta,\delta]^d}|\partial_J b(t)|\,dt .
\]
In one variable, averaging \(b(0)=b(t)-\int_0^t b'(s)\,ds\)
gives the derivative an integral kernel of absolute value at most
one. Applying this linear formula successively in the \(d\)
variables and then taking absolute values gives the display.
Apply it to \(b(t)=F(g\Phi(t))\). In a mixed derivative take
the variables in decreasing order of index. Each differentiation
inserts the right Lie derivative in direction
\(\operatorname{Ad}(h_{>j}^{-1})X_j\). Directions already inserted
depend only on later variables, so their coefficients are constant
in this differentiation. Each direction has Frobenius norm at
most two. Expansion in the fixed basis thus gives at most \(d^d\)
words with coefficients of magnitude at most \(2^d\).

For any one word, Cauchy--Schwarz on the cube, the lower Haar
Jacobian and the injective chart bound its integral in absolute
value by
\((2\delta)^{d/2}e^{Cd} I^{1/2}\|D_wF\|_2\).
Indeed the finite factor of the chart has measure \(1/I\), and
archimedean differentiation commutes with \(K_f\). Summing the
terms costs at most an additional
\(\exp(Cd\log(2d/\delta))\). This gives \eqref{eq:rest96}
with the claimed dependence.
\end{proof}

\begin{lemma}[Iterated derivatives of a fixed compact type]\label{lem:cas96}
At every real place use the invariant form
\(\operatorname{Re}\operatorname{tr}(XY)\) on \(\mathfrak u(p_+,p_-)\).
Let \(F\) be a smooth vector of a unitary representation, smooth
in Hilbert norm in all orders, on which the Casimir for this form
at \(\tau\) acts by the scalar \(\omega_\tau\) on the representation
generated by \(F\). The \(K_\infty\)-span of \(F\) is assumed
finite dimensional. Let \(L\ge1\) bound \(\sqrt{|\omega_\tau|}\)
and the operator norms on this span of \(dR(Z)\) for all compact
Lie directions of Frobenius norm one at every place. A word of
\(j\) real derivatives, each at one place and of Frobenius norm
at most two, satisfies
\begin{equation}
          \|D_wF\|_2\le \{C r(L+j+1)\}^j\|F\|_2 .          \label{eq:cas96}
\end{equation}
\end{lemma}
\begin{proof}
The span of words of length at most \(l\) on the initial compact
span is finite dimensional and compact-group stable. It is an
equivariant image of the sum of the tensor powers through \(l\)
of the complexified Lie algebra tensored with the initial span.
Use the adjoint action on the Lie algebra. For compact unit
\(Z\), its differential \(Y\mapsto[Z,Y]\) has operator norm at
most two for the Hilbert--Schmidt metric. On the tensor
representation the bound is therefore \(L+2l\). The same holds
on its image with the given unitary metric. Indeed the polar part
of an equivariant linear map of finite-dimensional unitary
representations gives a unitary identification of the orthogonal
complement of its kernel with its image.

At one place take Frobenius orthonormal bases for the compact
and noncompact parts \(\mathfrak k,\mathfrak p\). The invariant
form is minus the Frobenius form on the former and equals it on
the latter. The Casimir is \(\sum_{\mathfrak p}X^2-\sum_{\mathfrak k}X^2\).
Skew adjointness of the derived operators on smooth vectors gives
for any vector \(v\) in the preceding span
\[
 \sum_{\mathfrak k\cup\mathfrak p}\|dR(X)v\|^2
       =2\sum_{\mathfrak k}\|dR(X)v\|^2-\omega_\tau\|v\|^2
       \le \{2r^2(L+2l)^2+|\omega_\tau|\}\|v\|^2 .
\]
Here the sums index the indicated bases. By Cauchy--Schwarz this
bounds a single derivative in any real direction of norm at most
two by \(Cr(L+2l+1)\|v\|\). Induct on \(l\), enlarging the
absolute constant, to obtain \eqref{eq:cas96}.
\end{proof}

\begin{theorem}[Decay of the unscaled automorphic integral]\label{thm:decay96}
Suppose \(F\) is a square-integrable smooth automorphic vector
satisfying the level and compact-type assumptions above. For
any nonnegative orders of total \(A\), the last vector
\(\varpi e\) of norm \(-T\) satisfies
\begin{equation}
 \begin{split}
 |{\cal P}_{\varpi}(F,\phi)|
 &\le I^{1/2}\exp\{Cd_0\log(2d_0B)\}
   \frac{\{Cr(L+A+d_0+1)\}^{A+d_0}}{(4\pi\sqrt T)^A}\\[-2pt]
 &\hspace{27mm}{}\cdot\|F\|_{L^2([G])}\,\|\phi\|_{L^2([H])}.
 \end{split}                                                    \label{eq:decay96}
\end{equation}
The same absolute constant may be enlarged in both occurrences.
\end{theorem}
\begin{proof}
Apply Proposition \ref{prop:der95}
and bound the restriction of \(D_{\varpi}F\) by
Theorem~\ref{thm:rest96}. In the positive orthonormal basis
in the directions \(u_s,e,v_s\), the generator of
\(n_\tau(b\eta_\tau)\), \(|\eta_\tau|=1\), has two entries of
absolute value \(\sqrt2\) and all others zero. Its Frobenius
norm is two. Each word in the restriction bound therefore
satisfies Lemma~\ref{lem:cas96} with length at most \(A+d_0\).
Both period measures have total mass one. These observations give
\eqref{eq:decay96}.
\end{proof}

\section{Integral tests, both coefficient lattices and their minors}\label{nav:g-XII-1}

\subsection{Testing a definite endpoint at its original level}\label{nav:V-57}
Use the definite endpoint of the complex iteration in
Proposition~\ref{prop:adjtheta234}. Write \(N\) for its
rank, \(K_e\) for its finite level, and \(\chi_e\) for the
finite right level character, with the trivial character
allowed. Use its integral Hermitian lattice model. Let
\(\delta_e\) be the determinant of the model lattice,
which is the same in every class since the transports are
finite isometries. For the endpoint real type \(\sigma_e\)
of dimension \(d_e\), put \(e_e,a_e\) equal to the sums of
its polynomial degrees and its degrees with determinant
factors counted as in Proposition~\ref{prop:source305},
now using \(N\). Put \(L_e=\prod_\tau e_\tau!\) for these
endpoint polynomial degrees and set
\[
 D_e=2fN,\quad R_e=2D_e^{D_e/2}|\Delta_E|^{N/2}\delta_e,
 \qquad B_e=(\sqrt N R_e)^{Nf}.
\]
Let \(K'_e\) be the intersection of \(\ker\chi_e\) in \(K_e\)
with the kernel of reduction at one factor above the odd
split prime \(g_1\) in its lattice model. Set
\(m_e=[K_e:K'_e]\), let \(h_e\) be the number of fine
classes at \(K'_e\), and let \(w_e\) be the measure of one
such chart in the Petersson measure, using probability on
the real compact group within the chart.

\begin{lemma}[An integral endpoint test with the original character]\leavevmode
\label{lem:endpoint305}
For every endpoint form \(F\) of the indicated real type
and original finite level character, some integral
coefficient test \(Q_e\) of the same level and character
satisfies
\[
 |\langle F,Q_e\rangle|^2
    \ge \frac{m_e^2 w_e B_e^{-4a_e}}
                {d_e h_e N^{2e_e}}\,\|F\|^2,\qquad
 \|Q_e\|\le m_e\sqrt{w_e}\,L_e B_e^{2a_e}.
\]
\end{lemma}
\begin{proof}
The subgroup \(K'_e\) is normal in \(K_e\). A rational
stabilizer at this fine level is finite by positive
definiteness and preserves its class lattice. At the
chosen split place it lies in a conjugate of
\(1+g_1 M_N(\ZZ_{g_1})\), which is torsion-free by the
binomial valuation argument of Theorem~\ref{thm:unique221}.
Thus its stabilizer is trivial. All fine charts have the
same measure \(w_e\). On each one write
\(F(g)=\operatorname{tr}(M_c\sigma_e(g))\) for the real
compact argument \(g\). This representation of the
isotypic space by matrices is the compact group
orthogonality theorem. It gives
\[
                       \|F\|^2=(w_e/d_e)\sum_c\|M_c\|_{\rm HS}^2 .
\]

At each fine class apply the trace lattice argument of
Proposition~\ref{prop:source305} with rank \(N\) and determinant
\(\delta_e\). It gives an integral spanning frame in each
embedding and integral dual functionals with singular values
between \(B_e^{-1}\) and \(B_e\). The two Schur lists
of Proposition~\ref{prop:schurframe259} on this class have
length \(N^{e_e}\) each. Their lower frame bounds are
\(B_e^{-2a_e}\), and the norms of individual vectors
and covectors are at most \(L_eB_e^{a_e}\) and \(B_e^{a_e}\),
respectively. These conclusions use the determinant factors
in that proposition and apply to each transported lattice.

Let \(q_{cab}\) be supported on class \(c\) with matrix
coefficient \(\ell_a(\sigma_e(g)v_b)\) there. Compact
orthogonality makes the squared pairing with \(F\) equal
to \(w_e^2/d_e^2\) times the squared matrix entry of
\(M_c\) against the two frame members, with adjoints
according to the inner product convention. Summing over
both lists gives at least
\((w_e/d_e)^2 B_e^{-4a_e}\|M_c\|_{\rm HS}^2\) by applying
the two lower frame bounds in succession. The sum over
classes is at least \(w_e B_e^{-4a_e}\|F\|^2/d_e\).
One of the \(h_e N^{2e_e}\) terms attains at least its
average. Its test has norm at most
\(\sqrt{w_e}L_eB_e^{2a_e}\).

Apply to this test the unnormalized character sum
\(\sum_{u\in K_e/K'_e}\overline{\chi_e(u)}R(u)\) in the
unitary right action. This operator is \(m_e\) times
the orthogonal projector and acts on \(F\) as \(m_e\).
Its output \(Q_e\) has the desired pairing and norm.
The sum is independent of representatives because the
test has level \(K'_e\). Translation by \(K_e\) transports
the class lattices isometrically, and in translated
representatives it preserves the integral coefficient
model. The scalar character values are roots of unity.
The sum therefore gives an integral test at the original
level over the coefficient ring containing these values.
\end{proof}

\subsection{Integral tests on the compact highest lines}\label{nav:V-58}
At the last definite pair of Theorem~\ref{thm:last244}, let
\(R\) denote either rank. At embedding \(\tau\) write
\(e_{j,\tau}=\lambda_{\tau,j}-\lambda_{\tau,j+1}\), padding the
partition with zeros. Put \(l_\tau=|\lambda_\tau|\) and
\(d=\sum_\tau l_\tau\). Use the coordinate exterior tensor
product of Proposition~\ref{prop:cartan244} and its unit highest
vector \(t_R\). The span of the compact orbit of \(t_R\) is
\(V_R\), of dimension \(D_R\). For \(f\) square matrices
\(Z_\tau\) of order \(R\) put
\[
 v_R(Z)=\bigotimes_{\tau,j}
    (Z_\tau e_1\wedge\cdots\wedge Z_\tau e_j)^{\otimes e_{j,\tau}} .
\]
This is a polynomial map with integer coefficients into \(V_R\).
Indeed the compact-invariant subspace \(V_R\) is invariant under
the complexified Lie algebra and hence under the connected group
\(\prod_\tau\mathrm{GL}_R(\CC)\). Its closedness gives the assertion
at singular matrices as well. Determinant characters from the
oscillator splitting are included in the compact representation
on class functions. The polynomial above uses its partition
factor alone.

\begin{theorem}[Two finite families of highest-type tests]\leavevmode
\label{thm:frames311}
There is a finite family of vectors \(v=v_R(Z)\), with a
probability distribution, such that for every \(y\in V_R\)
\[
       {\mathbb E}|\langle y,v\rangle|^2
                         \ge \frac{\|y\|^2}{4D_R}.
\]
One family has algebraic integer coordinates in the exterior
tensor lattice and \(\|\sigma v\|\le R^{d/2}\) at every complex
embedding \(\sigma\). A second family has unit norm at every
embedding and coordinates integral away from two. In the latter
family \(2^{dR(R-1)/2}v\) has algebraic integer coordinates.
In either construction the logarithm of the number of tests is
\[
   O\big(fR^2\log(2+R+d+D_R)\big).
\]
The probabilities enter only in the proof of the bound.
\end{theorem}
\begin{proof}
For independent Haar probability matrices \(U_\tau\in U(R)\)
the covariance of \(v_R(U)\) is \(I/D_R\). It commutes with the
product unitary group, its image is in the irreducible \(V_R\),
and its trace is one. In degree zero the singleton \(t_R\) suffices.

For the integral construction let \(b=1/[2(d+1)]\) and
\(r=1-b\). Choose
\(K=\max(d,\lceil b^{-1}\log(8/b^2)\rceil)\) and \(M=2K+1\).
For any \(|w|\le r\) the real trigonometric polynomial
\[
        p_w(z)=1+2\real\sum_{j=1}^{K}w^j\bar z^{\,j},
                          \qquad |z|=1,
\]
is positive. The infinite sum with its constant term is
\((1-|w|^2)/|1-w\bar z|^2\ge b/(2-b)\), and its omitted tail
has absolute value at most \(2r^{K+1}/b\le b/4\).
At the \(M\)-th roots of unity the values \(p_w(z)/M\) sum to
one. Orthogonality of these roots gives
\(\sum_{z^M=1}z^a p_w(z)/M=w^a\) for \(0\le a\le d\).
Given \(U\), sample each entry of each \(Z_\tau\) independently
from this distribution with \(w=r(U_\tau)_{ij}\).
The conditional mean of every monomial is its value at \(rU\).
Thus the conditional mean of \(v_R(Z)\) is \(r^d v_R(U)\).
Jensen's inequality and \(r^d\ge1-db\ge1/2\) prove the bound.
Every conjugate matrix has columns of norm \(\sqrt R\).
The exterior product of \(j\) columns has norm at most
\(R^{j/2}\) by the Gram determinant inequality. Multiplication
over the slots gives the stated bound. There are at most
\(M^{fR^2}\) choices.

For the unitary construction, elimination of the entries below
the diagonal factors any \(U\in U(R)\) into at most
\(R(R-1)/2\) real rotations on two coordinates, diagonal phases,
and a final unitary diagonal. In each elimination multiply the
two affected coordinates by phases to make the two entries real
and then rotate one to zero. An upper triangular unitary is
diagonal by orthogonality of its columns. Reversing the
eliminations uses at most \(2R^2\) one-parameter factors,
including the diagonal factors. Round each angle to a multiple
of \(2\pi/M'\), where
\(M'\ge 4\pi R^2(1+d)\sqrt{D_R}\) is the least such integer.
The change of a single factor in operator norm is at most
\(\pi/M'\). Telescoping a product of unitaries bounds the error
in one matrix by \(2\pi R^2/M'\). In a \(j\)-fold exterior power
the error is at most \(j\) times as large, by restriction of
the \(j\)-fold tensor power. Telescoping again bounds the change
of the unit vector by \(2\pi R^2d/M'\le1/(2\sqrt{D_R})\).
The finite set of rounded products is therefore close to every
point of the orbit by this amount. Send a Haar orbit vector to
a closest vector in this finite set, resolving ties by a fixed
ordering. This map is measurable. The triangle inequality in
\(L^2\) bounds the change in the root mean square pairing by
\(\|y\|/(2\sqrt{D_R})\), giving the assertion.

The sine and cosine of a rounded angle lie in
\(\ZZ[\mu_{4M'},1/2]\), and twice each is integral. Complex
embeddings carry a rotation to a real orthogonal rotation and
a phase to a phase. Thus all conjugate products are unitary.
Clearing one factor of two per real rotation clears a matrix,
and raising this denominator to the tensor degree clears the
vector. The elimination order is fixed and there are at most
\((M')^{2fR^2}\) lists of angles. This proves the size bound.
\end{proof}

\subsection{Changing the product basis in a marked period}\label{nav:V-60}
At an even rank-product edge keep the physical matrix coordinates,
the discriminant-one elliptic differential, the finite unit vacuum
\(\phi_0\) on the standard trace module, and the
Schr\"odinger--Fock identification of Theorem~\ref{thm:vac242}.
Use two \(K\)-orthonormal bases \(B,B'\) of the trace Hermitian
space to construct the normalized CM sections \(G_\phi^B,G_\phi^{B'}\)
by Lemma~\ref{lem:basis59}. The finite Schwartz vector \(\phi\)
is the same vector in the fixed finite oscillator representation.

\begin{theorem}[The common scalar for all finite tests]\leavevmode
\label{thm:framechange314}
There is one algebraic scalar \(u(B',B)\) with
\[
 G_\phi^{B'}=u(B',B)G_\phi^B \quad\hbox{for every finite test},\qquad
                    |\sigma u(B',B)|=1
\]
at every complex embedding. For a family of arithmetic period
entries obtained from the same directed definite column, changing
these bases on its edges multiplies all entries by the same
algebraic scalar \(u\) of absolute value one at every embedding.
Every product \(\mathcal A_a\overline{\mathcal A_b}\) of two such
entries is identical in the two constructions. For a nonzero
entry the normalized rational norm valuation
\[
 \frac{\operatorname{ord}_p |N_{J/\QQ}\mathcal A_a|}{[J:\QQ]}
\]
is independent of the bases at every rational prime \(p\).
\end{theorem}
\begin{proof}
Write \(F_\phi(z)\) for the theta functional on the coherent
vector in the proof of Theorem~\ref{thm:vac242}. The physical
coordinates and the finite vector fix this function. The
theorem gives \(F_\phi=C_B G_\phi^B=C_{B'}G_\phi^{B'}\),
with both constants nonzero, for every test. Thus their ratio
is common to all vectors and all differential orders. For
\(\phi_0\) the two normalized sections are nonzero, so some
ordinary derivative of each at the origin in the same multi-index
is nonzero. These derivatives are algebraic by
Theorem~\ref{thm:transport59}. Their ratio proves algebraicity
of \(u(B',B)\).

At an embedding the identity between the two vacuum germs
persists coefficientwise, hence for the entire functions in
physical coordinates. Their squared metric norms on the conjugate
trace torus are both \(\nu_\sigma^{D}\), by
Theorem~\ref{thm:transport59}. To check that the tori in that
calculation coincide, let \(V\) have as columns the embedded
vectors of \(B\) and write \(B'=B U\) with \(U\in U(D,K)\).
The matrix changes from \(V\) to \(VU\), while the trace module
in product coordinates changes from \(M_B\) to \(U^{-1}M_B\).
Applying \(\sigma\) to both changes leaves their product fixed,
also after the common elliptic ideal and period are inserted.
Moreover \(\sigma(U)\) is unitary since conjugation on \(K\)
commutes with either embedding of \(K\). The metric and Haar
probability are therefore the same. Comparison of the two norms
gives \(|\sigma u(B',B)|=1\).

Each summand in the arithmetic column of
Theorem~\ref{thm:column242} contains one derivative at each
edge, with the fixed orientation of that edge. Its basis change
multiplies every term by the product of the edge scalars, with
conjugation at any conjugated edges. This gives a common \(u\)
also after restriction, compact averaging and the finite signed
sum. Since \(u\bar u=1\) at the original embedding, this is an
algebraic identity. The Gram products are unchanged. In a common
number field \(|N_{J/\QQ}u|=1\) by the equality at all embeddings,
which proves the norm assertion. Passing to a larger coefficient
field preserves the normalized valuation by transitivity of norms.
\end{proof}

The source gives a determinant formula up to a root of unity for
the square of this scalar in
\srcloc{Formalism/SharpRestriction/GlobalSchurBounds/StandardLatticesGL3}{646}.

\subsection{Finite tests before the choice of class representatives}\label{nav:V-61}
Fix an edge of the definite column in Theorem~\ref{thm:column242}.
Its rational matrix space is \(W\), its reference lattice is
\(M=\mathfrak b^{-1}\operatorname{Mat}_{s\times r}(\OO_E)\),
and its reference finite vacuum is \(\phi_0\).
Write every prescribed finite test, before evaluating it at a
pair of classes, in the form
\[
                \phi^0=\sum_b c_b W(b)\phi_0,\qquad b\in W/M .
\]
Include the finitely many tests obtained by moving the right
translations of a finite character sum onto the kernel. Choose
integers \(T^0,Q^0\ge1\) such that every \(T^0 c_b\) is an
algebraic integer and \(Q^0 b=0\) in \(W/M\) for every term.
The integers exist by the finite vacuum expansion of
Theorem~\ref{thm:vac242}. External scalar weights may be carried
outside this expansion throughout.

\begin{theorem}[Class representatives at a specified rational prime]\leavevmode
\label{thm:neutral315}
At every rational prime \(p\), all integrals for this column can
be evaluated on representatives equal to the identity over \(p\),
with their original class weights and levels. On each edge, after
the class actions, the section has a vacuum expansion whose
coefficients have valuations at least \(-v_p(T^0)\) at every
place over \(p\). Modulo \(M\otimes\ZZ_p\), each of its shifts
agrees with the \(p\)-component of an original \(b\).
This also holds in the product lattice coordinates of
Theorem~\ref{thm:adapted314} at \(p\).
\end{theorem}
\begin{proof}
At a vertex with original integration subgroup \(J_v\) and
representative \(u\), choose an open compact \(J'_v\subset J_v\)
in the stabilizer of the standard lattice, fixing the finite
tests and contained in the kernels of the finite level characters.
Apply Theorem~\ref{thm:class61} to the class of this particular
\(u\) at level \(J'_v\). We obtain \(u'=\gamma u j\) with
\(\gamma\) rational unitary, \(j\in J'_v\), and \(u'\) supported
at completely split rational primes outside \(p\) and outside
all marked and nonspherical places. Do this once per original
class, at each vertex. These representatives give precisely the
original list modulo \(J_v\). Since \(j\in J_v\), the stabilizer
at \(u'\) for \(J_v\) is the conjugate by \(\gamma\) of the
stabilizer at \(u\). The embedded endpoint representative is
\(u'\oplus1\), with rational multiplier \(\gamma\oplus1\).

For an automorphic integrand the class-coordinate substitution
is \(k\mapsto\gamma_\infty k\). Apply it simultaneously to
the incident factors at each vertex. The right level fixes the
integrand and the substitution preserves compact Haar measure.
This proves equality for each original class contribution to the
period, including the endpoint pairing and the signed sum.
The volume factor and the scalars \(C_i\) of
Theorem~\ref{thm:vac242} are independent of the representatives.
In a fixed marked basis the normalized scalar \(\mathcal A\)
is therefore the same.

For an edge, the pair of representatives acts by the identity at
every place over \(p\). At a moving split place over \(q\ne p\),
the split change of variables carries a Weyl translate to a
translate of the moved vacuum. The elementary lattice intertwiner
in the proof of Theorem~\ref{thm:one61} expands that vacuum in
translates of the reference vacuum. Its multipliers are roots of
unity times the inverse square root of a \(q\)-power index,
apart from the splitting character. Each is a unit over \(p\).
For the character assertion, take a power of the finite idele
at the moving places which is principal times a ray unit in
the kernel, as in that proof. The principal generator is a
unit at every place over \(p\). The algebraic archimedean
character, after squaring if necessary, is a product of integral
powers of conjugates of this generator and their inverses. Every
such factor is a unit over \(p\), so the original value is also
a unit there. The argument for a norm-one character uses its
norm-one principal generator in exactly this calculation.

At \(p\) the finite shift has retained its component \(b_p\),
and all shifts introduced at the moving places have zero class
in \(W_p/M_p\), where \(M_p=M\otimes\ZZ_p\). Passing to a rational
representative preserves these local classes. Indeed the map
\(W/M\to W(\mathbb A_f)/\widehat M\) is bijective by the Chinese
remainder theorem after clearing a common integer denominator.
The phases incurred by changing the representative are roots of
unity. Multiplication by these units and summation preserve the
claimed valuation bound. Finally the product basis of
Theorem~\ref{thm:adapted314} generates \(M_p\). Its vacuum
intertwiner has index prime to \(p\) and adds only shifts integral
at \(p\). Its \(K\)-linear change of coordinates identifies the
two local quotient modules. This proves the last assertion.
\end{proof}

\subsection{A resolution of the identity by definite class tests}\label{nav:V-64}
Use the definite fine charts of Theorem~\ref{thm:lastrow311}.
Let \(h\) be their number and \(w\) the volume of each chart.
On a side of rank \(R\), let \(A_R\) be the tensor product of
the prescribed exterior powers at the real places, with the
determinant lines and their indicated duals. Use the usual
tensor Hermitian metric. Its coordinate wedge tensors are
orthonormal. Fix the product \(t\) of the standard highest
tensors of norm one. Its compact cyclic span \(V_R\) has
dimension \(d_R\). The entire highest-line space at the fine
level consists on each class of the functions
\(g\mapsto\langle\rho_R(g)t,v\rangle\).
Take conjugate Hilbert spaces where necessary to make the
test maps linear throughout. Let \(i=[K:J]\), and impose
the original level character by the unnormalized right
character sum of Theorem~\ref{thm:lastrow311}.

From \(\mathcal A_R=\bigoplus_{a=1}^h A_R\) to this
original-character space define \(S_R\) by making the
function on its own class for each summand, adding the
functions, and applying the character sum. Put
\(s_R=i^2w/d_R\). For a signed definite operator \(T_u\)
of Theorem~\ref{thm:closedperiod318}, take such a full
source and a fixed full target of this kind. Use its actual
scalar \(\kappa_u\). Let the matrix of all normalized periods
on the coordinate wedge tests be
\[
     Z_u=\kappa_u^{-1} S_0^\dagger T_u S_u .
\]
Here the index \(0\) designates the same endpoint tests for
every word. These are finite matrices in orthonormal ambient
coordinates. Their entries include the prescribed signs.

\begin{theorem}[Physical normalization of closed products]\leavevmode
\label{thm:charttrace320}
One has \(S_R S_R^\dagger=s_R I\) on the entire stated
character space. Hence for every ordered list \(u_1,\ldots,u_k\)
the algebraic number
\[
\begin{split}
 C(u_1,\ldots,u_k)
   &=\operatorname{tr}\Big(\prod_j Z_{u_j}Z_{u_j}^*\Big)\\
   &=\left(\prod_j {s_0s_{u_j}\over|\kappa_{u_j}|^2}\right)
       \operatorname{tr}\Big(\prod_j T_{u_j}T_{u_j}^\dagger\Big)
\end{split}
\]
has, after multiplication by \(\prod_j D_{u_j}^2\), valuation
at least \(\sum_j b_{u_j,p}\) at every \(\mathfrak p\mid p\).
The integer \(D_u\) is the compact and coefficient clearing
of Theorem~\ref{thm:originalgram315} for the class tests and
their prescribed finite translates, and \(b_{u,p}\) includes
the split improvement of Theorem~\ref{thm:poscolumn316}.
All constants on the second line are the chart masses, type
dimensions, level indices and period normalizations.
\end{theorem}
\begin{proof}
On one chart, Schur orthogonality gives the squared
Petersson norm \(w\|P_Rv\|^2/d_R\), where \(P_R\) is the
compact orthogonal projector onto \(V_R\). Every vector
of the fine highest-line space occurs this way. Thus the
single-chart map composed with its adjoint is \(w/d_R\)
times its chart identity. Distinct charts have orthogonal
supports. The finite character sum is \(i\) times an
orthogonal projector, so composing the direct sum with
that character sum and its adjoint gives \(i^2 w/d_R\)
times the identity on the projected space. In particular
\(S_0^\dagger/\sqrt{s_0}\) is an isometric embedding.
Multiplication of the matrices \(Z_u Z_u^*\) now gives
the second line of the formula, including every scalar.

The coordinate tests have algebraic class functions,
so their normalized signed periods are algebraic by
Theorem~\ref{thm:column242}. Fix \(p\). Choose all the
chart representatives neutral at \(p\) as in
Theorem~\ref{thm:neutral315}, using one common refinement
for the finite list and the same endpoint chart choices.
For a coordinate wedge at such a class the left covector
is integral on the localized class lattice, which is
standard there. Its determinant duals have the same
property. On the right the highest vector is integral.
Hence the coefficient clearing of
Theorem~\ref{thm:endpoint315} applies to these tests,
including the finite level translates. No projection of
the individual wedge onto \(V_R\) is needed in this
integrality argument, since the compact integral itself
is still present with its specified clearing. Theorems
\ref{thm:originalgram315} and~\ref{thm:poscolumn316}
bound each product of a matrix entry and a conjugate
entry for its word. Expanding the trace proves the
asserted bound. For another prime the chart tests may
change. The second line of the formula proves equality
of the resulting traces in the original complex embedding
and hence equality of the algebraic numbers. Thus the
local estimates apply to one scalar on the given field.
\end{proof}

\subsection{Integral duality for the terminal polynomial tests}\label{nav:V-66}
At one embedding subtract the last entry of the terminal
highest weight to obtain a partition \(\lambda\) of \(d\)
with at most \(n\) rows. Use a generator for the separate
determinant line in the chosen frame. Write \(A\) for the
valuation ring at a finite place of the coefficient field,
with fraction field \(K\). On \(n\)-by-\(n\) matrices the
test space \(V\) over \(K\) has column multidegree
\(\lambda\) and is fixed by additions of earlier columns to
later ones. Write \(L_0\) for the \(A\)-span of products of
flag minors, one for each column of the Young diagram,
using the first \(\lambda'_j\) matrix columns for the
diagram column of height \(\lambda'_j\). All row choices
are allowed. Write \(L_1=V\cap A[X]\). The characteristic
zero Cauchy decomposition identifies this highest-weight
space with one Schur module for the matrix rows. The flag
minor products span it since their span is stable under row
changes by the minors formula and contains the product of
leading minors. This also treats \(d=0\) with empty product one.

For \(P=\sum P_\alpha X^\alpha,\ Q=\sum Q_\alpha X^\alpha\)
use the bilinear form
\[
 B(P,Q)=H^{-1}\sum_\alpha \alpha!\,P_\alpha Q_\alpha,
 \qquad H=\sum_\alpha\alpha!\,(P_{\lambda,\alpha})^2 ,
\]
where \(P_\lambda\) is the product of leading minors.
Thus \(H\) is the positive integral hook product of
Lemma~\ref{lem:hook86}. For a lattice \(L\) in \(V\) set
\(L^\#=\{v\in V:B(v,L)\subset A\}\).

\begin{theorem}[Detection by integral flag minors]\leavevmode
\label{thm:terminaldual325}
For every \(L_0\subseteq L\subseteq L_1\) one has
\[
                 H L\ \subseteq\ L^\#\ \subseteq\ L_0 .
\]
In particular, if every pairing of \(w\in V\) with the
flag minor products belongs to a fractional ideal
\(\mathfrak a\) of \(A\), then \(w\in\mathfrak a L_0\).
At primes not dividing \(H\) the three lattices
\(L_0,L,L_1\) agree and the pairing is perfect on them.
The same conclusions hold for the tensor of the terminal
weights over all embeddings, with \(H=\prod_\tau H_\tau\)
and products of the indicated tests.
\end{theorem}
\begin{proof}
Let \(T=(A^n)^{\otimes d}\), with the tensor basis
orthonormal for the bilinear dot product. Index its places
by the boxes of the diagram. In diagram row \(i\) send a
basis vector \(e_k\) to \(X_{ki}\) and multiply to obtain
a polynomial. Over \(K\) the inverse to this multiplication
on row-symmetric tensors is symmetrization within the rows.
Indeed the distinct monomials correspond to disjoint row
orbits of basis tensors. Put \(R=\prod_i\mathfrak S_{\lambda_i}\),
\(C=\prod_j\mathfrak S_{\lambda'_j}\),
\(a=\sum_{\sigma\in R}\sigma,\ b=a/|R|\) and
\(c=\sum_{\sigma\in C}\operatorname{sgn}(\sigma)\sigma\).
In the symmetric realization \(L_0=bcT\), by expansion of
each determinant. For a monomial \(X^\alpha\) its
symmetrized tensor has squared norm \(\alpha!/|R|\) by
counting its row orbit. Distinct monomials are orthogonal.
Thus \(B\) on \(V\) is \(|R|/H\) times the tensor dot product.

We record the scalar by which \(bc\) acts on \(V\).
For every permutation \(x\) of the diagram boxes, \(axc\)
is either zero or a signed copy of \(ac\). If a row contains
two boxes of a single moved column, their transposition
fixes the row sum and negates the moved column sum. If each
intersection contains at most one box, the first moved
column meets every nonempty row. After deleting it, the
next meets every row still having room, since its height
is exactly the number of those rows. Continuing, a single
permutation within the rows carries all moved columns to
their original column sets. Thus \(x\in RC\), proving
the claim. Expanding the middle two factors in \(acac\)
now gives \(acac=h\,ac\) for an integer \(h\).

Fill each diagram row \(i\) entirely with \(e_i\), giving
a tensor \(t_0\). The vector \(w_0=bc t_0\) represents
\(P_\lambda\) and has coefficient one against \(t_0\).
Indeed the row symmetrizer fixes \(t_0\), and only the identity
column permutation returns it to itself. On \(bcT\) the
operator \(bc\) acts as \(h/|R|\). Since \(b,c\) are
self-adjoint and \(bw_0=w_0\),
\[
 (w_0,w_0)=(t_0,cw_0)=(t_0,bcw_0)=h/|R|.
\]
The definition of \(H\) therefore gives \(h=H\).

If \(w\in L_0^\#\), its pairing with \(bcT\) shows that
\(t=(|R|/H)cw\) is in the integral self-dual tensor
lattice \(T\). It is alternating in each diagram column.
Every such integral tensor lies in \(cT\). On a basis
orbit with a repeated letter in a column, a transposition
forces its coefficient to vanish. On every remaining
orbit the column group acts freely, and the coefficients
are a single integral coefficient times the signs. Taking
one representative per orbit gives a preimage under \(c\).
Now \(bt=(|R|/H)bcw=w\), so \(w\in bcT=L_0\).
Consequently \(L^\#\subseteq L_0^\#\subseteq L_0\).
The integral coefficients of vectors in \(L_1\) give
\(H L\subseteq L^\#\) directly. Apply the two inclusions
also to \(L_1\) to get equality when \(H\) is a unit.
Scaling by a generator of \(\mathfrak a\) gives the
fractional ideal assertion. Over a discrete valuation
ring all these lattices are free, and their duals for a
product pairing tensor by taking dual bases. Each \(L_1\)
is saturated in its free module of integral monomials and
therefore extends by a free complement. Tensoring these
decompositions identifies the full integral polynomial
lattice on the product space with the tensor of the \(L_1\).
The inclusions proved above now give the product assertions.
\end{proof}

The coefficient convention, including the normalization by
the extremal vector, is specified in
\srcloc{Formalism/LocalToGlobal/Columns/CoefModB}{17--20, 55--66}.

\subsection{The original trace lattice at every finite place}\label{nav:VI-59}
Let \(E,F\) be the fields of Theorem~\ref{thm:first250}, let
\(K=\QQ(\sqrt{-d_0})\), and take the CM ideal \(\mathfrak c\)
of the given field. Write \(u_0=N_{K/\QQ}\mathfrak c\).
Use the ideal \(\mathfrak b_{\rm D}\) of Theorem~\ref{thm:pol49}
with \(\mathfrak b_{\rm D}\overline{\mathfrak b}_{\rm D}
=\mathfrak D_{E/K}\). Put
\(\mathfrak A=(\mathfrak c\mathcal O_E)\mathfrak b_{\rm D}^{-1}\)
and \(\xi_0=-1/(u_0\sqrt{-d_0})\). These are the ideal and scalar
in \srcloc{Formalism/LocalRingCardinality/DataHost}{147--161}.
For two Hermitian spaces with self-dual lattices
\(\Lambda_S,\Lambda_T\), give the space of \(E\)-linear maps
from \(S\) to \(T\) the Hermitian form
\(h(X,Y)=\operatorname{tr}(Y^\dagger X)\), where \(Y^\dagger\)
is the Hermitian adjoint. The sign of either Hermitian form may
be reversed. Define \(M=\operatorname{Hom}_{\mathcal O_E}
(\Lambda_S,\Lambda_T)\).

\begin{proposition}[Trace duality for the integral oscillator]\leavevmode
\label{prop:tracelat253}
At each finite place \(v\) of \(F\), the lattice
\((\mathfrak A M)_v\) is its own annihilator for the alternating
pairing given by the standard absolute additive character applied
to \(\operatorname{Tr}(\xi_0 h(X,Y))\). Here the trace runs from
\(E\otimes_F F_v\) to \(\QQ_p\), \(p\) below \(v\).
At a split place its two halves are exact duals. At every marked
place both components of \(\mathfrak A\) are units.
At an inert place this pairing and lattice admit coordinates with
a trace-zero unit and a conductor-\(\mathcal O_{F_v}\) character.
\end{proposition}
\begin{proof}
Self-duality of \(\Lambda_S\) identifies its Hermitian dual with
its \(\mathcal O_E\)-linear dual via conjugation. Thus the
rank-one maps \(x\mapsto t h_S(s,x)\), with the Hermitian form
here linear in \(x\), span \(M\). Local bases make this explicit:
the Gram matrices and their inverses are integral since their
determinants are units. Inserting them into the matrix trace
therefore gives a perfect Hermitian pairing on \(M_v\). In the
split algebra the pairing exchanges its factors. This identifies
the finite oscillator lattice specified by the rank-one generators
with \((\mathfrak A M)_v\).

The quadratic norm identity gives \(\mathfrak c\bar{\mathfrak c}
=(u_0)\). Since \(-d_0\) is the fundamental discriminant of
\(K\), transitivity of the different gives
\[
 (\xi_0)\mathfrak A\bar{\mathfrak A}
 =(\sqrt{-d_0})^{-1}\mathfrak D_{E/K}^{-1}
 =\mathfrak D_{E/\QQ}^{-1}.
\]
The trace dual of the integer ring is the inverse different.
Combining this fact with the perfect matrix pairing gives the
dual lattice \((\xi_0)^{-1}\mathfrak D_{E/\QQ}^{-1}
\bar{\mathfrak A}^{-1}M=\mathfrak A M\). Taking completions
gives the asserted annihilator for the additive character, whose
\(\QQ_p\)-kernel is \(\ZZ_p\). Conjugation negates \(\xi_0\),
so the trace pairing is alternating. At a split \(v\) it vanishes
on each half separately, and the annihilator identity identifies
the full dual of either half with the other.

The rational primes under the marked set split completely in
\(E\). For the two level primes, the two type primes and the
single auxiliary prime this follows from their defining
congruence \(p\equiv1\) modulo a cyclotomic conductor for \(E\):
Frobenius raises a root of unity to its \(p\)-th power.
In particular the different is a unit at each of them. The CM
ideal is chosen prime to the level and type primes. The
additional auxiliary prime is chosen prime to \(u_0\) as well.
Both components of \(\mathfrak b_{\rm D}\) and of
\(\mathfrak c\mathcal O_E\) are consequently units. The prime
choices are specified in
\srcloc{Development/AbelianVarietyTheory/Definitions/Foundations/AuxiliaryFrameNormalizedSpaces}{18--32}
and \srcloc{Development/AbelianVarietyTheory/Definitions/Divisors/LeastSplitAuxiliaryPrimeThetaChainSpec}{209--235}.

At an inert \(v\) let \(a_v\) generate \(\mathfrak A_w\),
and let \(b_v\in F_v^\times\) generate the inverse different
of \(F_v/\QQ_p\). Write \(X=a_v X'\) and \(Y=a_vY'\).
The extension \(E_w/F_v\) is unramified, so transitivity in
the displayed identity gives
\((\xi_0 a_v\bar a_v)=b_v\mathcal O_{E_w}\).
Hence \(\xi_0 a_v\bar a_v/b_v\) is a trace-zero unit. In
these coordinates the additive character on \(F_v\) is
\(z\mapsto\psi_v(b_v z)\), where \(\psi_v\) is the original
trace character. Its conductor is \(\mathcal O_{F_v}\) by
trace duality. The two scalar changes cancel in the argument
of the absolute character, so this is an identity for the
original pairing. Finally \(d_0\equiv7\bmod8\) makes \(2\)
split in \(K\), hence all dyadic places of \(F\) split in
\(E=FK\).
\end{proof}

\subsection{The spherical section for the original finite lattice}\label{nav:VI-60}
Use the self-dual map lattice \(L_v=(\mathfrak A M)_v\) of
Proposition~\ref{prop:tracelat253} for a step of ranks
\(u=n_i,\ N=n_i+d\). Outside \(S_0\) there is no congruence on
the displacement lattice. In an integral polarization write the
two halves of \(L_v\) as \(L_X,L_Y\). Use the position indicator
\(1_{L_X}\) and its matrix coefficient divided by its squared
norm. Equivalently the position Haar measure may give \(L_X\)
volume one. This latter convention retains the integer values
of the test. The determinant characters and the additive
character are those of the original pair.

\begin{proposition}[The unmarked local test in the doubling integral]\leavevmode
\label{prop:vacuumcore254}
At every finite \(v\notin S_0\) this test satisfies the depth
equations, is fixed by both prescribed lattice stabilizers and
gives the normalized spherical doubling integral on the spherical
source line. These assertions include places at which the absolute
trace character has nontrivial different. In particular the finite
set omitted in the partial product of Theorem~\ref{thm:coreeuler252}
may be exactly \(S_0\).
\end{proposition}
\begin{proof}
The two quadratic imaginary fields of the construction generate
the same extension of \(F\), since \(\sqrt{q d_0}\) lies in \(F\).
At every finite prime at least one of the two quadratic fields
is unramified or split. At \(2\) one is split. Thus a nonsplit
finite extension \(E_w/F_v\) is unramified and has odd residue
characteristic. Exact duality of \(L_X,L_Y\) gives invariance of
the indicator under the pure translations and modulations from
these lattices. For doubled mixed displacements the remaining
phase is trivial by the calculation in Theorem~\ref{thm:markedcore253}.

At a split place we can take the two idempotent halves. In
lattice bases \(L_X\) is a scalar multiple of the rectangular
integral matrix lattice. Dilation by that scalar, normalized on
\(L^2\), commutes with the two group substitutions and their
determinant factors in the split model. Consequently the entire
normalized matrix coefficient on the smaller group is identical
to that of the ordinary rectangular lattice. The unramified
splitting characters are trivial on unit determinants. This also
proves invariance under both maximal compact groups. This
calculation holds in residue characteristic two.

Suppose now \(E_w/F_v\) is a field and put \(Q=Nw\). Since
\(\kappa_i\) is unramified and trivial on \(F_v^\times\), it is
identically one on \(E_w^\times\), since the two fields have a common
uniformizer. Both Hermitian lattices have even rank and admit
integral hyperbolic bases. Indeed, over the quadratic residue
field the norm is surjective, so a nonsingular diagonal
Hermitian form of rank at least two has an isotropic line.
Lift its generator to a primitive isotropic vector over the
integer ring by Hensel's lemma. The derivative in a suitable
coordinate direction is a unit, since the vector is primitive,
the lattice pairing is perfect and the residue trace is onto.
A partner pairing to one can be made isotropic by subtracting
half its self-pairing times that vector. Splitting off the
resulting unimodular plane proves the claim by induction.

Polarize the larger factor into its two isotropic halves. The
Schr\"odinger variable then consists of \(N/2\) copies of the
smaller space, with position lattice given by \(N/2\) copies
of \(\mathfrak A_w\Lambda_u\) in these bases. For the trivial
splitting character the smaller unitary group acts by substitution
alone. In the polarization by the smaller factor the larger group
likewise acts by substitution. These are the two models of the
same Weil representation. Indeed switching the factors transfers
a trace-zero scalar from one form to the other, and changing the
polarization is a partial Fourier transform. On the changed
isotropic coordinates this transform sends a lattice indicator
to its volume times the indicator of its annihilator. By
Proposition~\ref{prop:tracelat253} that annihilator is exactly
the corresponding component of \(L_v\). This proves invariance
under the larger lattice stabilizer as well. The unitary action
by substitution and the switching convention are given in
Liu~\cite[Section~3]{XXIII-XVIII-liu328}
and Ichino~\cite[Section~2.2]{XXIII-XVIII-ichino328}.

For clarity one can identify the section on the integration
group itself. Write a Cartan representative of the smaller group
in the hyperbolic basis with entries \(\varpi^{t_j},\varpi^{-t_j}\)
for \(1\le j\le u/2\), \(t_j\ge0\). The normalized overlap of
the position lattice with its translate on each of the \(N/2\)
copies is \(Q^{-\sum_j t_j}\). Thus the coefficient at \(h\)
with these elementary divisors is \(Q^{-(N/2)\sum_jt_j}\).
The normalized spherical section at \(s_0=d/2\) has exactly this
restriction. To see it, take the primitive exterior generator of
the diagonal \(u\)-plane in the doubled integral lattice.
Applying \((h,1)^{-1}\) changes its lattice norm by
\(Q^{\sum_jt_j}\). On a diagonal representative the corresponding
wedged columns have disjoint coordinate pairs, so their maximum
minor has that absolute value. The integral compact groups
preserve the norm. Its reciprocal is the absolute determinant
of the Siegel Levi in the Iwasawa decomposition of \((h,1)\).
The normalized section raises it to \(s_0+u/2=N/2\).

For both kinds of place the equality on \((h,1)\) persists
for flat extensions of the sections, since both sides are then
multiplied by the same Iwasawa height to the power \(s-s_0\).
At a split place the ordinary unramified rectangular indicator
used in that comparison is the standard spherical test, as in
the split calculation for Theorem~\ref{thm:global175}. The
unramified doubling formula now gives precisely the local factors
there and in Theorem~\ref{thm:coreeuler252}. This proves the
assertion for the original fractional lattice.
\end{proof}

\subsection{Minors detected by incoming Bessel rows}\label{nav:VI-61}
Keep the two definite spaces, levels, measures, and families of
Theorem~\ref{thm:lastrow311}. Write \(\mathcal U\) for the
span of the actual incoming vectors with that word and initial
character. It is a complex subspace of the penultimate
highest-line space. For each source test \(U_a\) and target
test \(V_b\) form the matrix \(Z_{ba}=z(U_a,V_b)\).

\begin{theorem}[Nonzero arithmetic minors at the original levels]\leavevmode
\label{thm:minors311}
For every \(1\le j\le\dim\mathcal U\) there are \(j\) rows and
\(j\) columns of \(Z\) whose minor \(z_j\) satisfies
\[
                  |z_j|\ \ge\ j!\,(i_-i_+ L_*)^j .
\]
This holds with either family of Theorem~\ref{thm:frames311}
at either end. At least \(j=1\) occurs on every good orbit.
For \(j=1\) one may take the same target test as in the
lower bound on \(|\langle TX,V\rangle|\) for a given \(X\).
\end{theorem}
\begin{proof}
Move complex conjugation to the source space when necessary
so that the operator whose matrix is being formed is linear.
Order the two lists of tests and let \(p_a,q_b\) be their
probabilities, with a uniform choice of fine class included.
Let \(E_-\) have columns \(\sqrt{p_a}U_a\), and \(E_+\) have
columns \(\sqrt{q_b}V_b\). On the respective finite-level,
highest-line Hilbert spaces the mean square inequality in
Theorem~\ref{thm:lastrow311} gives
\[
        E_\pm E_\pm^*\ \ge\ \ell_\pm I,\qquad
        \ell_\pm=\frac{i_\pm^2 w_\pm}{4h_\pm D_\pm^2}.
\]
The last edge here includes the projection to the specified
target line and finite character. It carries the actual incoming
space to that target with squared norm at least \(a_*\) times
the squared source norm. By the minimum--maximum principle
its \(j\)-th singular value is therefore at least \(\sqrt{a_*}\).
Further,
\(TE_-E_-^*T^*\ge\ell_-TT^*\), and
\((TE_-)^*E_+E_+^*(TE_-)\ge\ell_+(TE_-)^*(TE_-)\).
Comparison of eigenvalues proves that the first \(j\) singular
values of \(E_+^*TE_-\) are at least
\(\sqrt{\ell_-\ell_+ a_*}\).

Multiply this matrix by the scalar used to define \(z\).
Its entries are now \(\sqrt{q_b p_a}\,Z_{ba}\), up to the
choice of conjugate coordinates, and its first \(j\) singular
values are at least \(i_-i_+L_*\). The squared Hilbert--Schmidt
norm of its \(j\)-th exterior power is the sum of the squared
absolute values of its \(j\)-rowed minors. In a singular value
basis this norm is at least \((i_-i_+L_*)^{2j}\). In the test
bases it is
\[
  \sum_{\substack{|A|=j\\ |B|=j}}
       \left(\prod_{a\in A}p_a\right)
       \left(\prod_{b\in B}q_b\right)
                      |\det Z_{B,A}|^2 .
\]
For any probability list the sum of its squarefree products
of size \(j\) is at most \(1/j!\), by expanding the \(j\)-th
power of the sum. Bounding the last display by its largest
squared minor times the total weight proves the result.
The assertion about a specified incoming vector is the
successive choice of the two tests in
Theorem~\ref{thm:lastrow311}.
\end{proof}

No rational basis of \(\mathcal U\) enters this calculation.
The rows and columns indexing the minor belong to the two
original projected compact spaces. The dimension in the
hypothesis is the dimension of the image of the preceding
theta lifts on the Bessel rows.

\subsection{Finite-place bounds along the entire definite column}\label{nav:VI-62}
Fix the signed period \(\mathcal P\) of Theorem~\ref{thm:column242},
with its prescribed projected endpoint and the algebraic data of
that theorem. Use its normalization
\(\mathcal P=V_J(\prod_i C_i)c^{-t/2}\mathcal A\), its finite
levels and its class representatives throughout. At edge \(i\)
write \(\ell_i=N_{i-1}N_i\), \(d_{i,\tau}\) for the surviving
degree at \(\tau\), and \(d_i=\sum_\tau d_{i,\tau}\). If no degree
profile survives, \(\mathcal A=0\). In the formulas below we use a
surviving profile. The scalar characters fix it throughout the signed
sum by Lemma~\ref{lem:degrees242}.

On each edge take the product presentations of its finite class
sections as in Theorem~\ref{thm:classden312}. For their derivatives
of block degrees \(d_{i,\tau}\) let \(\Gamma_{i,\mathfrak p}\)
be the maximum loss given by its finite calculation. Include all
terms and tuples of classes. For a conjugated derivative use the
place transported by conjugation. Use a common Galois coefficient
field \(J\) containing their algebraic data. At vertex \(v\)
let \(A_v\) be the least common multiple of its stabilizer orders.
Write \(n_v\) and \(L_{v,\tau}\) for its rank and the polynomial
degree bounds of Theorem~\ref{thm:column242}, including those at
the endpoint subgroup.

To specify the other finite scalar costs, expand the compact
integrand of that theorem using distinct formal symbols for the
derivatives of each edge and class section. Let \(B_0\ge1\) clear
the coefficients of these finitely many polynomials in the
symbols and the compact matrix entries and their conjugates.
This uses the endpoint coefficients, the real test coefficients,
the signed weights, and any finite projection scalars outside the
class sections. No derivative has been evaluated in forming \(B_0\).
Put
\[
 D_{\rm H}=\prod_{v,\tau}D(n_v,L_{v,\tau}),\qquad
 D_0=B_0D_{\rm H}\prod_v A_v.
\]

Let \(S\) be any finite set of good primes to which
Theorem~\ref{thm:direct76} applies to every edge and every class
section of these finite tests. Thus \(p\ge5\) is inert in \(K\),
unramified in \(E_x\), and the common local vacuum and integral
coefficients at \(p\) are as in that theorem. Choose the
orthonormal basis there for \(S\) on each edge. In an even cycle
\(C\) of the permutation \(\pi_p\) of the embeddings choose a half
\(H_{i,C}\) for edge \(i\), common to its \(\ell_i\) matrix entries.
With \([z]_+=\max(z,0)\) put
\[
 \begin{split}
 b_M(a,\ell)&=\max(0,\lceil(a-\ell(M-1))/M\rceil),\\
 g_{i,p}&=\sum_{C\ {\rm even}}\sum_{\tau\in H_{i,C}}
          \sum_{\substack{j\ge1\\ j\ {\rm odd}}}
   [\,b_{p^j}(d_{i,\tau},\ell_i)-d_{i,\pi_p^j\tau}\,]_+ .
 \end{split}
\]

\begin{theorem}[Divisibility after all compact integrations]\leavevmode
\label{thm:finitecolumn313}
For the scalar of this signed column, keeping its original
restriction and character projection, one has at every
\(\mathfrak p\mid p\)
\[
 v_{\mathfrak p}(D_0\mathcal A)\ \ge\
 \begin{cases}
    \sum_i g_{i,p} &p\in S,\\
    -\sum_i\Gamma_{i,\mathfrak p} &p\notin S.
 \end{cases}
\]
If \(A_i^{\rm jet}\) is the integer \(A\) of
Theorem~\ref{thm:periodden312} computed for the sections and degree
of edge \(i\), then \(D_0(\prod_i A_i^{\rm jet})\mathcal A\)
is an algebraic integer.
\end{theorem}
\begin{proof}
All column spaces are definite and their consecutive rank products
are even by Proposition~\ref{prop:ranks47}. Their section
presentations follow from the same normalized finite lattice
intertwiners and translated elliptic products used in
Theorem~\ref{thm:classden312}. This is precisely the construction
at each edge in the proof of Theorem~\ref{thm:column242}. The
chain rule there requires only the product expression, so the
bound \(-\Gamma_{i,\mathfrak p}\) and its integer clearing factor
hold at each edge in its own matrix coordinates.

For \(p\in S\), take any monomial derivative at edge \(i\). If
\(\alpha_{a,\tau}\) are its exponents in the \(\ell_i\) entries,
then \(\sum_a\alpha_{a,\tau}=d_{i,\tau}\).
For \(M=p^j\), division with remainder gives
\(\sum_a\lfloor\alpha_{a,\tau}/M\rfloor
\ge b_M(d_{i,\tau},\ell_i)\). The inequalities
\(\sum_a[u_a]_+\ge[\sum_a u_a]_+\) and
\(\sum_a\alpha_{a,\pi_p^j\tau}=d_{i,\pi_p^j\tau}\) now give
\(v_{\mathfrak p}(\partial^\alpha G_i(0))\ge g_{i,p}\) from
Theorem~\ref{thm:direct76}. This uses only the block degrees,
so it holds for every monomial produced by a compact matrix
substitution. The sum over \(j\) is finite since eventually
all the floors vanish.

Integrate the symbolic expansion over the compact variables
with probability Haar measure. Theorem~\ref{thm:compact313}
at one group applies separately to each coefficient of the
remaining variables and symbols. Iterating over \(v,\tau\)
shows that \(B_0D_{\rm H}\) clears every coefficient after
averaging. The integers \(A_v\) clear the stabilizer weights in
each tuple of finite classes. The expansion of
\(D_0\mathcal A\) is therefore an integral linear combination
over \(\OO_J\) of products of one derivative at each edge.
The original signs, finite tests and endpoint all occur in this
very expansion. Valuations add in each product, and taking
the sum preserves the common lower bound. This proves both
assertions, including when the signed sum vanishes.
\end{proof}

The compact cost at any subset of primes \(p\le Y\) obeys
\[
 \sum_{p\le Y} v_p(D_{\rm H})\log p
  \le \pi(Y)\sum_{v,\tau}\binom{n_v}{2}\log(n_v+L_{v,\tau})
\]
by Theorem~\ref{thm:compact313}. Thus for
\(Y\le\exp(x^\varepsilon/3)\) this cost is
\(o(\exp(x^\varepsilon/2))\) whenever the sum on the right
after \(\pi(Y)\) is bounded by a fixed power of \(x\), uniformly
over the interval. The inequality itself holds at all degrees.

\subsection{Integral class tensors detecting the incoming row}\label{nav:VI-63}
Keep the last edge, its finite characters, its fine classes,
the measures \(w_\pm,v_\pm\) and its actual incoming space
\(\mathcal U\) as in Theorem~\ref{thm:lastrow311}. For each
rank \(R=n,s\) and each fine class \(u\), take the \(E\)-basis
in \(\Lambda_u\) of Lemma~\ref{lem:terminalclass259}, applied
at that rank. Write \(X_{u,\tau}\) for its column matrix and
\(B_R\ge1\) for the bound of that lemma on its singular values
and inverse singular values at every embedding. Use only its
lattice assertion here; the fine class counts are \(h_\pm\)
of Theorem~\ref{thm:lastrow311}. Write \(\det^{j_\tau}\) for
the determinant twist of the partition model at \(\tau\).
Put
\[
       b_R=B_R^{d+R\sum_\tau |j_\tau|},\qquad
       b_-=b_n,\quad b_+=b_s,
\]
where the partition degree on either side is \(d\).
Use the last-edge normalization \(z(U,V)\), the constants
\(a_*,L_*,i_\pm,D_\pm\) of Theorem~\ref{thm:lastrow311}, and
the archimedean kernel bounds \(B_\sigma\) of
Theorem~\ref{thm:bothlast311}. The fine levels fix the tests
and have trivial stabilizers.

\begin{theorem}[A detected minor with denominators from the original finite tests]\leavevmode
\label{thm:intrinsicminor315}
There are finite lists of tests on these classes such that for
every \(0\ne X\in\mathcal U\) one terminal test \(V\) and one
source test \(U\) satisfy
\[
 |\langle TX,V\rangle|\ge
 \frac{i_+\sqrt{w_+a_*}}{2D_+\sqrt{h_+}\,b_+}\|X\|,
 \qquad |z(U,V)|\ge\frac{i_-i_+ L_*}{b_-b_+}.
\]
For \(1\le j\le\dim\mathcal U\) some \(j\)-rowed minor has
\(\displaystyle |z_j|\ge j!(i_-i_+L_*/(b_-b_+))^j\).
It is algebraic and satisfies, on the very same field,
\[
 \sum_{p\in S}g_p\log p
 \le B^0+\tfrac12\log j+\log(i_-i_+b_-b_+)
       +\tfrac d2\log(ns)+\frac1{[J:\QQ]}\sum_\sigma\log B_\sigma .
\]
Here \(S,g_p\) are as in Theorem~\ref{thm:bothlast311} at good
primes prime to \(T^0Q^0\), and \(B^0\) is the single last-edge
quantity of Theorem~\ref{thm:originalgram315} for the prescribed
finite input and its right-level translates.
\end{theorem}
\begin{proof}
On class \(u\) start with the integral exterior tensor list
\(v_R(Z)\) and its distribution from Theorem~\ref{thm:frames311}.
Apply the exterior tensor map of \(X_{u,\tau}\) on each
partition factor. On the determinant factor multiply by
\((\det X_{u,\tau})^{j_\tau}\) for \(j_\tau\ge0\), and by
\(\overline{\det X_{u,\tau}}^{\,|j_\tau|}\) for \(j_\tau<0\).
For the latter choice, contracting with the wedge of the
lattice basis gives an integral functional on the top
exterior power of \(\Lambda_u\) by Hermitian integrality.
Its coordinates in the dual line give exactly this conjugate
determinant. Thus the resulting vector belongs to the ambient
exterior tensor lattice of the class, including its determinant
lines. Since the original lattices are Hermitian unimodular
locally, their exterior, dual and tensor products carry
induced integral pairings. Pairing with the resulting vector
therefore gives an integral covector on that lattice. Use it in
\(k\mapsto\langle\rho_R(k)t_R,v\rangle\), keeping \(t_R\)
the fixed standard highest vector.

Write \(A_u\) for the map just applied to the test list.
The partition Cartan component is invariant under the
complexified general linear group, as proved in
Theorem~\ref{thm:frames311}. Hence \(A_u\) maps it onto
itself. The singular values on \(\wedge^h\) lie between
\(B_R^{-h}\) and \(B_R^h\), since they are the \(h\)-fold
products of singular values. The determinant scalars satisfy
the corresponding bounds with exponent \(R|j_\tau|\).
Consequently both \(\|A_u\|\) and \(\|A_u^{-1}\|\) on this
component are at most \(b_R\). The same norm bound holds
at all embeddings, by the lattice lemma, including on
the conjugate determinant factors. It follows that
\(\|\sigma v\|\le b_R R^{d/2}\). At the original embedding
the covariance is \(A_u\,{\mathbb E}(v_Rv_R^*)A_u^*\),
bounded below by \(I/(4D_R b_R^2)\).

Take the same uniform distribution on classes and apply the
same unnormalized finite character sums as in
Theorem~\ref{thm:lastrow311}. Compact orthogonality now gives
\(\mathbb E|\langle F,Q\rangle|^2\ge
i_R^2 w_R\|F\|^2/(4h_R D_R^2 b_R^2)\) at either end, where
\(Q\) denotes the resulting random test. Apply it first to
\(TX\), using \(\|TX\|\ge\sqrt{a_*}\|X\|\). For the selected
\(V\), apply it at the source to the Riesz vector of
\(F\mapsto\langle TF,V\rangle\), whose norm is at least
\(|\langle TX,V\rangle|/\|X\|\). This proves the simultaneous
assertions. In the exterior-power calculation of
Theorem~\ref{thm:minors311} the two covariance lower bounds
are now divided by \(b_-^2,b_+^2\). Each of the first \(j\)
singular values of the weighted normalized matrix is therefore
at least \(i_-i_+L_*/(b_-b_+)\). Summing the squared minors
with their probability weights, of total mass at most \(1/(j!)^2\),
proves the stated lower bound.

Move each finite right translation in the two character sums
to the kernel by invariance of Haar measure. This leaves
\(i_-i_+\) terms using the original one-class endpoint
functions, with root-of-unity weights. In the representatives
of Theorem~\ref{thm:neutral315} at \(p\), the transported
class vector in the exterior tensor model is integral in
standard local coordinates by construction and
Theorem~\ref{thm:endpoint315}. This follows also by applying
the rational unitary change to its lattice vectors and
contracted determinant covectors. The right highest vector
has unit highest coordinate. The two dimension factors in
\(z\) cancel the compact matrix-coefficient averages of
Theorem~\ref{thm:last244}. Thus each entry is a sum of
integral minor-polynomial derivatives of the actual finite
sections, with no stabilizer or compact denominator. These
expressions give algebraicity. Their valuations at a good
\(p\in S\) have the full-column lower bound \(g_p\) of
Theorem~\ref{thm:frob244}, since each monomial of a full
minor still uses every column.

At every other prime use the local proof of
Theorem~\ref{thm:originalgram315}. On expanding the two
determinants, all products have \(j\) factors of the given
degree in each orientation, with integral endpoint tensors.
It follows that \(v_{\mathfrak p}(z_j\bar z_j)\) is at least
\(2jg_p\) in \(S\) and at least \(-2j J^0_p\) elsewhere.
For the complex upper bound, the norm estimate on the two
new lists and the linearity of the minor contraction in each
vector multiply the entry bound of
Theorem~\ref{thm:bothlast311} by at most \(b_-b_+\).
Hadamard's inequality then gives
\[
 |\sigma z_j|\le j^{j/2}
       \big(i_-i_+ b_-b_+(ns)^{d/2} B_\sigma\big)^j .
\]
Apply the product formula to this nonzero diagonal Gram element.
The archimedean mean of its logarithm equals twice that for
\(z_j\), since conjugation permutes the embeddings of \(J\).
Division by \(2j\) gives the inequality.
\end{proof}

The tests just constructed probe the same incoming \(X\) in
the first inequality and the same \(V\) in the second. The
arithmetic scalar in the second still pairs \(V\) with \(TU\).
The interval argument requires a comparison with the pairing
against the full prescribed Bessel row and the strict cohomology
bound of Section~\ref{sec:missing}.

\subsection{The orthogonal projector onto a sum of theta images}\label{nav:VI-64}
In the common space \(E_0\) of Theorem~\ref{thm:closedperiod318},
choose any finite collection of the signed definite operators.
Let \(\mathcal R\) be the sum of their images with the original
endpoint metric. No equality of images for distinct words
is required. We give a denominator for projection onto
\(\mathcal R\) using only class integrals already computed on
this field. Keep \(A=\operatorname{adj}(d_0G)\) from that
theorem. Select a maximal independent list of columns
\(y_1,\ldots,y_r\) from the normalized period matrices
\(P_u\) of the collection. Write \(Y=(y_1,\ldots,y_r)\),
and let \(u_j\) be the word of column \(y_j\). If \(r=0\)
the projector is zero. Otherwise put
\[
 D_Y=\prod_j D_{u_j}^{\,2},\qquad
 d_Y=\det(Y^* A Y),\qquad b_{Y,p}=\sum_j b_{u_j,p}.
\]
The different heights and degree profiles of the words
enter through their own \(D_{u_j}\) and \(b_{u_j,p}\).
At an embedding \(\sigma:J\to\CC\) let \(L_{j,\sigma}\)
bound the Euclidean norm of \(\sigma y_j\). One can take
the square root of the sum of the entrywise squared bounds
from Theorem~\ref{thm:archcolumn313}. Let \(c_J\) denote
complex conjugation on \(J\) in the original embedding.

\begin{theorem}[Denominators for simultaneous orthogonal projection]\leavevmode
\label{thm:jointimage318}
The projector onto \(\mathcal R\) has the matrix
\[
 \Pi_{\mathcal R}
     ={A Y\operatorname{adj}(Y^* A Y)Y^*\over d_Y}.
\]
For each \(\mathfrak p\mid p\) the nonnegative number
\(\eta_{Y,\mathfrak p}=v_{\mathfrak p}(D_Yd_Y)-b_{Y,p}\)
bounds the negative valuations of its entries. One may
minimize this number over maximal independent lists at each
place. With \(n_{\mathfrak p}=[J_{\mathfrak p}:\QQ_p]\),
the total cost for each fixed list satisfies
\[
\begin{split}
 {1\over[J:\QQ]}\sum_{p,\mathfrak p\mid p}
         n_{\mathfrak p}\eta_{Y,\mathfrak p}\log p
 \ \le\ &\log D_Y-\sum_p b_{Y,p}\log p\\
 &+{1\over[J:\QQ]}\sum_\sigma\log\!\left(
 \|\sigma A\|_{\rm op}^{\,r}
       \prod_{j=1}^r L_{j,\sigma}L_{j,\sigma\circ c_J}\right).
\end{split}
\]
If another signed period on these endpoint tests has matrix
\(P_v\), the pairing after orthogonal projection of its
output onto \(\mathcal R\) has normalized matrix
\(P_v^{\mathcal R}=\Pi_{\mathcal R}^*P_v\).
Every product of one entry of \(P_v^{\mathcal R}\) with the
conjugate of another, times \(D_v^2\), has valuation at
least \(b_{v,p}-2\eta_{Y,\mathfrak p}\).
\end{theorem}
\begin{proof}
The physical outputs span \(G^{-1}\operatorname{im}Y\).
This equals \(\operatorname{im}(AY)\), since \(A\) is a
nonzero scalar multiple of \(G^{-1}\). The kernel of \(Y^*\)
is the metric orthogonal complement of this image. The matrix
\(A\) is positive definite at the original embedding, so
\(d_Y>0\). The displayed projector is the identity on
\(\operatorname{im}(AY)\) and zero on \(\ker Y^*\).
These two subspaces exhaust \(E_0\). In particular the
formula is unchanged if a different maximal list is used.

Expand the determinant into permutations and expand each
entry of \(Y^*AY\) in the coordinate basis of \(E_0\).
In each summand, column \(j\) appears once in a factor
\((y_j)_a\) and once in a conjugate factor
\(\overline{(y_j)_b}\). The same count holds in each entry
of the numerator. Indeed a term with the outside columns
\(y_i,y_j^*\) uses the cofactor obtained by deleting row
\(j\) and column \(i\) of \(Y^*AY\). The two outside
factors supply exactly the two deleted occurrences.
All the remaining coefficients are integers of \(J\).
The period estimate for each pair from column \(j\) then
gives at least \(b_{u_j,p}\) after clearing \(D_{u_j}^2\),
by Theorem~\ref{thm:originalgram315} with the improvement
of Theorem~\ref{thm:poscolumn316}. This proves the same
lower bound \(b_{Y,p}\) for \(D_Yd_Y\) and every entry of
the cleared numerator. Division gives the asserted
denominator for the projector.

The product formula for \(D_Y d_Y\) makes the left side
of the total cost estimate exactly
\[
 \log D_Y-\sum_p b_{Y,p}\log p+
                  [J:\QQ]^{-1}\sum_\sigma\log|\sigma d_Y|.
\]
At \(\sigma\), the determinant is the contraction of the
wedge of the columns \((\sigma\circ c_J)y_j\), transposed,
with \((\bigwedge^r\sigma A)(\bigwedge_j\sigma y_j)\).
The Euclidean norm of a simple wedge is at most the product
of the vector norms, by Gram--Schmidt. The norm of the
induced operator is at most \(\|\sigma A\|_{\rm op}^r\),
by singular value decomposition. Cauchy--Schwarz proves
the bound even when \(\sigma\) does not commute with \(c_J\).

Finally \(G\Pi_{\mathcal R}=\Pi_{\mathcal R}^*G\), which
gives the projected matrix formula in the period pairing.
Complex conjugation takes a place over \(p\) to another
place over \(p\). Since \(D_Yd_Y\) is fixed by \(c_J\),
the two places have the same \(\eta_{Y,\mathfrak p}\).
Expand a product of two projected entries and apply the
original pair bound to the two entries of \(P_v\) and
the projector bound to the two other factors. This gives
the stated loss. All operations took place at the original
endpoint and retained its metric and prescribed levels.
\end{proof}

For the final definite step one can take a full set of
independent columns of its period matrix. Its image contains
the \(k_0\)-dimensional image of the incoming subspace in
Theorem~\ref{thm:lastrow311}. Thus this projector has rank at
least \(k_0\). For the \(z(U,V)\) normalization one obtains
the same projector by using the entries of \(z\) in the
formula, since a common scalar on an entire matrix leaves
its image unchanged. The denominator estimates above use
the \(D_u,b_{u,p}\) belonging to the stated normalization.
Identifying one of these computed images with the row in
the Bessel pairing requires a common realization of the two
row spaces.

\subsection{Denominators from the finite class data}\label{nav:VIII-61}
Fix one \(x\) and the levels, representatives and right translates
of Theorem~\ref{thm:bothlast311}. Let \(d=\sum_\tau|\lambda_\tau|\)
be the differential degree of its last-edge minor polynomial and
\(D=fns\) the number of matrix coordinates \(Z_\mu\).
Use the marked \(K\)-orthonormal coordinates of
Lemma~\ref{lem:basis59}. Every one of the finitely many class sections
in the arithmetic matrix admits an expression
\[
 G(Z)=\sum_b c_b\prod_{h=1}^D
       (U_{w_{bh}}\Theta_{t_*})\big((T_bZ)_h\big),
       \qquad P_{bh}=[w_{bh}]\in C_{\mathrm{tors}} .
\]
Here \(T_b\) converts the matrix coordinates to the chosen product
coordinates. One may use \(T_b=V^*\) for a common orthonormal basis.
The coefficients \(c_b\) include all finite scalar factors from the
class construction. We describe below how to take this presentation
directly from that construction.

Choose positive integers \(N,H,m\) such that every \(Nc_b\) and
every entry of \(HT_b\) is integral, and \(mP_{bh}=0\) for every
term of every section. Use one coefficient field \(J\) for all of
them. Write \(\Phi_p(a,t)=at+\lfloor a/2\rfloor(\kappa_p-2t)_+\),
with \(\kappa_p\) as in Theorem~\ref{thm:elliptic312}.

\begin{theorem}[Arithmetic bounds for the specified class sections]\leavevmode
\label{thm:classden312}
For every matrix multi-index \(k\) with \(|k|=d\) and every
\(\mathfrak p\mid p\) one has
\[
 v_{\mathfrak p}(\partial^k G(0))
 \ge -v_{\mathfrak p}(N)-d\,v_{\mathfrak p}(H)
                      -\Phi_p(d,v_{\mathfrak p}(m)).
\]
There is a sharper finite calculation for any prescribed set of
multi-indices. For each contributing term choose nonnegative
integers \(n_{h\mu}\) with \(\sum_h n_{h\mu}=k_\mu\), requiring
\((T_b)_{h\mu}\ne0\) whenever \(n_{h\mu}>0\), and put
\(s_h=\sum_\mu n_{h\mu}\). Define its loss by
\[
 \sum_h\Phi_p(s_h,\delta_{\mathfrak p}(P_{bh}))
   -\!\sum_{n_{h\mu}>0} n_{h\mu}v_{\mathfrak p}((T_b)_{h\mu})
   -v_{\mathfrak p}(c_b).
\]
Take the maximum of zero and all such losses over the relevant
terms, classes and multi-indices. The resulting number
\(\Gamma_{\mathfrak p}\) bounds the negative valuation of every
integer linear combination of these derivatives.
\end{theorem}
\begin{proof}
Start with the product vacuum in Lemma~\ref{lem:basis59}. A normalized
intertwiner to a commensurable period lattice is the finite sum of
translations indexed by its quotient by the intersection, weighted
by signs or roots of unity and by the inverse square root of the
quotient order. The formula and its normalization are proved in
Corollary~\ref{cor:move55}. The class and level transports in
Theorem~\ref{thm:one61} move the vacuum to such a lattice and the
Schwartz test to finitely many further translates. For a smaller
split indicator the finite Fourier sum in
Theorem~\ref{thm:bothlast311} gives the additional coefficients,
including its reciprocal group order. Compose the translations
and multiply their coefficients. Their composition introduces a
root of unity since the alternating form is rational on rational
lattice vectors. A product Hermitian translation factors over its
elliptic coordinates. This gives precisely the displayed expression
for each class term. It retains the local splitting characters in
\(c_b\), including those supplied by Theorem~\ref{thm:one61}.

The integers in the statement can thus be read from finite data.
For example, take \(t\) so that each unscaled period module in
product coordinates lies in \(t^{-1}\OO_K^D\). Each lattice
intertwiner translates by a representative in such a module.
If all further prescribed translations have order dividing \(Q\)
relative to their period modules, their product coordinates
lie in \((Qt)^{-1}\OO_K^D\) before multiplying by the elliptic
period \(\delta\). Consequently \(m=Qt\) works even after
composing the translations. The coordinate change supplies \(H\).
For a term passing through quotient orders \(q_1,\ldots,q_a\),
the product \(\prod_i q_i\) clears the inverse square roots, since
each \(\sqrt{q_i}\) is an algebraic integer. Clearing also the
actual finite coefficients and character values and taking a common
multiple over the terms gives \(N\). Zero coefficients may be
discarded. None of these choices requires evaluating a theta
derivative.

Apply each matrix derivative as
\(\partial_\mu=\sum_h(T_b)_{h\mu}\partial_h\).
Expanding the commuting derivations gives exactly the indicated
arrays \(n_{h\mu}\), with integral multinomial coefficients.
The derivative in product coordinate \(h\) occurs \(s_h\) times
and acts only on its own factor. Theorem~\ref{thm:torsion312}
therefore bounds the negative valuation of the term by its
displayed loss. Sums and integral linear combinations obey the
same maximum bound.

For the first inequality, \(\delta_{\mathfrak p}(P_{bh})\le
v_{\mathfrak p}(m)\). For every \(a\) the function
\(\Phi_p(a,t)\) is increasing in the weak sense, as is seen
from its maximum formula in Theorem~\ref{thm:torsion312}.
Also \(\sum_h\Phi_p(s_h,t)\le\Phi_p(d,t)\) since
\(\sum_hs_h=d\) and
\(\sum_h\lfloor s_h/2\rfloor\le\lfloor d/2\rfloor\).
The two remaining losses are bounded by
\(d\,v_{\mathfrak p}(H)\) and \(v_{\mathfrak p}(N)\).
\end{proof}

\subsection{A trace lattice basis at an arbitrary prime}\label{nav:VIII-62}
On \(W=\operatorname{Mat}_{s\times r}(E)\) use the positive
trace Hermitian form of Lemma~\ref{lem:basis59}, with \(rs\)
even. Its standard trace module is
\(M=\mathfrak b^{-1}\operatorname{Mat}_{s\times r}(\OO_E)\),
where \(\mathfrak b\bar{\mathfrak b}=\mathfrak D_{E/K}\) by
Theorem~\ref{thm:pol49}. Use a finite list of the actual tests
on this edge. In the finite model of Theorem~\ref{thm:vac242}
write each in the basis of the original unit vacuum as
\[
           \phi=\sum_a c_{\phi,a} W(a)\phi_0 .
\]
Here \(a\) represents a rational class modulo \(M\).
Each sum is finite by smoothness as proved in that theorem.
The finite transports and splitting factors in
Theorem~\ref{thm:one61} and the finite Fourier expansions in
Theorem~\ref{thm:bothlast311} have algebraic coefficients.
Thus for the stipulated algebraic tests we can choose positive
integers \(Q,N\) with \(Qa\in M\) and all \(Nc_{\phi,a}\)
algebraic integers. These data are in the fixed finite model.

\begin{theorem}[Local product coordinates for the fixed finite model]\leavevmode
\label{thm:adapted314}
For every rational prime \(p\) there is a \(K\)-orthonormal basis
\(B_p\) of \(W\) generating \(M\otimes\ZZ_p\) over
\(\OO_K\otimes\ZZ_p\). In physical matrix coordinates take an
ordinary derivative of degree \(d_\tau\) in block \(\tau\) and
write \(d=\sum_\tau d_\tau\). At every \(\mathfrak p\mid p\),
in this basis its class section satisfies
\[
 v_{\mathfrak p}(\partial^\alpha G^{B_p}_\phi(0))
 \ge -v_p(N)-\Phi_p(d,v_p(Q))
           -\sum_\tau d_\tau v_{\mathfrak p}(\bar\tau\mathfrak b).
\]
For an ideal its valuation denotes the minimum valuation of
its elements after extension to the coefficient field, with
\(v_{\mathfrak p}(p)=1\). Conjugating an edge conjugates the
ideal in the last sum.
\end{theorem}
\begin{proof}
The module \(M\) is self-dual since the trace dual of
\(\OO_E\) is \(\mathfrak D_{E/K}^{-1}\) and
\(\mathfrak b\bar{\mathfrak b}=\mathfrak D_{E/K}\). It is
locally free over \(\OO_K\). If \(p\) splits in \(K\), the
two components of the local module are in perfect duality.
Choose a basis of the first and its dual in the second. The
sums of corresponding vectors form an orthonormal basis over
the product ring. This covers \(p=2\), since \(d_0\equiv7\bmod8\).

At a nonsplit prime \(p\) the residue characteristic is odd.
A local integral basis has a Hermitian Gram matrix with unit
determinant. At an unramified prime its reduction is equivalent
to the identity over the quadratic residue field, by orthogonal
diagonalization and surjectivity of the finite-field norm.
For the diagonalization, a nondegenerate Hermitian form has a
vector of nonzero length by polarization. Split off its line
and normalize it by the norm map, then continue on the
orthogonal complement. At a ramified prime the reduced form
is symmetric. Its determinant is a square in the residue field.
Indeed in a rational orthonormal basis, available by
Lemma~\ref{lem:basis59}, the Gram determinant is the norm of
the determinant of the local change of basis. This determinant
is a local unit since its norm is a unit. Reduction of the
norm is a square at a ramified prime. A nondegenerate symmetric
form over a finite field of odd order with square determinant
is equivalent to the identity. To see this, diagonalize and
pair the nonsquare entries after square rescalings make each
one the same nonsquare \(a\). There exist \(u,v\) with
\(u^2+v^2=a^{-1}\), by surjectivity of the norm in the finite
quadratic algebra obtained by adjoining a root of \(-1\),
including the split product algebra. The two orthogonal columns
\((u,v),(-v,u)\) turn \(aI_2\) into \(I_2\).

After lifting the residue change of basis the integral Gram
matrix \(G\) is the identity modulo the local maximal ideal.
The series \(G^{-1/2}=\sum_{n\ge0}\binom{-1/2}{n}(G-I)^n\)
converges integrally, because
\(\binom{-1/2}{n}=(-1)^n\binom{2n}{n}/4^n\) and \(p\) is odd.
It is self-adjoint and its square is \(G^{-1}\), giving a local
orthonormal basis. To obtain a rational one, apply the Cayley
approximation in the proof of Lemma~\ref{lem:basis63} at this
prime. The condition to generate the local lattice is open.
A norm-one scalar arbitrarily close to one moves any local
unitary matrix off the eigenvalue \(-1\). The inverse Cayley
transform then lies in the rational vector space of skew-Hermitian
matrices over \(\QQ_p\), where rational matrices in that vector
space are dense. Transform a sufficiently close approximation
back. This also works for the split product algebra.

Let \(M'\) be the \(\OO_K\)-span of the resulting basis.
The index \([M:M\cap M']\) is prime to \(p\) since the local
modules agree. In the normalized intertwiner from the product
lattice to \(M\) the inverse square root of this index is a
unit at every place above \(p\). Its translation vectors belong
to \(M\) and have order modulo \(M'\) dividing this same index.
A further translation by \(a\), with \(Qa\in M\), consequently
has order with \(p\)-part dividing the \(p\)-part of \(Q\) in
the product coordinates. Composing translations contributes
only roots of unity. The translated elliptic product in
Theorem~\ref{thm:torsion312} thus costs at most
\(\Phi_p(d,v_p(Q))\) in ordinary derivatives, by the product
rule of Theorem~\ref{thm:classden312}.

Finally every basis vector has its matrix entries locally in
\(\mathfrak b^{-1}\) at every prime of \(E\) above \(p\), by
the definition of \(M\). Each coefficient converting a physical
derivative in block \(\tau\) to a product derivative is a conjugate
entry of that embedded basis matrix. Its valuation is therefore
at least \(-v_{\mathfrak p}(\bar\tau\mathfrak b)\). Multiply these
bounds in the chain rule and use \(v_{\mathfrak p}(c_{\phi,a})
\ge-v_p(N)\). This proves the asserted inequality simultaneously
for every finite test on the list.
\end{proof}

\subsection{Lattices for the two slots of an endpoint coefficient}\label{nav:VIII-63}
At an endpoint of rank \(R\) write its class functions as sums
of \(f(k,u)=\ell_u(\rho(k)t)\). Use a number field containing
the specified algebraic coefficients and all the embeddings
in the tensor model. If \(L\) is a Hermitian \(\OO_E\)-lattice,
write \(W(L)\) for the prescribed integral coefficient module,
formed from its embedded copies by tensor products, exterior
powers, integral Young images and determinant lines, as in
Proposition~\ref{prop:terminalcore255}. For a negative
determinant power use the corresponding power of the dual
determinant line. Localizations and duals below are over
the coefficient integers, after extending the embedded lattices
there. Write \(\Lambda_u=E^R\cap u\widehat{\OO}_E^R\).

\begin{theorem}[A denominator for the endpoint independent of the representatives]\leavevmode
\label{thm:endpoint315}
Use \(W(\OO_E^R)\) for \(t\), and \(W(\Lambda_u)^\vee\) for
\(\ell_u\). For the algebraic ends of the signed period in
Theorem~\ref{thm:column242}, choose one integer \(B_{\rm end}\)
clearing the tensor products of these endpoint coefficients,
the prescribed real polynomial coefficients, the signed weights,
and any scalar of the projection kept outside the finite tests.
Then \(B_{\rm end}\) clears the coefficients of the formal
compact integrand over \(p\), in the representatives of
Theorem~\ref{thm:neutral315}, simultaneously for every term
and class. Here the physical derivatives of the translated finite
sections are formal indeterminates.
\end{theorem}
\begin{proof}
The lattices have full rank in their respective algebraic
representations. For the finite list of original classes one
may clear each coefficient in its indicated lattice and take
products for every monomial in the integrand, then take a common
multiple. This supplies the required integer before varying \(p\).
It also treats any fixed arithmetic vector in the image of the
terminal projection. If the endpoint pairing is written with an
explicit projection on its test vector, the coefficients used
here are those of the resulting vector.

For \(u'=\gamma u j\) as in Theorem~\ref{thm:neutral315},
preservation of the standard finite lattice by \(j\) gives
\(\Lambda_{u'}=\gamma\Lambda_u\). Since \(u'_p=1\), its
localization at \(p\) is the standard lattice. Automorphy
expresses the endpoint coefficient in the new chart as
\[
      f(k,u')=(\ell_u\circ\rho(\gamma_\infty^{-1}))(\rho(k)t).
\]
Indeed multiplying \((k,u')\) on the left by \(\gamma^{-1}\)
and on the right by \(j^{-1}\) in its finite component gives
\((\gamma_\infty^{-1}k,u)\). Thus \(t\) stays fixed, and
the indicated covector is integral on \(W(\Lambda_{u'})\)
whenever \(\ell_u\) was integral on \(W(\Lambda_u)\).
The formula is compatible with complex conjugation of a factor,
since integrality is imposed at all places above \(p\).

In the new chart, after the stipulated clearing, both slots are
integral for the standard local module. The polynomial giving
their matrix coefficient in the entries of \(k,\bar k\) has
integral coefficients. For an exterior power this follows from
the formula by minors. For an integral Young image write a
lattice vector as a sum of \(Yx\) for the integral Young
operator \(Y\). The variable matrix acts by
\(Y(k^{\otimes d}x)\); its coefficients belong to the Young
image because \(Y\) commutes with tensor powers of \(k\).
Applying an integral covector preserves coefficient integrality.
On \(U(R)\) inverse determinant is \(\det\bar k\), and the
dual standard action also uses conjugate entries. This proves
the polynomial assertion for every stated tensor construction.

The same polynomial substitutions in the physical coordinates
act on the incident edge symbols. Their coefficients consist of
the fixed real test coefficients and integral matrix operations.
The external finite order determinant values are roots of unity.
Any moved oscillator splitting scalars are already in the finite
section and are units at \(p\) by Theorem~\ref{thm:neutral315}.
Taking products, then the prescribed signed sums, proves the
claim with the single integer chosen at the outset.
\end{proof}

\subsection{Integral reverse tests in the last theta images}\label{nav:VIII-64}
Keep a finite list of matching complex iterations with the
same definite endpoint, highest line, finite level and character.
At each last definite step let \(T_u\) act on the full
penultimate coefficient space at its stated type and level.
Let \(\mathcal R=\sum_u\operatorname{im}T_u\) in the
common endpoint. Its projector \(\Pi\) is given by
Theorem~\ref{thm:jointimage318}. One can use the full
signed period normalization of its one-edge matrices
there, or their nonzero common scalar multiples for
computing the same image. Use the former normalization
for its denominator bounds. For an actual incoming
vector \(X_u\) of Theorem~\ref{thm:lastrow311} put
\(y_u=T_uX_u\).

\begin{theorem}[Retaining the actual Bessel values on an integral image]\leavevmode
\label{thm:imagereverse320}
For each endpoint test \(Q\) one has simultaneously
\(\langle y_u,\Pi Q\rangle=\langle y_u,Q\rangle\) for
all \(u\). In particular, given finitely many nonzero
incoming vectors and their respective integral detecting
tests \(Q_u\) from Theorem~\ref{thm:intrinsicminor315},
one can choose a common integral \(Q'\in\mathcal R\) and
an explicit positive integer \(m\) such that
\[
       |\langle y_u,Q'\rangle|/m
                 \ \ge\ \tfrac12|\langle y_u,Q_u\rangle|
                   \qquad\text{for every }u .
\]
Before the integer clearing, projection can only improve
the ratio of the absolute pairing to the norm of the common
test. Also, for a matching complex iteration
\(\Theta:\mathcal H\to E_0\) that factors through one of
these last operators on its prescribed row and satisfies
\(\Theta^\dagger\Theta\ge aI\), \(a>0\), the adjoint
restricted to \(\mathcal R\) maps onto \(\mathcal H\)
with a right inverse of norm at most \(a^{-1/2}\).
\end{theorem}
\begin{proof}
Every \(y_u\) belongs to \(\mathcal R\), hence is orthogonal
to \((I-\Pi)Q\). For \(l\) input vectors form
\(Q=\sum_{j=1}^l n_j Q_j\) with \(0\le n_j\le l\) integers.
For vector \(y_u\), fix the choices other than \(n_u\).
Its possible pairings have spacing
\(|\langle y_u,Q_u\rangle|>0\). At most one lies in the
open disk about zero of half that radius. The union of
the failing sets has fraction at most \(l/(l+1)<1\).
Choose a successful \(Q\). Then \(\Pi Q\ne0\), all its
pairings agree with those for \(Q\), and \(\|\Pi Q\|\le\|Q\|\).

Here is the integer clearing in the stated endpoint lattice.
Choose a basis from its integral class tests, let \(D_Q\)
be a positive integer clearing the coordinates of \(Q\)
in that basis, and use these coordinates for the projector
in Theorem~\ref{thm:jointimage318}. At \(\mathfrak p\mid p\)
put \(l_{\mathfrak p}=\max(0,-\min_{a,b}v_{\mathfrak p}
(\Pi_{ab}))\), with zero entries omitted. Write
\[
   N=\prod_p p^{\lceil\max_{\mathfrak p\mid p}
                                      l_{\mathfrak p}\rceil},
              \qquad m=D_QN .
\]
Only finitely many factors differ from one. For any of
the admissible determinant lists \(Y\) in that theorem,
\[
 \log N\le\sum_{p,\mathfrak p\mid p}
         [J_{\mathfrak p}:\QQ_p]\eta_{Y,\mathfrak p}\log p .
\]
Indeed \(l_{\mathfrak p}\le\eta_{Y,\mathfrak p}\) and
\(e(\mathfrak p/p)l_{\mathfrak p}\) is an integer. For a
place where the maximum is positive this integer is at
least its ceiling before multiplication by the ramification
index. This proves the bound prime by prime.
Now \(Q'=m\Pi Q\) lies in the integral span of the basis
and thus in the given lattice. Dividing its pairings by
the recorded \(m\) proves the assertion with the original
normalization.

Finally the range of the stated \(\Theta\) lies in
\(\mathcal R\). The vector
\(\Theta(\Theta^\dagger\Theta)^{-1}z\) lies there for every
\(z\in\mathcal H\), maps back to \(z\) under the adjoint,
and has squared norm
\(\langle z,(\Theta^\dagger\Theta)^{-1}z\rangle
 \le a^{-1}\|z\|^2\). On the matching data one can take
the span of the actual ordinary columns and the corresponding
bound of Theorems~\ref{thm:range239} and~\ref{thm:norm177}.
The Bessel values on the reverse tests are the forward
pairings by Proposition~\ref{prop:adjtheta234}.
\end{proof}

This gives integral tests inside the sum of the computed
last images, with a common endpoint for the finitely many
actual incoming vectors. Identifying that sum, or its
reverse image, with the independently specified row and
rational pairing in Section~\ref{sec:missing} remains an
additional assertion.

\subsection{Integral class tests with the original Haar measure}\label{nav:I-81}
Use the definite terminal row of Theorem~\ref{thm:terminalreading325}.
Write \(J=J_\xi\), \(\omega=\omega_\xi\), and fix the
coefficient ring of integers \(\mathcal O\), enlarged to
contain the values of the finite character. At class \(i\)
let \(L_i\) be the coefficient lattice of the Hermitian
lattice in that class, with the given determinant twists.
It can be taken to be the lattice of integral flag minor
products in its local frames. These lattices are projective
over \(\mathcal O\) and the class stabilizer preserves them.
Indeed formation of flag minors commutes with integral
basis change by expansion of minors, in both directions.
The local lattices agree with a fixed one almost everywhere
and define a global lattice by intersection. Put
\(\omega_i(\gamma)=\omega(g_i^{-1}\gamma g_i)\), and let
\(S_i\subset L_i\) be the submodule where
\(\rho(\gamma)y=\omega_i(\gamma)y\) for all \(\gamma\in
\Gamma_i\). The full integral class table is \(S=\bigoplus_iS_i\).
Use the inverse character and the algebraic dual representation
on testing forms.

\begin{theorem}[Integral class traces and change of level]\leavevmode
\label{thm:classtrace326}
For \(\ell\in L_i^\vee=\operatorname{Hom}_{\mathcal O}(L_i,\mathcal O)\)
let the test be supported on class \(i\) with value
\[
       t_i(\ell)=\sum_{\gamma\in\Gamma_i}
                        \omega_i(\gamma)\,\rho^\vee(\gamma)\ell .
\]
The lattice \(T_J\) of all sums of these tests is exactly
the full dual of \(S\) for the global bilinear integral
divided by \(\operatorname{vol}(J)\). In particular the
period with \(t_i(\ell)\), in the original Haar measure,
is \(\operatorname{vol}(J)\ell(y_i)\) on any table \(y\).
All these integral covectors can be realized locally in the
integral flag minor tests of Theorem~\ref{thm:terminaldual325}.

Let \(K\triangleleft J\) be a finite-index open subgroup
in the kernel of \(\omega\), including the common refinement
fixing the residue markings. For a dual testing form
\(z\) of level \(K\) put
\(\mathcal Tz=\sum_{a\in J/K}\omega(a)R_a z\),
where \(R_a z(g)=z(ga)\). Then
\[
  \mathcal T(T_K)=T_J,\qquad
  \frac{\langle y,\mathcal Tz\rangle}{\operatorname{vol}(J)}
       =\frac{\langle y,z\rangle}{\operatorname{vol}(K)}
\]
for any \(J\)-table \(y\) of character \(\omega\). The
pairings are the integrals of the natural vector and
covector contraction, with probability Haar at infinity.
\end{theorem}
\begin{proof}
The stabilizer acts on \(t_i(\ell)\) by the inverse character,
so the table gives an integral dual form. At \(y_i\in S_i\)
each term of its evaluation equals \(\ell(y_i)\). The class
has measure \(\operatorname{vol}(J)/|\Gamma_i|\) for this
invariant contraction, by Theorem~\ref{thm:terminalreading325}.
This proves the period formula. The submodule \(S_i\) is
saturated in \(L_i\), being the intersection of kernels
of integral linear maps to a torsion-free module. Thus
\(L_i/S_i\) is a finite projective module over the Dedekind
ring \(\mathcal O\), and every integral functional on
\(S_i\) extends to \(L_i\). The trace tests realize all
such functionals. The pairing is nondegenerate over the
coefficient field, since averaging there projects a dual
space onto the full dual of the corresponding character
space. This proves the equality of lattices. In a local
frame an integral covector on \(L_i\) is represented by
an element of \(L_i^\#\) under the normalized Fischer
pairing. Theorem~\ref{thm:terminaldual325} puts that
element in the integral span of flag minors.

Right invariance of Haar and \(y(ga^{-1})=\omega(a)^{-1}y(g)\)
give \(\langle y,\mathcal Tz\rangle=[J:K]\langle y,z\rangle\).
Since \(\operatorname{vol}(J)=[J:K]\operatorname{vol}(K)\),
the normalized identity follows. It puts \(\mathcal T(T_K)\)
inside the full dual just computed, because restriction to
level \(K\) preserves integrality of \(y\).

To prove surjectivity, keep the fine class represented by
\(g_i\) above a chosen coarse class. Put on it the fine
trace of \(\ell\) using \(\Gamma'_i=G(F)\cap g_iKg_i^{-1}\),
and zero elsewhere. In evaluating \(\mathcal T\) at \(g_i\),
the contributing \(a\)'s are precisely the image of
\(\Gamma_i\) in \(J/K\), with kernel \(\Gamma'_i\). For
a lift \(\gamma_a\) of such an \(a\), the value at \(g_i a\)
is the dual action of \(\gamma_a\) on the fine trace.
Multiplication by \(\omega(a)\) and summation now give
the displayed trace over \(\Gamma_i\), once per element.
All other coarse classes have zero value. This realizes
every \(t_i(\ell)\) and proves equality globally over
\(\mathcal O\), including at primes dividing either index.
\end{proof}

\section{The signed finite tests along the column}\label{nav:g-XII-2}

\subsection{The position lattice of a split finite test}\label{nav:V-62}
Let \(p\) be an odd rational prime split in \(K\) on an original
interval. Write \(F=E^+\). At each \(v\mid p\) the algebra
\(E\otimes_F F_v\) is split because \(E=FK\). Fix the actual
position half \(w_v\) and write \(\bar w_v\) for the other half.
On an edge of matrix dimension \(m=rs\) the local reference
lattice of Theorem~\ref{thm:neutral315} is
\(L_v\oplus L_v^\circ\), where \(L_v\) is in position space
and \(L_v^\circ\) is its frequency annihilator. A uniformizer
of \(F_v\) is denoted by \(\varpi_v\), and its residue field has
size \(q_v=p^{f_v}\). The position vacuum of norm one has
amplitude \(a_{0,v}\) on \(L_v\).

Expand the finite list of original edge tests, including all
the right translates that occur in the projection, as
\[
 \phi^0=\sum_{\lambda,\mathbf y} c_{\lambda,\mathbf y}
 W(b_\lambda^{(p)})\phi_0^{(p)}
       \otimes\bigotimes_{v\mid p}a_{0,v}{\bf1}_{y_v+P_v},
 \qquad P_v=\varpi_v^{h_v}L_v,\quad h_v\in\ZZ .
\]
The superscript \((p)\) denotes the components away from \(p\).
The sum exists by finite vacuum expansion there and local
constancy and compact support in the position variables. Use
distinct Weyl classes away from \(p\), with a fixed lift for
each, distinct cosets at \(p\), and combined coefficients.
The split Fourier formulas in Lemma~\ref{lem:split61} and
Theorem~\ref{thm:splitinput315} give algebraic relative
amplitudes for the stipulated algebraic finite tests. Define
\[
   d_{\lambda,\mathbf y}=c_{\lambda,\mathbf y}
                \prod_{v\mid p}q_v^{-m h_v/2}.
\]

\begin{theorem}[A product basis for the actual position supports]\leavevmode
\label{thm:posframe316}
There is a single rational orthonormal product basis for the
whole list of tests with local trace lattice
\(\bigoplus_{v\mid p}(P_v\oplus P_v^\circ)\).
In its product vacuum expansion the coefficient of each term,
before summing terms, is \(d_{\lambda,\mathbf y}\) times a
unit over \(p\), including after moving the vertex representatives
as in Theorem~\ref{thm:neutral315}. The local translation class
of this term at each \(v\mid p\) is \((y_v\bmod P_v,0)\) in
position and frequency coordinates. The physical fields,
coordinates and period normalizations are those of
Theorem~\ref{thm:column242}.
\end{theorem}
\begin{proof}
Let \(\mathfrak c\) be the fractional \(\OO_E\)-ideal with
exponents \(h_v,-h_v\) at \(w_v,\bar w_v\) and zero exponents
elsewhere. Its product with its conjugate is the unit ideal.
Hence \(\mathfrak c M\) is self-dual for the trace Hermitian form,
since \(M\) is self-dual. Its \(p\)-part is the lattice displayed
in the assertion. Indeed the dual of \(\varpi_v^{h_v}L_v\)
in frequency space is \(\varpi_v^{-h_v}L_v^\circ\).
The two \(K_p\)-idempotent components are in perfect duality.
Choosing dual bases and taking their pairwise sums gives a
local orthonormal basis. The rational Cayley approximation in
Theorem~\ref{thm:adapted314} gives a rational orthonormal
basis spanning this same lattice at \(p\), since generation
of the local lattice is an open condition.

In the split position model the lattice vacuum at \(v\) for
\(P_v\oplus P_v^\circ\) is constant on \(P_v\) and supported
there. These assertions follow respectively by invariance under
translations by \(P_v\) and modulations by \(P_v^\circ\). At
this odd prime the order-two characteristic on the product
period lattice is trivial on its \(p\)-part. Lemma~\ref{lem:split61}
fixes the relative normalization of the two unit vacua by their
nonzero positive inner product. In particular, writing
\(\varphi_{P_v}\) for the new unit position vacuum, one has in
this normalization
\[
       a_{0,v}{\bf1}_{y_v+P_v}
          = q_v^{-m h_v/2}\, W(y_v,0)\varphi_{P_v}.
\]
Indeed the scalar for \(y_v=0\) is the square root of the
volume of \(P_v\) divided by the square root of the volume
of \(L_v\). Translation gives the formula for every coset.
This identity holds for negative as well as positive \(h_v\).

To see that these are also the normalizations in the global
product expansion, use the normalized lattice intertwiner taking
the product vacuum to the original one. Its quotient of period
lattices is the product of the local quotients. Its factors are
the respective finite Weyl sums divided by the positive square
roots of their indices. The coset and cocycle description in
Theorem~\ref{thm:one61} identifies its \(p\)-factors with the
intertwiners of Lemma~\ref{lem:split61}. A different lift of a
Weyl class only introduces a root of unity. At a prime away
from \(p\) each index is a power of that other prime, so all
coefficients introduced there are units over \(p\). The same
holds for the moving split components and their characters by
Theorem~\ref{thm:neutral315}. They leave all \(p\)-components
fixed. Passing from the adelic shift to a rational representative
preserves its class in the new local lattice. This proves both
the coefficient and translation assertions with the marked
product realization. The change of rational basis affects the
signed entries on this directed edge by the common scalar of
Theorem~\ref{thm:framechange314}.
\end{proof}

\subsection{Traces of products of signed theta periods}\label{nav:V-63}
Fix one field and one value of \(x\) on the intervals of
Section~\ref{sec:missing}. Work at a fixed common definite
endpoint, with its prescribed subgroup, levels and tests.
Let \(E_0\) be a finite-dimensional endpoint space spanned
by algebraic class tests, and \(M_u\) a source space of this
kind for each word \(u\). Take bases from the given finite
spanning lists. A common endpoint can include a finite sum
of types by taking their actual span in the same space of
functions. The inner product on that span is the original
Petersson inner product. Pairings below are written \(y^*P x\),
where a star means conjugate transpose. For a bilinear
pairing use the conjugate of its row space. Apply this
convention to the entire common row.

Write \(w_0G\) and \(w_u H_u\) for the endpoint and source
metric matrices. Here \(w\) is the Haar volume of the
archimedean group times the common finite integration level
on the corresponding side. There are positive integers
\(d_0,d_u\) such that \(d_0G,d_uH_u\) have entries in the
integers of a common conjugation-stable Galois field \(J\).
Indeed the integral of the product of two class coefficients
is a finite sum of integrals with compact Haar probability
divided by the orders of the chart stabilizers, times \(w\).
Intersecting the finite levels and expanding the two signed
character sums gives this same formula on the projected
space. A class change acts on the compact coefficient by a
rational group element with algebraic matrix entries. Each
remaining compact integrand is polynomial in matrix entries
and their conjugates. This includes an inverse determinant
on a unitary group, where it equals the conjugate determinant.
Clear the finitely many algebraic polynomial coefficients by
a rational integer. Multiply by the least common multiple
of the stabilizer orders and by the product over the real
factors of the moment denominators of
Theorem~\ref{thm:compact313}, using the largest degrees
among these products of two coefficients. This gives an
admissible \(d\) for the chosen basis. Positive definiteness
follows from independence as functions and positivity of
the measure on every chart.

In these bases let \(P_u\) have entries \(\mathcal A_{ba}\)
of Theorem~\ref{thm:column242}, retaining the finite signed
sum and its restriction to \(U(Y_m)\). The physical pairing
is \(\kappa_u y^*P_u x\), with
\(\kappa_u=V_{J,u}(\prod_i C_i)c^{-t_u/2}\) for its fixed
degree profile. A word with no surviving profile gives the
zero matrix. Use positive-dimensional spaces below. A zero
source contributes the zero operator. Put
\[
 A=\operatorname{adj}(d_0G),\qquad
 A_u=\operatorname{adj}(d_u H_u),\qquad
 K_u=A P_u A_u P_u^* .
\]
All adjugates of one by one matrices are \(1\). For the
original finite tests write \(D_u=D_{{\rm intr},u}\) using
the maximum coefficient denominators and degree bounds
across the chosen lists. Write \(b_{u,p}=2\sum_i g_{i,p}\)
at a common good prime in \(S\), and
\(b_{u,p}=-2\sum_i J_{i,p}\) elsewhere, with the sharpened
loss of Theorem~\ref{thm:poscolumn316} at the selected split
primes. All valuations have \(v_{\mathfrak p}(p)=1\).

\begin{theorem}[Arithmetic of the output Gram operators]\leavevmode
\label{thm:closedperiod318}
At the original complex embedding, \(K_u\) is a positive
scalar multiple of the output Gram operator of the projected
signed definite lift in the stated metric. For every
\(\mathfrak p\mid p\), every entry of \(D_u^2 K_u\) has
valuation at least \(b_{u,p}\). In particular, for any list
\(u_1,\ldots,u_k\), allowing repetitions,
\[
 v_{\mathfrak p}\left(
    \Big(\prod_{j=1}^k D_{u_j}^2\Big)
              \operatorname{tr}(K_{u_1}\cdots K_{u_k})\right)
                  \ \ge\ \sum_{j=1}^k b_{u_j,p}.
\]
If the listed operators are nonzero, dividing this trace
by \(\prod_j\operatorname{tr}K_{u_j}\) gives the normalized
mixed trace in the physical Hilbert spaces. It is an element
of \(J\) independent of all the chosen bases, their integer
clearings and the nonzero scalars \(\kappa_u\).
\end{theorem}
\begin{proof}
The matrix of the Riesz representative \(T_u:M_u\to E_0\)
is \(\kappa_u(w_0G)^{-1}P_u\). The adjoint is
\(\overline{\kappa_u}(w_u H_u)^{-1}P_u^*\), since the two
metric matrices occur on opposite sides in the adjoint
identity. Consequently
\[
 \begin{aligned}
 T_uT_u^\dagger&=\gamma_u G^{-1}P_u H_u^{-1}P_u^*,
        &\gamma_u&=|\kappa_u|^2/(w_0w_u),\\
 K_u&=(\beta_u/\gamma_u)T_uT_u^\dagger,
        &\beta_u&=
    \det(d_0G)\det(d_uH_u)/(d_0d_u).
 \end{aligned}
\]
Both scalars are positive. This proves the assertions about
the normalized traces. In particular a change of basis in
the common endpoint is one change of coordinates for every
factor of the product.

All entries of \(A,A_u\) are algebraic integers. Each entry
of \(K_u\) is a sum of products of two such coefficients
with one entry of \(P_u\) and the conjugate of an entry of
the same \(P_u\). Theorems~\ref{thm:originalgram315}
and~\ref{thm:poscolumn316} give exactly the stated bound
on these pairs. Matrix multiplication and the ultrametric
inequality prove the trace estimate. The two period factors
from a word always occur together in this calculation.
Thus its local choice of arithmetic frame may be made
independently of the frames for the other words.
\end{proof}

The exterior traces also detect an entire incoming space at
the last definite step. Take the source and target to be the
full finite-level highest-line spaces with the original
characters in Theorem~\ref{thm:lastrow311}. The integral
lists of Theorem~\ref{thm:intrinsicminor315} span these
spaces, since their covariances are strictly positive there.
Use the last matrix \(P_*\) of normalized entries \(z(U,V)\)
and the corresponding physical lift. In this normalization
\(|\kappa_*|=w_-w_+|C|c^{-d/2}/(D_-D_+)\). For
\(k_0=\dim\mathcal U\) and \(1\le j\le k_0\) one has
\[
 \operatorname{tr}(\textstyle\bigwedge^j K_*)\ \ge\
        {k_0\choose j}\bigl(\beta_*a_*/\gamma_*\bigr)^j>0 .
\]
Indeed the squared norm of the physical lift of each vector
of \(\mathcal U\) is at least \(a_*\) times its squared norm
by Theorem~\ref{thm:lastrow311}. If fewer than \(k_0\)
eigenvalues of \(T_*^\dagger T_*\) were at least \(a_*\),
the orthogonal complement of their eigenspaces would meet
\(\mathcal U\) nontrivially. This contradicts the bound.
The positive eigenvalues on the output equal those on the
source since \(T_*\) identifies their eigenspaces. Taking
the elementary symmetric sum of order \(j\) gives the
inequality. It concerns the last definite operator on the
same field as the incoming Bessel vectors.

\subsection{Elementary divisors on the same split path}\label{nav:V-65}
Continue at the auxiliary place of Theorem~\ref{thm:splitpath322},
with residue characteristic \(\ell\). Divide out the explicit
unramified scalar \(\eta\) in that theorem, and keep all its
Haar and index factors. Let \(\chi_0\) have order \(N\) on
\(K^\times\). For a fixed tame character \(\xi\), take the
\(k\) consecutive twists \(\chi_j=\xi\chi_0^j\),
\(0\le j<k\le N\). Choose a generator \(g\) of \(K^\times\)
and the \(N\) terminal frequencies obtained by multiplying
the frequency of that theorem by lifts of \(g^i\),
\(0\le i<N\). Denote the resulting \(N\)-by-\(k\) matrix
of local values by \(H\), with frequency as row index.
Let \(p\ne\ell\), take a coefficient field containing the
finite character values, and work at a place over \(p\) with
valuation normalized by \(v(p)=1\).

\begin{theorem}[Integral minors for consecutive twists]\leavevmode
\label{thm:tamedivisors322}
The \(k\) elementary divisors of \(H\) are units for \(p\nmid N\).
For \(N=N_0p^e\) with \(e\ge1\) and \(p\nmid N_0\), their
valuations in increasing order, indexed by \(0\le j<k\), are
\[
 \lambda_j=\frac{m_j+(p-1)\sum_{a=1}^{e-1}
                   p^{a-1}\lfloor m_j/p^a\rfloor}
                   {p^{e-1}(p-1)},\qquad m_j=\lfloor j/N_0\rfloor .
\]
In particular the least valuation among minors of order \(h\)
is \(\sum_{j=0}^{h-1}\lambda_j\). The frequency with \(i=0\)
has a unit entry for every twist. With all \(N\) twists, the
largest valuation is \(v_p(N)\).
\end{theorem}
\begin{proof}
Theorem~\ref{thm:splitpath322} gives
\(H_{ij}=a_j\chi_j(g^i)^{-1}\) for one constant \(a_j\)
per column. Every \(a_j\) is a unit at the chosen place.
Indeed all its factors apart from the Gauss sums are roots
of unity and integral powers of \(q\). The Gauss sum of the
trivial character is \(-1\). For a nontrivial character it
is an algebraic integer and satisfies
\(G(\chi)\overline{G(\chi)}=q\), by the direct character
calculation in Theorem~\ref{thm:lastcolumn321}. Both factors
are integral, so each is a unit above \(p\ne\ell\). Removing
these column units and the row units \(\xi(g^i)^{-1}\)
leaves the evaluation matrix \(((\zeta^i)^j)_{i,j}\), with
\(\zeta=\chi_0(g)^{-1}\) primitive of order \(N\). This also
proves the statement at \(i=0\).

Use the monic polynomial of degree \(j\)
\(\prod_{b<j}(Z-\zeta^b)\) as the new column in place of
\(Z^j\). The change of columns is integral with determinant
one. Its entries above the leading diagonal are zero. At
\(i\ge j\) the entry is a unit times
\(\prod_{u=i-j+1}^{i}(1-\zeta^u)\).
If \(p\nmid N\), the distinct \(N\)-th roots remain distinct
in the residue field, which proves the first case.

Now \(v(1-\zeta^u)=0\) unless \(N_0\mid u\), since
reduction is injective on prime-to-\(p\) roots of unity.
For \(u=N_0v'\) with \(0<u<N\) its valuation is
\(p^{v_p(v')}/(p^{e-1}(p-1))\). To see this, a primitive
\(p^c\)-th root \(\alpha\) reduces to one. The factors
\(1-\alpha^l\), \(1\le l\le p^c\) and \(p\nmid l\), have
the same valuation because their quotients by \(1-\alpha\)
reduce to the nonzero residues of \(l\). Their product is
\(\Phi_{p^c}(1)=p\), giving the formula for every \(c\).
In any interval of \(j\) successive integers there are at
least \(\lfloor j/(N_0p^a)\rfloor\) multiples of \(N_0p^a\).
All these lower bounds are attained in \(1,\ldots,j\).
Writing \(p^h=1+(p-1)\sum_{a=1}^{h}p^{a-1}\) therefore
gives the minimum \(\lambda_j\) for column \(j\), at \(i=j\).

Divide each new column by its diagonal entry. The resulting
matrix is integral and its leading \(k\)-rowed square is
lower triangular with unit diagonal. Integral row operations
make that square the identity and then clear the remaining
rows. The original elementary divisors are consequently
the diagonal entries up to units, in their stated order.
For \(k=N\) the last diagonal valuation is the valuation
of \(\prod_{u=1}^{N-1}(1-\zeta^u)=N\). The assertion
about the minors follows by taking the diagonal form.
\end{proof}

The local calculation distinguishes a single common nonzero
test from the recovery of individual twists after summation.
For example, take the full character family and multiply
each column of \(H\) by \(a_j^{-1}\). Their sum has value
\(N\) at the first frequency and zero at the other \(N-1\)
frequencies, by orthogonality. Thus signed summation can
change the local content at a prime dividing \(N\), even
though each column is primitive there.

\subsection{The two components of the finite shift at a split prime}\label{nav:VII-64}
Suppose \(p\) splits in the marked imaginary quadratic field \(K\).
Denote the two idempotents of \(\OO_K\otimes\ZZ_p\) by
\(e_+,e_-\), in either order. For the list of finite tests of
Theorem~\ref{thm:neutral315}, let \(a_\epsilon\) be the least
nonnegative integer for which \(p^{a_\epsilon}e_\epsilon b_p\)
belongs to \(e_\epsilon M_p\) for every occurring \(b\), with
\(\epsilon=+,-\). Both are at most \(v_p(Q^0)\).

\begin{theorem}[A split-prime estimate in the marked CM differential]\leavevmode
\label{thm:splitinput315}
At a coefficient place restricting to \(+\) on \(K\), a
degree-\(d\) class derivative in the product basis loses at most
\(d a_-\) in valuation from translations. At a place over \(-\)
the corresponding loss is at most \(d a_+\). Thus the two
translation losses in a product of a class derivative with a
conjugate derivative of the same degree sum to at most
\(d(a_++a_-)\). The assertion applies after the class actions
at both ends of the edge.
\end{theorem}
\begin{proof}
The localized product basis of Theorem~\ref{thm:adapted314}
gives an isomorphism of \(\OO_K\otimes\ZZ_p\)-modules with
the standard product lattice. Its quotient map commutes with
\(e_+\) and \(e_-\). By Theorem~\ref{thm:neutral315} the class
actions and the vacuum intertwiners have introduced no further
class modulo this lattice at \(p\). For each component \(P\)
on the CM elliptic curve, passage from an additive shift to the
torsion point multiplies the coordinate and its period lattice
by the same elliptic period. This map on quotient modules is
\(\OO_K\)-linear. Hence \(p^{a_\epsilon}\) kills the
\(\epsilon\)-part of the \(p\)-primary component of \(P\).

Let \(I=\operatorname{Ann}_{\OO_K}(P)\).
The cyclic module \(\OO_KP\) is \(\OO_K/I\). Tensoring this
description with either local factor shows that the localized
ideal \(I\) is the annihilator of the localized point.
The exponent of \(I\) at \(-\) is therefore at most \(a_-\).
At \(+\) this is the exponent of \(\bar I\), which occurs in
Theorem~\ref{thm:torsion312}. The normalized valuation agrees
with this exponent since \(K\) is split at \(p\). There
\(\kappa_p=0\), so an ordinary derivative of order \(h\) of
this translated factor loses at most \(ha_-\). In total order
\(d\) the orders on the separate factors sum to \(d\).
This proves the bound above \(+\). Interchanging the local
factors proves the other bound. Complex conjugation exchanges
the two places on \(K\), which proves the product assertion.
\end{proof}

For the indicator in Proposition~\ref{prop:split173} the two
orders can be read directly from its position lattice. More
generally, let \(L_0\) be the position lattice of the vacuum
of amplitude \(a_0\), and let a test of amplitude \(a\) have
support \(y+P\) with \(P\subset L_0\) a full sublattice.
For \(x-y\in L_0\), orthogonality of characters gives
\[
 {\bf1}_P(x-y)=\frac1{[L_0:P]}
       \sum_{\xi\in P^\vee/L_0^\vee}\psi(\langle x-y,\xi\rangle).
\]
Multiplication by the translated vacuum supplies the condition
\(x-y\in L_0\). This is the Weyl expansion with coefficients
\((a/a_0)/[L_0:P]\) times roots of unity in the symmetric Weyl
convention. Its position classes have the order of \(y\) modulo
\(L_0\). Its frequency classes have maximum order the exponent
of \(L_0/P\), by the perfect character pairing.

In particular, at a rational prime completely split in \(E\),
take the position half over a chosen prime of \(K\), say \(+\).
At a place \(v\mid p\) the congruence indicator has
\(y=\binom{1_n}{0}\in L_0\), \(P=p^hL_0\) and
\([L_0:P]=p^{h ns}\) for \(s\)-by-\(n\) position matrices.
The finite Fourier denominator is still \(p^{h ns}\) for unit
relative amplitude. The contribution of this indicator to the
component exponents is \(0,h\). If it is the only nonspherical
factor above \(p\), the translation cost in the conjugate
product is at most \(dh\). The same holds for several such factors
with coherent orientation and maximum depth \(h\).
If the chosen position halves vary over \(v\mid p\), the two
maxima must be taken in their actual idempotent components.

\subsection{The local cost of the position cosets}\label{nav:VII-65}
Use one edge \(i\) and the expansion of
Theorem~\ref{thm:posframe316} at an odd split prime \(p\).
Write \(A_p\) for the maximum of \(0\) and the negative normalized
valuations of the nonzero \(d_{\lambda,\mathbf y}\) in this
list, at all places over \(p\). For each nonzero term let
\[
 k_v(y_v)=\min\{k\ge0:p^k y_v\in P_v\}.
\]
This is the exponent of \(p\) in the additive order of the
position class. For \(\epsilon=+,-\) let \(b_\epsilon\)
be the maximum of these integers over the terms and the
\(v\mid p\) whose position prime \(w_v\) restricts to
\(\epsilon\) on \(K\), using zero for an empty maximum. Put
\[
        J^{\rm pos}_{i,p}=A_p+
          \tfrac12d_i(b_++b_-)+\tfrac12d_i\rho_p .
\]

\begin{theorem}[Conjugate derivative bounds from the support of the finite test]\leavevmode
\label{thm:poscost316}
For two edge derivatives of the prescribed degree profile,
one with the orientation of the edge and the other conjugated,
the sum of their normalized valuations is at least
\(-2J^{\rm pos}_{i,p}\) in the product coordinates of
Theorem~\ref{thm:posframe316}, after converting derivatives
back to the physical matrix coordinates. This bound holds for
every pair of class sections on the input list. Among homothetic
position lattices \(P_v\) of the given shape at all \(v\mid p\)
the largest ones on whose cosets the list is constant minimize
this upper bound on the loss.
\end{theorem}
\begin{proof}
The coefficient losses are at most \(A_p\) in each copy.
The product basis identifies the local trace lattice with a
free \(\OO_K\otimes\ZZ_p\)-module and commutes with both CM
idempotents. In its coordinates the \(\epsilon\)-parts of all
elliptic torsion points have order dividing \(p^{b_\epsilon}\)
in their \(p\)-primary components, by
Theorem~\ref{thm:posframe316}. Localizing the conjugate of the
annihilator of such a point exchanges the two components.
Theorem~\ref{thm:torsion312} therefore costs at most \(d_i b_-\)
at a place over \(+\) and \(d_i b_+\) at a place over \(-\).
This is the proof of Theorem~\ref{thm:splitinput315} in the
new product lattice. The Leibniz expansion of a derivative of
order \(d_i\) has exactly \(d_i\) differentiations in total.

For the coordinate conversion, the physical entries of a product
basis vector belong locally to \(\mathfrak c\mathfrak b^{-1}\).
At block \(\tau\) the coefficients of a physical derivative in
the product variables are the conjugate entries, as in
Theorem~\ref{thm:adapted314}. Thus the lower bound per occurrence
is \(v_{\mathfrak p}(\bar\tau\mathfrak c)-
v_{\mathfrak p}(\bar\tau\mathfrak b)\). The conjugate copy gives
the same expression with \(\tau\) in place of \(\bar\tau\).
Their sum is \(-\rho_p\), because
\(\mathfrak c\bar{\mathfrak c}=\OO_E\) and
\(\mathfrak b\bar{\mathfrak b}=\mathfrak D_{E/K}\).
The two derivatives have the same number \(d_{i,\tau}\) of
occurrences at every block. Multiplying the bounds and summing
the terms of the coordinate expansions proves the assertion.

For the minimizing choice omit an identically zero list.
The translations of one position variable fixing every input
form a subgroup containing \(\varpi_v^hL_v\) for large \(h\)
by local constancy. Its elements lie in the difference of a
nonempty compact support with itself, by using a nonzero
coefficient function with the other variables fixed. Hence it
is bounded. There is a least integer \(h_v\) such that
\(\varpi_v^{h_v}L_v\) fixes the list. The separate subgroups
act in separate variables, so the simultaneous largest choice
is admissible. Shrinking any \(P_v\) subdivides each nonzero
coset with the same relative amplitude on every part. In the
normalized coefficients this multiplies each copy by the inverse
square root of the subdivision index. The maximum loss \(A_p\)
can only increase. Every new class maps to its old class, so
the two torsion exponents can only increase as well.
\end{proof}

For a single congruence test at a place of a rational prime
completely split in \(E\), take the position model with unit
relative amplitude of Proposition~\ref{prop:split173}. Its
support is \(y+p^hL_v\), where \(y=\binom{1_n}{0}\) has a unit
entry and \(m=ns\). With spherical unit vacua elsewhere over
\(p\), \(A_p=hns/2\) and \(b_++b_-=h\). The raw Fourier
expansion in \(L_v\) has every nonzero coefficient of valuation
\(-hns\) by Theorem~\ref{thm:splitinput315}. The support
calculation reduces the bound \(hns+dh/2\) to \(h(ns+d)/2\)
for degree \(d\). It remains valid for any list of translates
by integral invertible position matrices with unit amplitude
factors. Such matrices preserve \(p^hL_v\) and preserve
primitivity of \(y\). In particular this includes the two
congruence groups fixing the prescribed test.

There is a second case in the character transport of
Proposition~\ref{prop:split233}. In its Haar normalization
the transported spherical test at \(v_*\) has value
\(|\xi_w|_v^{m/2}a_{0,v}\) on
\(P_v=\xi_w^{-1}L_v\). If \(\operatorname{ord}_v(\xi_w)=s\),
then \(h_v=-s\). The relative amplitude \(q_v^{-ms/2}\)
multiplied by \(q_v^{-m h_v/2}\) is exactly one. Its position
class is zero. Thus both the normalization of this dilation
and its change of support have zero local cost. An additional
splitting scalar in a chosen realization remains in the actual
coefficient list. If it is common to the entries and has absolute
value one, its product with its conjugate is one. Distinct
phases in a signed sum are retained in that sum.

\subsection{The last column at a split twist place}\label{nav:VII-67}
Fix a split auxiliary place of a tuple on the interval of
Section~\ref{sec:missing}. Use the plain frames for its tame twist.
Write \(k\) for the local field, \(O\) for its integer ring,
\(\pi\) for a uniformizer and \(K\) for its residue field of
size \(q\). Use \(|\pi|=q^{-1}\), additive volume one on \(O\)
and an additive character \(\psi\) of conductor \(O\). Thus
\(\psi\) is trivial on \(O\) and nontrivial on \(\pi^{-1}O\).
Signs below use the Fourier kernel \(\psi(-zy)\).
In the two integral frames let \(E\subset K^{n+1}\) and
\(U\subset K^d\) be the distinguished coordinate spaces of the
same dimension \(r\ge1\). Keep their ordered bases.
For a matrix \(Z\) of size \(d\) by \(n+1\), the parahoric
Schr\"odinger test of the penultimate edge is
\[
 \Phi(Z)=
 \begin{cases}
 \chi(\det(\bar Z|_E:E\longrightarrow U)),
       & Z\text{ integral and }\bar Z|_E:E\overset{\sim}{\longrightarrow}U,\\
 0, &\text{otherwise}.
 \end{cases}
\]
Here \(\chi\) is the inverse of the tame splitting character
on residue units. Indeed its defining blocks require the
upper left block to be integral invertible, the lower left
block to vanish after reduction, and both right blocks to be
integral. These conditions and the displayed determinant give
the formula in the ordered frames.

Fix the actual source representative \(h\in{\rm GL}_{n+1}(k)\)
and write \([X\mid y]\) in the original last-column order,
inserting the frame permutation into \(h\). The source action
on the position function is
\(\mu(\det h)|\det h|^{d/2}\Phi([X\mid y]h)\), with the
corresponding fixed determinant phase if a permutation is
used to express the action. Denote this scalar by \(c_h\),
using the chosen original frames throughout. The other edge
supplies the unchanged top test \(\Phi^+(B)\). It is enough to
evaluate the following factor of the coupled kernel:
\[
 D_h(X,z)=\int_{k^d}\Phi([X\mid y]h)\psi(-zy)\,dy .
\]
The original kernel is \(c_h\Phi^+(B)D_h(X,z)\).

Write \(b\) for the last row of \(h\), \(W=[X\mid0]h\).
Choose \(i\) with \(|b_i|\) maximal and put
\[
 \beta=b_i,\quad y_0=-W_i/\beta,\quad
 W_0=W+y_0 b,\quad z'=\beta^{-1}z ,
\]
where \(W_i\) is the \(i\)-th column. If \(W_0\) is integral,
write \(A=\bar W_0|_E\) and \(\lambda=\overline{b/\beta}|_E\).
Fix a nontrivial additive character \(\theta\) of \(K\).
For \(z'\in\pi^{-1}O^{1\times d}\) let \(l:K^d\to K\)
be the unique linear form with
\(\psi(-z'u)=\theta(l(\bar u))\) on \(O^d\).

\begin{theorem}[Fourier integral of the marked last column]\leavevmode
\label{thm:lastcolumn321}
If \(W_0\) is nonintegral then \(D_h=0\). In the remaining
cases divide \(D_h\) by \(|\beta|^{-d}\psi(-zy_0)\).
For \(\lambda=0\) the result is zero unless \(A:E\to U\)
is an isomorphism and \(z'\) is integral. When both conditions
hold the result is \(\chi(\det A)\).

For \(\lambda\ne0\) it is zero unless \(A\) maps \(\ker\lambda\)
injectively into \(U\) and \(z'\in\pi^{-1}O^{1\times d}\).
Under these two conditions take \(e\in E\) with \(\lambda(e)=1\).
For \(v\in U\) define \(L(v)\) as the determinant, in the marked
bases of \(E,U\), of the map that agrees with \(A\) on
\(\ker\lambda\) and sends \(e\) to \(v\). The answer is zero
unless \(l|_U=tL\) for some \(t\in K\). In that event it is
\[
      q^{\,r-1-d}\theta(-l(Ae))G_t(\chi),\qquad
      G_t(\chi)=\sum_{a\in K^\times}\chi(a)\theta(ta).
\]
For nontrivial \(\chi\), \(G_0=0\) and every \(t\ne0\) gives
\(|G_t|=\sqrt q\). For trivial \(\chi\) the two values are
\(q-1\) at zero and \(-1\) elsewhere.
\end{theorem}
\begin{proof}
The \(i\)-th column of \(W+yb\) is \(\beta(y-y_0)\).
Hence integrality of that column requires
\(y=y_0+\beta^{-1}u\), \(u\in O^d\), with measure
\(|\beta|^{-d}du\). In this parametrization the full matrix
is \(W_0+u(b/\beta)\). The second term is integral, so
integrality of \(W_0\) is both necessary and sufficient.
On \(E\) the reduced matrix is \(A+\bar u\lambda\).
If \(\lambda=0\) its value is independent of \(u\).
Integration of a character on \(O^d\) gives one exactly
on the annihilator \(O^{1\times d}\) and zero elsewhere,
by translation invariance and the conductor convention.

In the second case the function of \(u\) is constant on
each coset modulo \(\pi O^d\). Translation by \(\pi O^d\)
therefore gives the asserted frequency restriction. On these
frequencies the additive character factors through \(K^d\).
The characters \(x\mapsto\theta(ax)\), \(a\in K\), are distinct
since multiplication by a nonzero element is onto. They
exhaust the additive dual by cardinality, proving the assertion
about \(l\) in each coordinate. Every residue class has measure
\(q^{-d}\). The restriction on \(\ker\lambda\) follows from
the invertibility required on \(E\).

Set \(U_0=A(\ker\lambda)\). When the restriction on the
kernel holds, \(L\) is a nonzero linear form on \(U\) with
kernel \(U_0\), by expansion of a determinant in its last
column. The admissible \(\bar u\) are precisely those for which
\(v=Ae+\bar u\) lies in \(U\) and \(L(v)\ne0\). Summing
\(\theta(l(v))\) on a fiber of \(L\) gives zero if \(l\) is
nonzero on \(U_0\), since the character is then nontrivial
there. Otherwise \(l|_U=tL\), and the fiber over \(a\ne0\)
contributes \(q^{r-1}\theta(ta)\chi(a)\), apart from the
common factor \(\theta(-l(Ae))\). This proves the evaluation.

For nontrivial \(\chi\), summing it over \(K^\times\) gives
zero. For \(t\ne0\) substitution gives
\(G_t=\chi(t)^{-1}G_1\). Additive character orthogonality gives
\(\sum_t |G_t|^2=q(q-1)\), by expanding the square and summing
over \(t\) first. This gives the claimed magnitude. The same
orthogonality gives the two elementary values for the trivial
character.
\end{proof}

For the identity representative \(h=1\) before permuting
bases, the separated column lies in the distinguished source
block of the plain frame. This follows directly by ordering
the extra coordinate among the \(s+1\) distinguished columns
at the penultimate stage. Thus \(\lambda\ne0\).
Choose integral \(X\) which sends the other \(r-1\) marked
columns to \(r-1\) distinct marked target basis vectors and
is zero on the other columns. These choices constrain only
the residue of \(X\). Choose \(B\) with identity
marked block and all other entries zero. Then
\(\Phi^+(B)=1\). As \(z\) runs through lifts of forms extending
\(tL\) for \(t\ne0\), the same positions and frequencies work
for every nontrivial tame twist. At an integral \(z\) the
factor vanishes for every such twist. To use the identity
representative in the global chain one must also show that
the preceding edges reach its double coset. The general
formula above applies to the representatives those edges
actually reach. The local definitions are
\srcloc{Development/AbelianVarietyTheory/Definitions/TwistedInnerProducts/StandardOrderExtendedBasisApplyLowBlock}{399, 650, 796, 889--920}.

\subsection{Compact integration along the supported path}\label{nav:VII-68}
Fix the split place, representatives \(h_j\), primitive core
vectors \(z_j\) for \(j>0\), and matrices of
Theorem~\ref{thm:splitpath322}. At vertex \(j\le s\) write
\(A_j\subset K^{n_j}\) for the marked core plane and \(P_j\)
for its stabilizer in \({\rm GL}_{n_j}(O)\) after reduction.
The full parahoric \(J_j\) preserves \(A_j\) together with
the flag coordinates in rank \(D_j\). Let \(\Gamma_j\)
be the image in the core group of the intersection of
\(h_jJ_jh_j^{-1}\) with the semidirect product of the core
general linear group and the flag unipotent group. This is
the compact group defining the middle weight
\(w_j={\rm vol}(\Gamma_j)^{-1}\).
Every full integral general linear group has volume one
in its own rank. Put \(d=n_{s+1}\), let \(U\subset K^d\)
be the \(r\)-dimensional marked target of the penultimate
test, and write \(J\subset{\rm GL}_d(O)\) for its parahoric.

\begin{theorem}[Middle weights and the terminal orbit]\leavevmode
\label{thm:splitcompact323}
One has \(w_0=1\). For \(1\le j\le s\) the weight is
\[
  w_j={n_j\brack a_j}_q\,q^{j n_j}(1-q^{-a_j}),
\]
where the Gaussian binomial coefficient counts planes.
Each compact integration over \(\Gamma_j\) cancels the
corresponding weight in the product of consecutive kernel
entries, with the core action at both adjacent edges.
At the initial endpoint this assertion uses the matching
source character, or equivalently holds with that endpoint
kept fixed.

For the final two entries, take the coset lattice
\(\pi^M O^d\) of Theorem~\ref{thm:adjacent321} with
\(M\ge s+1\). At the marked positions of
Theorem~\ref{thm:splitpath322} every contributing coset
belongs to the single \(J\)-orbit of the primitive integral
vectors with reduction in \(U\). This orbit has cardinality
\((q^r-1)q^{(M-1)d}\). Its representative, with weight
\([{\rm GL}_d(O):J]\) times the orbit size and with the
original Haar integral over \(J\), contributes exactly the
full signed last-column value. These assertions hold for
each final unit multiple of the frequency and for each tame
twist on the same levels.
\end{theorem}
\begin{proof}
At \(j=0\) the core is entirely marked and \(h_0=1\), so
its projected compact is the full integral group. At \(j>0\)
a lift of a core element \(g\) to the semidirect product
has diagonal blocks \(g,I+H\), lower left zero, and upper
right \(B\) with first column zero. Here \(H\) is strictly
upper triangular. Conjugating by \(h_j^{-1}\) gives upper
right block
\[
  \pi^{-j}\bigl(B+(gz_j-z_j)e_0^t-z_j e_0^tH\bigr).
\]
In column zero the free matrices \(B,H\) have no effect.
On the core one must have \(g\in P_j\). The additional
condition is that the marked coordinates of \(gz_j-z_j\)
belong to \(\pi^jO\) and the unmarked coordinates to
\(\pi^{j+1}O\). These conditions suffice, since we can
take \(B=H=0\). Thus \(\Gamma_j\) is exactly the stabilizer
of the residue of \(z_j\) modulo this mixed lattice in \(P_j\).

The orbit consists of a primitive marked column modulo
\(\pi^j\), together with arbitrary unmarked entries in
\(\pi O/\pi^{j+1}O\). To obtain such a column, complete its
marked lift to an invertible marked block over \(O\), use
the given unmarked entries in the lower left block, and set
the lower right block to the identity. This proves both
inclusions in the orbit description. Its size is
\((q^{ja_j}-q^{(j-1)a_j})q^{j(n_j-a_j)}\). Multiplication
by \([{\rm GL}_{n_j}(O):P_j]={n_j\brack a_j}_q\) proves
the weight formula.

For \(g\in\Gamma_j\), the lift
\(h_j^{-1}{\rm diag}(g,I)h_j\) belongs to the parahoric
and its determinant on the marked plane is
\(\det(\bar g|_{A_j})\). This follows from its block
triangular form and identity flag block. On replacing
\(X_{j-1}\) by \(gX_{j-1}\) in the preceding integral,
replace each upper free vector \(Y_i\) by \(gY_i\) as
well. The measure is preserved and the phase on the lower
free entries is fixed. The resulting multiplier on the
position test is \(\chi(\det(\bar g|_{A_j}))\).
At the next edge \(X_j\mapsto X_jg^{-1}\) amounts to
right multiplication of its affine matrix by
\({\rm diag}(g^{-1},I)\). This gives the inverse multiplier.
The calculation includes the penultimate test by keeping
the separated last column fixed. The two compact Weil
splittings also cancel, since consecutive \(\mu_i\) agree
on units. Thus the joined integrand is constant under each
of these simultaneous changes of core charts. An integral
over \(\Gamma_j\) contributes exactly its volume, which
proves the cancellation. At the input the same calculation
uses the inverse multiplier on a matching test, or one
simply omits the input average and its unit weight.

At the last variable the support condition is
\(y=-X_sz_s+\pi^s v\), where \(\bar v\) lies in \(U\)
outside the span of the \(r-1\) marked columns of \(X_s\).
For \(s\ge1\), reduction of \(y\) is the fixed nonzero
vector \(-\bar X_s\bar z_s\in U\). For \(s=0\) reduction
of \(y=v\) is again a nonzero vector in \(U\). The orbit
of all these primitive classes modulo \(\pi^M\) consists
of exactly the classes specified in the statement: the
marked block acts transitively on nonzero residue
vectors, and the principal congruence subgroup moves
any primitive lift to any other with the same reduction.
For the latter fact, extend a primitive vector to an
integral basis and change its first basis vector by the
specified multiple of \(\pi\). Counting the lifts gives
the cardinality. A scalar support lattice containing a
primitive position contains \(O^d\); hence the full orbit
is present in the coset space used for the decomposition.

Write \(P_a,T_a\) for the two pieces of one last coset.
Over probability Haar on \(J\), the integral of
\(P_a(kX_s)T_a(Bk^{-1},zk^{-1})\) is the same for
all \(a\) in the orbit by Theorem~\ref{thm:adjacent321}.
Their sum is the entire kernel just computed, since at
the marked positions the other orbits have zero summands.
It has the same value after simultaneous translation by
\(k\), by permutation of the cosets and cancellation of the
two characters. Therefore this probability average equals
that entire value divided by the orbit size. Restoring the
specified measure on \(J\) divides it by
\([{\rm GL}_d(O):J]\). The assigned weight cancels both
factors. The different core and final group actions commute
on the joined product since they multiply on opposite sides
of each shared position matrix. All these integrations can
therefore be performed together.
\end{proof}

The computation concerns the compact integrations on this
local path. Summing the other interior double cosets and
then projecting the global automorphic row are further
operations. The local projected group and weights appear in
\srcloc{Development/AbelianVarietyTheory/Definitions/TwistedInnerProducts/StandardOrderExtendedBasisApplyLowBlock}{675--701}
and \srcloc{Development/AbelianVarietyTheory/Definitions/LatticeFrames/CosInj}{110--118}.

\subsection{Reading every class of the standard terminal row}\label{nav:VII-69}
Take the standard definite terminal group \(G=U(E^n)\) over
\(F=E^+\) in its given orthonormal frame, on a field of
Section~\ref{sec:missing}. Let \(G_\infty=\prod_{v\mid\infty}G(F_v)\)
and fix a coefficient-field embedding \(\sigma:K\to\CC\).
Use the terminal polynomial module \(V_K\) of
Theorem~\ref{thm:terminaldual325}, including its prescribed
determinant character. Its coefficient map \(C:V_K\to
V_\CC\) applies \(\sigma\) in the same polynomial frame at
every tuple. Write \(\rho\) for the resulting compact
representation. Its invariant norm is the Hermitian Fischer
norm divided by the squared extremal norm \(H\). On choosing
the dual determinant generator, write \(B\) for the
corresponding normalized bilinear pairing as in that theorem.

At a tuple \(\xi\) use its actual finite level \(J_\xi\),
its finite character \(\omega_\xi\) in the right translation
convention, and the smaller level \(J_\xi^1\) that fixes
every marked residue frame. Given \(g\in G(\mathbb A_{F,f})\),
write \(D_\xi(g)\) for the lattice with local completions
\(g_v\Lambda_v\), marked at the indicated primes by the
reductions of the coordinates of \(g_v^{-1}\). Use the
genus of the original lattice. A table of class values
assigns \(w_i\in V_K\) to representatives \(g_i\) for
\(G(F)\backslash G(\mathbb A_{F,f})/J_\xi\).
Put \(\Gamma_i=G(F)\cap g_iJ_\xi g_i^{-1}\).
The table satisfies
\(\rho_K(\gamma)w_i=\omega_\xi(g_i^{-1}\gamma g_i)w_i\)
for \(\gamma\in\Gamma_i\), using \(K\)-valued finite
characters and the rational coefficient representation.

\begin{theorem}[Class values and the projected terminal norm]\leavevmode
\label{thm:terminalreading325}
For every such table there is a unique equivariant
coefficient function \(\phi\) on the decorated genus with
these values and this finite character. Its standard
adelic form, evaluated against any \(\ell\in V_\CC^*\),
is given at every finite adele by
\[
 f_{\phi,\ell}(k,g)
    =\ell\bigl(\rho(k)^{-1}C(\phi(D_\xi(g)))\bigr),
               \qquad k\in G_\infty .
\]
The same coefficient map works for every tuple. There is
one linear map \(j:V_\CC^*\to V_\CC\) with
\(\ell(Cw)=B(j\ell,Cw)\) for every \(w\), common
to all classes and tuples.

Let \(f_\ell\) be an independently given automorphic family
linear in the dual test, of the same compact and finite
types, after the prescribed projection commuting with right
translation at infinity. Write \(z_i\) for its coefficient
vector at \((1,g_i)\), so \(\ell(z_i)=f_\ell(1,g_i)\).
For any orthonormal basis \((\ell_a)\) of the compact dual,
with probability Haar at infinity, the original Haar at the
finite level and counting measure on \(G(F)\), one has
\[
 \sum_a\|f_{\ell_a}-f_{\phi,\ell_a}\|_{[G]}^2
    =\operatorname{vol}(J_\xi)\sum_i
              \frac{\|z_i-Cw_i\|^2}{|\Gamma_i|}.
\]
If the differences are algebraic in the given coefficient
coordinates, testing each difference by the integral flag
minors with \(B\) detects its divisibility as in
Theorem~\ref{thm:terminaldual325}, at every finite coefficient
place in the corresponding local frames.
\end{theorem}
\begin{proof}
Intersect \(E^n\) with the product of the local lattices
\(g_v\Lambda_v\). It contains a nonzero integer multiple
of the original lattice and is contained in the reciprocal
multiple, so it is a full lattice. Its completion gives
each of the specified lattices. Indeed
between these two bounds one works in a finite quotient
of the original lattice, which is the product of its local
quotients by the Chinese remainder theorem. The marked
coordinate reductions are surjective on each specified
residue frame because they are transported from the standard
one. Equality of two decorated lattices now says that their
ratio preserves the standard lattices locally and induces
the identity in the marked coordinates. This is precisely
membership in \(J_\xi^1\). Thus finite adeles modulo
\(J_\xi^1\) parametrize the decorated genus, with the
action of \(J_\xi/J_\xi^1\) by changes of marking.

If \(g=\gamma g_i u\), set
\(\phi(D_\xi(g))=\omega_\xi(u)\rho_K(\gamma)w_i\).
Two such expressions for \(g\) differ by a member of
\(\Gamma_i\). Its action on \(w_i\) is exactly the finite
character in the table, so the two values agree. This
extends restriction to the representatives inversely,
including when a stabilizer forces some values to vanish.
Scalar extension of the polynomial and determinant actions
intertwines \(\rho_K(\gamma)\) with
\(\rho(\gamma_\infty)\). Hence the displayed adelic
function is invariant under diagonal rational translation,
has the asserted finite character, and transforms on the
right at infinity by the dual test action. Conversely these
laws recover the table by evaluation at \(k=1,g=g_i\).
This gives the inverse-restriction realization of the
standard terminal row.

Coefficient conjugation in the chosen polynomial frame
preserves the complex space since its defining minor
products have integer coefficients. For a nonzero \(v\)
one has \(B(\bar v,v)>0\) by the monomial formula.
Consequently \(j\ell\) is the unique solution to
\(B(j\ell,v)=\ell(v)\) for every \(v\). This proves
linearity and independence of \(\xi\). More generally an
already specified bilinear pairing \(\operatorname{pr}(v,e)\)
with a terminal test \(e\) gives the solution
\(j(\operatorname{pr}(\,\cdot\,,e))\), with its given
scalar. Compact covariance follows by applying the defining
equality to \(\rho(k)^{-1}v\).

The class set is finite by finiteness of the genus of a
definite lattice and finiteness of the residue markings.
Each \(\Gamma_i\) is finite since it preserves the
definite trace lattice. The contribution of its class to
Haar integration is integration on \(G_\infty\times J_\xi\)
with factor \(1/|\Gamma_i|\), since diagonal rational
translation identifies two points of this parametrization
precisely by that finite group, acting freely. On this
piece the difference uses the vector \(z_i-Cw_i\)
transformed by a unitary compact operator and a finite
scalar of modulus one. The sum of the squared evaluations
over \(\ell_a\) is its squared norm by Parseval. Integration
gives the formula, including the stabilizer and level
volume. The local divisibility assertion follows by testing
in the dual frame at that class and applying the integral
duality theorem.
\end{proof}

The construction gives the actual coefficient reading for
the standard adelic realization of a decorated table. For
the independently projected theta family the right side is
an exact discrepancy at the prescribed level. Controlling
those class vectors, or proving their agreement with the
arithmetic ones in the physical normalization, remains a
global period problem. The table restriction and its inverse
are described in
\srcloc{Formalism/LocalToGlobal/Columns/RhoEqOfSmulEq}{157--193}.
The standard adelic realization uses that inverse in
\srcloc{Development/AbelianVarietyTheory/Definitions/AutomorphicMembership/IsWtOfParUnique}{327--346}.

\subsection{A nonzero path through the split twist kernels}\label{nav:VI-65}
At one auxiliary split place keep \(k,O,\pi,K,q,\psi,\chi\)
of Theorem~\ref{thm:lastcolumn321}. The residue character
\(\chi\) may be trivial. Write \(s=m-2\), with \(m\ge2\).
In the plain frames put
\[
 a_j=k_0+j,\quad n_j=a_j+jt,\quad r=k_0+s+1,\qquad
 D_j=n_j+s+1-j=r+jt \quad(0\le j\le s+1).
\]
Here \(k_0\) is the initial core dimension and \(t\) the
rank increment in the full space. The \(D_j\) coordinates
are the \(n_j\) core coordinates and \(s+1-j\) flag
coordinates. The marked \(r\)-plane consists of the first
\(a_j\) core coordinates and all the flag coordinates.
In these bases each edge uses the parahoric test \(\Phi\)
of Theorem~\ref{thm:lastcolumn321}. The splitting characters
are \(\mu_i=\nu_i\delta\), with \(\nu_i\) unramified and
\(\chi\) the inverse of \(\delta\) on residue units. Put
\(\theta(v)=\psi(-\pi^{-1}\widetilde v)\) for \(v\in K\), where
\(\widetilde v\) is any integral lift. Set
\(G(\chi)=\sum_{v\ne0}\chi(v)\theta(v)\), summing in \(K\).

We describe precisely the interior entry used below. At
edge \(j<s\) put \(c=s-j\). With flag indices starting at
zero, the affine matrix has core block \(X_j\), lower left
block zero, and \(c+1\) remaining columns. Column zero is
\((0,e_0)\). For \(1\le i\le c\) its upper vector \(Y_i\)
is free. Its lower entry in row \(l\), \(0\le l<c\), is
\(b_{li}\) for \(l<i\), one for \(l=i\), and zero for \(l>i\).
Call this matrix \(A\). The entry at representatives
\(h_j,h_{j+1}\) integrates
\(\Phi(h_{j+1}^{-1}Ah_j)\psi(-\sum_{l<c}b_{l,l+1})\)
over the free entries with product additive Haar. Multiply
by the original Weil scalar
\[
 C_j=\frac{\mu_{j+1}(\det h_j)}{\mu_{j+1}(\det h_{j+1})}
       |\det h_j|^{D_{j+1}/2}|\det h_{j+1}|^{-D_j/2}
\]
and by the reciprocal of the relative unipotent volume.
The latter is the volume of \(N\cap h_{j+1}J h_{j+1}^{-1}\)
divided by the volume of \(N(O)\), before taking the
reciprocal. Here \(J\) is the target parahoric and \(N\)
has identity diagonal, free upper flag entries, and free
core entries over flag columns \(1,\ldots,c-1\).
At \(j=s\) use the Fourier integral in the separated last
column with its source Weil scalar and the unchanged top test.
Take \(h_{s+1}=1\) in \(C_s\).

\begin{theorem}[One supported path for every tame twist]\leavevmode
\label{thm:splitpath322}
There are representatives starting at \(h_0=1\), core
positions and a terminal frequency, all independent of
\(\delta\), for which every consecutive local entry is
nonzero. Their product, with the preceding normalizations, is
\[
       \eta\,\chi(u)\,q^{\,r-1-(s+1)n_{s+1}}\,
                  G(\chi)^{\,1+s(s+1)/2},
 \qquad
 \eta=\prod_{j=1}^{s}(\nu_{j+1}/\nu_j)(\det h_j).
\]
Here \(u\in K^\times\) comes from the ordered marked bases.
Multiplying the last frequency by an integral unit of residue
\(b\) multiplies the product by \(\chi(b)^{-1}\).
The representatives define a successive choice of supported
double cosets in the local unfolding, on the given levels.
\end{theorem}
\begin{proof}
Put \(z_0=0\). For \(1\le j\le s\) take \(z_j\) to be a
marked core basis vector. With \(e_0^t\) extracting the
first flag coordinate, choose
\[
 h_j=
 \begin{pmatrix}
 I&\pi^{-j}z_j e_0^t\\0&\pi^{-j}I
 \end{pmatrix}\quad(0\le j\le s).
\]
At \(j=0<s\) take \(X_0\) to identify the marked core with
a coordinate complement of \(z_1\) in the next marked core.
For \(1\le j<s\) take a further marked basis vector
\(w\ne z_{j+1}\) in the next core.
Send \(z_j\) to \(z_{j+1}+\pi^jw\) and the other marked
basis vectors bijectively to the remaining marked basis
vectors other than \(z_{j+1},w\). Send unmarked vectors
to zero. In both cases the marked columns of \(X_j\),
together with \(v_j=\pi^{-j}(X_jz_j-z_{j+1})\), give an
integral basis of the next marked core. At \(j=0\) this last
vector is \(-z_1\).

In \(h_{j+1}^{-1}Ah_j\) the extra column zero is
\((v_j,\pi e_0)\). In each extra column \(i\ge1\) the
upper free vector becomes
\(Y_i'=\pi^{-j}(Y_i-z_{j+1}b_{0i})\). The lower free
entries become \(v_{li}=\pi b_{li}\); each original lower
diagonal one becomes \(\pi\). The Jacobian from these new
variables back to the old measure is
\(q^{c(c+1)/2-jn_{j+1}c}\), by triangular change of
variables. Integrality requires the new variables to be
integral. The \(n_{j+1}-a_{j+1}\) unmarked coordinates
in each \(Y_i'\) must vanish modulo \(\pi\), at volume cost
\(q^{-(n_{j+1}-a_{j+1})c}\).

Reduce the marked determinant. Its first \(a_j+1\) columns
form the basis just found in the upper rows, with zero lower
entries. The lower square on columns \(1,\ldots,c\) is upper
triangular with diagonal \(\bar v_{l,l+1}\), \(0\le l<c\).
The determinant is a fixed residue unit \(u_j\) times the
product of these \(c\) entries. The additive character
is the product of \(\theta(\bar v_{l,l+1})\) on the same
entries. Each gives the one-dimensional integral
\(q^{-1}G(\chi)\), and all remaining unrestricted integral
coordinates have volume one. Thus before the two exterior
scalars this edge contributes
\[
   \chi(u_j)G(\chi)^c\,
   q^{\,\binom{c}{2}-jn_{j+1}c-(n_{j+1}-a_{j+1})c}.
\]

For its unipotent volume write a group element of \(N\) as
\(\left(\begin{smallmatrix}I&B\\0&I+H\end{smallmatrix}\right)\),
with \(H\) strictly upper triangular and \(B_0=0\).
Product measure in the free entries is Haar. Left translation
is affine triangular with unit diagonal in a column ordering.
Conjugating by \(h_{j+1}^{-1}\) gives the upper block
\(\pi^{-j-1}(B-z_{j+1}e_0^tH)\) and lower block \(I+H\).
Membership in \(J\) requires integral \(H\). For each
\(1\le i<c\), the translated upper column lies in
\(\pi^{j+1}O^{n_{j+1}}\) before the division, with a further
factor \(\pi\) required in its unmarked rows. These conditions
are sufficient because \(J\) is the inverse image of the
marked plane stabilizer. The reciprocal relative volume is
therefore exactly
\[
       q^{(c-1)((j+1)n_{j+1}+n_{j+1}-a_{j+1})}.
\]

At the final split pair the source matrix gives
\([X_s\mid \pi^{-s}(X_s z_s+y)]\). Choose \(X_s\) to
identify its marked core with the first \(r-1\) marked
target vectors and to vanish on the other columns. Use an
integral covector \(l_0\) vanishing on this image whose
residue restricts to the oriented determinant form \(L\)
of Theorem~\ref{thm:lastcolumn321} on the marked target.
Take frequency \(\pi^{-s-1}l_0\) and top test with identity
marked block and all other entries zero. Substitution
\(y=-X_sz_s+\pi^s v\) contributes \(q^{-sn_{s+1}}\), with
unit additive phase since \(l_0X_s=0\). The last-column
formula contributes \(q^{r-1-n_{s+1}}G(\chi)\) in the
oriented frames. Its dependence on the last frequency is
the stated inverse character. This calculation also covers
\(s=0\) with its empty set of interior edges.

The character factors in \(\prod C_j\) telescope to
\(\eta\), since the common \(\delta\) cancels at every
intermediate determinant. The absolute determinant factors
give \(\prod_{j=1}^s|\det h_j|^t\), since
\(D_{j+1}-D_{j-1}=2t\). Their power of \(q\) is
\(t\sum_{j=1}^s j(s+1-j)\). It cancels the sum of
\(-(n_{j+1}-a_{j+1})c\) in the interior integrals.
Adding the remaining interior exponents gives
\[
 \binom{s+1}{3}
  +\sum_{j=0}^{s-1}(s-1-2j)n_{j+1}
  +t\sum_{j=0}^{s-1}(s-j-1)(j+1)=0 .
\]
Indeed the last sum without \(t\) is \(\binom{s+1}{3}\).
In the middle one substitute \(n_{j+1}=k_0+(j+1)(t+1)\);
the constant coefficient sums to zero and the linear term
gives \(-(t+1)\binom{s+1}{3}\). This leaves exactly the
last power of \(q\) in the theorem. The units \(u_j\) and
any fixed basis signs give \(u\). Nonvanishing of \(G(\chi)\),
including \(G(1)=-1\), was proved in
Theorem~\ref{thm:lastcolumn321}. By definition a successor
double coset is supported if the interior entry is nonzero
for some core position. The computed entries establish this
at each step starting with the required identity.
\end{proof}

Products over the distinct split places give compatible
nonzero local entries for a fixed tuple. The global period
also integrates the core positions, sums over paths, and
projects to its prescribed row. The calculation retains a
single path before those operations. Its affine matrices,
phase and relative volume are the ones in
\srcloc{Development/AbelianVarietyTheory/Definitions/Foundations/SimpleSum}{53--68, 91--93}
and \srcloc{Development/AbelianVarietyTheory/Definitions/TwistedInnerProducts/StandardOrderExtendedBasisApplyLowBlock}{718--732}.

\subsection{Integration over the flag columns at a ramified twist}\label{nav:VI-66}
Fix a split twist place of one of the tuples in
Section~\ref{sec:missing}. We use the plain frames, Haar
measures and conductor of Theorem~\ref{thm:splitpath322}.
Assume that the residue character \(\chi:K^\times\to\CC^\times\)
is nontrivial. Let \(0\le j<s\), \(c=s-j\), \(n=n_{j+1}\)
and \(a=a_{j+1}=r-c\). Use the representative \(h_j\) of
that theorem at the source. At the target consider a matrix
\[
 h=H k_1,\qquad
 H=\begin{pmatrix}I&\pi^{-e_0}z f_0^t\\
                 0&\operatorname{diag}(\pi^{-e_l})_{0\le l<c}
   \end{pmatrix},\qquad k_1\in{\rm GL}_{n+c}(O).
\]
Here \(z\in k^n\), the \(e_l\) are arbitrary integers and
\(f_l\) are the flag basis vectors. In the coordinates of
\(H\), let \(W\) be the image under \(\bar k_1\) of the
marked target plane. Give \(W\) its transported determinant
form. Put \(H_i^0=K^n\oplus\langle f_0,\ldots,f_{i-1}\rangle\)
for \(0\le i\le c\). The entry below is the integral over the
free columns of the affine matrix \(A\) in
Theorem~\ref{thm:splitpath322}, before the Weil scalar and
the reciprocal unipotent volume.

\begin{theorem}[Support and value of the flag integral]\leavevmode
\label{thm:ramifiedflag324}
At a core matrix \(X\) the entry vanishes unless all
\(e_l=j+1\), \(X\) and \(v=\pi^{-j}(Xz_j-z)\) are integral,
and \(W\) projects onto the full flag space \(K^c\).
A further necessary condition is that the reduced marked
columns of \(X\), followed by \(\bar v\), form a basis of
\(W\cap K^n\). These conditions together are sufficient.
Choose \(w_i\in W\) lifting \(f_{i-1}\) under projection,
for \(1\le i\le c\). Let \(u\) be the determinant of the
specified core basis followed by the \(w_i\), in the
transported frames and with the fixed marked column ordering.
The integral equals
\[
     \chi(u)\,G(\chi)^c\,
       q^{\,\binom c2-jnc-(n-a)c}.
\]
\end{theorem}
\begin{proof}
Integrality of \(k_1^{-1}H^{-1}Ah_j\) is equivalent to
integrality before multiplication by \(k_1^{-1}\). The core
columns there are \((X,0)\). The first extra column is
\((v,\pi^{d_0} f_0)\), where \(d_l=e_l-j\).
In extra column \(i\ge1\) the upper vector is
\(\pi^{-j}(Y_i-z b_{0i})\). Its free lower entries are
\(v_{li}=\pi^{d_l} b_{li}\) for \(l<i\). The fixed entry
in row \(i<c\) is \(\pi^{d_i}\). Thus integrality requires
\(d_l\ge0\) for every \(l\). The new free variables range
independently over integral vectors. The marked test now
depends only on their residues and uses the determinant
form on \(W\). If \(d_l>1\), integration in
\(v_{l,l+1}\) over any coset of \(\pi O\) gives zero.
Indeed the phase has factor \(\psi(-\pi^{-d_l}v_{l,l+1})\),
which is nontrivial on this subgroup. Hence every \(d_l\)
is zero or one.

Integrate the remaining columns in the order \(c,c-1,\ldots,1\).
Write \(W_i=W\cap H_i^0\). Initially the last column ranges
over all of \(H_c^0\) after reduction. At step \(i\), the
integrand from the later columns is a constant times the
character of the determinant of the columns still present,
required to give a basis of \(W_i\). The current column
ranges over all of \(H_i^0\). This description holds at
\(i=c\) by definition. Fix the earlier columns. If they
do not span a hyperplane in \(W_i\), their contribution
vanishes. Otherwise the determinant, as a function of the
current column in \(W_i\), is a nonzero linear functional
with this hyperplane as kernel.

If \(d_{i-1}=0\), the additive factor of this column is
one. Summing over its allowable vectors gives zero, since
each nonzero determinant occurs equally often and
\(\sum_{t\in K^\times}\chi(t)=0\). Consequently a surviving
step has \(d_{i-1}=1\). The earlier columns then belong to
\(H_{i-1}^0\), since the only possible fixed entry in row
\(i-1\) has zero reduction. If \(W_i\) projects trivially
on that row the same character sum is zero. Otherwise
\(W_{i-1}\) has codimension one in \(W_i\). The earlier
columns must form a basis of \(W_{i-1}\). Choose an element
of \(W_i\) with last coordinate one and use it to induce
the determinant form on \(W_{i-1}\). At coordinate \(t\ne0\)
there are \(|W_{i-1}|\) choices for the current column,
the determinant acquires a factor \(t\), and the phase is
\(\theta(t)=\psi(-\pi^{-1}\widetilde t)\).
The normalized residue integral therefore multiplies the
remaining determinant character by
\(|W_{i-1}|\,G(\chi)/|H_i^0|\).
If \(i>1\), the next column ranges over all of \(H_{i-1}^0\)
because its fixed entry \(\pi^{d_{i-1}}\) reduces to zero.
This proves the induction assertion.

All projections in the flag have rank one at their
successive quotients exactly when \(W\) projects onto
\(K^c\). In that event \(\dim W_{i-1}=a+i-1\).
The remaining determinant is the one in the statement.
Adding an element of \(W\cap K^n\) to any chosen lift
leaves it unchanged. The product of the residue factors
without this determinant is
\((q^{a-1-n}G(\chi))^c\). The additive change of variables
back to the original Haar measure has Jacobian
\(q^{c(c+1)/2-jnc}\), since there are \(c(c+1)/2\)
lower free entries, each divided by \(\pi\), and \(c\)
upper vectors multiplied by \(\pi^j\). This proves the
formula. The Gauss sum has absolute value \(\sqrt q\) by
Theorem~\ref{thm:lastcolumn321}, proving sufficiency.
\end{proof}

\subsection{Integral cost of the prescribed scalar orbit projections}\label{nav:VI-67}
Use the common split prime \(T\), the words \(c\), and the
refined finite column spaces of Corollary~\ref{cor:rank111}.
Write \(f\) for the number of places above \(T\) and
\(b=T-1\). At each of those places the scalar group acts
freely on the nonzero residue covectors. On a product orbit
the integral coefficient arrays are therefore copies of
\((\mathcal O^b)^{\otimes f}\) with the same underlying
class coefficient module. This follows also by transporting
a base covector as in Theorem~\ref{thm:dyad234}. Tensoring
with the class module counts one copy per coefficient rank.
Let \(B_c=R_c/b^f\) denote the total number of copies
counted with these ranks. On these columns the class trace
and projection arguments of Theorems~\ref{thm:classtrace326}
and~\ref{thm:projclass326} apply to their definite group.
Keep the normalized pairings with their specified level volumes.

\begin{theorem}[Elementary divisors of the zero-sum column pairing]\leavevmode
\label{thm:scalarcost326}
For the projection onto arrays of zero sum in each scalar
coordinate, the paired lattices have elementary divisors
\[
         b^j\quad\text{with multiplicity}\quad
             B_c\binom fj(b-2)^{f-j},\qquad 0\le j\le f .
\]
At a coefficient place these give exponents \(jv_{\mathfrak p}(b)\).
The globally normalized logarithmic index per retained rank
is exactly
\[
 \frac{\log|C_P|}{[K:\QQ]\,R_c^\circ}
     =\frac{f\log(T-1)}{T-2}\le \exp(-x^{1/4})
\]
for all sufficiently large \(x\) on the common intervals.
Here \(C_P\) is the finite integral quotient over \(\mathcal O_K\).
A direct summand of relative rank
\((1-1/(T-2))^f\) has unit elementary divisors at every
finite prime simultaneously. For the separate projection
onto one specified scalar character, each of the \(B_c\)
retained directions has elementary divisor \(b^f\) over
a coefficient ring containing the character values.
\end{theorem}
\begin{proof}
For one scalar orbit put \(M=\mathcal O^b\), use its
coordinate basis \(e_1,\ldots,e_b\), and let \(A\) be its
augmentation kernel, with basis \(a_i=e_i-e_b\), \(1\le i<b\).
Equal class measures on this free orbit give the orthogonal
projection \(P_1 z=z-(\sum_i z_i)\mathbf 1/b\).
Set \(v=P_1e_b\). Then \(P_1M=A+\mathcal O v\) and
\(b v=-\sum_{i<b}a_i\). It follows that
\(a_1,\ldots,a_{b-2},v\) form a basis of \(P_1M\) and that
replacing \(v\) by \(b v\) gives a basis of \(A\).
Both choices work over the rational integers.

On the product orbit \(P=P_1^{\otimes f}\). Splitting
\(M=A\oplus\mathcal O e_b\) in each coordinate shows that
the saturated zero-sum module is \(A^{\otimes f}\), since
the coordinate-sum kernel omits all direct summands using
\(e_b\) in that coordinate. The projection of the full
lattice is \((P_1M)^{\otimes f}\). Tensor the bases just
constructed. A tensor uses the last basis element in
exactly \(j\) coordinates in
\(\binom fj(b-2)^{f-j}\) ways. Its inclusion multiplier
is \(b^j\). Together with the class module this gives the
stated elementary divisors by Theorem~\ref{thm:projclass326}.
The tensors using only the \(a_i\) for \(i\le b-2\) give
the claimed direct summand, including for a projective
base coefficient lattice.

There are \(B_c(b-1)^f=R_c^\circ\) retained directions.
Differentiating the binomial identity gives
\(\sum_j j\binom fj(b-2)^{f-j}=f(b-1)^{f-1}\).
The absolute norm of \(b\mathcal O_K\) is \(b^{[K:\QQ]}\).
These equalities prove the mean index formula.
In Theorem~\ref{thm:common111} the common prime satisfies
\(\log T\ge(16-o(1))m x^\eps\), with
\(m\asymp x^{1/2-\eps/16}\), and the field degree \(f\)
is polynomially bounded in \(x\). Since \(\log y/y\)
is decreasing for \(y>e\), these estimates give the
displayed inequality uniformly for the selected words.

Finally the character projector on the regular product
orbit sends a coordinate vector to \(b^{-f}\) times its
character vector, up to a unit. That vector has unit
identity coordinate. It spans the saturated integral
character line, while the projected full lattice is the
same line divided by \(b^f\). This yields the last assertion.
\end{proof}

The exponential bound concerns the actual scalar-orbit
projection on the column side. It also bounds the loss
per row in any exterior degree \(k\ge
R_c^\circ\exp(-x^{\eps/8})\) by
\(\exp(-x^{1/4}/2)\) for all sufficiently large \(x\),
after summing the finite valuations with their norm weights
and dividing by \([K:\QQ]\). Indeed the largest \(k\)
nonnegative exponents at each prime have sum at most the
total, and \(R_c^\circ/k\le\exp(x^{\eps/8})\).
The size of the projection
onto the designated automorphic row is a separate term in
the period comparison.

\subsection{The two adjacent level volumes}\label{nav:VIII-65}
We compute the last two split local factors from the tests of
Theorem~\ref{thm:lastcolumn321}. Work first at one auxiliary
twist place \(v\). Let \(J\subset K_0={\rm GL}_d(O)\) be the
middle parahoric. In the marked basis it preserves the subspace
\(U\) modulo \(\pi\). Its character is the tame determinant
character \(\delta\) on the lower diagonal block. We use Haar
volume one on \(K_0\).

Write \(\Psi(X,y)=c_h\Phi([X\mid y]h)\) with the fixed
frame convention of that theorem. Take the two scalar
lattices \(M\subset S\subset k^d\) used to record the last
variable. Thus \(\Psi\) is supported in \(S\) as a function
of \(y\), and is invariant by translations by \(M\).
Such bounds exist for each finite collection of the marked
tests. Indeed their common compact support is contained in
a scalar box and compactness gives a common neighborhood
of translation invariance for the finitely many locally
constant functions. Both lattices are preserved by \(K_0\).
Write \(Y=S/M\) and define \(L^\perp=\{z:\psi(zy)=1
\text{ for all }y\in L\}\) for a lattice \(L\). For
\(a\in Y\) put
\[
 P_a(X)=\Psi(X,a),\qquad
 T_a(B,z)=\Phi^+(B)\,{\rm vol}(M)
                \mathbf1_{M^\perp}(z)\psi(-za).
\]
The character in this formula is used only on \(M^\perp\),
where it is independent of the lift of \(a\).

\begin{theorem}[Coset weights and original Haar measures]\leavevmode
\label{thm:adjacent321}
The coupled last kernel equals \(\sum_{a\in Y}P_aT_a\)
pointwise. Its contraction over the middle definite group
can be formed with one summand for each \(J\)-orbit, with
coefficient \(|Ja|\). The other variables, including the
two outer tests, may be kept fixed. The same formula holds
at all auxiliary twist places together, with products of
the orbit sizes.

At each vertex \(i\) let \(J_{i,v}\subset K_{i,v}\) be the
tuple level inside the base level at its auxiliary places.
In the transported frames these are the parahoric and the
integral general linear group. Put
\[
 I_i(t)=\prod_{v\in Q(t)}[K_{i,v}:J_{i,v}],
 \qquad a_j(t)=H\kappa_{j+1}\,{\rm vol}(K_{j+1,t}).
\]
Here \(K_{i,t}\) is the global finite level at vertex \(i\),
\(H\) is the order of the common finite character group, and
\(\kappa_i\) is the fixed measure scalar of that vertex.
Keep the same Haar measure at every tuple, including the
base tuple \(0\) with \(Q(0)=\varnothing\).
For an adjacent pair let \(q_j(t)\) be the product of the
two local level volumes over \(Q(t)\), each measured relative
to its base level. Then
\[
 \frac{q_j(t)}{a_j(t)a_{j+1}(t)}
             = \frac{1}{a_j(0)a_{j+1}(0)} .
\]
The equality holds also at coefficient primes dividing the
orders of the residue general linear groups.
\end{theorem}
\begin{proof}
Partition the defining integral in \(y\) into the disjoint
sets \(a+M\). On each such set \(\Psi\) is constant and
the integral of its character is
\({\rm vol}(M)\psi(-za)\mathbf1_{M^\perp}(z)\), by translation.
This proves the pointwise identity with the full top test.
In particular the support and the phases of a coset are
kept in both factors. The unnormalized character table
between \(Y\) and \(M^\perp/S^\perp\) has inverse its
conjugate transpose divided by \(|Y|\). Indeed the pairing
is perfect by the scalar lattice description, and distinct
columns are orthogonal by character summation. Thus refining
the additive cosets before summing introduces only powers
of the residue characteristic into the denominators of this
Fourier inversion.

For \(g\in J\), let \(g_U\) act on \(U\) after reduction
and put \(c(g)=\delta(\det g_U)^{-1}\). The parahoric test
at the first edge satisfies
\(\Psi(gX,gy)=c(g)\Psi(X,y)\), since \(g\) preserves the
integral target lattice and multiplies the marked determinant
by \(\det g_U\). On the next edge it is the source which acts
by \(g\), whence \(\Phi^+(Bg)=c(g)\Phi^+(B)\).
The two statements use the same inverse tame character.
Unramified splitting factors are one on these determinant
units. From \(zga=z(ga)\) and preservation of \(M\) we get
\[
 P_{ga}(X)=c(g)P_a(g^{-1}X),\qquad
 T_{ga}(B,z)=c(g)^{-1}T_a(Bg,zg).
\]
These identities apply to arbitrary \(h\), since left
multiplication on the first test commutes with its right
source translation. The scalar phases cancel in the product.

Give \(J\) probability measure \(dk\). The middle average is
\[
       \int_J P_a(kX)T_a(Bk^{-1},zk^{-1})\,dk .
\]
The last two identities leave this average unchanged when
\(a\) is replaced by \(ga\), by left translation of \(k\).
In the adelic definite contraction one may insert this
probability average by right translation of the intermediate
integration variable. The two adjacent Weil actions on that
variable have inverse changes of coordinates. Their scalar
splittings on \(J\) are inverse there because the common
tame character is the same and the absolute determinant
is one. Thus the inserted average is exactly the compact
average just computed in these frames. All other local
factors are fixed under this translation. Compactness of
the definite quotient and finiteness of the displayed sum
justify the interchange. It follows that the contraction is
constant on each orbit. Grouping its terms gives \(|Ja|\).
Distinct places act on distinct coordinates. Taking their
product gives the asserted simultaneous grouping, including
the empty product for the base tuple.

Let \(I=[K_0:J]\). Multiplicative Haar measure gives
\({\rm vol}(J)=1/I\) by the coset partition. In the convention
where this volume is placed outside the middle sum, the
weight of the representative is accordingly \(I|Ja|\).
At the last vertex its singleton index has weight \(I'\)
and local volume \(1/I'\). More explicitly a split local
index for a distinguished block of rank \(r\) in rank \(d\)
is the number of \(r\)-planes in \(K^d\). Reduction maps
\({\rm GL}_d(O)\) onto \({\rm GL}_d(K)\), by lifting a
matrix with unit determinant, and \(J\) is the inverse
image of a plane stabilizer. Counting ordered bases gives
the integer \(\prod_{i=0}^{r-1}(q^d-q^i)/(q^r-q^i)\).

Globally the levels agree outside \(Q(t)\), and the local
coordinates in its distinct factors can be chosen independently.
Hence \({\rm vol}(K_{i,t})={\rm vol}(K_{i,0})/I_i(t)\).
Transport by a common frame preserves the index and this
volume ratio. Rational stabilizers cause no correction in
passing to definite classes. Over a parent class with finite
stabilizer \(\Gamma\), the fine classes are the \(\Gamma\)-orbits
in the \(I_i(t)\)-element set of level cosets. Summing the
reciprocal orders of their stabilizers gives
\(I_i(t)/|\Gamma|\) by orbit--stabilizer, which is exactly
offset by the new level volume. Finally
\[
 a_j(t)=a_j(0)/I_{j+1}(t),\qquad
 q_j(t)=\bigl(I_{j+1}(t)I_{j+2}(t)\bigr)^{-1}.
\]
Substitution gives the claimed equality over the rationals
for the index factors and with the original measure scalars.
\end{proof}

The last identity concerns the two volume factors in the
seesaw quotient. The orbit sizes and the signed values
\(P_aT_a\) inside a period still enter its global integral.
The relevant normalizations are recorded in
\srcloc{Development/AbelianVarietyTheory/Definitions/LatticeFrames/CosInj}{110--150}
and \srcloc{Development/AbelianVarietyTheory/Definitions/LatticeFrames/NotBadPrimeBase}{103--107}.

\subsection{All intermediate double cosets at a ramified twist}\label{nav:VIII-66}
At a split twist place use the frames and kernels of
Theorem~\ref{thm:splitpath322}. At a vertex with core
dimension \(n\) and \(c\ge1\) flag coordinates write
\(J\subset{\rm GL}_{n+c}(O)\) for the inverse image of
the marked plane stabilizer. The group on the left of the
double quotient consists of the matrices
\[
 R=\left\{\begin{pmatrix}g&B\\0&T\end{pmatrix}:
   g\in{\rm GL}_{n}(k),\quad B f_0=0,\quad
   T\ \text{upper unitriangular}\right\}.
\]
For a fixed predecessor, a double coset is supported if
some representative gives a nonzero affine integral for
some core matrix. The initial representative is the identity.

\begin{theorem}[Uniqueness of all intermediate double cosets]\leavevmode
\label{thm:ramifiedcoset324}
If the tame splitting character is nontrivial on residue
units, there is exactly one supported successor double
coset at every interior step, beginning at the identity.
The matrices \(h_1,\ldots,h_s\) of
Theorem~\ref{thm:splitpath322} represent them. Thus the sum
over all interior double cosets at this place, in these
representatives, has the value of that single path as a
function of its core positions and terminal variables. Its
integral over the common position set and the terminal
frequency ball of Theorem~\ref{thm:splitmass323}, with the
compact averages there, is exactly the value in that theorem.
For an admissible global twist tuple this applies at every
place of its auxiliary twist set \(Q(t)\). Its whole
contribution from these places on the product of the stated
sets is the product of the local values, with empty product
one for the base tuple.
\end{theorem}
\begin{proof}
Right target multiplication by \(J\) changes the marked
test by a nonzero constant from its determinant character.
Left target multiplication by an element of \(R\) changes
the affine variables and the core position. Indeed, acting
on \(A\) on the left by its inverse fixes extra column zero,
sends \(X\) to \(g^{-1}X\), and makes an invertible affine
change of the other free columns. In the first lower
superdiagonal each changed entry differs from its old
value by a constant coming from the fixed diagonal.
The change of measure has positive constant Jacobian
and the additive phase changes by a constant of modulus
one. Acting on \(A\) on the right by the predecessor's
corresponding group changes \(X\) to \(Xg\), fixes extra
column zero and has the same properties. Right multiplication
of a predecessor representative by its parahoric again gives
a determinant character. The Weil factors never vanish.
Thus the existence of a nonzero entry is independent of
representatives on either side.

We describe representatives of the target quotient. For a
matrix with column lattice \(\Lambda\), choose a basis of
\(\Lambda\cap k^n\) and send it by the core group to the
standard basis. For \(0\le l<c\), intersect \(\Lambda\)
with the core and the first \(l+1\) flag lines. Its image
in the last of these lines is a fractional ideal
\(\pi^{-e_l}O\), since \(\Lambda\) is a full lattice.
Lift its generator. Successive subtraction of such lifts
expresses every lattice vector in the resulting basis.
Each lift, divided by its pivot, has the form \(f_l\)
plus a vector in the preceding coordinates. The upper
unitriangular flag operations and the free core columns
of \(R\) remove these earlier entries for \(l\ge1\):
working in increasing \(l\), subtract multiples of
coordinate \(l\) from the earlier coordinates. This
fixes all preceding basis columns. In column zero the
normalized lift is \((z,f_0)\). Hence the representative
has the form \(H k_1\) of Theorem~\ref{thm:ramifiedflag324}.
For this fixed column lattice the remaining right cosets
modulo \(J\) correspond exactly to the residue planes
\(W=\bar k_1 E\), where \(E\) is the marked plane in
lattice coordinates. Indeed two integral bases give the
same plane precisely when their ratio belongs to \(J\).

Suppose the predecessor is \(h_j\). The preceding theorem
forces \(e_l=j+1\) for every \(l\). It also forces \(W\)
to project onto the flags, with \(W\cap K^n\) spanned by
the marked columns of \(\bar X\) and
\(\bar v=\overline{\pi^{-j}(Xz_j-z)}\).
For \(j=0\) we have \(v=-z\), so \(z\) is primitive and
\(\bar z\in W\cap K^n\). For \(j>0\) integrality gives
\(\bar z=\bar X\bar z_j\). The marked columns of \(\bar X\)
are independent, and \(\bar z_j\) is a marked basis
vector. This again makes \(z\) primitive and puts its
reduction in \(W\cap K^n\).

An integral core basis change sends \(z\) to the chosen
\(z_{j+1}\), since a primitive vector extends to an
integral basis. Call this vector \(z_*\). For the resulting
lattice with \(e_l=j+1\) take \(T=I\). In lattice
coordinates a core action \(g\) has upper flag column zero
equal to \(\pi^{-j-1}(gz_*-z_*)\); the other upper flag
columns are \(\pi^{-j-1}B f_i\). The reductions of these
columns can be prescribed arbitrarily, independently of
any prescribed \(\bar g\) fixing \(\bar z_*\). To see this,
lift the other columns of \(\bar g\) in a basis starting
with \(z_*\) and choose \(gz_*=z_*+\pi^{j+1}u\) with
the desired \(\bar u\). This gives an integral invertible
\(g\). Choose \(B f_i\) for \(i>0\) by multiplying the
desired lifts by \(\pi^{j+1}\). The resulting element of
\(R\) stabilizes the lattice.

The stabilizer of \(\bar z_*\) in the residue core group
is transitive on \(a_{j+1}\)-planes containing it, by
extending bases beginning with that vector. Send \(W\cap K^n\)
to the standard marked core plane. Since \(W\) projects
onto every flag coordinate together, the independent
upper flag columns just constructed can then cancel the
core coordinates of lifts of every \(f_i\) in \(W\).
The resulting plane is the standard marked plane.
There is therefore at most one supported double coset.
The nonzero entry constructed in
Theorem~\ref{thm:splitpath322} supplies one at each step.
Induction gives uniqueness along the entire chain, including
the empty list when \(s=0\). The terminal orbit and its Haar
weight are evaluated in Theorem~\ref{thm:splitcompact323}.
This proves the assertions at one place.

For the last assertion, recall how the fields of the tuple
specify its auxiliary twists. At a finite place \(v\) of
the real subfield the tuple character kills the principal
units, and away from the allowed split places it kills all
local units. The support used for \(Q(t)\) consists of
the \(v\) where the restriction to local units is nontrivial,
choosing one prime of the CM field over each such \(v\).
At a split \(v\) the character is trivial on the diagonal
copy of the base field's multiplicative group. Its two
local characters there are therefore inverses. Nontriviality
on the product of the two unit groups implies nontriviality
on the units of either chosen factor. Since principal units
are in the kernel, the corresponding residue character and
its inverse \(\chi\) are nontrivial. This checks the
hypothesis simultaneously at every place of \(Q(t)\).
The integrals over their product are finite products under
product Haar measure. The base character is unramified
everywhere and has empty auxiliary support.
\end{proof}

The definitions of the auxiliary support and the affine
successor integral are in
\srcloc{Formalism/LocalToGlobal/Divisibility/Tuples}{50--67, 110--128},
\srcloc{Development/AbelianVarietyTheory/Definitions/TwistedInnerProducts/StandardOrderExtendedBasisApplyLowBlock}{718--732, 922--938}.
The projection onto the prescribed global automorphic row
still involves the global form and its arithmetic normalization.

\subsection{The exact integral loss in a projected terminal period}\label{nav:VIII-67}
At a finite coefficient place localize the terminal table
of Theorem~\ref{thm:classtrace326} over the discrete
valuation ring \(O\) with uniformizer \(\varpi\) and
fraction field \(K\). Write \(S,T\) for its integral row
and its lattice of trace tests, paired in the original
measure divided by the finite level volume. These are
mutual duals. Fix a \(K\)-rational idempotent \(P\) on
\(S_K\). This includes an orthogonal projection when
its matrix is defined over \(K\), as for the theta images
of Theorem~\ref{thm:jointimage318}. Write \(P^t\) for the
transpose on covectors. Define
\[
 S_P=S\cap P S_K,\qquad S_Q=S\cap(1-P)S_K,\qquad
 T_P=T\cap P^t T_K,\qquad C_P=S/(S_P+S_Q).
\]

\begin{theorem}[Elementary divisors of the projected period]\leavevmode
\label{thm:projclass326}
The nonunit elementary divisors of the normalized global
pairing on \(S_P\times T_P\) are exactly those of the finite
module \(C_P\). Write the exponents as
\(c_1\ge\cdots\ge c_d\ge0\), padding by zeros to
\(d=\dim_K P S_K\). Then \(T_P\) is the full dual of
\(P S\) and \(P S/S_P\simeq C_P\). In particular
\(\varpi^{c_1}\) clears the projection on the original
lattice with the least possible valuation when \(d>0\).

Let an operator with values in \(P S_K\) be tested on an
integral source lattice. If all its \(k\)-row determinants
against projected integral trace tests have valuation at
least \(b_k\), its exterior coefficients on the saturated
target \(S_P\) have valuations at least
\(b_k-\sum_{i=1}^k c_i\), for \(1\le k\le d\).
This is the sharp universal loss for that projection. For
the original Haar measure restore \(\operatorname{vol}(J)^k\)
in the tested determinant.
\end{theorem}
\begin{proof}
Projection induces an isomorphism \(S/S_Q\simeq P S\).
Its restriction to \(S_P\) is the identity. Both lattices
span the same space, and the quotient is \(C_P\).
A covector is fixed by \(P^t\) precisely when it annihilates
the complementary space. Such a covector is integral on
\(S\) precisely when it is integral on \(P S\), by the
identity \(\ell(s)=\ell(Ps)\). Theorem~\ref{thm:classtrace326}
identifies all of them with actual integral trace tests
in \(T_P\).

Choose a basis \(e_i\) of \(P S\) with
\(\varpi^{c_i}e_i\) a basis of \(S_P\), by the elementary
divisor theorem over \(O\). The dual basis to \(e_i\)
is a basis of \(T_P\). The tested pairing matrix is
therefore diagonal with entries \(\varpi^{c_i}\).
These formulas also give the least scalar that sends
\(P S\) into \(S_P\), equivalently into \(S\).

For a wedge indexed by \(i_1,\ldots,i_k\) the tested
coefficient is its coefficient in the saturated basis
multiplied by \(\varpi^{c_{i_1}+\cdots+c_{i_k}}\).
The largest such sum is \(c_1+\cdots+c_k\). The coordinate
wedge on the first \(k\) basis vectors attains the loss.
It can itself be realized as the wedge of images of \(k\)
source basis vectors. Finally each unnormalized scalar
test has one factor of the level volume by the class
trace formula, proving the assertion for determinants.
\end{proof}

The result applies to any arithmetic or analytic column
known in the stated rational projected row. At a prescribed
projection, the size of \(C_P\) measures a genuine integral
comparison that remains after the class and level factors
have been evaluated.

\subsection{The two finite type projections on the actual terminal classes}\label{nav:IV-63}
Fix the standard definite terminal group on \(E_x\), the coefficient
lattice and the embedding of Theorem~\ref{thm:terminalreading325}.
At each of the two completely split primes \(g_1,g_2\), distinguish
the \(r\)-dimensional block in each marked frame. Write \(f=[E_x^+:\QQ]\)
and let
\[
 \mathcal G=\prod_{i=1}^2\prod_{v\mid g_i,\ v\text{ in }E_x^+}
                \mathrm{GL}_r(\mathbf F_{g_i}),\qquad m_G=|\mathcal G|.
\]
Use the original levels on the other blocks and at all other
places. On each distinguished block first fix its full residue
marking, then allow the prescribed subgroup \(P\subseteq\mathcal G\)
with its finite character \(\omega\). Write \(S\) for the original
integral table at the latter level. A splitting number field for
the finite group and its characters may be used throughout.
Let \(\mathcal X\) be the set of irreducible characters in the
prescribed central idempotent at the two primes. The argument
works for any such set, including the one obtained by summing
the selected products of finite series idempotents. Put
\[
 k_{\mathcal X}(a)=\sum_{\chi\in\mathcal X}\chi(1)\chi(a^{-1}).
\]

\begin{theorem}[Class matrix and uniform integral loss for the finite types]\leavevmode
\label{thm:terminaltypes327}
The harmonic projection \(P_{\mathcal X}\) onto these finite
types, on the \(P,\omega\) level, is the same endomorphism in
decorated class coordinates and in the standard adelic row.
It is computed on each integral table \(w\) by
\[
       (m_G P_{\mathcal X} w)(x_i)
           =\sum_{a\in\mathcal G}k_{\mathcal X}(a)\,w(x_i\widehat a),
\]
where \(x_i\) ranges over the original class representatives
and \(\widehat a\) changes their distinguished markings by \(a\).
In particular \(m_G P_{\mathcal X}S\subseteq S\).

At a finite coefficient place choose bases for the integral
class components and let \(M\) be the matrix of this displayed
sum. For \(1\le l\le\operatorname{rank}M\), let \(d_l(M)\)
be the least valuation of an \(l\)-rowed minor. The largest
possible loss for an \(l\)-row determinant of the projected
integral trace pairing of Theorem~\ref{thm:projclass326} is
exactly \(l\,v(m_G)-d_l(M)\) in the level-volume normalization.
Over all finite coefficient places the normalized logarithmic
loss per exterior degree is at most
\[
 \log m_G \ \le\ f r^2\log(g_1g_2)
       \ll_{\psi,\varepsilon} f r^2 x^\varepsilon.
\]
This bound holds at every \(x\) on the same intervals, with
the same scalar for every word and twist tuple at that \(x\).
For every fixed \(A>0\) it is \(o_{\psi,\varepsilon,A}(b x^{-A})\)
there, where \(b=\exp(x^\varepsilon)\).
\end{theorem}
\begin{proof}
At either split prime a matrix in the distinguished general linear
group lifts to an integral invertible matrix there: lift its
entries, observing that its determinant is a unit. Extend the lift
by the identity on the other blocks. It normalizes the level with
the distinguished residue marking fixed. In the block parabolic
description the unipotent rectangle is preserved by the change
of basis, the distinguished diagonal block is the identity
modulo the prime, and the characters of the other diagonal
blocks are unchanged. Two lifts induce the same change of
marking and hence the same translation on these tables. Right
translations \(R_a w(x)=w(x\widehat a)\) give a representation
of \(\mathcal G\); they preserve the integral table on this finer
level. Indeed the underlying lattice and its coefficient fiber
are unchanged when its residue marking varies. They are unitary
in the adelic norm by right invariance of Haar.

The central sum \(\sum_a k_{\mathcal X}(a)R_a/m_G\) acts by a
scalar in an irreducible representation of character \(\eta\),
by Schur's lemma. Its trace is
\(\sum_{\chi\in\mathcal X}\chi(1)m_G^{-1}
   \sum_a\chi(a^{-1})\eta(a)\).
Character orthogonality gives \(\eta(1)\) if
\(\eta\in\mathcal X\), and zero otherwise. Thus this central
sum is the orthogonal type projector on the fine space. It
commutes with the \(P,\omega\) projection, so it restricts to
the given coarse level. Each \(k_{\mathcal X}(a)\) is integral
because eigenvalues of a finite order matrix are roots of
unity. The unnormalized sum belongs to the fine integral
table and once again satisfies the coarse transformation law.
That law cuts out the coarse integral table by equality on
the same marked lattices. No factor of \(|P|\) occurs here.

For a class representative \(x_i\), write
\(x_i\widehat a=\lambda x_j u\) with \(\lambda\) rational and
\(u\) in the original level. The summand is read as
\(k_{\mathcal X}(a)\omega_\xi(u)\rho_K(\lambda)w_j\), with
\(\omega_\xi\) the whole level character of the tuple.
The class stabilizer condition makes it independent of that
decomposition. Apply the compact coefficient functional of
Theorem~\ref{thm:terminalreading325} to each summand. The
result is the same finite sum of right translates of the
standard adelic form at every class and every archimedean
argument, with no adjustment of Haar measure.

There is also an explicit verification for the two coordinate
lists used at this step. Let \(L\) and \(V\) have as columns
any bases of the coarse character space and its retained type
space, respectively, in fine table coordinates. If \(H_0\)
is the positive definite fine Gram matrix, set
\[
 R=(L^*H_0L)^{-1}L^*H_0 V .
\]
The larger space contains the smaller. Since
\((L^*H_0L)^{-1}L^*H_0\) is a left inverse to \(L\),
\(LR=V\). This holds also for a zero space by the unique
empty map. Thus coefficients \(F\) in the retained basis
and \(RF\) in the coarse basis give the identical table on
every marked class. Inverse restriction of either list gives
the same standard adelic form, using the same coefficient
map at every tuple. This proves the coefficient equality
without imposing integrality on either chosen basis.

Finally the saturation of the image of \(M=m_GP_{\mathcal X}\)
is \(S\cap P_{\mathcal X}S_K\), locally over the coefficient
ring. In a basis of \(P_{\mathcal X}S\) the saturated
submodule is generated by multiples of valuations
\(c_1\ge\cdots\ge c_d\ge0\), as in
Theorem~\ref{thm:projclass326}. Hence the nonzero invariant
factors of \(M\), measured in its saturated target, have
valuations \(v(m_G)-c_i\) in increasing order.
The saturation is a direct summand locally, so the minimum
valuation in the full matrix of \(l\)-rowed minors is
\(l\,v(m_G)-\sum_{i\le l}c_i\). The class trace theorem
restores the original level volume once per row.
Summing \(v_w(m_G)\log\mathrm Nw/[K:\QQ]\) gives \(\log m_G\).
There are \(f\) factors above each split prime, and
\(|\mathrm{GL}_r(\mathbf F_g)|\le g^{r^2}\) by counting all
matrices. Proposition~\ref{prop:index90} gives
\(g_i\le2(kb)^8\), with \(\log k=O_{\psi,\varepsilon}(\log x)\).
The ranks and \(f\) grow at most polynomially in \(x\).
This proves the uniform bound and the last assertion.
\end{proof}

The two coordinate maps occur in
\srcloc{Formalism/LocalToGlobal/Divisibility/RowSideTie}{76}
and the finite type selector in
\srcloc{Formalism/SharpRestriction/LevelLetterCocycles/DataN}{1001}.
The equality above concerns both realizations of this finite
type projection. Comparing the full signed iterated theta
column with the independently specified Bessel period still
requires identifying their values on the common target.

\subsection{Integral separation of the marked finite types}\label{nav:typeblocks329}
Use the same field \(E_x\), pair of marked primes, coefficient maps and
Haar measure as in Theorem~\ref{thm:terminaltypes327}.
Let \(G=\mathcal G\) be its product of finite general linear groups,
and let \(S\) be the integral table at the prescribed level and
character. At a finite place of a splitting field for \(G\) and the
coefficients, let \(O\) be the local ring, \(K\) its field of fractions
and \(v\) the additive valuation normalized by \(v(\pi)=1\).
Take \(S\) over \(O\). For a conjugacy class \(D\) write
\(z_D=\sum_{g\in D}g\). On an irreducible character \(\chi\), the
class sum acts by the scalar
\[
 \lambda_\chi(D)=\frac{|D|\chi(g)}{\chi(1)}
 \qquad(g\in D).
\]
This follows by taking the trace of a central operator.
The scalar lies in \(O\) because it is an eigenvalue of the integral
matrix of multiplication by \(z_D\) on the regular representation,
and \(O\) is integrally closed. Group the characters together
when all their \(\lambda_\chi(D)\) have the same residues modulo
the maximal ideal. Denote one such class of characters by \(B\).
For a subset \(A\) of characters let \(e_A\) be the sum of its
ordinary central character idempotents.

\begin{theorem}[The integral cost within each residue class of characters]\leavevmode
\label{thm:typeblocks329}
Every \(e_B\) belongs to \(O[G]\). A sum \(e_A\) belongs to \(O[G]\)
exactly when \(A\) is a union of the classes \(B\).
In particular \(S=\bigoplus_B S_B\), where \(S_B=e_BS\),
also at coefficient primes dividing \(|G|\).

Let \(\mathcal X\) be the actual set of retained finite types, and
write \(S_K=S\otimes_O K\). The discrepancy of its two integral
lattices has the decomposition
\begin{equation}\label{eq:typeblocks329}
 \frac{e_{\mathcal X}S}{S\cap e_{\mathcal X}S_K}
 \simeq
 \bigoplus_{\substack{B:\ \varnothing\ne B\cap\mathcal X\ne B}}
 \frac{e_{B\cap\mathcal X}S_B}
      {S_B\cap e_{B\cap\mathcal X}(S_B\otimes_O K)} .
\end{equation}
For each class in the sum, every elementary divisor exponent is
at most
\[
 d_B=\left(v(|G|)-\min_{g\in G}
 v\!\left(\sum_{\chi\in B\cap\mathcal X}
                   \chi(1)\chi(g^{-1})\right)\right)_+.
\]
Here \(v(0)=+\infty\), and \(z_+=\max(z,0)\).
The same \(d_B\) is obtained using the omitted characters in \(B\).
In particular, complete residue classes contribute zero loss in
every exterior degree of the projected class pairing in the
original measure. This statement applies to every \(x\) and every
auxiliary tuple on the common interval.
\end{theorem}

\begin{proof}
For each \(\eta\notin B\), choose a conjugacy class \(D_\eta\)
on which its residue differs from that of \(B\). Such a class
exists by the definition of \(B\). In the center of \(O[G]\) form
\[
 T=\prod_{\eta\notin B}
       \bigl(z_{D_\eta}-\lambda_\eta(D_\eta)\,1\bigr).
\]
It vanishes on all irreducibles outside \(B\). On each
\(\chi\in B\) it acts by a unit \(t_\chi\), since each factor
has nonzero residue there. Thus
\[
 1-\prod_{\chi\in B}(1-t_\chi^{-1}T)
\]
is integral. It acts by one on \(B\) and by zero outside.
The regular representation in characteristic zero contains all
the irreducibles and is faithful, so this element is \(e_B\).
Conversely an integral central element has an expansion
\(\sum_D c_Dz_D\), \(c_D\in O\), because centrality makes its
coefficients constant on conjugacy classes. Its eigenvalues on
two characters in \(B\) must be congruent. For an idempotent these
eigenvalues are zero or one. The residues distinguish those two
values, proving the asserted converse for \(O[G]\).

Apply these elements first on the tables with full residue
markings, where each right translation is integral. Their
centrality makes them commute with every element of the level
subgroup. Hence they preserve its prescribed character. The integral
coarse tables are exactly those fine tables satisfying this
equivariance, by restriction to the original class representatives
in Theorem~\ref{thm:terminaltypes327}. Thus \(e_BS\subseteq S\).
Their sum is the identity and distinct idempotents have product
zero. This gives the integral direct sum, with no division by the
order of the level subgroup. For any idempotent \(e\), each element
of \(S\cap eS_K\) equals its own image under \(e\), so belongs to
\(eS\). Taking images and intersections on the direct summands
gives \eqref{eq:typeblocks329}. When a whole \(B\) is retained its
map on \(S_B\) is the identity. When it is omitted its image is zero.

For \(C=B\cap\mathcal X\), the coefficient of \(g\) in \(e_C\)
is \(|G|^{-1}\sum_{\chi\in C}\chi(1)\chi(g^{-1})\), by character
orthogonality. Consequently \(d_B\) is precisely the least
nonnegative exponent with \(\pi^{d_B}e_C\in O[G]\).
Since \(e_B\) is integral and \(e_{B\setminus C}=e_B-e_C\),
the same test, at each exponent, works for \(B\setminus C\).
The integral multiple preserves the coarse table by the same
equivariance argument and sends it into \(S_B\cap e_C S_K\).
It therefore annihilates the quotient in the corresponding
summand. Finally the elementary divisors of a direct sum over
a discrete valuation ring are obtained by putting the two lists
together and ordering them. The largest loss in exterior degree
\(l\) is the sum of the largest \(l\) exponents, padded with zeros,
by Theorem~\ref{thm:projclass326}. The original Haar and stabilizer
weights enter that pairing just as in Theorem~\ref{thm:classtrace326}.
\end{proof}

This identifies the finite types that can contribute an integral
denominator to the prescribed projection. The equality with the
independently constructed Bessel functional still requires the
comparison of their coefficients on the surviving classes.

\clearpage
\section{References}\label{nav:bibliography-XXIII}
\begingroup\def\section*#1{}\begin{thebibliography}{99}\small
\bibitem{XXIII-XVIII-fulton351} W. Fulton, \emph{Young Tableaux}, London Mathematical Society Student Texts 35, Cambridge University Press, 1997, Chapters 1 and 8. \href{https://www.cambridge.org/core/books/young-tableaux/A7570B10D82AE7233E25E5D6F70A07B6}{Monograph}.
\bibitem{XXIII-XVIII-IV-kudla29} S. S. Kudla, \emph{Notes on the local theta correspondence}, European School of Group Theory, 1996. \href{https://www.math.toronto.edu/skudla/castle.pdf}{Chapter III, Section 6, pp.~53--54}.
\bibitem{XXIII-XVIII-IV-minima32} J.-H. Evertse, \emph{Notes on Diophantine Approximation}, Chapter 6, Theorem 6.3.3. \href{https://www.math.leidenuniv.nl/~evertsejh/Minkowski.pdf}{Geometry of numbers}.
\bibitem{XXIII-XVIII-liu328} Y. Liu, \emph{Theta correspondence for almost unramified representations of unitary groups}, 2021. \href{https://arxiv.org/abs/2105.14223}{Preprint}.
\bibitem{XXIII-XVIII-ichino328} A. Ichino, \emph{Theta lifting for discrete series representations of real unitary groups}, 2020. \href{https://arxiv.org/abs/2002.09148}{Preprint}.
\end{thebibliography}\endgroup
\end{document}
